% concatenated from Higher_intertwings_Prsimtou.tex

\documentclass[10pt]{article}

\usepackage{fullpage, amsmath}
\usepackage{amsthm,amsfonts,url,amssymb}
\usepackage{mathrsfs}
\usepackage{amscd}
\usepackage{color}
\usepackage{chemarrow}
\usepackage{pictexwd,dcpic}
\usepackage[textwidth=18cm,centering]{geometry}
\usepackage{extarrows}
\usepackage{enumitem}
\usepackage{lipsum}
\usepackage[all,cmtip]{xy}
\usepackage{leftidx}
\usepackage[version=3]{mhchem}
\usepackage{makeidx}
\usepackage{bm} 
\usepackage{stmaryrd}
\usepackage[linkcolor=red,anchorcolor=green,citecolor=blue]{hyperref}

\makeindex
\topmargin=0.23in
\oddsidemargin=0.15in
\evensidemargin=0.15in
\textwidth=6.4in
\textheight=9.0in
\setlength{\parskip}{2mm}


\newtheorem{thm}{Theorem}[section]
\newtheorem{cor}[thm]{Corollary}
\newtheorem{lem}[thm]{Lemma}
\newtheorem{pro}[thm]{Proposition}
\newtheorem{dfn}[thm]{Definition}
\newtheorem{hypothesis}[thm]{Hypothesis}
\newtheorem{rmk}[thm]{Remark}
\newtheorem{notation}[thm]{Notation}
\newtheorem{exmp}[thm]{Example}
\newtheorem{conjecture}[thm]{Conjecture}
\newtheorem{condition}[thm]{Condition}
\newtheorem{expect}[thm]{Expectation}

\newcommand{\hooklongrightarrow}{\lhook\joinrel\longrightarrow}
\newcommand{\twoheadlongrightarrow}{\relbar\joinrel\twoheadrightarrow}

\DeclareMathOperator{\Aut}{Aut}
\DeclareMathOperator{\Cl}{Cl}
\DeclareMathOperator{\Cov}{Cov}
\DeclareMathOperator{\Det}{Det}
\DeclareMathOperator{\Div}{Div}
\DeclareMathOperator{\End}{End}
\DeclareMathOperator{\Et}{\acute{E}t}
\DeclareMathOperator{\Ext}{Ext}
\DeclareMathOperator{\Frob}{Frob}
\DeclareMathOperator{\Gal}{Gal}
\DeclareMathOperator{\GL}{GL}
\DeclareMathOperator{\Grass}{Grass}
\DeclareMathOperator{\Koh}{H}
\DeclareMathOperator{\CKoh}{\check H}
\DeclareMathOperator{\Hom}{Hom}
\newcommand{\CHom}{{\cal H}om}
\DeclareMathOperator{\Id}{Id}
\DeclareMathOperator{\Isom}{Isom}
\DeclareMathOperator{\Lie}{Lie}
\DeclareMathOperator{\Mod}{Mod}
\DeclareMathOperator{\Nor}{Nor}
\DeclareMathOperator{\NS}{NS}
\DeclareMathOperator{\Pic}{Pic}
\DeclareMathOperator{\PGL}{PGL}
\DeclareMathOperator{\Proj}{Proj}
\DeclareMathOperator{\Quot}{Frac}
\newcommand{\Rkoh}{{\rm R}}
\DeclareMathOperator{\Rep}{\mathrm{Rep}}
\DeclareMathOperator{\Res}{Res}
\DeclareMathOperator{\Rig}{Rig}
\DeclareMathOperator{\SL}{SL}
\DeclareMathOperator{\Spm}{Sp}
\DeclareMathOperator{\Spec}{Spec}
\DeclareMathOperator{\Spf}{Spf}
\DeclareMathOperator{\Stab}{Stab}
\DeclareMathOperator{\Supp}{Supp}
\DeclareMathOperator{\Tor}{Tor}
\DeclareMathOperator{\Tr}{Tr}
\DeclareMathOperator{\Var}{V}
\newcommand{\pdr}{{\rm pdR}}

\newcommand{\ad}{{\rm ad}}
\newcommand{\alg}{{\rm alg}}
\newcommand{\an}{{\rm la}}
\newcommand{\ccyc}{{\epsilon}}
\DeclareMathOperator{\charakt}{char}
\DeclareMathOperator{\card}{card}
\DeclareMathOperator{\codim}{codim}
\DeclareMathOperator{\coker}{coker}
%\newcommand{\con}{{\rm con}}
\newcommand{\cris}{{\rm cris}}
%\DeclareMathOperator{\deg}{deg}
\DeclareMathOperator{\di}{div}
\DeclareMathOperator{\diag}{diag}
\newcommand{\dR}{{\rm dR}}
\newcommand{\et}{{\rm\acute{e}t}}
%\DeclareMathOperator{\gcd}{gcd}
\DeclareMathOperator{\ggT}{ggT}
\DeclareMathOperator{\grad}{grad}
\DeclareMathOperator{\id}{\,id}
\DeclareMathOperator{\im}{Im}

%\DeclareMathOperator{\ker}{ker}
\DeclareMathOperator{\length}{length}
\renewcommand{\mod}{{\rm\,mod\,}}
\DeclareMathOperator{\ord}{ord}
\newcommand{\perf}{{\rm perf}}
\DeclareMathOperator{\rank}{rank}
\newcommand{\red}{{\rm red}}
\newcommand{\recc}{{\rm rec}}
\DeclareMathOperator{\reg}{reg}
\newcommand{\rig}{{\rm rig}}
\DeclareMathOperator{\rk}{rk}
\newcommand{\sep}{{\rm sep}}
\newcommand{\sh}{{\rm sh}}
\newcommand{\temp}{{\rm temp}}
\newcommand{\topol}{{\rm top}}
\newcommand{\tors}{{\rm tors}}
\newcommand{\norm}{{\rm Norm}}
\newcommand{\gal}{{\rm Gal}}
\newcommand{\pcr}{{\rm pcr}}
\newcommand{\ur}{{\rm ur}}
\newcommand{\spa}{{\rm Spa}}
\newcommand{\criss}{{\rm cris}}
\newcommand{\stt}{{\rm st}}
\newcommand{\fil}{{\rm Fil}}
\newcommand{\proj}{{\rm Proj}}
\newcommand{\ind}{{\rm Ind}}
\newcommand{\homo}{{\rm Hom}}
\newcommand{\EndO}{{\rm End}}
\newcommand{\ext}{{\rm Ext}}
\newcommand{\GLN}{{\rm GL}}
\newcommand{\PGLN}{{\rm PGL}}
\newcommand{\leg}{{\rm lg}}
\newcommand{\ana}{{\rm an}}
\newcommand{\st}{{\rm St}}
\newcommand{\pol}{{\rm pol}}
\newcommand{\op}{{\overline{{\mathbf{P}}}}}
\newcommand{\lrr}{{\langle r\rangle}}
\newcommand{\opp}{{\overline{{\mathbf{P}}}}}
\newcommand{\ob}{{\overline{{\mathbf{B}}}}}
\newcommand{\on}{{\overline{{\mathbf{N}}}}}
\DeclareMathOperator{\weight}{wt}
\newcommand{\cEo}{{\rm \overline{\cE}}}
\newcommand{\unr}{{\rm unr}}
\newcommand{\dett}{{\rm det}}
%\newcommand{\bI}{{\mathbb{I}}}
\newcommand{\hH}{{\mathrm{H}}}
\newcommand{\regu}{{\mathrm{reg}}}
\newcommand{\ra}{{\longrightarrow}}
\DeclareMathOperator{\val}{val}
%%%%%      Greek letters      %%%%%%%%

\newcommand{\soc}{{\rm soc}}
\newcommand{\opki}{{\overline{\mathcal{P}}_{\Delta_{k,i}}}}
\newcommand{\jsm}{{\mathcal{J}^{\infty}}}
\newcommand{\cten}{{\rm \widehat{\otimes}_E}}
\newcommand{\ten}{{{\otimes}_E}}
\newcommand{\Modo}{{\textrm{mod}}}
\newcommand{\fge}{{\rm fg}}
\newcommand{\adm}{{\rm ad}}
\newcommand{\sm}{{\rm sm}}
\newcommand{\spr}{{\rm Sp}}
\newcommand{\spf}{{\rm Spf}}
\newcommand{\df}{{\mathrm{DF}}}
\newcommand{\wdre}{{\textbf{r}}}
\newcommand{\padelc}{{(\Delta_{\Dpik}^z)}}
\newcommand{\padel}{{(\Delta_{\Dpik}^0)}}
\newcommand{\rec}{{\rm rec}_L}
\newcommand{\lalg}{{\rm lalg}}
\newcommand{\lra}{\longrightarrow}
\newcommand{\obp}{\overline{\mathbb{Q}}_p^\times}
\newcommand{\obpo}{\overline{\mathbb{Q}}_p}
\newcommand{\EX}{{E^\times}}
\newcommand{\LX}{{L^\times}}
\DeclareMathOperator{\usl}{\mathfrak sl}
\DeclareMathOperator{\ugl}{\mathfrak gl}
\DeclareMathOperator{\Ima}{Im}
\newcommand{\MSpec}{{\rm MaxSpec}}
\newcommand{\unra}{{\rm unr}}
\newcommand{\lamoverx}{{L({\bm\lambda}_\bh)}}
\newcommand{\recp}{{{\rm rec}_p}}
\newcommand{\univer}{{\rm univ}}
\newcommand{\fsl}{{\rm fs}}
\newcommand{\crys}{{\rm crys}}
\newcommand{\cpvar}{{\mathbf{{C}}_p}}
\newcommand{\lM}{{\overline{M}}}
\newcommand{\lL}{{\overline{L}}}
\newcommand{\cyc}{{\rm cyc}}
\newcommand{\Art}{{\rm Art}}
\newcommand{\tr}{{\rm tr}}
\newcommand{\Ker}{{\rm Ker}}
\newcommand{\ab}{{\rm ab}}
\newcommand{\wt}{{\rm wt}}
\newcommand{\semi}{{\rm st}}
\newcommand{\pr}{{\rm pr}}
\newcommand{\dr}{{\rm dR}}
\newcommand{\alge}{{\rm alg}}
\newcommand{\res}{{\rm Res}}
\newcommand{\tee}{{{\otimes}_{\cR_{E,L}}}}
\newcommand{\te}{{\otimes_E}}
\newcommand{\wte}{{\widehat{\otimes}}}
\newcommand{\tene}{{\widehat{\otimes}_{E}}}
\newcommand{\tenp}{{\widehat{\otimes}_{E,\pi}}}
\newcommand{\teni}{{\widehat{\otimes}_{E,\iota}}}
\newcommand{\undelram}{{\delta^0}}
\newcommand{\undel}{{\delta}}
\newcommand{\gr}{{\rm gr}}
\newcommand{\hpi}{{{\mathbf{h}}}}
\newcommand{\Dpik}{{\mathbf{D}}}
\newcommand{\pdR}{{\rm pdR}}
\newcommand{\lrrtwo}{{\rm \langle 2\rangle}}
\newcommand{\ag}{{\rm Ad}}


\newcommand{\sbanpik}{{\big(\mathrm{Spec}\hspace{2pt}\FZ_{\Omega_{S_0}}\big)^{\mathrm{rig}}}}

\newcommand{\sbanpiku}{{\big(\mathrm{Spec}\hspace{2pt}\FZ_{\Omega_{S^u_0}}\big)^{\mathrm{rig}}}}
\newcommand{\sbanpi}{{\left(\Spec\mathfrak{Z}_{\Omega_{r}}\right)^{\mathrm{rig}}}}


\newcommand{\rigchu}{{\mathcal{Z}_{\bL_{S^u_0},\mathcal{O}_L}}}
\newcommand{\rigchlu}{{\mathcal{Z}_{\bL_{S^u_0},L}}}



\newcommand{\vg}{{(\varphi,\Gamma)}}
\newcommand{\wdr}{{W^+_{\mathrm{dR}}}}
\newcommand{\bdr}{{B^+_{\mathrm{dR}}}}
\newcommand{\univ}{{\rm univ}}
\newcommand{\loc}{{\rm loc}}
\newcommand{\ct}{{\rm ct}}
\newcommand{\frk}{{\rm rk}}
\newcommand{\lp}{{\rm lp}}
\newcommand{\WD}{{\rm WD}}
%\newcommand{\pcr}{{\rm pcr}}
\newcommand{\ver}{{\rm ver}}
\newcommand{\gen}{{\rm gen}}
\newcommand{\jh}{{\rm JH}}
\newcommand{\stx}{{x_{\rm st}}}
\newcommand{\bmdel}{{ \bm{\delta}_\bh}}
\newcommand{\ul}{\underline}
\newcommand{\omepik}{\Omega_{S^u_0}}

\newcommand{\parb}{{\rm par}}

\newcommand{\bersteineigenvarpik}{{\mathcal{E}_{\omepik,\fp,\bm{\lambda}_\bh}^\infty(\overline{\rho})}}

\newcommand{\defvarnon}{{U^\Box_{\omega,\underline{\bh}}(\overline{r})}}
\newcommand{\defvarr}{{\mathfrak{X}_{\overline{r}}^\Box}}
\newcommand{\defvarring}{{R_{\overline{r}}^\Box}}


\newcommand{\defvaru}{{X^\Box_{\omepik,\mathbf{{h}}}(\overline{r})}}
\newcommand{\defvarrho}{{X^\Box_{\omepik,\mathbf{{h}}}(\overline{r})}}
\newcommand{\defvarrhonon}{{U^\Box_{\omepik,\mathbf{{h}}_{\Dpik}}(\overline{r})}}

\newcommand{\complocalbersteineigenvarpik}{{\widehat{\mathcal{E}_{\omepik,\fp,\bm{\lambda}_\pi}^\infty}(\overline{\rho})_x}}
\newcommand{\complocaldefvarrho}{{ \widehat{X^\Box_{\omepik,\mathbf{h}}(\overline{r})}_{x_L}}}

\newcommand{\complocaldefvarrhosub}{{ \widehat{\xfpxxl}}}
\newcommand{\complocaldefvarrhop}{{\widehat{X_{\mathrm{tri}}(\overline{r})}_{x}}}
\newcommand{\complocaldefvarrhox}{{ \widehat{X^\Box_{\omepik,\mathbf{{h}}}}(\overline{r})_{x}}}


%\newcommand{\complocaldefvarrhosub}{{ \widehat{\xfpxxl}}}
%\newcommand{\complocaldefvarrhop}{{\widehat{X_{\mathrm{tri}}(\overline{r})}_{x}}}
%\newcommand{\complocaldefvarrhop}{{ 


%\renewcommand{\Gamma}{\varGamma}
\usepackage{amsfonts}
\newcommand{\BOne} {{\mathchoice{\hbox{\rm1\kern-2.7pt l\kern.9pt}}
		{\hbox{\rm1\kern-2.7pt l\kern.9pt}}
		{\hbox{\scriptsize\rm1\kern-2.3pt l\kern.4pt}}
		{\hbox{\scriptsize\rm1\kern-2.4pt l\kern.5pt}}}}

\newcommand{\BA}{{\mathbb{A}}}
\newcommand{\BB}{{\mathbb{B}}}
\newcommand{\BC}{{\mathbb{C}}}
\newcommand{\BD}{{\mathbb{D}}}
\newcommand{\BE}{{\mathbb{E}}}
\newcommand{\BF}{{\mathbb{F}}}
\newcommand{\BG}{{\mathbb{G}}}
\newcommand{\BH}{{\mathbb{H}}}
\newcommand{\BI}{{\mathbb{I}}}
\newcommand{\BJ}{{\mathbb{J}}}
\newcommand{\BK}{{\mathbb{K}}}
\newcommand{\BL}{{\mathbb{L}}}
\newcommand{\BM}{{\mathbb{M}}}
\newcommand{\BN}{{\mathbb{N}}}
\newcommand{\BO}{{\mathbb{O}}}
\newcommand{\BP}{{\mathbb{P}}}
\newcommand{\BQ}{{\mathbb{Q}}}
\newcommand{\BR}{{\mathbb{R}}}
\newcommand{\BS}{{\mathbb{S}}}
\newcommand{\BT}{{\mathbb{T}}}
\newcommand{\BU}{{\mathbb{U}}}
\newcommand{\BV}{{\mathbb{V}}}
\newcommand{\BW}{{\mathbb{W}}}
\newcommand{\BX}{{\mathbb{X}}}
\newcommand{\BY}{{\mathbb{Y}}}
\newcommand{\BZ}{{\mathbb{Z}}}

\newcommand{\ba}{{\mathbf{a}}}
\newcommand{\bA}{{\mathbf{A}}}
\newcommand{\bB}{{\mathbf{B}}}
\newcommand{\bC}{{\mathbf{C}}}
\newcommand{\bD}{{\mathbf{D}}}
\newcommand{\bE}{{\mathbf{E}}}
\newcommand{\be}{{\mathbf{e}}}
\newcommand{\bF}{{\mathbf{F}}}
\newcommand{\bG}{{\mathbf{G}}}
\newcommand{\bH}{{\mathbf{H}}}
\newcommand{\bI}{{\mathbf{I}}}
\newcommand{\bJ}{{\mathbf{J}}}
\newcommand{\bK}{{\mathbf{K}}}
\newcommand{\bL}{{\mathbf{L}}}
\newcommand{\bM}{{\mathbf{M}}}
\newcommand{\bN}{{\mathbf{N}}}
\newcommand{\bO}{{\mathbf{O}}}
\newcommand{\bP}{{\mathbf{P}}}
\newcommand{\bQ}{{\mathbf{Q}}}
\newcommand{\bR}{{\mathbf{R}}}
\newcommand{\bS}{{\mathbf{S}}}
\newcommand{\bT}{{\mathbf{T}}}
\newcommand{\bU}{{\mathbf{U}}}
\newcommand{\bV}{{\mathbf{V}}}
\newcommand{\bW}{{\mathbf{W}}}
\newcommand{\bX}{{\mathbf{X}}}
\newcommand{\bY}{{\mathbf{Y}}}
\newcommand{\bZ}{{\mathbf{Z}}}
\newcommand{\bz}{{\mathbf{z}}}
\newcommand{\bc}{{\mathbf{c}}}
\newcommand{\bx}{{\mathbf{x}}}
\newcommand{\bv}{{\mathbf{v}}}
\newcommand{\bh}{{\mathbf{h}}}
\newcommand{\bk}{{\mathbf{k}}}
\newcommand{\bd}{{\mathbf{d}}}
\newcommand{\bj}{{\mathbf{j}}}
\newcommand{\bq}{{\mathbf{q}}}

\newcommand{\cN}{\mathcal N}
\newcommand{\cL}{\mathcal L}
\newcommand{\co}{\mathcal O}
\newcommand{\cA}{\mathcal A}
\newcommand{\cB}{\mathcal B}
\newcommand{\cR}{\mathcal R}
\newcommand{\cH}{\mathcal H}
\newcommand{\cC}{\mathcal C}
\newcommand{\cS}{\mathcal S}
\newcommand{\cD}{\mathcal D}
\newcommand{\cI}{\mathcal I}
\newcommand{\cW}{\mathcal W}
\newcommand{\cT}{\mathcal T}
\newcommand{\cM}{\mathcal M}
\newcommand{\cF}{\mathcal F}
\newcommand{\cV}{\mathcal V}
\newcommand{\cE}{\mathcal E}
\newcommand{\cG}{\mathcal G}
\newcommand{\cJ}{\mathcal J}
\newcommand{\cO}{\mathcal O}
\newcommand{\cP}{\mathcal P}
\newcommand{\cQ}{\mathcal Q}
\newcommand{\cU}{\mathcal U}
\newcommand{\cZ}{\mathcal Z}
\newcommand{\cY}{\mathcal Y}
\newcommand{\cX}{\mathcal X}


\newcommand{\FA}{{\mathfrak{A}}}
\newcommand{\FB}{{\mathfrak{B}}}
\newcommand{\FC}{{\mathfrak{C}}}
\newcommand{\FD}{{\mathfrak{D}}}
\newcommand{\FE}{{\mathfrak{E}}}
\newcommand{\FF}{{\mathfrak{F}}}
\newcommand{\FG}{{\mathfrak{G}}}
\newcommand{\FH}{{\mathfrak{H}}}
\newcommand{\FI}{{\mathfrak{I}}}
\newcommand{\FJ}{{\mathfrak{J}}}
\newcommand{\FK}{{\mathfrak{K}}}
\newcommand{\FL}{{\mathfrak{L}}}
\newcommand{\FM}{{\mathfrak{M}}}
\newcommand{\FN}{{\mathfrak{N}}}
\newcommand{\FO}{{\mathfrak{O}}}
\newcommand{\FP}{{\mathfrak{P}}}
\newcommand{\FQ}{{\mathfrak{Q}}}
\newcommand{\FR}{{\mathfrak{R}}}
\newcommand{\FS}{{\mathfrak{S}}}
\newcommand{\FT}{{\mathfrak{T}}}
\newcommand{\FU}{{\mathfrak{U}}}
\newcommand{\FV}{{\mathfrak{V}}}
\newcommand{\FW}{{\mathfrak{W}}}
\newcommand{\FX}{{\mathfrak{X}}}
\newcommand{\FY}{{\mathfrak{Y}}}
\newcommand{\FZ}{{\mathfrak{Z}}}


\newcommand{\fa}{{\mathfrak{a}}}
\newcommand{\fb}{{\mathfrak{b}}}
\newcommand{\fc}{{\mathfrak{c}}}
\newcommand{\fd}{{\mathfrak{d}}}
\newcommand{\fe}{{\mathfrak{e}}}
\newcommand{\ff}{{\mathfrak{f}}}
\newcommand{\fg}{{\mathfrak{g}}}
\newcommand{\fh}{{\mathfrak{h}}}
\newcommand{\Fi}{{\mathfrak{i}}}
\newcommand{\fj}{{\mathfrak{j}}}
\newcommand{\fk}{{\mathfrak{k}}}
\newcommand{\fl}{{\mathfrak{l}}}
\newcommand{\fm}{{\mathfrak{m}}}
\newcommand{\fn}{{\mathfrak{n}}}
\newcommand{\fo}{{\mathfrak{o}}}
\newcommand{\fp}{{\mathfrak{p}}}
\newcommand{\fq}{{\mathfrak{q}}}
\newcommand{\fr}{{\mathfrak{r}}}
\newcommand{\fs}{{\mathfrak{s}}}
\newcommand{\ft}{{\mathfrak{t}}}
\newcommand{\fu}{{\mathfrak{u}}}
\newcommand{\fv}{{\mathfrak{v}}}
\newcommand{\fw}{{\mathfrak{w}}}
\newcommand{\fx}{{\mathfrak{x}}}
\newcommand{\fy}{{\mathfrak{y}}}
\newcommand{\fz}{{\mathfrak{z}}}

\newcommand{\sE}{\mathscr E}
\newcommand{\sB}{\mathscr B}
\newcommand{\sH}{\mathscr H}
\newcommand{\sC}{\mathscr C}
\newcommand{\sF}{\mathscr F}
\newcommand{\sD}{\mathscr D}
\newcommand{\sI}{\mathscr I}
\newcommand{\sA}{\mathscr A}
\newcommand{\sW}{\mathscr W}
\newcommand{\sL}{\mathscr L}
\newcommand{\sK}{\mathscr K}
\newcommand{\sN}{\mathscr N}
\newcommand{\sX}{\mathscr X}
\newcommand{\sV}{\mathscr V}


		

\begin{document}	
	
\title{\textbf{\textsc{Toward $p$-adic Hodge parameters for potentially crystalline representations of $\GLN_n$}}}


\author{Yiqin He
\thanks{Morningside Center of Mathematics, Chinese Academy of Sciences,\;No. 55, Zhongguancun East Road, Haidian District, Beijing 100190, P.R. China; e-mail address: \texttt{heyiqin@amss.ac.cn}
}}

\date{}
\maketitle


\begin{abstract}
Let $p$ be a prime number,\;$n\geq 2$,\;and let $L$ be a finite extension of $\bQ_p$.\;Let $\rho_L$ be an $n$-dimensional non-critical generic potentially crystalline $p$-adic representation of the absolute Galois group of $L$ with regular Hodge--Tate weights.\;Building on Ding's results and strategy in the crystabelline case \cite{ParaDing2024},\;and on the recent work of Breuil--Ding \cite{BDcritical25},\;we construct an explicit locally analytic representation $\pi_{1}(\rho_L)$ and describe explicitly the Hodge-filtration data of $\rho_L$ that it determines.\;When $\rho_L$ arises from a patched $p$-adic automorphic representation,\;we show, under mild hypotheses, that $\pi_{1}(\rho_L)$ is a subrepresentation of the $\GLN_n(L)$-representation globally associated with $\rho_L$ by using the framework of Bernstein eigenvarieties that were developed by Breuil-Ding \cite{Ding2021}. \;
\end{abstract}

	
{\hypersetup{linkcolor=black}
\tableofcontents}

\numberwithin{equation}{section}
	
\numberwithin{thm}{section}	



\setlength{\baselineskip}{12pt}

\section{Introduction}

This paper aims to extend the discussion of crystabelline (resp.,\;crystalline) $p$-adic Galois representations in \cite{ParaDing2024} (resp.,\;the recent work \cite{BDcritical25}) to potentially crystalline case.\;For simplicity,\;we assume in the introduction that $L=\bQ_p$.\;Let $\rho_p:\gal_{\bQ_p}\rightarrow \GLN_n(E)$ be a de Rham $p$-adic Galois representation,\;where $\gal_{\bQ_p}$ is the absolute Galois group of $\bQ_p$ and $E$ is a finite extension of $\bQ_p$.\;By Fontaine's theory,\;we can attach an $n$-dimensional Weil--Deligne representation $\wdre(\rho_p)$ and thus an irreducible smooth representation $\pi_{\mathrm{sm}}(\rho_p)$ of $\GLN_n(\bQ_p)$ over $E$ (via the classical local Langlands correspondence).\;Assume that $\rho_p$ has regular Hodge--Tate weights $\bh=(\bh_1>\cdots>\bh_n)$.\;Then the locally algebraic representation $\pi_{\lalg}(\rho_p):=\pi_{\mathrm{sm}}(\rho_p)\otimes_EL(\bh-\theta)$ is expected to be the locally algebraic subrepresentation of the conjectural locally analytic representation $\pi_{\ana}(\rho_p)$ via the $p$-adic local Langlands correspondence,\;where $\theta=(0,-1,\cdots,1-n)$ and $L(\bh-\theta)$ is the algebraic representation of $\GLN_n(\bQ_p)$ with highest weight $\bh-\theta$.\;The representation $\pi_{\lalg}(\rho_p)$ only encodes the information in the $F$-semisimplification of $\wdre(\rho_p)$ and $\bh$,\;and therefore misses the Hodge-filtration information of $\rho_p$.\;A basic problem (and starting point) in the locally analytic $p$-adic local Langlands program is to recover the Hodge filtration from locally ${\BQ}_p$-analytic representations of $\GLN_n({\bQ}_p)$.\;

After the pioneering work of Breuil, Colmez, and others, this question was completely settled for the $\GLN_2(\bQ_p)$ case.\;For general $\GLN_n(\bQ_p)$, however, it remained mysterious for a long time.\;Recently, Yiwen Ding made a breakthrough in understanding the $p$-adic Hodge parameters of \textit{non-critical generic} crystabelline representations of $G$, namely representations that become crystalline after restriction to the absolute Galois group of a suitable abelian extension of $\bQ_p$ \cite{ParaDing2024}.\;More precisely, Ding constructed an explicit locally $\bQ_p$-analytic representation $\pi_{1}(\rho_p)$ which determines $\rho_p$ uniquely and has the following form (i.e.,\;$\pi_1(\rho_p)$ is a non-split extension of $\pi_{\lalg}(\rho_p)^{\oplus(2^n-\frac{n(n+1)}{2}-1)}$ by $\pi_1(\wdre(\rho_p),\bh)$):
\[\pi_1(\rho_p)=\left[\pi_1(\wdre(\rho_p),\bh)-\pi_{\lalg}(\rho_p)^{\oplus(2^n-\frac{n(n+1)}{2}-1)}\right],\]
where $\pi_1(\wdre(\rho_p),\bh)$ is obtained from the amalgamated sum of some locally analytic principal series and $\soc_{G}\pi_1(\rho_p)=\pi_{\lalg}(\rho_p)$.\;Moreover, when $\rho_p$ comes from a $p$-adic automorphic representation, $\pi_{1}(\rho_p)$ is a subrepresentation of the globally associated $\GLN_n(\bQ_p)$-representation $\widehat{\pi}(\rho_p)^{\ana}$.\;


The appearance, with large multiplicity, of extra locally algebraic constituents in the cosocle of $\pi_{1}(\rho_p)$ is a new phenomenon.\;Such constituents seem to occur widely for general de Rham $p$-adic Galois representations and may encode information about the Hodge filtration of $\rho_p$.\;In this article, we study this question for \textit{non-critical generic potentially crystalline} $p$-adic Galois representations.\;This extends the crystabelline picture of \cite{ParaDing2024} to the potentially crystalline setting and connects it with the local--global compatibility framework of the patched Bernstein eigenvariety developed in \cite{Ding2021}.\;

Although we impose the genericity assumption throughout this article in order to keep the exposition and the statements as clean as possible, the methods developed here are not expected to rely on it in an essential way.\;After replacing the generic smooth representation by the corresponding finite-length object in its Bernstein block and keeping track of possible coincidences among Jordan--H\"older factors and extension classes,  \textit{the same constructions should apply without essential difficulty to non-critical potentially crystalline representations without imposing the genericity assumption on $\wdre(\rho_p)$, or equivalently on $\pi_{\mathrm{sm}}(\rho_p)$}.\;We do not pursue this greater generality here only to avoid additional notation.\;In our work \cite{HEparaforsemitable}, we discuss the corresponding question in the \textit{semistable} case, which provides further evidence for the role of $p$-adic Hodge parameters.\;That work reduces the general semistable case to the crystalline case and to the so-called Steinberg case, the opposite extreme in which the Weil--Deligne representation $\wdre(\rho_p)$  has maximal monodromy rank.\;This reduction is, in spirit, parallel to the Bernstein--Zelevinsky classification theorem for smooth representations.\;

%Some of the results in \cite{HEparaforsemitable} also apply to more general potentially semistable case.\;



In the sequel, let $\Delta:=\{1,\cdots,n-1\}$ be the set of simple roots of $\GLN_n$.\;Fix $S_0\subseteq \Delta$, and let $\bL_{S_0}:=\prod_{i=1}^s\GLN_{r_i}\subseteq \GLN_n$, for integers $r_1,\cdots,r_s$ satisfying $r_1+\cdots+r_s=n$, be the standard Levi subgroup associated with $S_0$.\;For $u\in \sW_s$, let $S_0^u$ be the subset of $\Delta$ such that the associated Levi subgroup satisfies $\bL_{S_0^u}=\prod_{i=1}^s\GLN_{r_{u^{-1}(i)}}$.\;For $i=1,\cdots,s$ and $u\in\sW_s$, put $t_0^u=0$,\;$t_{s}^u=n$ and
\[t_i^u=\sum_{j=1}^i r_{u^{-1}(j)}.\]
Thus $\Delta\backslash S_0^u=\{t_1^u,\cdots,t_{s-1}^u\}$.\;Let $\bB$ be the Borel subgroup of $\GLN_n$ consisting of upper triangular matrices, and let $\bP_{S_0^u}\supseteq \bL_{S_0^u}\bB$ be the standard parabolic subgroup associated with $S_0^u$.\;For $1\leq i\leq s$, fix a cuspidal Bernstein component $\Omega_{i}$ of $\GLN_{r_i}(\bQ_p)$.\;Then, for $u\in \sW_s$, $\Omega_{S^u_0}:=\prod_{i=1}^s\Omega_{u^{-1}(i)}$ is a cuspidal Bernstein component of the Levi subgroup $\bL_{S^u_0}$.\;Put $G:=\GLN_n(\bQ_p)$.\;

Let $\cR_{E}$ be the Robba ring over $\bQ_p$ with coefficients in $E$.\;We write $\cR_{E}(\delta)$ for the $(\varphi,\Gamma)$-module over $\cR_{E}$ of rank one associated to a continuous character $\delta:\bQ_p^\times\rightarrow E$.\;For $\alpha\in E$,\;let $\unr(p)$ be the unramified character of $\bQ_p^\times$ sending uniformizers $p$ to $\alpha$.\;For integer $\bh$,\;let $z^{h}$ be the character sending $z\in \bQ^\times_p$ to $z^{h}$.\;


For a closed point $x_i\in \Spec\FZ_{\Omega_{i}}(E)$ ($1\leq i\leq r$), let $\pi_{x_i}$ be the associated irreducible cuspidal sm ooth representation over $E$ of type $\Omega_i$ and let $\wdre_{x_i}:=\mathrm{rec}(\pi_{x_i})$ be the associated  $r$-dimensional absolutely irreducible Weil--Deligne representation via the normalized classical local Langlands correspondence (see \cite{scholze2013local}).\;$\wdre_{x_i}$ corresponds to a Deligne--Fontaine module $\df_{x_i}$, by Fontaine's equivalence of categories as in \cite[Proposition 4.1]{breuil2007first}, and a $p$-adic differential equation ${\Delta_{x_i}}$ over $\cR_{E}$.\;Here ${\Delta_{x_i}}$ is a de Rham $(\varphi,\Gamma)$-module of rank $r$ over $\cR_{k(x),L}$ with constant Hodge--Tate weights $0$ such that, after forgetting the Hodge filtration, $D_{\mathrm{pst}}(\Delta_{x_i})$ is isomorphic to $\df_{x_i}$, by Berger's theory \cite[Theorem A]{berger2008equations}.\;




Fix  a potentially crystalline $p$-adic Galois representation $\rho_p$, and assume that the associated potentially crystalline $(\varphi,\Gamma)$-module $\Dpik$ over $\cR_{E}$ of rank $n$ admits, for each $u\in \sW_s$, a \textit{non-critical} $\Omega_{S^u_0}$-filtration $\cF_u$ in the sense of \cite{Ding2021}; this is a parabolic filtration by saturated $(\varphi,\Gamma)$-submodules of $\Dpik$ over $\cR_{E}$:
\[\cF_u=\fil_{\bullet}^{\cF_u}\Dpik: \ 0 =\fil_0^{\cF_u}\Dpik \subsetneq \fil_1^{\cF_u}\Dpik \subsetneq \cdots \subsetneq \fil_{s}^{\cF_u} \Dpik=\Dpik\]
such that there is an injection of $(\varphi,\Gamma)$-modules over $\cR_{E,L}$ of rank $r_{u^{-1}(i)}$:
\begin{equation}
	\mathbf{I}^u_{i}:E_{u,i}:=\gr_{i}^{\cF_u} \Dpik \hookrightarrow \Delta_{x_{u^{-1}(i)}}\otimes_{\cR_E} \cR_{E}(z^{t^u_i}),\quad(\text{note that}\;E_{u,i}[1/t]=\Delta_{x_{u^{-1}(i)}}[1/t])
\end{equation}
for some  $\underline{x}:=(x_{i})_{1\leq i\leq s}\in \Spec\FZ_{\Omega_{S_0}}(E)$ (so that $\underline{x}^u=(x_{u^{-1}(i)})_{1\leq i\leq s}\in \Spec\FZ_{\Omega_{S^u_0}}(E)$)
and the non-critical assumption means that the Hodge--Tate weights of $\fil_{i}^{\cF_u}\Dpik$ (resp.,\;$\gr_{i}^{\cF_u} \Dpik$) are
\[(\hpi_{1},\hpi_{2},\cdots,\hpi_{t_i^u}),\quad (\text{resp.,\;}(\hpi_{t^u_{i-1}+1},\hpi_{t^u_{i-1}+2},\cdots,\hpi_{t_i^u})),\]
In this paper,\;we usually use $\bullet-\bullet$ to denote an extension of two objects, such as $(\varphi,\Gamma)$-modules or $G$-representations, where the first (resp., second) object is the subobject (resp., quotient)).\;Then $\cF_u$ corresponds to the following successive  extension of $\{E_{u,i}\}$ :
\[\Dpik=[{E}_{u,1}-{E}_{u,2}-\cdots-{E}_{u,s}]=[{E}_{1}-{E}_{2}-\cdots-{E}_{s}],\;\]
where we write $E_i:=E_{1,i}$ for simplicity.\;In particular, if $S_0=\emptyset$, then $\Dpik$ is crystabelline and $\cF_u$ becomes the classical triangulation (or the refinement associated to $u$).\;

Let $L'$ be a finite Galois extension of $\bQ_p$ such that $\Dpik|_{L'}$ is a crystalline $(\varphi,\Gamma)$-module over $\cR_{E,L'}$ of rank $n$, where $\cR_{E,L'}$ is the Robba ring over $L'$ with coefficients in $E$.\;By Fontaine's theory, we consider the Deligne--Fontaine module attached to $\Dpik$:
$$\df_{\Dpik}=(\varphi,\Gal(L'/\bQ_p),D_{\pcr}(\Dpik))$$
where $D_{\pcr}(\Dpik)=D_{\mathrm{cr}}^{L'}(\Dpik\otimes_{\cR_{E}}\cR_{E,L'})$ is a finite free $L'_0\otimes_{\bQ_p}E$-module of rank $n$, and $L'_0$ is the maximal unramified subextension of $L'$ over $\bQ_p$.\;Since $\Dpik$ is potentially crystalline, hence de Rham, $D:=D_{\dr}(\Dpik)\cong (D_{\pcr}(\Dpik)\otimes_{L_0'}L')^{\Gal(L'/\bQ_p)}$ is a free $E$-module of rank $n$.\;The module $D$ carries a natural Hodge filtration, expressed by the following complete flag :
\[\fil^{\bullet}_{H}(D): \ 0 \subsetneq \fil^{-\hpi_{n}}_{H}(D) \subsetneq \fil^{-\hpi_{n-1}}_{H}(D) \subsetneq \cdots \subsetneq \fil^{{-\hpi_{1}}}_{H}(D)=D.\]
For $1\leq i\leq s$, fix an $E$-basis $\{e_{i,j}\}_{1\leq j\leq r_i}$ of $D_{\dr}(E_{i})$.\;Then $\{e_{1,j}\}_{1\leq j\leq r_1},\cdots,\{e_{s,j}\}_{1\leq j\leq r_s}$ form an $E$-basis of $D$; we relabel them as $e_{1},e_{2},\cdots,e_{n}$ using lexicographic order.\;Under such basis,\;$\fil^{\bullet}_{H}(D)$ corresponds to a point of $\bZ_{S_0}\backslash \GLN_n/\bB$,\;which we call the \textit{$p$-adic Hodge parameters} of $\rho_p$.\;Let $\homo_{\fil}(D,D)$ be the subspace of endomorphisms of $D$ that preserve the filtration $\fil^{H}_{\bullet}(D)$.\;
%For $1\leq i\leq n$,\;there exist $E$-lines $\cL_{l}\in E\langle e_{l},e_{l-1},\cdots,e_{1}\rangle$ such that\[\fil^{H}_{-\bh_{n-i}}(D)=\oplus_{l=n-i+1}^n\cL_l.\]

Throughout this paper,\;we assume that $\underline{x}$ is \textit{regular} and \textit{generic}, namely $\Delta_{x_i}\neq \Delta_{x_j}$ and $\Delta_{x_i}\neq \Delta_{x_j}\otimes_{\cR_{E}}\cR_{E}(\unr(p))$ for $i\neq j$.\;Equivalently, $\pi_{\mathrm{sm}}(\rho_p)$ is a generic smooth representation of $G$, and $\wdre(\rho_p)$ is a generic Weil--Deligne representation.\;Recall $\lambda=\bh-\theta$.\;Consider the following locally $\bQ_p$-analytic principal series of $G$:
\[\mathrm{PS}_{u}(\underline{x},\bh):=\Big(\ind^G_{\op_{S_0^u}(\bQ_p)}\pi_{\sm}(\underline{x}^u)\eta_{S_0^u}\otimes_E L_{S_0^u}(\lambda)\Big)^{\ana},\pi_{\sm}(\underline{x}^u):=\boxtimes_{j=1}^s\pi_{x_{u^{-1}(j)}},\]
where $L_{S_0^u}(\lambda)$ is the unique irreducible algebraic representation of $\bL_{S_0^u}(\bQ_p)$ with highest weight $\lambda$, $\op_{S_0^u}$ is the parabolic subgroup opposite to $\bP_{S_0^u}$, and $\eta_{S_0^u}$ is the square root of the modulus character of $\op_{S_0^u}$.\;The explicit structure of $\mathrm{PS}_{u}(\underline{x},\bh)$ is well-understand by Orlik-Strauch \cite[The main theorem]{orlik2015jordan}.\;By the generic assumption,\;$\soc_G\mathrm{PS}_{u}(\underline{x},\bh)\cong \pi_{\lalg}(\rho_p)\cong \mathrm{PS}^{\lalg}_{u}(\underline{x},\bh)$,\;where  $\mathrm{PS}^{\lalg}_{u}(\underline{x},\bh)$ is the space of locally $\bQ_p$-algebraic vectors in $\mathrm{PS}_{u}(\underline{x},\bh)$.\;For simplicity, put $\pi_{\lalg}(\underline{x},\bh):=\mathrm{PS}^{\lalg}_{u}(\underline{x},\bh)\cong \pi_{\lalg}(\rho_p)$.\;




Let $\mathrm{PS}_{u,1}(\underline{x},\bh)$ be the first two layers of the socle filtration of $\mathrm{PS}_u(\underline{x},\bh)$.\;The amalgamated sum 
\[\bigoplus^{u\in \sW_{s}}_{\pi_{\lalg}(\underline{x},\bh)}\mathrm{PS}_{u,1}(\underline{x},\bh)\]
 has the unique quotient $\pi^{\flat}_{1}(\underline{x},\bh)$ of socle $\pi_{\lalg}(\underline{x},\bh)$,\;which lies in the following short exact sequence (here $\sW_{s}^{\widehat{j},\emptyset}=\sW_{\widehat{j}}\backslash\sW_s$, $\sW_{\widehat{j}}$ is the subgroup of $\sW_s$ generated by the simple reflections $\{s_i\}_{i\in \widehat{j}}$, and $\widehat{j}=\Delta_s\backslash \{j\}$):
\begin{equation}\label{firstwholeintro}
	0\rightarrow \pi_{\lalg}(\underline{x},\bh) \rightarrow \pi^{\flat}_{1}(\underline{x},\bh)\rightarrow \oplus_{\substack{1\leq j\leq s-1\\u\in \sW_{s}^{\widehat{j},\emptyset}}}C_{t_j^u,u}\rightarrow 0,
\end{equation}
where $C_{t_j^u,u}$ are explicit pure irreducible locally analytic representations given by the Orlik--Strauch functor (see \cite[The main theorem]{orlik2015jordan}).\;We state our three main theorems and several remarks as follows.\;



\begin{thm}\label{thmintor1}(Proposition \ref{construext1}) There exists a map, depending only on the choice of $\log_p(p)\in E$,
	\begin{equation}\label{reconsfortDpikcirc}
{t}^{\flat}_{D}:\ext^1_{G}\big(\pi_{\lalg},\pi^{\flat}_{1}(\underline{x},\bh)\big)\twoheadrightarrow
		\Big(\prod_{i=1}^s\ext^{1}(E_i[1/t],E_i[1/t])\Big)\oplus\homo^{\flat}_{\fil}(D,D),
	\end{equation}
and $\homo^{\flat}_{\fil}(D,D)\subseteq \homo_{\fil}(D,D)$ is described in Proposition \ref{introforhomoflat}.\;
\end{thm}

Define $\Delta'=\bigcup_{u\in \sW_s}\Delta\backslash S_0^u=\bigcup_{u\in \sW_s}\{t_1^u,\cdots,t_{s-1}^u\}\subseteq\Delta$, and write $\Delta':=\{i_{1}<i_{2}<\cdots<i_{|\Delta'|}\}$.\;For $i\in\Delta'$, put $\cI_{i}^{\flat}:=\{u\in \sW_s:i\notin S^u_0 \}$.\;For any $u\in \cI_{i}^{\flat}$,\;suppose that $i=t^u_{l_i(u)}$ for some $1\leq l_i(u)\leq s-1$.\;Define
\[D^{\flat}_{(i)}:=E\langle e_{u^{-1}(q),j}:l_i(u)+1\leq q\leq s,1\leq j\leq r_{u^{-1}(q)},u\in \cI_{i}^{\flat}\rangle\subseteq D,\]
and set $\bd_i=n-\dim_ED^{\flat}_{(i)}$.\;Note that $\bd_i=0$ if and only if $D^{\flat}_{(i)}=D$.\;

\begin{thm}\label{Hodgeparadeterintro}(Theorem \ref{Hodgeparadeter} \& Corollary \ref{Hodgeparadetercor}) For $1\leq t\leq |\Delta'|$, $\ker(t^{\flat}_{D})$ determines
	\begin{equation}
		\begin{aligned}
			\Big\{\fil_H^{-\bh_{{\max\{i_{{j}},\bd_{i_{t}}\}}+1}}(D_{\sigma})\oplus D_{\sigma}/\fil_H^{-\bh_{\bd_{i_{t}}}}(D_{\sigma})\Big\}_{1\leq j\leq t}
		\end{aligned}
	\end{equation}	
	In particular, if $\bd_{i_t}=0$, then $\ker(t^{\flat}_{D_{\sigma}})$ determines $\{\fil_H^{-\bh_{i_{{j}}+1}}(D_{\sigma})\}_{1\leq j\leq t}$.\;
\end{thm}
The tautological extension of $\pi_{\lalg}(\underline{x},\bh)\otimes_E\ker(t^{\flat}_{D})$ by $\pi^{\flat}_{1}(\underline{x},\bh)$ gives a locally analytic representation $\pi^{\flat}_{\min}(\Dpik)$.\;
\begin{thm}\label{thmintor11} (Theorem \ref{mainthmglobal}:\;local--global compatibility in the patched setting) Let $R_{\infty}$ be the patched Galois deformation ring and $\Pi_{\infty}$ be the patched $R_{\infty}$-admissible unitary representation.\;If $\rho_p$ comes from a maximal ideal $\fm_{\rho}$ of $R_{\infty}[1/p]$, then $\pi^{\flat}_{\min}(\Dpik)\hookrightarrow \Pi_{\infty}[\fm_{\rho}]$.\;
\end{thm}




\begin{rmk}In general,\;for a de Rham $\Dpik$,\;when a suitable locally analytic representation $\pi_{1}(D)$ is available,\;\cite[(6)]{BDcritical25} predicts a surjection
 \begin{equation}
	\begin{aligned}
t_{D}:\ext^1_{G}\big(\pi_{\lalg}(D),\pi_{1}(D)\big)\twoheadrightarrow\ext^1_{\varphi}(D,D)\oplus\homo_{\fil}(D,D),\;
	\end{aligned}
\end{equation}
where $\ext^1_{\varphi}(D,D)$ means extensions as $\varphi$-modules.\;In our case, we expect $\pi^{\flat}_{1}(\underline{x},\bh)$ to be a subrepresentation of $\pi_{1}(D)$, and (\ref{reconsfortDpikcirc}) gives $t_{D}|_{\ext^1_{G}\left(\pi_{\lalg},\pi^{\flat}_{1}(\underline{x},\bh)\right)}$ and its image.\;For $u\in \sW_s$, recall that $\Dpik=[{E}_{u,1}-{E}_{u,2}-\cdots-{E}_{u,s}]$.\;Assume that for each $1\leq i\leq s$, one can associate a locally analytic representation $\pi_{\ana}({E}_{u,i})$ that determines ${E}_{u,i}$.\;Consider the locally analytic parabolic induction
\begin{equation*}
	\mathrm{PS}_{\cF_{u}
		}(\underline{x},\bh):=\big(\ind^G_{\op_{S^u_0}(L)}\pi_{\ana}(\underline{x}^u)\eta_{S_0^u}\big)^{\ana},\;\pi_{\ana}(\underline{x}^u):=\boxtimes_{i=1}^s\pi_{\ana}({E}_{u,i}).\;
\end{equation*}
We have a injection of locally analytic $G$-representations $\mathrm{PS}_{S_0,u}(\underline{x},\bh)\hookrightarrow\mathrm{PS}_{\cF_{u}}(\underline{x},\bh)$.\;Then the  amalgamated sum (if possible) of $\{\mathrm{PS}_{\cF_{u}}(\underline{x},\bh)\}_{u\in \sW_{s}}$ has a unique quotient $\pi_{1}^+(\underline{x},\bh)$ of socle $\pi_{\lalg}(\underline{x},\bh)$  such that $\pi^{\flat}_{1}(\underline{x},\bh)\hookrightarrow\pi_{1}^+(\underline{x},\bh)$; we expect $\pi_{1}^+(\underline{x},\bh)\hookrightarrow \pi_{1}(D)$.\;Then (\ref{reconsfortDpikcirc}) can be extended to a morphism
\begin{equation}\label{introt+D}
	t^+_{D}:\ext^1_{G}\big(\pi_{\lalg}(\underline{x},\bh),\pi^+_{1}(\underline{x},\bh)\big)\rightarrow\ext^1_{\varphi}(D,D)\oplus\homo_{\fil}(D,D).
\end{equation}
See  Expectation \ref{expectintro} for a further comment.\;
\end{rmk}

Let $\fg$ (resp.,\;$\fb$) be the $E$-Lie algebra of $\GLN_n$ (resp.,\;$\bB$).\;Moreover,\;for $u\in \sW_s$,\;let $\fz_{S^u_0}$ (resp.,\;$\fn_{S^u_0}$) be the $E$-Lie algebra of the center $\bZ_{S^u_0}$ of $\bL_{S_0^u}$ (resp.,\;of the unipotent radical of $\bP_{S^u_0}$).\;Put $\tau_{S^u_0}=\fz_{S^u_0}\ltimes\fn_{S^u_0}$ the full radical.\;Under the basis $\{e_{i}\}_{1\leq i\leq n}$, we identify $\homo_{E}(D,D)$ with $
\fg$.\;If we choose $g\in \GLN_n(E)$ such that $g\bB\in \GLN_n/\bB$ gives the coordinates of the complete Hodge flag $\fil^{\bullet}_{H}(D)$ in this basis, then $\homo_{\fil}(D,D)$ is identified with $\mathrm{Ad}_g(\fb)$.\;Let $u^{\sharp}\in \sW_n$ be the unique element sending $e_{i,j}$ for $1\leq j\leq r_i$ to $e_{u^{-1}(i),j}$ for $1\leq j\leq r_i$,\;which thus induces a permutation on $e_{1},\cdots,e_{n}$.\;

\begin{pro}\label{introforhomoflat}(Propositions \ref{dfnisoforflatfilhomo} \& \ref{comparesharpflat})
	$\homo^{\flat}_{\fil}(D,D)\cong \mathrm{Ad}_{g}(\fb)_{S_0}^{\flat}:=\sum_{u\in \sW_s}\mathrm{Ad}_{(u^{\sharp})^{-1}}(\tau_{S^u_0})\cap \mathrm{Ad}_{g}(\fb)$ and 
$\dim_E\ker(t^{\flat}_{D})=2^s-1-\dim_E\mathrm{Ad}_{g}(\fb)_{S_0}^{\flat}$.\;
\end{pro}
\begin{rmk}In particular,\;if $S_0=\emptyset$,\;then $\mathrm{Ad}_{g}(\fb)_{S_0}^{\flat}=\mathrm{Ad}_{g}(\fb)$ and $\dim_E\ker(t^{\flat}_{D})=2^n-\frac{n(n+1)}{2}-1$.\;We come back to \cite{ParaDing2024}.\;
\end{rmk}	 











\begin{rmk}\label{rmkintroII} (An example for Theorem \ref{Hodgeparadeterintro}) A typical example is $\bL_{S_0}\cong \GL_1^{\oplus (n-m)}\times \GLN_m$.\;As $E$-vector spaces, $D=D^1\oplus\cdots\oplus D^{n-m}\oplus D^{n-m+1}$, with $\dim_ED^j=1$ for $1\leq j\leq n-m$ and $\dim_ED
		_{\sigma}^{n-m+1}=m$.\;For $i\in \Delta'$,\;we have
		\begin{equation}
			D_{(i)}=\left\{
			\begin{array}{ll}
				D/D^{n-m+1},\;&n-i<m\\
				D ,\;&n-i\geq m
			\end{array}
			\right.,\;\bd_{i}=\left\{
			\begin{array}{ll}
				m,\;&n-i<m\\
				0,\;&n-i\geq m
			\end{array}
			\right.,\;
		\end{equation}
		If $m\leq n-m$ (resp.,\;$m>n-m$), then $\Delta'=\Delta$ (resp.,\;$\Delta'=\Delta\backslash \{n-m+1,\cdots,m-1\}$).\;For $n-i\geq m$, $\ker(t^{\flat}_{D})$ determines $\{\fil_H^{-\bh_{j+1}}(D)\}_{1\leq j\leq i}$.\;If $n-i<m$ (resp.,\;$i>m-1$), then $\ker(t^{\flat}_{D})$ determines
		\begin{equation}
			\begin{aligned}
				\Big\{\fil_H^{-\bh_{j+1}}(D)\oplus D/\fil_H^{-\bh_{m}}(D)\Big\}_{m \leq j\leq i}\cup \Big\{\fil_H^{-\bh_{m+1}}(D)\oplus D/\fil_H^{-\bh_{m}}(D)\Big\}.
			\end{aligned}
		\end{equation}	
\end{rmk}

We finally  relate Theorem \ref{thmintor1} to certain parabolic deformations.\;For $u\in \sW_{s}$, note that $\bZ_{S_0^u}(\bQ_p)\cong \bQ_p^s$ is the center of $\bL_{S^u_0}(\bQ_p)$.\;Let $\homo(\bZ_{S_0^u}(\bQ_p),E)$ (resp.,\;$\homo(\bQ_p^{\times},E)$) be the $E$-vector space of locally analytic $E$-valued characters on $\bZ_{S_0^u}(\bQ_p)$ (resp.,\;$\bQ_p^{\times}$).\;Let $\ext^1_u(\Dpik,\Dpik)\subseteq \ext^1(\Dpik,\Dpik)$ be the space of $\bP_{S_0^u}$-parabolic deformations of $\Dpik$ that preserve $\cF_u$.\;There are several deformation subspaces $\ext^{1,\flat}_{u}(\Dpik,\Dpik)\subseteq \ext^{1,0}_{u}(\Dpik,\Dpik)$ of $\ext^{1}_{{u}}(\Dpik,\Dpik)$ which appear on the Bernstein eigenvariety. After identifying elements of $\ext^1(\Dpik,\Dpik)$ with deformations of $\Dpik$ over $\cR_{E[\epsilon]/\epsilon^2}$ (where $\cR_{E[\epsilon]/\epsilon^2}$ is the Robba ring over $\bQ_p$ with coefficients in $E[\epsilon]/\epsilon^2$), an element $\widetilde{D}\in \ext^1(\Dpik,\Dpik)$ belongs to $\ext^{1,\flat}_u(\Dpik,\Dpik)$ (resp.,\;$\ext^{1,0}_u(\Dpik,\Dpik)$) if
\[\widetilde{D}=[\widetilde{E}_{u,1}-\widetilde{E}_{u,2}-\cdots-\widetilde{E}_{u,s}]\]
and $\widetilde{E}_{u,i}\cong {E}_{u,i}\otimes_{\cR_{E}}\cR_{E[\epsilon]/\epsilon^2}(1+\psi_i\epsilon)$ (resp.,\;$\widetilde{E}_{u,i}\hookrightarrow  \Delta_{x_{u^{-1}(i)}}\otimes_{\cR_{E}}\cR_{E[\epsilon]/\epsilon^2}(z^{t^u_i}(1+\psi_i\epsilon))$) for some $\psi_i\in \homo(\bQ_p^{\times},E)$ and $1\leq i\leq s$.\;For $?\in\{\flat,0\}$,\;we have a  natural surjection (by the generic assumption)
\begin{equation}
	\kappa_{u}^{?}:\ext^{1,?}_{{u}}(\Dpik,\Dpik)\rightarrow \homo(\bZ_{S_0^u}(\bQ_p),E),
\end{equation}
which sends $\widetilde{D}$ to its deformation parameters $(\psi_i)_{1\leq i\leq s}$.\;Put $\ext^{1,\flat}_0(\Dpik,\Dpik)=\ker \kappa^{\flat}_1$.\;For any subspace $V\subseteq \ext^{1}(\Dpik,\Dpik)$, put $\overline{V}:=V/(V\cap \ext^{1,\flat}_0(\Dpik,\Dpik))$.\;In particular, we obtain an isomorphism
$\overline{\ext}^{1,\flat}_{u}(\Dpik,\Dpik)\xrightarrow{\sim }\homo(\bZ_{S_0^u}(\bQ_p),E)$.\;Moreover,\;for $?\in\{\flat,0\}$, we have a natural map
\begin{equation}\label{dfnforgDpikintro}
	\overline{g}^{?}_{\Dpik}:\bigoplus_{u\in\sW_{s}}\overline{\ext}^{1,?}_u(\Dpik,\Dpik)\rightarrow\overline{\ext}^{1}(\Dpik,\Dpik).
\end{equation}
We have a natural morphism:
\begin{equation}
	\begin{aligned}
		\homo(\bZ_{S^u_0}(\bQ_p),E)&\rightarrow\ext^1_{G}\big(\mathrm{PS}_{u}(\underline{x},\bh),\mathrm{PS}_{u}(\underline{x},\bh)\big),\;\\
		\psi&\mapsto \Big(\ind^G_{\op_{S_0^u}(\bQ_p)}\pi_{\sm}(\underline{x}^u)\eta_{S_0^u}\otimes_E L_{S_0^u}(\lambda)\otimes_E(1+(\psi\circ\mathrm{det}_{\bL_{S_0^u}})\epsilon)\Big)^{\ana},\;
	\end{aligned}
\end{equation}
which leads to a map, obtained by pushing forward via $\mathrm{PS}_{u}(\underline{x},\bh)\hookrightarrow \pi^{\flat}_{1}(\underline{x},\bh)$ and pulling back via the natural map $\pi_{\lalg}(\underline{x},\bh)\cong \mathrm{PS}^{\lalg}_{u}(\underline{x},\bh)\hookrightarrow \mathrm{PS}_{u}(\underline{x},\bh)$,
\[\zeta_{u}:\homo(\bZ_{S^u_0}(\bQ_p),E)\rightarrow  \ext^1_{G}\left(\pi_{\lalg}(\underline{x},\bh),\pi^{\flat}_{1}(\underline{x},\bh)\right).\]
Then the composition $\oplus_{u\in\sW_{s}}(\zeta_u\circ \kappa^{\flat}_u)$ gives
\begin{equation}\label{gammDpikcirc}
	\begin{aligned}
		\gamma^{\flat}_{\Dpik}:\bigoplus_{u\in\sW_{s}}\overline{\ext}^{1,\flat}_{u}(\Dpik,\Dpik)\rightarrow \ext^1_{G}\left(\pi_{\lalg}(\underline{x},\bh),\pi^{\flat}_{1}(\underline{x},\bh)\right).\;
	\end{aligned}
\end{equation}
The theory of almost de Rham representations induces a natural morphism $\nu:{\ext}^{1}(\Dpik,\Dpik)\rightarrow \homo_{\fil}(D,D)$. Its kernel is ${\ext}^{1}_g(\Dpik,\Dpik)$, the subspace of de Rham deformations of $\Dpik$.\;We have the following short exact sequence:
\begin{equation}\label{exactseqforhomfilanext1}
	0\rightarrow \prod_{i=1}^s\ext^1_g(E_i,E_i)\rightarrow  \overline{\ext}^{1}(\Dpik,\Dpik) \rightarrow \homo_{\fil}(D,D)\rightarrow 0.
\end{equation}


\begin{pro}\label{thmintor2}
\begin{itemize}
	\item[(1)] (Proposition \ref{splitingforext1gps})	Put ${\ext}^{1,\flat}(\Dpik,\Dpik)=\mathrm{Im}(g^{\flat}_{\Dpik})$ the image (\ref{dfnforgDpikintro}).\;The restriction of (\ref{exactseqforhomfilanext1}) to $\overline{\ext}^{1,\flat}(\Dpik,\Dpik)$ admits a splitting, depending only on the choice of $\log_p(p)\in E$:
	\[f_{\Dpik}^{\flat}:\overline{\ext}^{1,\flat}(\Dpik,\Dpik)\xrightarrow{\sim} \ext^{1,\flat}_{\varphi^f}(\Dpik[1/t],\Dpik[1/t]) \oplus\mathrm{Im}(\nu|_{\overline{\ext}^{1,\flat}(\Dpik,\Dpik)}),\]
	where $\ext^{1,\flat}_{\varphi^f}(\Dpik[1/t],\Dpik[1/t]):=\prod_{i=1}^s\ext^{1,\flat}_g(E_i,E_i)\cong \homo(\bZ_{S_0}(\bQ_p),E)$.\;
	\item[(2)]	(Proposition \ref{mainthmforfullrefine}) $f_{\Dpik}^{\flat}\circ \overline{g}^{\flat}_{\Dpik}=t^{\flat}_{D}\circ\gamma^{\flat}_{\Dpik}$.\;
\end{itemize}	
\end{pro}



\begin{rmk}\label{rmkintroIII}(Infinite fern for potentially crystalline cases;\;Theorem \ref{infinitePoten},\;Theorem \ref{Inifinitefernpoten}) The natural map $g_{\Dpik}:\bigoplus_{u\in \sW_{s}}{\ext}^{1}_{u}(\Dpik,\Dpik)\rightarrow {\ext}^1(\Dpik,\Dpik)$  is surjective.\;
\end{rmk}

\begin{rmk}\label{expectintro}
  (\ref{gammDpikcirc}) can be extended  to $\gamma^+_{\Dpik}:\bigoplus\nolimits_{u}\overline{\ext}^{1}_{u}(\Dpik,\Dpik)\rightarrow\ext^1_{G}\big(\pi_{\lalg}(\underline{x},\bh),\pi^+_{1}(\underline{x},\bh)\big)$.\;We  suspect that there are $p$-adic Hodge parameters that are contributed by the conjectural supersingular constituents in $\pi_1(D)/\pi^+_{1}(\underline{x},\bh)$.\;Thus  $\ker({t}^{+}_{D})$ still not determines $\{\fil_H^{-\bh_{i+1}}(D_{\sigma})\}_{i\in \Delta}$.\;
\end{rmk}









%A typical example is $n=4$,\;$S_0=\{1,3\}$ and $\Dpik=[E_1-E_2]=[E_2'-E_1']$ .\;Let $e_1,e_2$ (resp.,\;$e_3,e_4$) be a basis of $D_{\dr}(E_1)$ (resp.,\;$D_{\dr}(E_2)$).\;We list the Hodge filtration $\fil_{\bullet}^{H}D$  of $D$ under the basis $e_1,e_2,e_3,e_4$ (which is determined by the $p$-adic Hodge parameters $\cL_{12},\;\cL_{23},\;\cL_{13},\;\cL_{34}$ and $\cL_{14}$):
%\begin{equation*}\label{fil1}
%	\fil_{i}^{H}(D)=\left\{
%	\begin{array}{ll}
%		D,\;&i\leq -h_1\\
%		E(e_2+\cL_{12}e_1)\oplus \fil^H_{-h_3}(D),\;&-h_1<i\leq -h_2\\
%		E(e_3+\cL_{23}e_2+\cL_{13}e_1)\oplus \fil^H_{-h_4}(D),\;&-h_2<i\leq -h_3\\
%		E(e_4+\cL_{34}e_3+e_2+\cL_{14}e_1),\;& -h_3<i\leq -h_4,\\				0,\;& i>-h_4.
%	\end{array}
%	\right.
%\end{equation*}
%Thus the Hodge parameter $\cL_{12}$ (resp.,\;$\cL_{34}$) is contained in the block $E_1$ (resp.,\;$E_2$) and  $\cL_{14},\cL_{23},\cL_{13}$ are the so-called higher $\cL$-invariants.\;Secondly,\;under the basis $e_3,e_4,e_1,e_2$,\;the Hodge filtration can be rewritten by 
%\begin{equation*}\label{fil2}
%	\fil_{i}^{H}(D)=\left\{
%	\begin{array}{ll}
%		D,\;&i\leq -h_1\\
%		E(e_4-(\frac{(1-\cL_{23}\cL_{34})(\cL_{12}-\cL_{14})+\cL_{13}-\cL_{14}\cL_{23}}{\cL_{23}\cL_{12}-\cL_{13}})e_3)\oplus \fil^H_{-h_3}(D),\;&-h_1<i\leq -h_2\\
%		E(e_1-\frac{\cL_{23}}{\cL_{13}-\cL_{14}\cL_{23}}e_4+\frac{1-\cL_{23}\cL_{34}}{\cL_{13}-\cL_{14}\cL_{23}}e_3)\oplus \fil^H_{-h_4}(D),\;&-h_2<i\leq -h_3\\
%		E(e_2+\cL_{14}e_1+e_4+\cL_{34}e_3),\;& -h_3<i\leq -h_4,\\				0,\;& i>-h_4.
%	\end{array}
%	\right.
%\end{equation*}
%Therefore,\;in $E_1'$ (resp.,\;$E_2'$),\;we see the higher parameter $\cL_{14}$ (resp.,\;a relationship between higher parameters $\cL_{14}$,\;$\cL_{23}$ and $\cL_{13}$).\;In this case,\;$\pi^+_{1}(\underline{x},\bh)$ already determines $\Dpik$ except another relationship between higher parameters $\cL_{14}$,\;$\cL_{23}$ and $\cL_{13}$,\;the author suspects that the remainder one relationship is related to the unknown so-called supersingular consistents.\; 









\section*{Acknowledgment}
The author thank Yiwen Ding and Zicheng Qian for discussions or answers to questions.\;

\section{Preliminaries}

\subsection{General notation}


\noindent Let $L$ (resp.\;$E$) be a finite extension of $\bQ_p$ with $\co_L$ (resp.\;$\co_E$) as its ring of integers and $\varpi_L$ (resp.\;$\varpi_E$) a uniformizer.\;Suppose $E$ is sufficiently large containing all the embeddings of $L$ in $\overline{\BQ}_p$.\;Put $\Sigma_L:=\{\sigma: L\hookrightarrow \overline{\BQ}_p\} =\{\sigma: L\hookrightarrow E\}$.\;Let $\val_L$ (resp. $\val_p$) be the $p$-adic valuation on $\overline{\bQ}_p$ normalized by sending uniformizers of $\co_L$ (resp.,\;the ring $\bZ_p$ of integers in $\bQ_p$) to $1$.\;Let $d_L:=[L:\bQ_p]=|\Sigma_L|$ and $q_L:=p^{f_L}=|\co_L/\varpi_L|$.\;Let $\gal_{L}:=\gal(\overline{L}/L)$ be the absolute Galois group of $L$.\;


Let $\cR_{L}:=\bB_{\rig,L}^{\dagger}$ be the Robba ring associated to $L$.\;For $\bQ_p$-affinoid algebra (for example $A\in\{E,E[\epsilon]/\epsilon^2\}$),\;let $\cR_{A,L}:=\cR_{L}\widehat{\otimes}_{\bQ_p} A$ for the Robba ring associated to $L$ with $A$-coefficient.\;We write $\cR_{A,L}(\delta_A)$ for the $(\varphi,\Gamma)$-module over $\cR_{A,L}$ of rank one associated to a continuous character $\delta_A:L^\times\rightarrow A^\times$.\;Let $D$ be a $(\varphi,\Gamma)$-module over $\cR_{A,L}$,\;we write $D(\delta_A):=D\otimes_{\cR_{A,L}}\cR_{A,L}(\delta_A)$ for simplicity.\;

For a Lie algebra $\fg$ over $L$,\;and $\sigma\in \Sigma_L$,\;let $\fg_{\sigma}:=\fg\otimes_{L,\sigma} E$ (which is  a Lie algebra over $E$).\;


Let $\GL_n$ be the general linear group over $L$.\;Let $\Delta:=\Delta_n$ be the set of simple roots of $\GL_n$ (with respect to the Borel subgroup $\bB$ of upper triangular matrices), and we identify the set $\Delta$ with $\{1,\cdots, n-1\}$ such that $i\in \{1,\cdots, n-1\}$ corresponds to the simple root $\alpha_{i}:\;(x_1,\cdots, x_n)\in \ft \mapsto x_{i}-x_{i+1}$, where $\ft$ denotes the $L$-Lie algebra of the torus $\bT$ of diagonal matrices.\;Each $I\subseteq \Delta$ gives rise to a (standard) Levi subgroup $\bL_I\subseteq \GL_n$ such that $\bT\subseteq \bL_I$ and $I$ is equal to the set of simple roots of $\mathbf{L}_{I}$.\;Let $\mathbf{P}_{I}:=\mathbf{L}_{I}\bB$ be the (standard) parabolic subgroup of
$\GL_n$ containing $\bB$.\;In particular,\;we have $\mathbf{P}_{\Delta}=\GL_n$, $\mathbf{P}_{\emptyset}=\mathbf{B}$.\;Let $\overline{\mathbf{P}}_{I}$ be the parabolic subgroup opposite to $\mathbf{P}_{I}$.\;Let $\bN_{I}$ (resp.\;$\overline{\bN}_{I}$) be the unipotent radical of $\bP_{I}$ (resp.\;$\overline{\bP}_{I}$).\;We have Levi decompositions $\bP_I=\bL_I\bN_I$,\;$\overline{\bP}_I=\bL_I\overline{\bN}_I$ and $\Delta \backslash  I$ are precisely the simple roots of the unipotent radical $\bN_{I}$.\;Let $\bZ$ (resp.,\;$\bZ_I$) be the center of $\GL_n$ (resp.,\;$\mathbf{L}_{I}$).\;Let $\fg$,\;$\fb$,\;$\fp_{I}$,\;$\fl_{I}$,\;$\fn_I$ and $\fz_I$ be the $L$-Lie algebras of $\GL_n$,\;$\bB$,\;$\mathbf{P}_{I}$, $\mathbf{L}_{I}$,\;$\bN_I$ and $\bZ_I$ respectively.\;We put $G:=\GL_n(L)$.\;



Let $m\in\BZ_{\geq 1}$,\;and $\pi$ be an irreducible smooth admissible representation of $\GLN_m(L)$,\;let $\rec(\pi)$ be the $m$-dimensional absolutely irreducible $F$-semi-simple Weil-Deligne representation of the Weil group $W_L$ via the normalized classical local Langlands correspondence (normalized in \cite{scholze2013local}).\;We normalize the reciprocity isomorphism $\rec:L^\times\rightarrow W_L^{\mathrm{ab}}$ of local class theory such that the uniformizer $\varpi_{L}$ is mapped to a geometric Frobenius morphism,\;where $W_L^{\mathrm{ab}}$ is the abelization of the Weil group $W_L\subset \gal_L$.\;

For a group $A$ and $a\in A$,\;denote by $\unr(\alpha)$ the unramified character of $L^\times$ sending uniformizers to $\alpha$.\;If $\bk:=(\bk_{\tau})_{\tau\in \Sigma_L}\in \BZ^{\Sigma_L}:=\prod_{\sigma\in \Sigma_L}\BZ$,\;we denote $z^{\bk}:=\prod_{\tau\in  \Sigma_L}\tau(z)^{\bk_{\tau}}$.\;For a character $\chi$ of $\cO_L^{\times}$,\;denoted by $\chi_{\varpi_L}$ the character of $L^{\times}$ such that $\chi_{\varpi_L}|_{\cO_L^{\times}}=\chi$ and $\chi_{\varpi_L}(\varpi_L)=1$.\;Let ${\ccyc}:\gal_L\rightarrow \bZ_p^\times$ be the $p$-adic cyclotomic character.\;Then we have ${\ccyc}\circ\rec=\unr(q_L^{-1})\prod_{\tau\in \Sigma_L}\tau$ by local class theory.\;Let $A$ be an affinoid $E$-algebra.\;A locally $\bQ$-analytic character $\delta:L^\times\rightarrow A^\times$ induces a $\bQ_p$-linear map $L\rightarrow A$,\;$x\mapsto \frac{d}{dt}\delta(\exp(tx))|_{t=0}$ and hence it induces an $E$-linear map $L\otimes_{\bQ_p} E=\prod_{\tau\in \Sigma_L}E\rightarrow A$.\;There exist $\wt(\delta):=(\wt_{\tau}(\delta))_{\tau\in \Sigma_L}$ such that the latter map is given by $(a_{\tau})_{\tau\in \Sigma_L}\mapsto\sum_{\tau\in \Sigma_L}a_{\tau}\wt_{\tau}(\delta)$.\;We call $\wt(\delta)$ the weight of $\delta$.\;




Let $\ul{\lambda}:=(\lambda_{\sigma,1}, \cdots, \lambda_{\sigma,n})_{\sigma\in \Sigma_L}$ be a weight of $\ft_{L}$.\;For $I\subseteq \Delta$,\;we call that $\ul{\lambda}$ is $I$-dominant (resp.,\;strictly $I$-dominant)  if $\lambda_{\sigma,i}\geq \lambda_{\sigma,i+1}$ (resp.,\;$\lambda_{\sigma,i}> \lambda_{\sigma,i+1}$)   for all $i\in I$ and $\sigma\in \Sigma_L$.\;We denote by $X_{I}^+$ the set of $I$-dominant integral weights  of $\ft_{\Sigma_L}$.\;For $\ul{\lambda}\in X_{I}^+$, there exists a unique irreducible algebraic representation, denoted by $L_{I}(\ul{\lambda})$, of $\res_{L/
\bQ_p}\bL_{I}$ with highest weight $\ul{\lambda}$ with respect to $\res_{L/
\bQ_p}\bL_{I}\cap \bB$,\;so that $\overline{L}_{I}(-\ul{\lambda}):=(L_{I}(\ul{\lambda}))^\vee$ is the irreducible algebraic representation of $(\bL_{I})_{L}$ with highest weight $-\ul{\lambda}$ with respect to $\res_{L/
\bQ_p}\bL_{I}\cap \ob$.\;If $\ul{\lambda}\in X_{\Delta}^+$, let $L(\ul{\lambda}):=L(\ul{\lambda})_{\Delta}$.\;A $\bQ_p$-algebraic representation of $G$ over $E$ is the induced action of $G\subset \GL_{n,L}(E)$ on an algebraic representation of $\res_{L/
\bQ_p}\GL_{n}$.\;Let $\ul{\lambda}\in X_I^+$ be an integral weight,\;denote by $M_I(\ul{\lambda}):=\text{U}(\fg_{L})\otimes_{\text{U}(\fp_{I,L})} L_{I}(\ul{\lambda})$ (resp.\;$\overline{M}_I(\ul{\lambda}):=\text{U}(\fg_{L})\otimes_{\text{U}(\overline{\fp}_{I,L})}L_{I}(-\ul{\lambda})$),\;the corresponding Verma module with respect to $\fp_{I,L}$ (resp.\;$\overline{\fp}_{I,L}$).\;Let $L(\ul{\lambda})$ (resp. $\overline{L}(-\ul{\lambda})$) be the unique simple quotient of $M(\ul{\lambda}):=M_{\emptyset}(\ul{\lambda})$ (resp. of $\overline{M}(-\ul{\lambda}):=\overline{M}_{\emptyset}(-\ul{\lambda})$).\;

%Actually, when $\ul{\lambda}\in X_{\Delta}^+$ (i.e.\;$-\ul{\lambda}\in X_{\Delta}^-$),\;$L(\ul{\lambda})$ is finite dimensional and isomorphic to the algebraic representation $L(\ul{\lambda})$ introduced above (hence there is no conflict of notation).\;We have $\overline{L}(-\ul{\lambda})\cong L(\ul{\lambda})^{\vee}$.\;

%In general,\;for any subset $I$ of $\Delta$,\;and $\ul{\lambda}\in X_{I}^+$,\;we define the generalized parabolic Verma module $M_{I}(\ul{\lambda}):=\text{U}(\fg_{\Sigma_L})\otimes_{\text{U}(\fp_{I,L})} L(\ul{\lambda})_{I}$ (resp. $\overline{M}_{I}(-\ul{\lambda}):=\text{U}(\fg_{\Sigma_L})\otimes_{\text{U}(\overline{\fp}_{I,L})} \overline{L}(-\ul{\lambda})_{I})$ with respect to $\fp_{I,L}$ (resp.\;$\overline{\fp}_{I,L}$).
%see \cite[Chapter.\;9]{humphreysBGG} for precise definition.\;


Denote by $\sW_n$ ($\cong S_n$) the Weyl group of $\GL_n$, and denote by $s_{i}$ the simple reflection corresponding to $i\in \Delta$.\;For any $I\subset \Delta$,\;define $\sW_{I}$ to be the subgroup of $\sW_{n}$ generated by simple reflections $s_{i}$ with $i\in I$ (so that $\sW_I$ is the Weyl group of $\bL_I$).\;Let $I,J\subseteq\Delta$,\;recall that $\sW_{I}\backslash \sW_n/\sW_{J}$ has a canonical set of representatives,\;which we will denote by $\sW^{I,J}_n$ (resp.,\;$\sW^{I,J}_{n,\max}$),\;given by taking in each double coset the elements of minimal (resp.,\;maximal) length.\;Let $w_0:=w_{0,n}$  be the longest elements in $\sW_{n}$.\;


%More general,\;for subset $S\subseteq \Sigma_L$ and  $\underline{I}=\prod_{\sigma\in S}I_{\tau}\subseteq \Delta^{S}$ and $\underline{J}=\prod_{\sigma\in S}J_{\tau}\subseteq \Delta^{S}$,\;we put $\sW_{\underline{I},S}:=\prod_{\sigma\in S} \sW_{I_\sigma}$,\;$\sW^{\underline{I},\underline{J}}_{n,S}:=\prod_{\sigma\in S}\sW^{I_{\sigma},J_{\sigma}}_{n}$  and $\sW^{\underline{I},\underline{J}}_{n,S,\max}:=\prod_{\sigma\in S}\sW^{I_{\sigma},J_{\sigma}}_{n,\max}$.\;In particular,\;if $\underline{I}=\underline{I}_S$ and $\underline{J}=\underline{J}_S$ for $I,J\subseteq\Delta$,\;we write $\sW_{{I},S}:=\sW_{\underline{I}_S,S}$ $\sW^{{I},{J}}_{n,S}:=\sW^{\underline{I}_S,\underline{J}_S}_{n,S}$ and  $\sW^{{I},{J}}_{n,S,\max}:=\sW^{\underline{I}_S,\underline{J}_S}_{n,S,\max}$ for simplicity.\;If $S=\Sigma_L$,\;we replace the subscript $S$ by $L$.\;
%For $i\in \Delta$ and $\sigma\in \Sigma_L$,\;let $s_{i,\sigma}\in \sW_{n,\sigma}$ be the simple reflection corresponding to $i\in \Delta$.\;For any $I\subset \Delta$ ,\;define $\sW_{I,\sigma}$ to be the $\sigma$-copy of $\sW_{I}$,\;i.e.,\;to be the subgroup of $\sW_{n,\sigma}$ generated by simple reflections $s_{i,\sigma}$ with $i\in I$.\;For $S\subseteq \Sigma_L$ and  $I\subset {\Delta}$, we put $\sW_{I,S}:=\prod_{\sigma\in S} \sW_{I,\sigma}$.\;If $\ul{i}_{S}:=(i_{\sigma})_{\sigma\in S}\in {\Delta}^{|S|}$,\;denote by $s_{\ul{i}_S}:=\prod_{\sigma\in S}s_{i_\sigma,\sigma}$.\;




%When $J=\emptyset$,\;we put $^{I}\sW_{n}=[\sW_{I}\backslash\sW_{n}/\sW_\emptyset]^{\min}$.\;For $S\subseteq \Sigma_L$ and  $I\subset {\Delta}$,\;we put $^{I}\sW_{n,S}:=\prod_{\sigma\in S}{^{I}\sW_{n,\sigma}}$.\;

%If $H$ is a closed subgroup of $G$ and $(\rho,V)$ is an abstract representation of $H$ on $E$-vector space,\;we put $H^w:=w^{-1}Hw$ (resp.,\;${}^wH:=wHw^{-1}$),\;and $\rho^w(h):=\rho(whw^{-1})$.\;Then $(\rho^w,V)$ forms an abstract representation of $H^w$ on the same $E$-vector space $V$.;Moreover,\;we have relative Bruhat decompositions\[G=\coprod_{w\in [\sW_{r,I}\backslash \sW_n/\sW_{r,J}]}\bP_{{I}}(L)w\bP_{{J}}(L)=\coprod_{w\in [\sW_{r,I}\backslash \sW_n/\sW_{r,J}]}\op_{{I}}(L)w\op_{{J}}(L).\]




If $V$ is a continuous representation of $G$ over $E$, we denote by $V^{\bQ_p-\ana}$ its locally $\bQ_p$-analytic vectors.\;If $V$ is a locally $\bQ_p$-analytic representation of $G$, we denote by $V^{\mathrm{sm}}$ (resp.\;$V^{\mathrm{lalg}}$) the smooth (resp., locally $\bQ_p$-algebraic) subrepresentation consisting of its smooth (resp., locally $\bQ_p$-algebraic) vectors; see \cite{Emerton2007summary} for details.\;Let $\pi_P$ be either a locally $\bQ_p$-analytic representation of $P$ on a locally convex $E$-vector space of compact type or a smooth representation of $P$ over $E$.\;We denote by
\begin{equation}\label{smoothadj2}
	\begin{aligned}
		&(\mathrm{Ind}_{P}^{G}\pi_P)^{\bQ_p-\ana}:=\{f:G\rightarrow \pi_P \text{\;locally $\bQ_p$-analytic},\;f(pg)=pf(g)\},\\
		&\text{resp.,\;}(\mathrm{Ind}_{P}^{G}\pi_P)^{\infty}=\{f:G\rightarrow \pi_P \text{\;smooth},\;f(pg)=pf(g)\}
	\end{aligned}
\end{equation}
the locally $\bQ_p$-analytic parabolic induction (resp., the unnormalized smooth parabolic induction) to $G$.\;It becomes a locally $\bQ_p$-analytic representation (resp., a smooth representation) of $G$ over $E$ by equipping it with the left action of $G$ by right translation on functions: $(gf)(g')=f(g'g)$.\;


%Let $\mathbf{WD}_{L'/L,E}$ be the category of representations $(r,N,V)$ of $W_L$,\;on an $E$-vector space $V$ of finite dimension such that $r$ is unramified when restricted to the $W$.\;Let $\mathbf{DF}_{L'/L,E}$ be the category of Deligne-Fontaine modules,\;i.e.,\;the category of quadruples $(\varphi,N,\gal(L'/L),D)$ where $D$ is an $L'_0\otimes_{\bQ_p}E$-module free of finite rank,\;which is endowed with a Frobenius $\varphi:D\rightarrow D$ (resp.,\;an $L'_0\otimes_{\bQ_p}E$-linear endomorphism $N:D\rightarrow D$) such that $N\varphi=p\varphi N$ and an action of $\gal(L'/L)$ commuting with $\varphi$ and $N$ such that $g((l\otimes e)d)=(g(L)\otimes e)d$ for $g\in \gal(L'/L),\;l\in L'_0,\;e\in E,\;d\in D$.\;Then the  Fontaine's theory asserts that there is a functor $\mathrm{WD}:\mathbf{WD}_{L'/L,E}\rightarrow \mathbf{DF}_{L'/L,E}$ gives an equivalence of categories (\cite[Proposition 4.1]{breuil2007first}).\;


Let $r$ be an integer, and let $\Omega_r$ be a cuspidal Bernstein component of $\GLN_r(L)$.\;Let $\FZ_{\Omega_{r}}$ be the rational Bernstein centre over $E$ (see \cite[Section 3.2]{PATCHING2016}).\;For a closed point $x\in \Spec\FZ_{\Omega_{r}}$, let $\pi_x$ be the associated irreducible cuspidal smooth representation over $k(x)$.\;Then, via the normalized classical local Langlands correspondence (see \cite{scholze2013local}), $\pi_x$ gives rise to an $r$-dimensional absolutely irreducible Weil--Deligne representation $\wdre_x:=\rec(\pi_x)$ of $W_L$ over $k(x)$.\;$\wdre_x$ corresponds to an absolutely irreducible Deligne--Fontaine module $\df_x$, by Fontaine's equivalence of categories as in \cite[Proposition 4.1]{breuil2007first}, and a $p$-adic differential equation ${\Delta_x}$ over $\cR_{k(x),L}$.\;Here ${\Delta_x}$ is a de Rham $(\varphi,\Gamma)$-module of rank $r$ over $\cR_{k(x),L}$ with constant Hodge--Tate weights $0$ such that, after forgetting the Hodge filtration, $D_{\mathrm{pst}}(\Delta_x)$ is isomorphic to $\df_x$, by Berger's theory \cite[Theorem A]{berger2008equations}.\;Assume that $\wdre_x$ is unramified when restricted to $W_{L'}$ for some finite Galois extension $L'$ of $L$.\;Then $\df_x=(\varphi_x,N=0,\Gal(L'/L),\df_x)$, where
$\varphi_x$ is the Frobenius semilinear operator on $\df_x$.\;

Throughout the paper, we use $\bullet-\bullet$ to denote an extension of two objects, such as $(\varphi,\Gamma)$-modules or $G$-representations, where the first (resp., second) object is the subobject (resp., quotient).\;
 






\subsection{\texorpdfstring{Potentially crystalline $(\varphi,\Gamma)$-modules over $\cR_{E,L}$}{Potentially crystalline (phi,Gamma)-modules over R E,L}}\label{Omegafil}


We give a quick review, following \cite[Section 2.3]{Ding2021}, of the structure of \textit{non-critical and generic} potentially crystalline $(\varphi,\Gamma)$-modules over $\cR_{E,L}$.\;


Throughout this article, fix $S_0\subseteq\Delta$ and write $\bL_{S_0}:=\prod_{i=1}^s\GLN_{r_i}$ for integers $r_1,\cdots,r_s$ such that $n=r_1+\cdots+r_s$.\;For $1\leq i\leq s$, fix a cuspidal Bernstein component $\Omega_{i}$ of $\GLN_{r_i}$.\;Then $\Omega_{S_0}:=\prod_{i=1}^s\Omega_{i}$ is a cuspidal Bernstein component of $\bL_{S_0}$.\;For $u\in \sW_s$, let $S_0^u$ be the subset of $\Delta$ such that $\bL_{S_0^u}=\prod_{i=1}^s\GLN_{r_{u^{-1}(i)}}$.\;Put $\Omega_{S^u_0}:=\prod_{i=1}^s\Omega_{u^{-1}(i)}$, which is a cuspidal Bernstein component of $\bL_{S^u_0}$.\;For $1\leq i\leq s$, put $t_0^u=0$ and $t_i^u=\sum_{j=1}^ir_{u^{-1}(j)}$, so that $t^u_s=n$.\;Then $\Delta\backslash S_0^u=\{t_i^u:i\in \Delta_s\}$, where $\Delta_s=\{1,\cdots,s-1\}$.\;

%We also identify $\Delta\backslash S_0^u$ with $\Delta_s$ via the map, denoted by $J_u$, sending $i\in \Delta_s$ to $t^u_i\in \Delta\backslash S_0$; thus $\Delta\backslash S_0^u=J_u(\Delta_s)$.\;

%In this paper, we further assume that $r_1\leq r_2\leq\cdots\leq r_s$.\;



Let $\Dpik$ be a potentially crystalline $(\varphi,\Gamma)$-module over $\cR_{E,L}$ of rank $n$.\;Let $L'$ be a finite Galois extension of $L$ such that $\Dpik|_{L'}:=\Dpik\otimes_{\cR_{E,L}}\cR_{E,L'}$ is a crystalline $(\varphi,\Gamma)$-module over $\cR_{E,L'}$ of rank $n$.\;We consider the filtered  Deligne--Fontaine module associated with $\Dpik$:
\[\df_{\Dpik}=(\varphi,\Gal(L'/L),D_{\pcr}(\Dpik),\fil_H^{\bullet}(D_{L'})),\]
where $D_{\pcr}(\Dpik)=D_{\mathrm{cr}}^{L'}(\Dpik|_{L'})$ is a finite free $L'_0\otimes_{\bQ_p}E$-module of rank $n$, and $L'_0$ is the maximal unramified subextension of $L'$ over $\bQ_p$.\;The $\varphi$-action on $D_{\pcr}(\Dpik)$ is induced by the $\varphi$-action on $B_{\mathrm{cris}}$, and the $\Gal(L'/L)$-action on $D_{\pcr}(\Dpik)$ is the residual action of $\gal_L$.\;Moreover, $D_{L'}:=D_{\pcr}(\Dpik)\otimes_{L_0'}L'$ and the  filtered structure $\fil_H^{\bullet}(D_{L'})$ is  induced by the filtration on $B_{\mathrm{cris}}$.\;Note that $D_{L'}=\prod_{\sigma\in \Sigma_L}D_{L',\sigma}$ with $D_{L',\sigma}:=D_{L'}\otimes_{L,\sigma}E$ a $L'\otimes_{L,\sigma}E$-module  and $\fil_H^{\bullet}(D_{L'})=\prod_{\sigma\in \Sigma_L}\fil_H^{\bullet}(D_{L',\sigma})$.\;Since $\Dpik$ is potentially crystalline, it is de Rham.\;Hence $D_{\dr}(\Dpik)\cong D_{L'}^{\Gal(L'/L)}$ is a free $L\otimes_{\bQ_p}E$-module of rank $n$,\;which  is equipped with a natural Hodge filtration  $\fil^{\bullet}_{H}D_{\dr}(\Dpik):=(\fil^{\bullet}_{H}(D_{L'})^{\Gal(L'/L)}$.\;In this article, we define the Hodge--Tate weights of a de Rham representation as the opposites of the jumps of the filtration on the covariant de Rham functor, so that the Hodge--Tate weight of ${\ccyc}$ is $1$.\;Assume that $D_{\dr}(\Dpik)$ has distinct Hodge--Tate weights $\bh:=(\hpi_{1,\sigma}>\hpi_{2,\sigma}>\cdots>\hpi_{n,\sigma} )_{\sigma\in \Sigma_L}$.\;Denote $\hpi_{i}=(\hpi_{i,\sigma})_{\sigma\in \Sigma_L}$ for $1\leq i\leq n$.\;

Let $\wdre_{\Dpik}$ be the Weil--Deligne representation associated with $\df_{\Dpik}$ by Fontaine's equivalence of categories; see \cite[Proposition 4.1]{breuil2007first}.\;We say that $\wdre_{\Dpik}$ admits an $\Omega_{S_0}$-filtration $\cF$ if $\wdre_{\Dpik}$ admits an increasing filtration $\cF$ by Weil--Deligne subrepresentations:
\[\cF=\fil_{\bullet}^{\cF} \wdre_{\Dpik}: \ 0 =\fil_0^{\cF} \wdre_{\Dpik} \subsetneq \fil_1^{\cF} \wdre_{\Dpik} \subsetneq \cdots \subsetneq \fil_{s}^{\cF} \wdre_{\Dpik}=\wdre_{\Dpik},\]
such that the irreducible smooth cuspidal representation associated with $\gr^{\cF}_{i}\wdre_{\Dpik}$ lies in the cuspidal Bernstein component $\Omega_{i}$ for all $1\leqslant i\leqslant s$.\;By \cite[Proposition 4.1]{breuil2007first}, the $\Omega_{S_0}$-filtration $\cF$ on $\wdre_{\Dpik}$ corresponds to an increasing $\Omega_{S_0}$-filtration on $\df_{\Dpik}$, still denoted by $\cF$, by Deligne--Fontaine submodules
\[\cF=\fil_{\bullet}^{\cF}\df_{\Dpik}: \ 0 =\fil_0^{\cF} \df_{\Dpik} \subsetneq \fil_1^{\cF} \df_{\Dpik} \subsetneq \cdots \subsetneq \fil_{s}^{\cF} \df_{\Dpik}=\df_{\Dpik},\]
such that $\fil_{i}^{\cF} \df_{\Dpik}$ corresponds to $\fil_{i}^{\cF} \wdre_{\Dpik}$ via \cite[Proposition 4.1]{breuil2007first}.\;Thereofer,\;for $u\in\sW_s$, we also have an $\omepik$-filtration $\cF_u$ on $\wdre_{\Dpik}$ (resp.,\;$\df_{\Dpik}$) by Weil--Deligne subrepresentations (resp.,\;Deligne--Fontaine submodules):
\begin{equation}
	\begin{aligned}
		\cF_u=\fil_{\bullet}^{\cF_u} \wdre_{\Dpik}: &\ 0 =\fil_0^{\cF_u} \wdre_{\Dpik} \subsetneq \fil_1^{\cF_u}\wdre_{\Dpik} \subsetneq \cdots \subsetneq \fil_{s}^{\cF_u} \wdre_{\Dpik}=\wdre_{\Dpik},\\
		\cF_u=\fil_{\bullet}^{\cF_u}\df_{\Dpik}: &\ 0 =\fil_0^{\cF_u} \df_{\Dpik} \subsetneq \fil_1^{\cF_u} \df_{\Dpik} \subsetneq \cdots \subsetneq \fil_{s}^{\cF_u} \df_{\Dpik}=\df_{\Dpik}
	\end{aligned}
\end{equation}
such that $\gr_{i}^{\cF_u} \df_{\Dpik}$ is associated with $\gr_{i}^{\cF_u} \wdre_{\Dpik}=\gr_{u^{-1}(i)}^{\cF}\wdre_{\Dpik}$.\;


Let ${\Delta_\Dpik}$ be the $p$-adic differential equation over $\cR_{E,L}$ associated with $\df_{\Dpik}$.\;The $\Omega_{S^u_0}$-filtration on $\df_{\Dpik}$ induces an $\Omega_{S^u_0}$-filtration $\fil_{\bullet}^{\cF_u}{\Delta_\Dpik}=\{\fil_{i}^{\cF_u}{\Delta_\Dpik}\}$ on ${\Delta_\Dpik}$ by saturated $(\varphi,\Gamma)$-submodules over $\cR_{E,L}$, such that $\fil_{i}^{\cF_u}{\Delta_\Dpik}$ is the $p$-adic differential equation over $\cR_{E,L}$ associated with $\fil_{i}^{\cF_u} \df_{\Dpik}$.\;Consider the $(\varphi,\Gamma)$-module $\cM_\Dpik=\Dpik[1/t]\cong \Delta_\Dpik [1/t]$ over $\cR_{E,L}[1/t]$.\;After inverting $t$, the filtration $\cF_u$ on $\Delta_\Dpik$ induces an increasing filtration $\fil_{i}^{\cF_u}\cM_\Dpik:=\fil_{i}^{\cF_u}{\Delta_\Dpik}[1/t]$ on $\cM_\Dpik$ by $(\varphi,\Gamma)$-submodules over $\cR_{E,L}\big[\frac{1}{t}\big]$.\;Therefore, the filtration $\cF_u$ on $\cM_\Dpik$ induces an increasing $\bP_{S^u_0}$-filtration on $\Dpik$:
\[\cF_u=\fil_{\bullet}^{\cF_u}\Dpik: \ 0 =\fil_0^{\cF_u}\Dpik \subsetneq \fil_1^{\cF_u}\Dpik \subsetneq \cdots \subsetneq \fil_{s}^{\cF_u} \Dpik=\Dpik,\;\fil_{i}^{\cF_u}\Dpik=(\fil_{i}^{\cF_u}\cM_\Dpik)\cap \Dpik,\]
by saturated $(\varphi,\Gamma)$-submodules of $\Dpik$ over $\cR_{E,L}$.\;

For each $\sigma\in \Sigma_L$,\;the natural Hodge filtration on $D_{\dr}(\Dpik)_{\sigma}:=D_{\dr}(\Dpik)\otimes_{L,\sigma}E\cong D_{L',\sigma}^{\Gal(L'/L)}$ can be expressed by the following complete flag :
\[\fil^{\bullet}_{H}D_{\dr}(\Dpik)_{\sigma}: \ 0 \subsetneq \fil^{-\hpi_{n,\sigma}}_{H} D_{\dr}(\Dpik)_{\sigma} \subsetneq \fil^{-\hpi_{n-1,\sigma}}_{H}D_{\dr}(\Dpik)_{\sigma} \subsetneq \cdots \subsetneq \fil^{{-\hpi_{1,\sigma}}}_{H}D_{\dr}(\Dpik)_{\sigma}=D_{\dr}(\Dpik)_{\sigma}.\]
On the other hand, the $\Omega_{S^u_0}$-filtration $\cF_u$ on $\df_{\Dpik}$ induces an $\Omega_{S^u_0}$-filtration $\cF_u$ on $D_{\dr}(\Dpik)$ by free $L\otimes_{\bQ_p}E$-submodules $\fil_{\bullet}^{\cF_u}D_{\dr}(\Dpik):=(\fil_{\bullet}^{\cF_u}\df_{\Dpik}\otimes_{L_0'}L')^{\Gal(L'/L)}$.\;By Berger's equivalence of categories, $\fil_{i}^{\cF_u}\Dpik$ corresponds to the filtered Deligne--Fontaine module $\fil_{i}^{\cF_u} \df_{\Dpik}$ equipped with the filtration induced from the Hodge filtration on $D_{\pcr}(\Dpik)$.\;The \textit{non-critical} assumption means that the Hodge--Tate weights of $\fil_{i}^{\cF_u}\Dpik$ (resp.,\;$\gr_{i}^{\cF_u} \Dpik$) are given by
\begin{equation}
	\begin{aligned}
		&\{\hpi_{1,\sigma},\hpi_{2,\sigma},\cdots,\hpi_{t_i^u,\sigma}\}_{\sigma\in \Sigma_L},\;
		\text{resp.,\;}&\{\hpi_{t^u_{i-1}+1,\sigma},\hpi_{t^u_{i-1}+2,\sigma},\cdots,\hpi_{t^u_i,\sigma}\}_{\sigma\in \Sigma_L}.\;
	\end{aligned}
\end{equation}
Using Berger's equivalence of categories \cite[Theorem A]{berger2008equations} and comparing the weights (or see \cite[(2.4)]{Ding2021}), we obtain an injection of $(\varphi,\Gamma)$-modules over $\cR_{E,L}$ of rank $r_i$:
\begin{equation}\label{Dpikinjectionu}
	\mathbf{I}^u_{i}:\gr_{i}^{\cF_u} \Dpik \hookrightarrow \gr^{\cF_u}_{i}{\Delta_\Dpik}\tee \cR_{E,L}(z^{\bh_{t^u_i}}),
\end{equation}
for $i=1,\cdots,s$.\;Put $E_{u,i}:=\gr_{i}^{\cF_u}\Dpik$ and $E_{i}:=E_{1,i}$.\;Then $\cF_u$ (resp.,\;$\cF:=\cF_1$) corresponds to the following successive  extension of $E_{u,i}$ (resp.,\;$E_{i}:=E_{1,i}$) for $1\leq i\leq s$ :
\begin{equation}\label{omegafilforu}
\Dpik=[E_{u,1}-E_{u,2}-\cdots-E_{u,s}],\;\quad (\text{resp.,\;}\Dpik=[E_{1}-E_{2}-\cdots-E_{s}]).\;
\end{equation}
For $1\leq i\leq s$, let $\df_{E_i}:=(\varphi_i,\Gal(L'/L),D_{\pcr}(E_i))$ (resp.,\;$\Delta_{E_i}$) be the underlying Deligne--Fontaine module (resp.,\;$p$-adic differential equation) associated with $E_i$.\;Then $\gr_{i}^{\cF_u} \df_{\Dpik}\cong \df_{E_{u^{-1}(i)}}$ and $\gr^{\cF_u}_{i}{\Delta_\Dpik}\cong \Delta_{E_{u^{-1}(i)}}$ for $1\leq i\leq s$.\;

In the sequel, let $\mathcal{Z}_{\bL_{S^u_0},L}$ (resp.,\;$\mathcal{Z}_{\bL_{S^u_0},\cO_L}$) be the rigid space over $E$ parametrizing continuous characters of $\bL_{S^u_0}(L)$ (resp.,\;$\bL_{S^u_0}(\cO_L)$).\;

\begin{dfn}\label{noncriticalassumption}
\begin{itemize}
	\item[(1)] We say that $\Dpik$ is \textit{regular and generic} if :
	\[\df_{E_i}\neq \df_{E_j},(\varphi_i,\Gal(L'/L),D_{\pcr}(E_i))\neq (p\varphi_j,\Gal(L'/L),D_{\pcr}(E_j))\]
	for $i\neq j$; equivalently, $\Delta_{E_i}\neq \Delta_{E_j}(\unr(q_L))=\Delta_{E_j}\otimes_{\cR_{E}}\cR_{E}(\unr(q_L))$ for $i\neq j$.\;
	\item[(2)] Suppose that $\Delta_{E_i}\cong \Delta_{x_i}$ for some $x_i\in \Spec\FZ_{\Omega_{i}}(E)$.\;For $u\in \sW_s$, put
	\begin{equation}
		\begin{aligned}
			&\underline{x}^u:=(x_{u^{-1}(i)})\in \sbanpiku,\underline{\delta}^u=(z^{\bh_{t^u_i}})_{1\leq i\leq s}\in \mathcal{Z}_{\bL_{S^u_0},L},\\
			&\underline{x}^u_+:=(x'_{u^{-1}(i)})\in \sbanpiku,\underline{\delta}^{0,u}:=\underline{\delta}^u|_{\bL_{S^u_0}(\cO_L)}\in \mathcal{Z}_{\bL_{S^u_0},\cO_L},\;
		\end{aligned}
	\end{equation}
	where $x'_{u^{-1}(i)}\in \Spec\FZ_{\Omega_{u^{-1}(i)}}$ satisfies $\pi_{x'_{u^{-1}(i)}}\cong \unr(\varpi_L^{\bh_{t^u_i}})\pi_{x_{u^{-1}(i)}}$ for $1\leq i\leq s$.\;In the language of \cite[Definition 4.1.6]{Ding2021},\;The $\Omega_{S^u_0}$-filtration $\cF_u$ has parameters  $(\underline{x}^u,\underline{\delta}^u)\in \sbanpiku\times\rigchlu$ (resp.,\;$(\underline{x}^u_+,\underline{\delta}^{0,u})\in \sbanpiku\times \mathcal{Z}_{\bL_{S^u_0},\cO_L}$).\;	
\end{itemize}
	
	
	
	

	
\end{dfn}


%
%For each $\sigma\in \Sigma_L$, fix a basis of $D_{\dr}(\Dpik)_{\sigma}$ over $E$.\;Then $\fil^{\bullet}_{H}$ (resp.,\;$\fil_{\bullet}^{\cF}$) corresponds to an $E$-point $(g_{2,\sigma}\bB(E))_{\sigma\in \Sigma_L}\in \res_{L/\bQ_p}\GLN_{n}/\res_{L/\bQ_p}\bB$ (resp.,\;$(g_{1,\sigma}\bP_{S_0}(E))_{\sigma\in \Sigma_L}\in \res_{L/\bQ_p}\GLN_{n}/\res_{L/\bQ_p}\bP_{S_0}$).\;For each $\sigma\in \Sigma_L$, there is therefore a unique $w_{\cF,\sigma}\in \sW^{S_0,\emptyset}_{n,\max}$ such that
%\[(g_{1,\sigma}\bP_{S_0}(E),g_{2,\sigma}\bB(E))\in\GLN_n(E)(1,w_{\cF,\sigma})(\bP_{S_0}\times\bB)(E)\subset (\GLN_n/\bP_{S_0}\times\GLN_n/\bB)(E).\]





\subsection{(Parabolic) deformations of $(\varphi,\Gamma)$-modules}\label{genedeformations}


For a locally $L$-analytic group $H$, let $\homo(H,E)$ (resp.,\;$\homo_{\mathrm{sm}}(H,E)$) be the $E$-vector space of locally $\bQ_p$-analytic (resp., locally constant) $E$-valued characters on $H$.\;For $\sigma\in \Sigma_L$, let $\homo_{\sigma}(H,E)\subseteq \homo(H,E)$ be the subspace of locally $\sigma$-analytic $E$-valued characters on $H$.\;In particular, if $H=L^{\times}$, then by \cite[Section 1.3.1]{2015Ding}, $\homo(L^{\times},E)$ is generated by $\val_L$ and $\{\log_{p,\sigma}\}_{\sigma\in \Sigma_L}$, where $\log_p:L^{\times}\rightarrow L$ is the $p$-adic logarithm restricted to $\cO_L^{\times}$.\;Then $\homo_{\mathrm{sm}}(L^{\times},E)\cong E\val_L$, $\homo_{\sigma}(L^{\times},E)=E\langle\val_L,\log_{p,\sigma}\rangle$, and $\homo(L^{\times},E)=E\langle\val_L,\{\log_{p,\sigma}\}_{\sigma\in \Sigma_L}\rangle$.\;In particular, $\dim_E\homo(L^{\times},E)=1+d_L$, $\dim_E\homo_{\mathrm{sm}}(L^{\times},E)=1$, and $\dim_E\homo_{\sigma}(L^{\times},E)=2$.\;Moreover, a choice of $\log_p(p)\in E$ (we may choose $\log_p(p)=0$) defines a section $\homo_{\sigma}(\cO_L^{\times},E)\hookrightarrow \homo_{\sigma}(L^{\times},E)$ of the restriction map $\homo_{\sigma}(L^{\times},E)\rightarrow \homo_{\sigma}(\cO_L^{\times},E)$ by sending $\log_{p,\sigma}$ to its unique extension to $L^{\times}$ sending $p$ to $\log_p(p)$.\;

Let $\cG$ be an increasing $\bP_{\cG}$-parabolic filtration on $\Dpik$ by saturated $(\varphi,\Gamma)$-submodules of $\Dpik$ over $\cR_{E,L}$:
\[\cG=\fil_{\bullet}^{\cG}\Dpik: \ 0 =\fil_0^{\cG}\Dpik \subsetneq \fil_1^{\cG}\Dpik \subsetneq \cdots \subsetneq \fil_{h}^{\cG} \Dpik=\Dpik.\]
Put $M_i:=\gr_i^{\cG}\Dpik$; then $\Dpik=[M_1-M_2-\cdots-M_h]$.\;Let $\bL_{\cG}$ be the standard Levi subgroup of $\bP_{\cG}$, and let $\bZ_{\cG}\cong \BG_m^h$ be the center of $\bL_{\cG}$.\;Denote by $F_{\Dpik,\cG}$ the deformation functor
\[F_{\Dpik,\cG} :\Art(E):=\{\text{Artinian local $E$-algebra with residue field $E$}\}\longrightarrow \{\text{sets}\}\]
which sends $A$ to the set of isomorphism classes $\{(D_A,\pi_A,\cG_A)\}/\sim$, where
\begin{itemize}
	\item[(1)] $D_A$ is a $(\varphi,\Gamma)$-module over $\cR_{A,L}$ of rank $n$ with $\pi_A:D_A\otimes_AE\cong \Dpik$;
	\item[(2)] $\cG_A=\fil^{\cG_A}_{\bullet}D_A$ is an increasing filtration of $D_A$ by $(\varphi,\Gamma)$-submodules over $\cR_{A,L}$, such that the $\fil^{\cG_A}_{i}D_A$ are direct summands of $D_A$ as $\cR_{A,L}$-modules and $\pi_A(\fil^{\cG_A}_{i}D_A)\cong \fil_{i}^{\cG} \Dpik=M_i$.
\end{itemize}
Moreover, we identify $F_{\Dpik,\cG}(E[\epsilon]/\epsilon^2)$ with a subspace $\ext^1_{\cG}(\Dpik,\Dpik)\subseteq \ext^1(\Dpik,\Dpik)$ of $\cG$-parabolic deformations of $\Dpik$; namely, $\widetilde{D}\in \ext^1(\Dpik,\Dpik)$ belongs to $\ext^1_{\cG}(\Dpik,\Dpik)$ if and only if 
\[\widetilde{D}=[\widetilde{M}_{1}-\widetilde{M}_{2}-\cdots-\widetilde{M}_{h}],\;\]
where $\widetilde{M}_{i}$ is a deformation of $M_{i}$ over $E[\epsilon]/\epsilon^2$ for $1\leq i\leq h$.\;


We next recall the content of \cite[Section 4.1.1]{Ding2021}.\;For $1\leq i\leq h$, let $\Delta_{M_i}$ be the differential equation associated with $M_i$; it is a de Rham $(\varphi,\Gamma)$-module of rank $h_i$ over $\cR_{E,L}$ with constant Hodge--Tate weight $0$.\;Suppose that we have an injection $M_i\hooklongrightarrow \Delta_{M_i}\otimes_{\cR_{E,L}}\cR_{E,L}(\delta_i)$ of $(\varphi,\Gamma)$-modules over $\cR_{E,L}$ for some character $\delta_i:L^\times\rightarrow E^\times$.\;For each $\sigma\in \Sigma_L$, we usually assume that $0$ is the minimal $\sigma$-Hodge--Tate weight of $M_i\otimes_{\cR_{E,L}}\cR_{E,L}(\delta^{-1}_i)$.\;In \cite[Sections 4.1.1 \& 4.1.2]{Ding2021}, the authors consider the following functor:
\[F_{M_{i}}^0 (\text{resp.},\;F_{M_{i}}^{\flat}):\Art(E)\longrightarrow \{\text{sets}\}\]
which sends $A$ to the set of isomorphism classes $\{(D_A,\pi_A,\delta_A)\}/\sim$, where
\begin{description}
	\item[(1)] $D_A$ is a $(\varphi,\Gamma)$-module of rank $r$ over $\cR_{A,L}$ with $\pi_A:D_A\otimes_AE\cong M_i$;
	\item[(2)] ${\delta}_{A}:L^\times\rightarrow A^\times$ is a continuous character such that ${\delta}_{A}\equiv {\delta}_i\mod \mathfrak{m}_A$;
	\item[(3)] there exists an injection (resp., an isomorphism) $D_A\hooklongrightarrow \Delta_{M_i}\otimes_{\cR_{E,L}}\cR_{A,L}({\delta}_{A})$ (resp.,\;$D_A\cong M_i\otimes_{\cR_{E,L}}\cR_{A,L}({\delta}_{A})$) of $(\varphi,\Gamma)$-modules over $\cR_{A,L}$.
\end{description}


The map sending the data $\{(D_A,\pi_A,\cG_A)\}$ to $\{(\gr^{\cG_A}_iD_A,\pi_A|_{\gr_i^{\cG_A}D_A})\}$ induces a natural morphism $F_{\Dpik,\cF}\longrightarrow \prod_{i=1}^hF_{M_i}$.\;For $?\in\{\flat,0\}$, put
\[F_{\Dpik,\cG}^{?}=F_{\Dpik,\cG}\times_{\prod_{i=1}^h F_{M_i}}\prod_{i=1}^h F_{M_i}^{?},\]
i.e.\;$F_{\Dpik,\cG}^0$ (resp.,\;$F_{\Dpik,\cG}^{\flat}$) sends $A$ to the set of isomorphism classes $\{(D_A,\pi_A,\cG_A,\delta_A)\}/\sim$ satisfying the above conditions $(1)$--$(3)$ and
\begin{description}
	\item[(4)] For $1\leq i\leq h$, there is an injection (resp., an isomorphism) of $(\varphi,\Gamma)$-modules over $\cR_{A,L}$:
	\begin{equation}\label{injectiondefF0deforma}
		\gr^{\cG_A}_iD_A\hooklongrightarrow \Delta_{M_i}\otimes_{\cR_{E,L}}\cR_{A,L}({\delta}_{A,i}),\;\text{resp.\;}\gr^{\cG_A}_iD_A\cong M_{i}\otimes_{\cR_{E,L}}\cR_{A,L}({\delta}_{A,i}).\;
	\end{equation}
\end{description}
We have inclusions of functors $F_{\Dpik,\cG}^{\flat}\subseteq F_{\Dpik,\cG}^{0}\subseteq F_{\Dpik,\cG}$, and their $E[\epsilon]/\epsilon^2$-points correspond to the inclusions of $\ext^1$-spaces
\[\ext^{1,\flat}_{\cG}(\Dpik,\Dpik)\subseteq \ext^{1,0}_{\cG}(\Dpik,\Dpik)\subseteq \ext^1_{\cG}(\Dpik,\Dpik)\] 
such that $\widetilde{D}=[\widetilde{M}_{1}-\widetilde{M}_{2}-\cdots-\widetilde{M}_{h}]$ belongs to $\ext^{1,0}_{\cG}(\Dpik,\Dpik)$ (resp.,\;$\ext^{1,\flat}_{\cG}(\Dpik,\Dpik)$) if $\widetilde{M}_{i}\hookrightarrow  \Delta_{M_i}\otimes_{\cR_{E,L}}\cR_{E[\epsilon]/\epsilon^2,L}(\delta_i(1+\psi_i\epsilon))$ (resp.,\;$\widetilde{M}_{i}\cong M_{i}\otimes_{\cR_{E,L}}\cR_{E[\epsilon]/\epsilon^2}(\delta_i(1+\psi_i\epsilon))$) for some $\psi_i\in \homo(L^{\times},E)$.\;We thus have natural maps
\begin{equation}
	\begin{aligned}
		\kappa^{?}_{u}:\ext^{1,?}_{\cG}(\Dpik,\Dpik)\rightarrow\homo(\bZ_{\cG}(L),E),\;\widetilde{D}\mapsto (\psi_i)_{1\leq i\leq h}.\;
	\end{aligned}
\end{equation}
for $?\in\{\flat,0\}$.\;In particular, if $\rank(M_i)=1$ for all $1\leq i\leq h$, the functors $\{F_{\Dpik,\cG}^{?}\}_{?\in\{\emptyset,\flat,0\}}$ and the extension groups $\{\ext^{1,?}_{\cG}(\Dpik,\Dpik)\}_{?\in\{\emptyset,\flat,0\}}$ coincide, and we drop the symbols $\{\flat,0\}$ from the subscripts.\;


\section{Deformations of potentially crystalline $(\varphi,\Gamma)$-modules}

Throughout this section, we fix a non-critical generic potentially crystalline $(\varphi,\Gamma)$-module $\Dpik$ with Hodge--Tate weights $\bh$ in Section \ref{Omegafil}  and keep the notation of Section \ref{Omegafil}.\;





\subsection{Computations on extension groups}




For $u\in\sW_s$, recall the $\omepik$-filtration $\cF_{u}$, written as $\Dpik:=[E_{u,1}-E_{u,2}-\cdots-E_{u,s}]$ in (\ref{omegafilforu}).\;We apply the discussion in Section \ref{genedeformations} to $\cG=\cF_{u}$.\;For $?\in\{\emptyset,\flat,0\}$, write $\ext^{1,?}_{u}(\Dpik,\Dpik):=\ext^{1,?}_{\cF_{u}}(\Dpik,\Dpik)$ for simplicity.\;We have the following inclusions of extension groups:
$\ext^{1,\flat}_{u}(\Dpik,\Dpik)\subseteq\ext^{1,0}_{u}(\Dpik,\Dpik)\subseteq \ext^{1}_{u}(\Dpik,\Dpik)$ and the natural maps
\begin{equation}
	\begin{aligned}
		\kappa^{?}_{u}:\ext^{1,?}_{u}(\Dpik,\Dpik)\rightarrow\homo(\bZ_{S^u_0}(L),E)
		\end{aligned}
\end{equation}
where $\bZ_{S^u_0}\cong \BG_m^s$ is the center of $\bL_{S^u_0}$ and $?\in\{\flat,0\}$.\;By the genericity assumption, $\kappa^{\flat}_{u}$ and $\kappa^{0}_{u}$ are surjective.\;By \cite[Proposition 4.1.17,\;Corollary 4.1.8]{Ding2021} and its proof, we have
\begin{pro}\label{dimesionforParafilgeneric}
\begin{itemize}
	\item[(1)] $\dim_E\ext^1(\Dpik,\Dpik)=1+d_Ln^2$ and $\dim_E\ext^1_u(\Dpik,\Dpik)=1+d_L\dim \bP_{S^u_0}$.\; 
	\item[(2)]  $\dim_E\ext^{1,\flat}_{u}(\Dpik,\Dpik)=1+d_L(\dim\bN_{S^u_0}+s) $ and $\dim_E\ext^{1,0}_{u}(\Dpik,\Dpik)=1+d_L(\frac{n(n-1)}{2}+s)$.\;
\end{itemize}
\end{pro}	

%Thus,\;$\dim_E\homo_{w_{0,s}}(\bZ_{S^-_0}(L),E)=s(1+d_L)$ and $\dim_E\homo_1(\bZ_{S_0}(L),E)=s(1+d_L)$,\;where ${S^-_0}:=S^{w_{0,s}}_0$ and $\bZ_{S^-_0}(L)=\bZ^{w_{0,s}}_{S_0}(L)$.\;Put $\homo_{g'}(\bZ_{S^u_0}(L),E):=\{\psi\in \homo(\bZ_{S^u_0}(L),E):\exists\;\psi_0\in \homo(L^{\times},E),\text{such that\;}\psi-\psi_0\flat\mathrm{det}_{\bZ_{S^u_0}}\in \homo_{\sm}(\bZ_{S^u_0}(L),E)\}$.\;Note that $\dim_E\homo_{g'}(\bZ_{S^u_0}(L),E)=s+d_L$.\;It is clear that 
%\[\homo_{\sm}(\bZ_{S^u_0}(L),E)\subseteq \homo_{g'}(\bZ_{S^u_0}(L),E)\subseteq 	\homo_u(\bZ_{S^u_0}(L),E).\]



For $u\in \sW_{s}$ and $?\in\{\flat,0\}$,\;let  $\ext^{1,?}_{g,u}(\Dpik,\Dpik)\subseteq\ext^{1,?}_{u}(\Dpik,\Dpik)$  be the preimage of $\homo_{\sm}(\bZ_{S^u_0}(L),E)$ of via $\kappa^{?}_u$.\;Let $\ext^1_g(\Dpik,\Dpik)$ be the subspace of de Rham deformations of $\Dpik$.\;Note that $\dim_E{\ext}^{1}_g(\Dpik,\Dpik)=1+d_L\frac{n(n-1)}{2}$.\;
\begin{pro}\label{studyforkerkappau}
\begin{itemize}
	\item[(1)] $\ext^1_g(\Dpik,\Dpik)\subseteq \ext^{1}_{u}(\Dpik,\Dpik)$ and $\ext^{1,0}_{g,u}(\Dpik,\Dpik)=\ext^{1}_{g}(\Dpik,\Dpik)$.\;
	\item[(2)] For $u\in \sW_{s}$,\;we have $\ext^{1,\flat}_{g,u}(\Dpik,\Dpik)=\ext^{1,\flat}_{u}(\Dpik,\Dpik)\cap \ext^1_g(\Dpik,\Dpik)$.\;
\end{itemize}
	\end{pro}
\begin{proof}By \cite[Appendix,\;Proposition A.3]{Dingsocle},\;for any $i<j$,\;we deduce that $\dim_E\hH^1_g(\gal_L,W(E_{u^{-1}(i)}^{\vee}\otimes_{\cR_{E,L}}E_{u^{-1}(j)}))=0$,\;where $W(-)$ means the associated $E$-$B$ pair.\;The d\'{e}vissage argument shows the inclusion $\ext^1_g(\Dpik,\Dpik)\subseteq \ext^{1}_u(\Dpik,\Dpik)$.\;Since $\cR_{E[\epsilon]/\epsilon^2}(1+\psi_i\epsilon)$ is de Rham iff $\psi\in \homo_{\sm}(L^{\times},E)$,\;we deduce that $\ext^{1,\flat}_{u}(\Dpik,\Dpik)\cap \ext^1_g(\Dpik,\Dpik)\subseteq \ext^{1,\flat}_{g,u}(\Dpik,\Dpik)$.\;On the other hand,\;for $(\psi_i)\in \homo_{\sm}(\bZ_{S^u_0}(L),E)$,\;we show that \[\widetilde{D}=[	\widetilde{E}_{u,1}-\widetilde{E}_{u,2}-\cdots-	\widetilde{E}_{u,s}],\;\]
with $\widetilde{E}_{u,i}\cong \widetilde{E}_{u,i}\otimes_{\cR_{E,L}}\cR_{E[\epsilon]/\epsilon^2}(1+\psi_i\epsilon)$ and $\psi\in \homo_{\sm}(\bQ^{\times}_p,E)$ is de Rham.\;It suffices to show the de Rhamness of $\widetilde{D}_1^j:=[\widetilde{E}_{u,1}-\widetilde{E}_{u,2}-\cdots-\widetilde{E}_{u,j}]$ (with $j$ increasing) step by step.\;This is true for $j=1$ obviously.\;Suppose that the $\widetilde{D}_1^j$ is de Rham and we show the de Rhamness of $\widetilde{D}_1^{j+1}$.\;It suffices to show that each element in
\[\hH^1\left(\gal_L,W(\widetilde{D}')\right),\;\widetilde{D}':=\widetilde{D}_1^j\otimes_{\cR_{E,L}}\widetilde{E}_{u,j+1}^{\vee}\otimes_{\cR_{E,L}}\cR_{E[\epsilon]/\epsilon^2,L}(1-\psi_{j+1}\epsilon)\]
is de Rham.\;However,\;since all the Hodge-Tate weights of $\widetilde{D}'$ are positive,\;we deduce from \cite[Proposition A.3]{Dingsocle}  that $\hH_g^1(\gal_L,W(\widetilde{D}'))=\hH^1(\gal_L,W(\widetilde{D}'))$ (note that $\widetilde{H}^2_{\Sigma_L}(\gal_L,W(\widetilde{D}'))={H}^2(\gal_L,W(\widetilde{D}'))$).\;We thus obtain $\ext^{1,\flat}_{g,u}(\Dpik,\Dpik)\subseteq\ext^{1,\flat}_{u}(\Dpik,\Dpik)\cap \ext^1_g(\Dpik,\Dpik)$.\;Since $\ext^1_g(E_{u,i},E_{u,i})\subseteq \ext^{1,0}(E_{u,i},E_{u,i})$,\;we get that $\ext^{1,0}_{g,u}(\Dpik,\Dpik)=\ext^{1}_{g}(\Dpik,\Dpik)$. 
\end{proof}

The following lemma is an analogue of \cite[Lemma 2.11]{ParaDing2024}.\;For any $u_1,u_2\in \sW_s$,\;under the isomophism $\ext^{1,0}_{g,u_1}(\Dpik,\Dpik)\cong \ext^{1}_{g}(\Dpik,\Dpik)\cong \ext^{1,0}_{g,u_2}(\Dpik,\Dpik)$,\;we can identity $\ext^{1,\flat}_{g,u_1}(\Dpik,\Dpik)$ with $ \ext^{1,\flat}_{g,u_2}(\Dpik,\Dpik)$.\;Put
$\ext^{1,\flat}_{g,\{u_1,u_2\}}(\Dpik,\Dpik):=\ext^{1,\flat}_{g,u_1}(\Dpik,\Dpik)\cap \ext^{1,\flat}_{g,u_2}(\Dpik,\Dpik)\cong \ext^{1,\flat}_{g,u_1}(\Dpik,\Dpik)\cong \ext^{1,\flat}_{g,u_2}(\Dpik,\Dpik)$.\;


\begin{lem}\label{smhomoindepent}For any $u_1,u_2\in \sW_s$, the following diagram commutes:
\begin{equation}
	\xymatrix{ \ext^{1,\flat}_{g,\{u_1,u_2\}}(\Dpik,\Dpik)\ar@{=}[d] \ar[r]_{\kappa_{u_1}} & \homo_{\sm}(\bZ^{u_1}_{S_0}(L),E) \ar[d]^{\sim}_{u_2u_1^{-1}}\\
		\ext^{1,\flat}_{g,\{u_1,u_2\}}(\Dpik,\Dpik) \ar[r]_{\kappa_{u_2}} & \homo_{\sm}(\bZ^{u_2}_{S_0}(L),E) }.
\end{equation}
\end{lem}
\begin{proof}
It suffices to prove the statement when $u_2u_1^{-1}=s_k$ is a simple reflection.\;Let $\widetilde{D}\in  \ext^{1,\flat}_{g,\{u_1,u_2\}}(\Dpik,\Dpik)$, and suppose that
\[\widetilde{D}=\widetilde{E}_{u_1,1}-\widetilde{E}_{u_1,2}-\cdots-\widetilde{E}_{u_1,s},\]
with $\widetilde{E}_{u_1,i}\cong E_{u_1,i}\otimes_{\cR_{E,L}}\cR_{E[\epsilon]/\epsilon^2}(1+\psi_i\epsilon)$.\;By assumption, $u^{-1}_1(j)=u^{-1}_2(j)$ for $j\neq k,k+1$.\;Thus, for $j<k$ or $j>k+1$, we have $\fil^j_{\cF_{u_1}}\widetilde{D}=\fil^j_{\cF_{u_2}}\widetilde{D}$.\;Since $\homo(\widetilde{E}_{u_1,1},\widetilde{D})\cong \homo(\widetilde{E}_{u_1,1},\widetilde{E}_{u_1,1})\cong E[\epsilon]/\epsilon^2$, d\'{e}vissage for $\fil^j_{\cF_{u_2}}\widetilde{D}$ gives $\homo(\widetilde{E}_{u_1,1},\widetilde{E}_{u_2,1})\cong E[\epsilon]/\epsilon^2$, hence $\widetilde{E}_{u_1,1}=\widetilde{E}_{u_2,1}$.\;Now consider $\widetilde{D}/\widetilde{E}_{u_1,1}$ and repeat the same argument.\;We obtain $\widetilde{E}_{u_1,j}=\widetilde{E}_{u_2,j}$ for $j<k$.\;For $j=k$, we see that $\homo(\widetilde{E}_{u_1,k},\widetilde{E}_{u_2,k+1})\cong E[\epsilon]/\epsilon^2$, so $\widetilde{E}_{u_1,k}=\widetilde{E}_{u_2,k+1}$.\;Exchanging $\cF_{u_1}$ and $\cF_{u_2}$ gives $\widetilde{E}_{u_1,k+1}=\widetilde{E}_{u_2,k}$.\;For $j>k+1$, we consider $\widetilde{D}/\fil^{k+1}_{\cF_{u_1}}$ and see that $\widetilde{E}_{u_1,j}=\widetilde{E}_{u_2,j}$.\;This completes the proof.\;
\end{proof}

For any $u\in\sW_s$, let $\ext^{1,\flat}_{0,u}(\Dpik,\Dpik):=\ker\kappa^{\flat}_{u}\subseteq\ext^{1,\flat}_{g,u}(\Dpik,\Dpik)$.\;The above lemma shows that $\ker\kappa^{\flat}_{u}$ is independent of the choice of $u$; we denote it by $\ext^{1,\flat}_0(\Dpik,\Dpik)$.\;Note that 
\[ \dim_E{\ext}^{1}_0(\Dpik,\Dpik)=1+d_L\dim\bN_{S^u_0}-s=1+d_L\dim\bN_{S_0}-s.\]
For any subspace $V\subseteq \ext^1(\Dpik,\Dpik)$, let $\overline{V}:=V/(V\cap \ext^{1,\flat}_0(\Dpik,\Dpik))$.\;Therefore, $\kappa_{u}^{\flat}$ induces isomorphisms
\[\overline{\ext}^{1,\flat}_{u}(\Dpik,\Dpik)\xrightarrow{\sim }\homo(\bZ_{S^u_0}(L),E),\;\overline{\ext}^{1,\flat}_{g,u}(\Dpik,\Dpik)\xrightarrow{\sim }\homo_{\sm}(\bZ_{S^u_0}(L),E)\]
and (comparing with the short exact sequence in (\ref{commutativefortriandwhole}))
\[\dim_E\overline{\ext}^{1}(\Dpik,\Dpik)=d_L\dim_E\bP_{S_0}+s.\]







\subsection{Reinterpretation and supplements for deformations}\label{reinterfor33}



We recall briefly Fontaine's theory of almost de Rham representations.\;Let $B_{\pdr}^+:=B_{\dr}^+[\log t]$ and $B_{\pdr}:=B_{\pdr}^+[1/t]$.\;The $\gal_L$-action on $B_{\dr}$ extends uniquely to an action of $\gal_L$ on $B_{\pdr}$ with $g(\log t)=\log t+\log(\ccyc(g))$.\;Let $v_{\pdr}$ denote the unique $B_{\dr}$-linear derivation of $B_{\pdr}$ such that $v_{\pdr}(\log t)=-1$.\;Note that $v_{\pdr}$ and $\gal_L$ commute, and both preserve $B_{\pdr}^+$.\;


Let $\mathrm{Rep}_{B_{\dr}}(\gal_L)$ (resp.,\;$\mathrm{Rep}_{B^+_{\dr}}(\gal_L)$) be the category of finite free $B_{\dr}$-representations (resp., $B^+_{\dr}$-representations) of $\gal_L$.\;If $W\in \mathrm{Rep}_{B_{\dr}}(\gal_L)$, let $D_{\pdr}(W):=(B_{\pdr}\otimes_{B_{\dr}}W)^{\gal_L}$; this is a finite-dimensional $L$-vector space of dimension at most $\dim_{B_{\dr}}W$.\;The $B_{\dr}$-representation $W$ is called \textit{almost de Rham} if $\dim_LD_{\pdr}(W)=\dim_{B_{\dr}}W$.\;The $B^+_{\dr}$-representation $W^+$ is called \textit{almost de Rham} if $W^+[1/t]$ is almost de Rham.\;Let $\mathrm{Rep}_{\pdr}(\gal_L)$ be the category of almost de Rham $B_{\dr}$-representations of $\gal_L$.\;

Keep the notation of Section \ref{Omegafil}.\;In this section, we use the language of $E$-$B$-pairs (see \cite{nakamura2009classification}).\;Recall that the $(\varphi,\Gamma)$-module $\cM_\Dpik:=\Dpik[1/t]$ over $\cR_{E,L}[1/t]$ admits an $\Omega_{S_0^u}$-filtration $\cF_u$ with graded pieces $E_{u^{-1}(i)}[1/t]$ for $1\leq i\leq s$ and $\Dpik=[E_{1}-E_{2}-\cdots-E_{s}]$.\;Let $\bW_{\Dpik}=W_{\dr}(\cM_\Dpik)$ (resp.,\;$\bW^+_{\Dpik}:=W_{\dr}^+(\Dpik)$) be the $\bB_\dr\otimes_{\bQ_p}E$-representation (resp.,\;$\bB_\dr^+\otimes_{\bQ_p}E$-representation) of $\gal_L$ associated with $\cM_\Dpik$.\;Note that $D:=D_{\pdr}(\bW_\Dpik)\cong (L\otimes_{\bQ_p}E)^{\oplus n}$ for simplicity.\;For $\sigma\in \Sigma_L$, put $D_{\sigma}:=D_{\pdr,\sigma}(\bW_\Dpik):=D_{\pdr}(\bW_\Dpik)\otimes_{L\otimes_{\bQ_p}E}(L\otimes_{L,\sigma}E)$.\;Then $D=\prod_{\sigma\in \Sigma_L}D_{\sigma}$.\;For $1\leq i\leq s$, put $D^i:=D_{\pdr}(\bW_{E_i})$ and $D^i_{\sigma}:=D_{\pdr,\sigma}(\bW_{E_i})$.\;Then $D^i=\prod_{\sigma\in \Sigma_L}D^i_{\sigma}$.\;



For $u\in \sW_s$,\;the filtration $\cF_u$ on $\cM_{\Dpik}$ induces a $\bP_{S^u_0}$-filtration:
\[\bF_{u}=\fil_{\bullet}^{\bF_{u}}\bW_\Dpik: \ 0 =\fil_0^{\bF_{u}}\bW_\Dpik \subsetneq \fil_1^{\bF_{u}}\bW_\Dpik \subsetneq \cdots \subsetneq \fil_{s}^{\bF_{u}}\bW_\Dpik=\bW_\Dpik\]
on $\bW_\Dpik$ by $(\bB_\dr\otimes_{\bQ_p}E)$-subrepresentations of $\bW_{\Dpik}$.\;For $1\leq i\leq s$, $\bF_{u}$ has graded pieces $\gr_i^{\bF_u}\bW_{\Dpik}\cong \bW_{E_{u^{-1}(i)}}\cong (\bB_\dr\otimes_{\bQ_p}E)^{\oplus r_{u^{-1}(i)}}$.\;We thus get a $\bP_{S^u_0}$-filtration on $D$ and on each $D_{\sigma}$:
\begin{equation}
	\begin{aligned}
		\bF_{u}=\fil_{\bullet}^{\bF_{u}}(D)&: \ 0 =\fil_0^{\bF_{u}}(D) \subsetneq \fil_1^{\bF_{u}}(D) \subsetneq \cdots \subsetneq \fil_{s}^{\bF_{u}}(D)=D,\\
		\bF_{u,\sigma}=\fil_{\bullet}^{\bF_{u,\sigma}}(D_{\sigma})&: \ 0 =\fil_0^{\bF_{u,\sigma}}(D_{\sigma}) \subsetneq \fil_1^{\bF_{u,\sigma}}(D_{\sigma}) \subsetneq \cdots \subsetneq \fil_{s}^{\bF_{u,\sigma}}(D_{\sigma})=D_{\sigma}.
	\end{aligned}
\end{equation}
Note that $\gr_i^{\bF_u}D_{\pdr}(\bW_\Dpik)\cong D^i\cong (L\otimes_{\bQ_p}E)^{\oplus r_{u^{-1}(i)}}$.\;We have $\bF_{u}=\prod_{\sigma\in \Sigma_L}\bF_{u,\sigma}$.\;On the other hand, the $B_{\dr}^+$-lattice $\bW^+_\Dpik$ induces complete flags of $D_{\sigma}$ and $D$, respectively (see \cite[Section 6.3]{Ding2021}):
\begin{equation}
	\begin{aligned}
		\fil^{\bullet}_{H}(D_{\sigma})&: \ 0 \subsetneq \fil^{-\hpi_{n,\sigma}}_{H} (D_{\sigma}) \subsetneq \fil^{-\hpi_{n-1,\sigma}}_{H}(D_{\sigma}) \subsetneq \cdots \subsetneq \fil^{{-\hpi_{1,\sigma}}}_{H}(D_{\sigma})=D_{\sigma},\\
	%	\fil^{\bullet}_{H}(D)&: \ 0 \subsetneq \fil^{-\hpi_{n}}_{H}(D) \subsetneq \fil^{-\hpi_{n-1}}_{H}(D)  \subsetneq \cdots \subsetneq \fil^{{-\hpi_{1}}}_{H}(D)=D,\;
	\end{aligned}
\end{equation}
where $\fil^{-\bh_{n+1-i,\sigma}}_{H}(D_{\sigma}):=\big(t^{-\bh_{n+1-i,\sigma}}D^+_{\sigma}\big)^{\gal_L}$ and $D^+_{\sigma}:=D_{\pdr}(\bW^+_\Dpik)\otimes_{L\otimes_{\bQ_p}E}(L\otimes_{L,\sigma}E)$.\;
%\begin{equation}
%	\begin{aligned}
%		&\fil^{-\bh_{n+1-i,\sigma}}_{H}(D):=\bigoplus_{\sigma\in\Sigma_L}\fil^{-\bh_{n+1-i,\sigma}}_{H}(D_{\sigma}),
%		&
%	\end{aligned}
%\end{equation}
%with Note that $\fil_H^{\bullet}(D)=\prod_{\sigma\in \Sigma_L}\fil_{H}^{\bullet}(D_{\sigma})$.\;

In the sequel,\;for $1\leq i\leq s$,\;let $\{e_{i,j,\sigma}\}_{1\leq j\leq r_i}$ be a fixed $E$-basis of $D^i_{\sigma}$.\;Under the $\bP_{S_0}$-filtration $\bF_1$ on $D_{\sigma}$,\;we use the following ordering of the basis $\{e_{1,j,\sigma}\}_{1\leq j\leq r_1},\cdots,\{e_{s,j,\sigma}\}_{1\leq j\leq r_s}$ of $D_{\sigma}$ (i.e.,\;the lexicographical order),\;and relabel  them with $e_{1,\sigma},\cdots,e_{n,\sigma}$.\;



Let $\homo_{\fil}(D_{\sigma},D_{\sigma})$ be the subspace of $E$-linear endomorphisms of $D_{\sigma}$ that preserve the filtration $\fil_{H}^{\bullet}(D_{\sigma})$.\;By the same argument as in the discussion before \cite[(112)]{BDcritical25}, we obtain canonical $E$-linear morphisms:
\begin{equation}\label{maptofilDsigma}
	\begin{aligned}
	%	\nu_{\Sigma_L,\sigma}:&\;\ext^{1}(\Dpik,\Dpik)\rightarrow \homo_{\fil}(D_{\sigma},D_{\sigma}),\\
		\nu:	&\;\ext^{1}(\Dpik,\Dpik)\rightarrow \homo_{\fil}(D,D):=\bigoplus_{\sigma\in\Sigma_L}\homo_{\fil}(D_{\sigma},D_{\sigma}),\\
		\nu_{\sigma}:	&\;\ext^{1}_{\sigma}(\Dpik,\Dpik)\rightarrow \homo_{\fil}(D_{\sigma},D_{\sigma}).\;
	\end{aligned}
\end{equation}
where $\ext^{1}_{\sigma}(\Dpik,\Dpik)\subseteq \ext^{1}(\Dpik,\Dpik)$ is the subspace consisting of deformations that are $\Sigma_L\backslash \{\sigma\}$-de Rham.\;
By \cite[Lemma 2.4.1]{BDcritical25}, $\nu$ and $\nu_{\sigma}$ has kernel $\ext^{1}_g(\Dpik,\Dpik)$.\;


For $u\in \sW_s$, let $\homo_{\bF_u}(D_{\sigma},D_{\sigma})$ be the subspace of endomorphisms of $D_{\sigma}$ that preserve the filtration $\bF_{u,\sigma}$.\;For $?\in \{\flat,0\}$,\;put
\[\homo^{?}_{\bF_u}(D_{\sigma},D_{\sigma}):=\left\{f\in \homo_{\bF_u}(D_{\sigma},D_{\sigma}):\;f|_{\gr_{i}^{\bF_u}(D_{\sigma})} \text{\;scalar,\;}\;\forall \;1\leq i\leq s\right\}\subseteq \homo_{\bF_u}(D_{\sigma},D_{\sigma}).\]
For $?\in \{\emptyset,\flat,0\}$,\;put
\begin{equation}
	\begin{aligned}
		\homo^{?}_{\fil,\bF_u}(D_{\sigma},D_{\sigma}):=&\;\homo_{\fil}(D_{\sigma},D_{\sigma})\cap \homo^{?}_{\bF_u}(D_{\sigma},D_{\sigma}),\\
		\homo^{?}_{\fil,\bF_u}(D,D):=&\;\bigoplus_{\sigma\in\Sigma_L}\homo^{?}_{\fil,\bF_u}(D_{\sigma},D_{\sigma}).\;
	\end{aligned}
\end{equation}
For $?\in\{\flat,0\}$, put $\ext^{1,?}_{\sigma,u}(\Dpik,\Dpik):=\ext^{1,?}_{u}(\Dpik,\Dpik)\cap \ext^{1}_{\sigma}(\Dpik,\Dpik)$.\;Therefore,\;for $\star\in\{\emptyset,\sigma\}$, (\ref{maptofilDsigma}) induces
\begin{equation}\label{sendtofilD}
	\begin{aligned}
		\nu^?_{\star,u}:&\;\ext^{1,?}_{\star,u}(\Dpik,\Dpik)/\ext^{1,?}_{g,u}(\Dpik,\Dpik)\rightarrow\homo^{?}_{\fil,\bF_u}(D_{\star},D_{\star}),\\
		\nu_{\star,u}:&\;\ext^{1}_{\star,u}(\Dpik,\Dpik)/\ext^{1,?}_{g,u}(\Dpik,\Dpik)\rightarrow\homo_{\fil,\bF_u}(D_{\star},D_{\star}).\;
	\end{aligned}
\end{equation}



\begin{lem}
\begin{itemize}
	\item[(1)]
	For $1\leq i\leq s$,\;there is a  commutative diagram of short exact sequences of $E$-vector spaces: :
	\begin{equation}
		\xymatrix{
		0   \ar[r]  & \ext^{1,\flat}_g(E_i,E_i) \ar@{^(->}[d]\ar[r]  & \ext^{1,\flat}(E_i,E_i) \ar@{^(->}[d]\ar[r]  &  \bigoplus\limits_{\sigma\in\Sigma_L}E \ar@{=}[d] \ar[r] & 0  \\
		0    \ar[r] & \ext^{1,0}_g(E_i,E_i) \ar@{=}[d] \ar[r] & \ext^{1,0}(E_i,E_i)\ar[r] \ar@{^(->}[d]  & \bigoplus\limits_{\sigma\in\Sigma_L}E \ar[r]\ar@{^(->}[d] & 0 \\
		0    \ar[r] & \ext^1_g(E_i,E_i)  \ar[r] & \ext^{1}(E_i,E_i)\ar[r]   & \bigoplus\limits_{\sigma\in\Sigma_L}\homo_{\fil}(D^i_{\sigma},D^i_{\sigma}) \ar[r] & 0 }  
	\end{equation}	
		where $\ext^{1,\flat}_g(E_i,E_i)\cong \homo_{\sm}(L^{\times},E)$, and we have a natural inclusion $E\hookrightarrow \homo_{\fil}(D^i_{\sigma},D^i_{\sigma})$.\;There is also a $\sigma$-version obtained by dropping the symbol $\oplus_{\sigma\in\Sigma_L}$ and restricting to $\ext^{1,?}_{\sigma}(E_i,E_i)$.\;
	\item[(2)] For $1\leq i\leq s$,\;$\ext^{1}_g(E_i[1/t],E_i[1/t])\cong \homo_{\sm}(L^{\times},E)$.\;We have a natural map (by inverting $t$):
	\begin{equation}\label{mapblockinverset}
		\prod_{i=1}^s\ext^{1}_g(E_i,E_i)\rightarrow \prod_{i=1}^s\ext^{1}_g(E_i[1/t],E_i[1/t]).
	\end{equation}
		This map induces an isomorphism $\prod_{i=1}^s\ext^{1,\flat}_g(E_i,E_i)\xrightarrow{\sim} \prod_{i=1}^s\ext^{1}_g(E_i[1/t],E_i[1/t])$.\;
\end{itemize}
\end{lem}
In particular, (\ref{mapblockinverset}) is not injective if $S_0\neq\emptyset$ (resp., is an isomorphism if $S_0=\emptyset$).\;In the sequel,\;put 
\begin{equation}
	\begin{aligned}
		&\ext^{1,\flat}_{\gr}({\Dpik},{\Dpik}):=\prod_{i=1}^s\ext^{1,\flat}_g(E_i,E_i)\xrightarrow{\sim}\prod_{i=1}^s\ext^{1}_g(E_i[1/t],E_i[1/t]),\;\\
		&\ext^{1}_{\gr}({\Dpik},{\Dpik}):=\ext^{1,0}_{\gr}({\Dpik},{\Dpik}):=\prod_{i=1}^s\ext^{1}_g(E_i,E_i).
	\end{aligned}
\end{equation}
For any $1\leq i\leq s$,\;we have a decomposition $\ext^{1}(E_{i},E_{i})=\ext^{1,\flat}(E_{i},E_{i})\oplus \hH^1_{(\varphi,\Gamma)}(\EndO^0(E_{i}))$,\;where $\EndO^0(E_{i})=\EndO(E_{i})/\cR_{E,L}$.\;Therefore,\;$\ext^{1,\flat}_{\gr}({\Dpik},{\Dpik})$ is a direct summand of $\ext^{1}_{\gr}({\Dpik},{\Dpik})$ and thus we have a natural projection
\begin{equation}\label{projtocenterdefor}
	\ext^{1}_{\gr}({\Dpik},{\Dpik})\twoheadrightarrow \ext^{1,\flat}_{\gr}({\Dpik},{\Dpik}).\;
\end{equation}



For $?\in\{\emptyset,\flat,0\}$ and $\star\in\{\emptyset,\sigma\}$ for some $\sigma\in \Sigma_L$,\;consider the natural morphism
\begin{equation}\label{natmorphismforGDpikfern}
	g_{\Dpik,\star}^{?}:\bigoplus_{u\in \sW_{s}}{\ext}^{1,?}_{\star,u}(\Dpik,\Dpik)\rightarrow {\ext}_{\star}^1(\Dpik,\Dpik).\;
\end{equation}
For $?\in\{\flat,0\}$ and $\star\in\{\emptyset,\sigma\}$,\;let
${\ext}_{\star}^{1,?}(\Dpik,\Dpik)$ be the image of $g_{\Dpik,\star}^{?}$ and put $\nu_{\star}^{?}=\nu|_{{\ext}^{1,?}_{\star}(\Dpik,\Dpik)}$ and $\homo^{?}_{\fil}(D_{\star},D_{\star}):=\mathrm{Im}(\nu_{\star}^{?})$.\;


\begin{lem}\label{lemforexactforextgropus} 
\begin{itemize}
	\item[(1)]	There is an isomorphism
	\begin{equation}\label{expianextphi}
		{\ext}^{1,0}_g(\Dpik,\Dpik)/{\ext}^{1,\flat}_0(\Dpik,\Dpik)\xrightarrow{\sim}\ext^{1}_{\gr}({\Dpik},{\Dpik}).
	\end{equation}
	\item[(2)] For $u\in \sW_s$,\;$?\in\{\emptyset,\flat,0\}$ and $\star\in\{\emptyset,\sigma\}$,\;there exists a commutative diagrm of short exact sequences of $E$-vector spaces: 
	\begin{equation}\label{commutativefortriandwhole}
		\xymatrix{
			0   \ar[r]  & \ext^{1,?}_{\gr}({\Dpik},{\Dpik}) \ar@{=}[d]\ar[r]  & {\ext}^{1,?}_{\star,u}(\Dpik,\Dpik)/{\ext}^{1,\flat}_{0}(\Dpik,\Dpik) \ar@{^(->}[d]\ar[r]  &  \homo^{?}_{\fil,\bF_u}(D_{\star},D_{\star}) \ar@{^(->}[d] \ar[r] & 0  \\
			0    \ar[r] & \ext^{1,?}_{\gr}({\Dpik},{\Dpik}) \ar[r] & {\ext}_{\star}^{1,?}(\Dpik,\Dpik)/{\ext}^{1,\flat}_{0}(\Dpik,\Dpik)\ar[r]   &  \homo^{?}_{\fil}(D_{\star},D_{\star}) \ar[r] & 0  }.
	\end{equation}
%		\[0\rightarrow \ext^{1,?}_{\gr}({\Dpik},{\Dpik})\rightarrow {\ext}^{1,?}_{\star,u}(\Dpik,\Dpik)/{\ext}^{1,\flat}_{0}(\Dpik,\Dpik)\rightarrow  \homo^{?}_{\fil,\bF_u}(D_{\star},D_{\star})\rightarrow 0.\]
%	\item[(3)] There exists a short exact sequence of $E$-vector spaces :
%		\begin{equation}\label{extforextfilhomo}
%		0\rightarrow \ext^{1,?}_{\gr}({\Dpik},{\Dpik})\rightarrow {\ext}_{\star}^{1,?}(\Dpik,\Dpik)/{\ext}^{1,\flat}_0(\Dpik,\Dpik)\rightarrow  \homo^{?}_{\fil}(D_{\star},D_{\star})\rightarrow 0.
%	\end{equation}
\end{itemize}
\end{lem}
\begin{proof}Recall the dimensions of ${\ext}^{1,\flat}_0(\Dpik,\Dpik)$ and ${\ext}^{1,0}_g(\Dpik,\Dpik)={\ext}^{1}_g(\Dpik,\Dpik)$.\;These subspaces lead to a natural morphism ${\ext}^{1,0}_g(\Dpik,\Dpik)\rightarrow \ext^{1}_{\gr}({\Dpik},{\Dpik})$.\;Thus,\;we get $\dim_E{\ext}^{1,0}_g(\Dpik,\Dpik)/{\ext}^{1,\flat}_0(\Dpik,\Dpik)=s+d_L\dim_E\bL_{S_0}\cap\bN_{\emptyset}$,\;which is equal to the $E$-dimension of the right-hand side of (\ref{expianextphi});\;hence we get the isomorphism in $(1)$.\;
\end{proof}

Similar to \cite[(121)]{BDcritical25} and the commutative diagram below  \cite[(123)]{BDcritical25},\;we have maps
\[\ext^{1}(\Dpik,\Dpik)\rightarrow \ext^{1}(\cM_{\Dpik},\cM_{\Dpik})\rightarrow \bigoplus_{\sigma\in\Sigma_L}\homo_{E}(D_{\sigma},D_{\sigma}),\]
and thus obtain a commutative diagram of short exact sequences: 
\begin{equation}\label{diagDtoM}
	\begin{aligned}
		\xymatrix{
		0   \ar[r]  &  \ext^{1}_{\gr}({\Dpik},{\Dpik}) \ar[d]^{(\ref{projtocenterdefor})} \ar[r] & {\ext}^{1}(\Dpik,\Dpik)/{\ext}^{1,\flat}_{0}(\Dpik,\Dpik) \ar[d]\ar[r]  &  \bigoplus\limits_{\sigma\in\Sigma_L}\homo_{\fil}(D_{\sigma},D_{\sigma})\ar@{^(->}[d] \ar[r] & 0  \\
		0    \ar[r] &  \ext^{1,\flat}_{\gr}({\Dpik},{\Dpik}) \ar[r] & \ext^{1}(\cM_{\Dpik},\cM_{\Dpik})/{\ext}^{1,\flat}_{0}(\cM_{\Dpik},\cM_{\Dpik})\ar[r]   & \bigoplus\limits_{\sigma\in\Sigma_L}\homo_{E}(D_{\sigma},D_{\sigma})   \ar[r] & 0 }
	\end{aligned}
\end{equation}
where ${\ext}^{1,\flat}_{0}(\cM_{\Dpik},\cM_{\Dpik})$ is the image of ${\ext}^{1,\flat}_{0}(\Dpik,\Dpik)$ via the natural map $\ext^{1}(\Dpik,\Dpik)\rightarrow \ext^{1}(\cM_{\Dpik},\cM_{\Dpik})$.\;


Recall that $L'$ is a finite Galois extension of $L$ such that $\Dpik|_{L'}$ is a crystalline $(\varphi,\Gamma)$-module over $\cR_{E,L'}$ of rank $n$.\;Let $\cM_{\Dpik}|_{L'}:=\cM_{\Dpik}\otimes_{\cR_{E,L}[1/t]}\cR_{E,L'}[1/t]$ and $E_i[1/t]|_{L'}:=E_i[1/t]\otimes_{\cR_{E,L}[1/t]}\cR_{E,L'}[1/t]$ for $1\leq i\le s$.\;Note that $E_i[1/t]|_{L'}$ is crystalline for all $1\leq i\leq s$.\;For all $1\leq i\neq j\leq s$,\;we deduce from \cite[Proposition 2.7]{nakamura2009classification} that $\ext^1_e(E_i[1/t]|_{L'},E_j[1/t]|_{L'})=0$,\;and thus obtain
\[\cM_{\Dpik}|_{L'}\cong \oplus_{i=1}^sE_i[1/t]|_{L'}.\]
Similar to the argument in \cite[(122)-(123)]{BDcritical25},\;we have a commutative diagram of short exact sequences: 
\begin{equation}\label{exacttoL'}
	\begin{aligned}
		\xymatrix{
			0   \ar[r]  &  \ext^{1,\flat}_{\gr}({\Dpik},{\Dpik}) \ar[d]^{\sim}\ar[r]  & \ext^{1}(\cM_{\Dpik},\cM_{\Dpik})/{\ext}^{1,\flat}_{0}(\cM_{\Dpik},\cM_{\Dpik}) \ar[d]\ar[r]  &  \bigoplus\limits_{\sigma\in\Sigma_L}\homo_{E}(D_{\sigma},D_{\sigma}) \ar[d]^{\sim}    \\
			0    \ar[r] &  \ext^{1,\flat}_{\gr}({\Dpik}|_{L'},{\Dpik}|_{L'}) \ar[r] & \ext^{1}(\cM_{\Dpik}|_{L'},\cM_{\Dpik}|_{L'}) \ar[r]   & \bigoplus\limits_{\sigma\in\Sigma_L}\homo_{E}(D_{L',\sigma},D_{L',\sigma})    }
	\end{aligned}
\end{equation}
where $D_{\sigma}\cong D_{L',\sigma}^{\gal(L'/L)}$ and $\ext^{1,\flat}_{\gr}({\Dpik}|_{L'},{\Dpik}|_{L'})\cong \prod_{i=1}^s\ext^{1}_g(E_i[1/t]|_{L'},E_i[1/t]|_{L'})$.\;The first vertical map in (\ref{exacttoL'}) is an isomorphism since $\ext^{1}_g(E_i[1/t],E_i[1/t])\cong \ext^{1}_g(E_i[1/t]|_{L'},E_i[1/t]|_{L'})\cong \homo_{\mathrm{sm}}(L^{\times},E)$. By comparing dimensions,\;the last morphism in the bottom exact sequence of (\ref{exacttoL'}) is thus surjective,\;since $\dim_E\ext^{1}(\cM_{\Dpik}|_{L'},\cM_{\Dpik}|_{L'})=s+d_Ln^2$.\;Combining (\ref{diagDtoM}) with (\ref{exacttoL'}),\;we actually obtain the following commutative diagram (an analogue, in the potentially crystalline case, of the commutative diagram before \cite[Proposition 2.4.4]{BDcritical25}):
\begin{equation}\label{diagDtoML'}
	\begin{aligned}
		&\xymatrix{
			0   \ar[r]  &  \ext^{1}_{\gr}({\Dpik},{\Dpik}) \ar[d]^{\sim}\ar[r]  & {\ext}^{1}(\Dpik,\Dpik)/{\ext}^{1,\flat}_{0}(\Dpik,\Dpik) \ar[d]\ar[r]  &  \bigoplus\limits_{\sigma\in\Sigma_L}\homo_{\fil}(D_{\sigma},D_{\sigma})\ar@{^(->}[d] \ar[r] & 0  \\
			0    \ar[r] &   \ext^{1,\flat}_{\gr}({\Dpik}|_{L'},{\Dpik}|_{L'})\ar[r] & \ext^{1}(\cM_{\Dpik}|_{L'},\cM_{\Dpik}|_{L'})\ar[r]   & \bigoplus\limits_{\sigma\in\Sigma_L}\homo_{E}(D_{L',\sigma},D_{L',\sigma}) \ar[r]  & 0 }
	\end{aligned}
\end{equation}
Similar to \cite[Proposition 2.4.4]{BDcritical25},\;we see that
\begin{pro}\label{splitingforext1gps}
	There is a splitting of the bottom exact sequence in (\ref{exacttoL'}) which depends only on the choice of $\log_p(p)\in E$.\;Therefore,\;for $\ast\in \{\emptyset,u\}$ and $\star\in\{\emptyset,\sigma\}$,\;we have 
	\begin{equation}\label{splittingforextgps}
		\begin{aligned}
%			&f^{\flat}:{\ext}^{1,\flat}_{\ast}(\Dpik,\Dpik)/{\ext}^{1,\flat}_{0,\ast}(\Dpik,\Dpik)\xrightarrow{\sim}\ext^{1,\flat}_{\gr}({\Dpik},{\Dpik})\oplus\mathrm{Im}(\nu^{\flat}_{\ast}),\\
			&f^{\flat}_{\Dpik,\star}:{\ext}^{1,\flat}_{\star,\ast}(\Dpik,\Dpik)/{\ext}^{1,\flat}_{0,\ast}(\Dpik,\Dpik)\xrightarrow{\sim}\ext^{1,\flat}_{\gr}({\Dpik},{\Dpik})\oplus\mathrm{Im}(\nu^{\flat}_{\star,\ast}).\;
		\end{aligned}
	\end{equation}
\end{pro}
\begin{proof}By the same argument as in the proof of \cite[Proposition 2.4.4]{BDcritical25},\;the following map is also surjective:
\[\ext^{1}(\cM_{\Dpik}|_{L'},\cM_{\Dpik}|_{L'})\rightarrow \bigoplus_{i=1}^s\ext^{1}(E_i[1/t]|_{L'},E_i[1/t]|_{L'})\twoheadrightarrow \bigoplus_{i=1}^s\ext^{1}_g(E_i[1/t]|_{L'},E_i[1/t]|_{L'}).\]
	This gives the first assertion.\;The assertion in (\ref{splittingforextgps}) follows by replacing the top exact sequence in (\ref{diagDtoML'}) with the first row of  (\ref{commutativefortriandwhole}) (put $?=\flat$ and $\star\in\{\emptyset,\sigma\}$).\;
\end{proof}
\begin{rmk}The bottom exact sequence in (\ref{diagDtoM})
splits.\;For $1\leq i,j\leq s$,\;the natural injection $E_i[1/t]\hookrightarrow \cM_{\Dpik}$ and the surjection $\cM_{\Dpik}\twoheadrightarrow E_j[1/t]$ induce a natural map 
\[\ext^{1}(\cM_{\Dpik},\cM_{\Dpik})\rightarrow \oplus_{i,j=1}^s\ext^{1}(E_i[1/t],E_j[1/t]).\]	
This gives a splitting surjection $\ext^{1}(\cM_{\Dpik},\cM_{\Dpik})\rightarrow \oplus_{i=1}^s\ext^{1}(E_i[1/t],E_i[1/t])\twoheadrightarrow \oplus_{i=1}^s\ext^{1}_g(E_i[1/t],E_i[1/t])$. In general,\;the top exact sequence in (\ref{diagDtoML'}) may not admit a splitting.\;
\end{rmk}

%For $u\in \sW_s$,\;let $\ext^{1}_u(\cM_{\Dpik},\cM_{\Dpik})$ be the subspace of $\cF_u$-parabolic deformations of $\cM_{\Dpik}$ over $\cR_{E[\epsilon]/\epsilon^2,L}$;\;then we have a similar commutative diagram:
%\begin{equation}\label{diagDtoM2}
%	\xymatrix{
%		0   \ar[r]  & \ext^{1}_{\gr}({\Dpik},{\Dpik}) \ar[d]\ar[r]  & {\ext}^{1}_u(\Dpik,\Dpik)/{\ext}^{1,\flat}_{0,u}(\Dpik,\Dpik) \ar[d] \ar[r]  &  \homo_{\fil}(D,D) \ar@{^(->}[d]   \ar[r] & 0\\
%		0    \ar[r] & \ext^{1,\flat}_{\gr}({\Dpik},{\Dpik}) \ar[r] & \ext^{1}_u(\cM_{\Dpik},\cM_{\Dpik})/{\ext}^{1,\flat}_{0,u}(\cM_{\Dpik},\cM_{\Dpik})\ar[r]   & \bigoplus\limits_{\sigma\in\Sigma_L}\homo_{\bF_u}(D_{\sigma},D_{\sigma})   \ar[r] & 0}  
%\end{equation}
%where ${\ext}^{1,\flat}_{0,u}(\cM_{\Dpik},\cM_{\Dpik})$ is the image of ${\ext}^{1,\flat}_{0,u}(\Dpik,\Dpik)$ via the natural map $\ext^{1}(\Dpik,\Dpik)\rightarrow \ext^{1}(\cM_{\Dpik},\cM_{\Dpik})$. We have a composition of canonical surjective maps
%\begin{equation}
%	\begin{aligned}
%		\ext^{1}_u(\cM_{\Dpik},\cM_{\Dpik})\twoheadrightarrow \ext^{1}_u(\cM_{\Dpik},\cM_{\Dpik})/{\ext}^{1,\flat}_{0,u}(\cM_{\Dpik},\cM_{\Dpik})\twoheadrightarrow \prod_{i=1}^s\ext^1(E_{u^{-1}(i)}[1/t],E_{u^{-1}(i)}[1/t]).
%	\end{aligned}
%\end{equation}
%We have a natural surjection 
%\[\bigoplus_{\sigma\in\Sigma_L}\homo_{\fil,\bF_u}(D_{\sigma},D_{\sigma})\rightarrow \bigoplus_{i=1}^s\bigoplus_{\sigma\in\Sigma_L}\homo_{E}(D^{u^{-1}(i)}_{\sigma},D^{u^{-1}(i)}_{\sigma}).\]
%In conclusion,\;we have a commutative diagram of maps:
%\begin{equation}\label{diagDtoM3}
%	\begin{aligned}
%		\xymatrix{
%			\ext^{1}_u(\cM_{\Dpik},\cM_{\Dpik}) \ar[d]\ar@{->>}[r]  & \prod_{i=1}^s\ext^1(E_{u^{-1}(i)}[1/t],E_{u^{-1}(i)}[1/t])\ar@{->>}[d]   \\
%			\bigoplus_{\sigma\in\Sigma_L}\homo_{\bF_u}(D_{\sigma},D_{\sigma}) \ar@{->>}[r]  & \bigoplus_{i=1}^s\bigoplus_{\sigma\in\Sigma_L}\homo_{E}(D^{u^{-1}(i)}_{\sigma},D^{u^{-1}(i)}_{\sigma}) }.
%	\end{aligned}
%\end{equation}



Let $u^{\sharp}\in \sW_n$ be the unique element sending $e_{i,j,\sigma}$ for $1\leq j\leq r_i$ to $e_{u^{-1}(i),j,\sigma}$ for $1\leq j\leq r_i$,\;which thus induces a permutation on $e_{1,\sigma},\cdots,e_{n,\sigma}$.\;Therefore,\;we identify the space $\homo_{E}(D_{\sigma},D_{\sigma})$ with $\fg_{\sigma}$,\;so that
 $\homo_{\bF_u}(D_{\sigma},D_{\sigma})$ (resp.,\;$\homo^{\flat}_{\bF_u}(D_{\sigma},D_{\sigma})$) is identified with $\mathrm{Ad}_{(u^{\sharp})^{-1}}(\fp_{S^u_0,\sigma})$ (resp.,\;$\mathrm{Ad}_{(u^{\sharp})^{-1}}(\tau_{S^u_0,\sigma})$),\;where $\tau_{S^u_0,\sigma}=\fz_{S^u_0,\sigma}\ltimes\fn_{S^u_0,\sigma}$.\;Then we have 
\begin{equation}
	\begin{aligned}
		&\homo_{\fil}(D_{\sigma},D_{\sigma})=\mathrm{Ad}_{g_{\sigma}}(\fb_{\sigma}),\\
		&\homo_{\fil,\bF_u}(D_{\sigma},D_{\sigma})=\mathrm{Ad}_{(u^{\sharp})^{-1}}(\fp_{S^u_0,\sigma})\cap \mathrm{Ad}_{g_{\sigma}}(\fb_{\sigma}),\\
		&\homo^{\flat}_{\fil,\bF_u}(D_{\sigma},D_{\sigma})=\mathrm{Ad}_{(u^{\sharp})^{-1}}(\tau_{S^u_0,\sigma})\cap \mathrm{Ad}_{g_{\sigma}}(\fb_{\sigma}).\;
	\end{aligned}
\end{equation}
\begin{rmk}\label{explainHomasLiealg}For each $u\in \sW_s$,\;we deduce from the non-critical assumption that $g_{\sigma}\bB_{\sigma}=(u^{\sharp})^{-1}b_{\sigma}(u)w_{0,\sigma}\bB_{\sigma}$ for some $b_{\sigma}(u)\in \bB_{\sigma}$ (of course $b_{\sigma}(u)$ depends on $u$).\;Note that $\overline{\fb}_{\sigma}=\mathrm{Ad}_{w_0}(\fb_{\sigma})$ coincides with the Borel algebra of lower triangular matrices.\;We get that 
	\begin{equation}\label{noncriticalintersection}
		\begin{aligned}
			&\mathrm{Ad}_{(u^{\sharp})^{-1}}(\fp_{S^u_0,\sigma})\cap \mathrm{Ad}_{g_{\sigma}}(\fb_{\sigma})\cong \mathrm{Ad}_{(u^{\sharp})^{-1}b_{\sigma}(u)}(\fp_{S^u_0,\sigma}\cap \mathrm{Ad}_{w_0}(\fb_{\sigma}))=\mathrm{Ad}_{(u^{\sharp})^{-1}b_{\sigma}(u)}(\fl_{S^u_0,\sigma}\cap \overline{\fb}_{\sigma}),\\
			&\mathrm{Ad}_{(u^{\sharp})^{-1}}(\tau_{S^u_0,\sigma})\cap \mathrm{Ad}_{g_{\sigma}}(\fb_{\sigma})=\mathrm{Ad}_{(u^{\sharp})^{-1}b_{\sigma}(u)}(\fz_{S^u_0,\sigma}).
		\end{aligned}
	\end{equation}
\end{rmk}
By the proof of Theorem \ref{Inifinitefernpoten},\;we get that $\sum_{u\in \sW_s}\mathrm{Ad}_{(u^{\sharp})^{-1}}(\fp_{S^u_0,\sigma})\cap \mathrm{Ad}_{g_{\sigma}}(\fb_{\sigma})=\mathrm{Ad}_{g_{\sigma}}(\fb_{\sigma})$.\;Put
\begin{equation}
	\mathrm{Ad}_{g_{\sigma}}(\fb_{\sigma})_{S_0}^{\flat}:=\sum_{u\in \sW_s}\mathrm{Ad}_{(u^{\sharp})^{-1}}(\tau_{S^u_0,\sigma})\cap \mathrm{Ad}_{g_{\sigma}}(\fb_{\sigma})\subseteq \mathrm{Ad}_{g_{\sigma}}(\fb_{\sigma}).\;
\end{equation}
Note that $\dim_E\mathrm{Ad}_{g_{\sigma}}(\fb_{\sigma})_{S_0}^{\flat}$ depends on  the sharp of $\bP_{S_0}$.\;
\begin{thm} (Infinite fern in the non-critical potentially crystalline case)\label{infinitePoten} 
	\begin{itemize} 
		\item[(1)] For $\star\in\{\emptyset,\sigma\}$,\;the natural map  $g_{\Dpik,\star}:\bigoplus_{u\in \sW_{s}}{\ext}^{1}_{\star,u}(\Dpik,\Dpik)\rightarrow {\ext}^1_{\star}(\Dpik,\Dpik)$  is surjective.\;
		\item[(2)] For $?\in\{\flat,0\}$,\;${\ext}^{1,?}(\Dpik,\Dpik)/\ext^{1,?}_{g}({\Dpik},{\Dpik})\cong \bigoplus_{\sigma\in \Sigma_L}\mathrm{Ad}_{g_{\sigma}}(\fb_{\sigma})_{S_0}^{\flat}$.\;
	\end{itemize}
\end{thm}
\begin{proof}By commutative diagram (\ref{commutativefortriandwhole}),\;it suffices to show the following assertion for Lie-algebras :
		\[\sum_{u\in \sW_s}\mathrm{Ad}_{(u^{\sharp})^{-1}}(\fp_{S^u_0,\sigma})\cap \mathrm{Ad}_{g_{\sigma}}(\fb_{\sigma})=\mathrm{Ad}_{g_{\sigma}}(\fb_{\sigma}).\]
This is Lemma \ref{Inifinitefernpoten} .\;
\end{proof}


%where $g_L=(g_{\sigma})_{\sigma\in \Sigma_L}$ in Remark \ref{explainHomasLiealg})





	


\section{Locally analytic representations in the potentially crystalline case}

Throughout this section,\;we fix a non-critical generic potentially crystalline $(\varphi,\Gamma)$-module $\Dpik$ with Hodge--Tate weights $\bh$ in Section \ref{Omegafil}.\;Let $\theta=(0,-1,\cdots,1-n)\in X_{\Delta}^+$ and put $\lambda=\bh-\theta=(\lambda_i=\bh_i+i-1)_{1\leq i\leq n}\in X_{\Delta}^+$.\;






For admissible locally $\bQ_p$-analytic representations $V_1,V_2$ (\cite{Coadmissible}),\;let $\ext^1_{G}(V_1,V_2):=\ext^1_{D(G,E)}(V_2^{\vee},V_1^{\vee})$, where $D(G,E)$ is the locally $\bQ_p$-analytic distribution algebra (so that the continuous strong duals $V_1^{\vee}$ and $V_2^{\vee}$ are coadmissible modules over $D(G,E)$),\;and this latter extension group is defined in the abelian category of abstract $D(G,E)$-modules.\;By \cite[Lemma 2.1.1]{breuil2019ext1},\;$\ext^1_{G}(V_1,V_2)$ is equal to the extension group of admissible locally $\bQ_p$-analytic representations.\;Suppose that $V_1,V_2$ have the same central character $\chi$.\;Let $\ext^1_{G,Z}(V_1,V_2)$ be the subspace of $\ext^1_{G}(V_1,V_2)$ of locally $\bQ_p$-analytic extensions with central character $\chi$.\;For $\sigma\in \Sigma_L$ and any $U(\fg_{\Sigma_L\backslash\{\sigma\}})$-finite representations $V_1,V_2$,\;we denote by $\ext^1_{G,\sigma}(V_1,V_2)\subseteq \ext^1_{G}(V_1,V_2)$ the subspace of extensions which are 
$U(\fg_{\Sigma_L\backslash\{\sigma\}})$-finite.\;Moreover,\;we put $\ext^1_{G,\sigma,Z}(V_1,V_2):=\ext^1_{G,\sigma}(V_1,V_2)\cap \ext^1_{G,Z}(V_1,V_2)$.\;


Denote by $\sW_{\Sigma_L}\cong \prod_{\sigma\in \Sigma_L}\sW_{n,\sigma}$  the Weyl group of $\res_{L/\bQ_p}\GLN_n$,\;where $\sW_{n,\sigma}\cong \sW_n$.\;For $i\in \Delta$ and $\sigma\in \Sigma_L$,\;let $s_{i,\sigma}\in \sW_{n,\sigma}$ be the simple reflection corresponding to $i\in \Delta$.\;



For any $I\subseteq \Delta$,\;write $\bL_I=\GLN_{h_1}\times\cdots\times \GLN_{h_r}$.\;Then
\[\delta_{\op_I}:=\prod_{i=1}^r|\mathrm{det}_{\GLN_{h_i}}|^{-\sum_{j=i+1}^rh_j+\sum_{j=1}^{i-1}h_j}=\prod_{i=1}^r|\mathrm{det}_{\GLN_{h_i}}|^{-n-h_i+2\sum_{j=1}^{i}h_j}\]
is the  modulus character of $\op_I$.\;Let $\eta_I=\delta_{\op_I}^{1/2}$.\;


For $u\in \sW_s$,\;recall the $E$-point $\underline{x}^u=(x_{u^{-1}(i)})_{1\leq i\leq s}\in \sbanpiku$.\;Note that $\pi_{x_i}$ is the smooth representation $\pi_{\mathrm{sm}}(E_{i})$ of $\GLN_{r_i}(L)$ over $E$ associated with $E_{i}$.\;For $u\in \sW_{s}$,\;put 
\[\pi_{\sm}(\underline{x}^u):=\boxtimes_{j=1}^s\pi_{\mathrm{sm}}(E_{u^{-1}(j)})=\boxtimes_{j=1}^s\pi_{x_{u^{-1}(j)}}.\]
For $S_0^u\subseteq I\subseteq \Delta$,\;consider the smooth principal series of $\bL_{I}(L)$ (if $I=\Delta$,\;we drop the symbol $I=\Delta$):
\begin{equation}
	\begin{aligned}
		\mathrm{PS}^{\infty}_{I,u}(\underline{x})&:=\Big(\ind^{\bL_I(L)}_{\op_{S^u_0}(L)\cap \bL_I(L) }\pi_{\sm}(\underline{x}^u)\eta_{S^u_0}\Big)^{\infty}
	\end{aligned}
\end{equation}
which is irreducible by the generic assumption and is independent on the choice of $u\in \sW_s$ if $I=\Delta$.\;


For $i\in \Delta$,\;put $\widehat{i}=\Delta\backslash\{i\}$.\;For $u\in \sW_s$,\;$i\in \Delta\backslash S_0^u$ and $\sigma\in \Sigma_L$,\;put (see \cite[The main theorem]{orlik2015jordan})
\[C_{i,u,\sigma}:=\cF^{G}_{\op_{S^u_0}}\Big(\overline{L}(-s_{i,\sigma}\cdot\underline{\lambda}),\pi_{\sm}(\underline{x}^u)\eta_{S^u_0}\Big)\cong \cF^{G}_{\op_{\widehat{i}}}\Big(\overline{L}(-s_{i,\sigma}\cdot\underline{\lambda}),\mathrm{PS}^{\infty}_{\widehat{i},u}(\underline{x})\Big),\]
which is an irreducible admissible locally $\sigma$-analytic representation of $G$.\;




%Recall $ \Delta_s=\{1,\cdots,s-1\}$ and  $\Delta\backslash S_0^u=\{t_{1}^u,\cdots,t_{s-1}^u\}=J_u(\Delta_s)$.\;


\begin{lem}For $u,u'\in \sW_s$,\;$1\leq j,j'\leq s$,\;and $\sigma,\sigma'\in \Sigma_L$,\;we have $C_{t^u_j,u,\sigma}\cong C_{t^{u'}_{j'},u',\sigma'}$ if and only if $t^u_j=t^{u'}_{j'}$,\;$\sigma=\sigma'$,\;$j=j'$,\;and $\{u^{-1}(1),\cdots,u^{-1}(j)\}=\{(u')^{-1}(1),\cdots,(u')^{-1}(j')\}$.\;
\end{lem}
%From now on,\;put
%\[\mathcal{OS}^1_{\gen}(\underline{x},\bh)=\Big\{C_{j,u,\sigma}:u\in \sW_{s},j\in \Delta\backslash S_0^u,\sigma\in \Sigma_L\Big\}=\Big\{C_{{t_j^u},u,\sigma}: j\in \Delta_s,u\in\sW_{s}^{\widehat{j},\emptyset},\sigma\in \Sigma_L\Big\}.\;\]
%Note that $|\mathcal{OS}^1_{\gen}(\underline{x},\bh)|=(2^s-2)d_L$.\;




%The following lemma is the same as \cite[Lemma 2.1.31]{BQ24} essentially.\;
%\begin{lem}For $u\in \sW_s$,\;we have $I_u^+:=p_1(D_R(w_{\widehat{j}}u)\cap I_0')$ and
%	$I_u^-:=p_1(D_R(u)\cap I_0')$.\;
%\end{lem}
%\begin{proof}
%If $I_0'=\Delta_s$ or equivalently $I_0=\Delta$,\;this is the same as \cite[Lemma 2.1.31]{BQ24}.\;For general case,\;it only suffices to take intersection with $I_0'$.\;
%\end{proof}






\subsection{Preliminary on extension groups}


%Write ${S^-_0}:=S^{w_{0,s}}_0$ and $\underline{x}^{-}:=\underline{x}^{w_{0,s}}$ for simplicity.\;

The following discussion generalizes the results in  \cite[Section 2.1]{BDcritical25} to potentially crystalline case (or cuspidal case).\;In the sequel,\;write
\[\pi_{\lalg}(\underline{x},\bh)=\mathrm{PS}^{\infty}_{1}(\underline{x})\otimes_E L(\lambda)\cong \mathrm{PS}^{\infty}_{u}(\underline{x})\otimes_E L(\lambda).\]
for all $u\in \sW_s$.\;Similar to \cite[Lemma 2.1.6]{BDcritical25},\;we have
\begin{lem}\label{parameterline} Let $\sigma\in \Sigma_L$.\;We have isomorphisms of finite-dimensional $E$-vector spaces:
	\begin{equation}
		\begin{aligned}
			&\homo_{\sm}(\bZ_{S^{u}_0}(L),E)\oplus_{\homo_{\sm}(L^{\times},E)}\homo(L^{\times},E)\xrightarrow{\sim} \ext^1_{G}\big(\pi_{\lalg}(\underline{x},\bh),\pi_{\lalg}(\underline{x},\bh)\big),\\
\text{resp.,\;}	&\homo_{\sm}(\bZ_{S^{u}_0}(L),E)\oplus_{\homo_{\sm}(L^{\times},E)}\homo_{\sigma}(L^{\times},E)\xrightarrow{\sim} \ext^1_{G,\sigma}\big(\pi_{\lalg}(\underline{x},\bh),\pi_{\lalg}(\underline{x},\bh)\big).\;
		\end{aligned}
	\end{equation}
which have $E$-dimension $s+d_L$ (resp.,\;$s+1$).\;
\end{lem}
\begin{proof}Write $\pi_{\lalg}:=\pi_{\lalg}(\underline{x},\bh)$ for simplicity.\;We prove that $\dim_E\ext^1_{G}\left(\pi_{\lalg},\pi_{\lalg}\right)$ is $s+d_L$ and define the maps.\;We define a canonical injection:
	\begin{equation}\label{lalginj1}
		\begin{aligned}
			\homo_{\sm}(\bZ_{S^u_0}(L),E)&\rightarrow\ext^1_{G}\big(\pi_{\lalg},\pi_{\lalg}\big),\;\\
			\psi&\mapsto \Big(\ind^G_{\op_{S_0^u}(L)}\pi_{\sm}(\underline{x}^u)\eta_{S_0^u}\otimes_E\big(1+\big(\psi\circ\mathrm{det}_{\bL_{S_0^u}}\big)\epsilon\big)\Big)^{\mathrm{sm}}\otimes_EL(\lambda).
		\end{aligned}
	\end{equation}
On the other hand,\;we define a canonical injection
\[\homo(L^{\times},E)\hookrightarrow \ext^1_{G}\left(\pi_{\lalg},\pi_{\lalg}\right),\;\psi\mapsto \pi_{\lalg}\otimes_E\big(1+\big(\psi\circ\mathrm{det}_{\GLN_n}\big)\epsilon\big).\]
These two injections are coincides on $\homo_{\mathrm{sm}}(L^{\times},E)$.\;By the same argument as in \cite[Lemma 3.16]{HigherLinvariantsGL3(Qp)},\;we have
\[\dim_E\ext^1_{G}\left(\pi_{\lalg},\pi_{\lalg}\right)=s+d_L,\;\dim_E\ext^1_{G,\sigma}\left(\pi_{\lalg},\pi_{\lalg}\right)=s+1.\]
Then we get the first isomorphism in the lemma,\;and its restriction on $\homo_{\sigma}(L^{\times},E)$ gives the second isomorphism.\;
\end{proof}


\begin{lem}\label{parameterlineextgp}
	Let $j\in \Delta_s$,\;$u\in \sW_{s}^{\widehat{j},\emptyset}$ and $\sigma\in \Sigma_L$.\;We have 
	\begin{equation}
		\begin{aligned}
			\dim_E\ext^1_G\left(C_{t^u_j,u,\sigma},\pi_{\lalg}(\underline{x},\bh)\right)=1,\;
			\dim_E\ext^1_G\left(\pi_{\lalg}(\underline{x},\bh),C_{t^u_j,u,\sigma}\right)=1.
		\end{aligned}
	\end{equation}
%	There exists  a unique (up to a scalar) locally analytic $\sigma$-analytic representations of the form $[\pi_{\lalg}(\underline{x},\bh)-C_{t^u_j,u,\sigma}]$ (resp.,\;$[C_{t^u_j,u,\sigma}-\pi_{\lalg}(\underline{x},\bh)]$). 
\end{lem}
\begin{proof}The assertion follows from \cite[Proposition 5.1.14]{BQ24}.\;
\end{proof}

Let $\sigma\in \Sigma_L$,\;$u\in \sW_s$ and  $i=t^u_j\in \Delta\backslash S_0^u$ for some $j\in \Delta_s$,\;fix $V_{i,u,\sigma}$ a locally $\sigma$-analytic representation of $G$,\;which is isomorphic to a non-split extension of $C_{i,u,\sigma}$ by $\pi_{\lalg}(\underline{x},\bh)$ (see Lemma \ref{parameterlineextgp},\;$V_{i,u,\sigma}$ is unique up to a scalar),\;i.e.,\;
\[V_{i,u,\sigma}:=\big[\pi_{\lalg}(\underline{x},\bh)-C_{i,u,\sigma}\big],\]
and fix an injection $\iota_{i,u,\sigma}:\pi_{\lalg}(\underline{x},\bh)\hookrightarrow V_{i,u,\sigma}$.\;We have   $C_{i,u,\sigma}\cong V_{i,u,\sigma}/\pi_{\lalg}(\underline{x},\bh)$ and  $\bL_{S^u_0}\subseteq \bL_{\widehat{i}}\cong \GLN_i\times\GLN_{n-i}$.\;


\begin{pro}\label{reinterforextgp}There is a canonical isomorphism of $(s+2)$-dimensional $E$-vector spaces associated to $(V_{i,u,\sigma},\iota_{i,u,\sigma})$:
	\begin{equation}\label{isoforfullextgp}
		\begin{aligned}
			\homo_{\sm}(\bZ_{S^u_0}(L),E)\oplus_{\homo_{\sm}(\bL_{\widehat{i}}(L),E)}\homo_{\sigma}(\bL_{\widehat{i}}(L),E)\xrightarrow{\sim}\ext^1_{G,\sigma}\left(\pi_{\lalg}(\underline{x},\bh),V_{i,u,\sigma}\right).
					\end{aligned}
	\end{equation}
and the restriction of (\ref{isoforfullextgp}) on $\homo_{\sigma}(\bL_{\widehat{i}}(L),E)$ only depends on $(V_{i,u,\sigma},\iota_{i,u,\sigma})$.\;Moreover,\;there is a short exact sequence:
\begin{equation}
	\begin{aligned}
		0\rightarrow \ext^1_{G,\sigma}\Big(\pi_{\lalg}(\underline{x},\bh),\pi_{\lalg}(\underline{x},\bh)\Big) &\xrightarrow{\iota_{i,u,\sigma}} \ext^1_{G,\sigma}\Big(\pi_{\lalg}(\underline{x},\bh),V_{i,u,\sigma}\Big) \\&\rightarrow \ext^1_{G,\sigma}\Big(\pi_{\lalg}(\underline{x},\bh),V_{i,u,\sigma}/\pi_{\lalg}(\underline{x},\bh)\Big)\rightarrow 0.
	\end{aligned}
\end{equation}
\end{pro}
\begin{proof}It follows from Lemma \ref{parameterline} and Lemma \ref{parameterlineextgp} and 
	the  exact sequence $0\rightarrow \ext^1_{G,\sigma}\Big(\pi_{\lalg},\pi_{\lalg}\Big) \rightarrow \ext^1_{G,\sigma}\Big(\pi_{\lalg},V_{i,u,\sigma}\Big) \rightarrow \ext^1_{G,\sigma}\Big(\pi_{\lalg},V_{i,u,\sigma}/\pi_{\lalg}\Big)$ that $\dim_E\ext^1_{G,\sigma}\Big(\pi_{\lalg},V_{i,u,\sigma}\Big)\leq s+2$.\;We construct the desired map and show that the last map in such exact sequence is surjective.\;We first define the injection (by Lemma \ref{parameterline}):
\begin{equation}
		\homo_{\sm}(\bZ_{S^u_0}(L),E)\hookrightarrow \ext^1_{G,\sigma}\Big(\pi_{\lalg},\pi_{\lalg}\Big)\hookrightarrow \ext^1_{G,\sigma}\Big(\pi_{\lalg},V_{i,u,\sigma}\Big).\;
\end{equation}
We next construct a canonical injection $\homo_{\sigma}(\bL_{\widehat{i}}(L),E)\hookrightarrow\ext^1_{G,\sigma}\left(\pi_{\lalg}(\underline{x},\bh),V_{i,u,\sigma}\right)$.\;Put
	\begin{equation}
		R_{u}:=\left(\ind^G_{\op_{\widehat{i}}(L)}\mathrm{PS}^{\infty}_{\widehat{i},u}(\underline{x})\otimes L_{\widehat{i}}(\lambda_{\sigma})\right)^{\sigma-\ana}.\;
	\end{equation}
Consider the following injection:
	\begin{equation}
		\begin{aligned}
			\homo_{\sigma}(\bL_{\widehat{i}}(L),E)&\hookrightarrow \ext^1_{G,\sigma}\left(R_{u},R_{u}\right)\\
			\psi&\mapsto \left(\ind^G_{\op_{\widehat{i}}(L)}\mathrm{PS}^{\infty}_{\widehat{i},u}(\underline{x})\otimes L_{\widehat{i}}(\lambda_{\sigma})\otimes(1+(\psi\circ\det\nolimits_{\bL_{\widehat{i}}})\epsilon)\right)^{\sigma-\ana}.
		\end{aligned}
	\end{equation}
	Composed with the pull-back map for the natural map $\pi_{\lalg}\hookrightarrow R_{u}$,\;we get a map $f_u:\homo_{\sigma}(\bL_{\widehat{i}}(L),E)\hookrightarrow \ext^1_{G,\sigma}\big(\pi_{\lalg},R_{u}\big)\rightarrow  \ext^1_{G,\sigma}\big(\pi_{\lalg},V_{i,u,\sigma}\big)
	$ (it is still an injection and we can prove that it only depends on $(V_{i,u,\sigma},\iota_{i,u,\sigma})$.\;The map $f_u$ induces an injection
	\[\homo_{\sm}(\bL_{\widehat{i}}(L),E)\rightarrow \ext^{1}_{G,\sigma}\big(\pi_{\lalg},\pi_{\lalg}\big),\]
	which is compatible with the injection $\homo_{\sm}(\bZ_{S^u_0}(L),E)\rightarrow\ext^1_{G,\sigma}\big(\pi_{\lalg},\pi_{\lalg}\big)$.\;Thus,\;we get the desired map.\;On the other hand,\;by the second isomorphism in Lemma \ref{reinterforextgp},\;we see that the image of $\homo_{\sigma}(\bL_{\widehat{i}}(L),E)$ in $\ext^1_{G,\sigma}\big(\pi_{\lalg},V_{i,u,\sigma}\big)$ is not contained in $\ext^1_{G,\sigma}\left(\pi_{\lalg},\pi_{\lalg}\right)$,\;so that the second map in (\ref{isoforfullextgp}) is an surjection and thus an isomorphism by comparing the dimensions.\;
\end{proof}





\begin{pro}
	We have a perfect paring of $1$-dimensional $E$-vector spaces:
	\begin{equation}\label{dualpairing}
		\ext^1_{G,\sigma}\left(\pi_{\lalg}(\underline{x},\bh),C_{i,u,\sigma}\right)\times\ext^1_{G,\sigma}\left(C_{i,u,\sigma},\pi_{\lalg}(\underline{x},\bh)\right)\xrightarrow{\sim}\homo_{\sigma}(\GLN_{n-i}(\cO_L),E).\;
	\end{equation}
	and a canonical isomorphism of $1$-dimensional $E$-vector space $E\xrightarrow{\sim }\homo_{\sigma}(\GLN_{n-i}(\cO_L),E)$,\;$1\mapsto \log_{p,\sigma}\circ \det$.\;
\end{pro}
\begin{proof}
	Similar to the proof of \cite[Proposition 2.1.10]{BDcritical25},\;we have the following (formal) isomorphism of $1$-dimensional $E$-vector spaces:
	\[\ext^1_{G,\sigma}\left(C_{i,u,\sigma},\pi_{\lalg}\right)\otimes_E\ext^1_{G,\sigma}\left(\pi_{\lalg},C_{i,u,\sigma}\right)^{\vee}\xrightarrow{\sim}\ext^1_{G,\sigma}\left(C_{i,u,\sigma}\otimes\ext^1_{G,\sigma}\left(\pi_{\lalg},C_{i,u,\sigma}\right),\pi_{\lalg}\right),\]
	where the $G$-representation $C_{i,u,\sigma}\otimes\ext^1_{G,\sigma}\left(\pi_{\lalg},C_{i,u,\sigma}\right)$ has the trivial action of $G$ on $\ext^1_{G,\sigma}\left(\pi_{\lalg},C_{i,u,\sigma}\right)$. Any 
	representative $V'_{i,u,\sigma}$ of the image of the canonical vector of the left hand side is isomorphic to a non-split extension of $C_{i,u,\sigma}$ by $\pi_{\lalg}$.\;Then $V'_{i,u,\sigma}$ corresponds to an injection $\iota:\pi_{\lalg}\hookrightarrow V'_{i,u,\sigma}$ and an isomorphism $\kappa:V'_{i,u,\sigma}/\pi_{\lalg}\xrightarrow{\sim}C_{i,u,\sigma}\otimes_E\ext^1_{G,\sigma}\left(\pi_{\lalg},C_{i,u,\sigma}\right)$.\;By the same strategy as in the \cite[(37)]{BDcritical25},\;we have a canonical isomorphism of $1$-dimensional $E$-vector spaces (where the embedding $\homo_{\sigma}(\cO_L^{\times},E)\hookrightarrow \homo_{\sigma}(\bL_{\widehat{i}}(\cO_L^{\times}),E)$ is induced by $\det:\bL_{\widehat{i}}(\cO_L^{\times})\rightarrow \cO_L^{\times}$):
	\[\homo_{\sigma}(\GLN_{n-i}(\cO_L),E)\xrightarrow{\sim}\homo_{\sigma}(\bL_{\widehat{i}}(\cO_L^{\times}),E)/\homo_{\sigma}(\cO_L^{\times},E).\]
	By the proof of Proposition \ref{reinterforextgp} (similar to \cite[(34)]{BDcritical25}),\;the latter space is isomorphic to $\ext^1_{G,\sigma}\left(\pi_{\lalg},V/\pi_{\lalg}\right)$ and thus isomorphic to (via $\kappa$) $\ext^1_{G,\sigma}\left(\pi_{\lalg},C_{i,u,\sigma}\otimes_E\ext^1_{G,\sigma}\left(\pi_{\lalg},C_{i,u,\sigma}\right)\right)\xleftarrow{\sim} \ext^1_{G,\sigma}\left(\pi_{\lalg},C_{i,u,\sigma}\right)\otimes_E\ext^1_{G,\sigma}\left(\pi_{\lalg},C_{i,u,\sigma}\right)$,\;the result follows.\;
\end{proof}
\begin{rmk}\label{rmkforextgroupsigma}Similar to the Step $1$ in the proof of \cite[Proposition 2.2.4]{BDcritical25},\;the surjection $\bL_{\widehat{i}}(L)\twoheadrightarrow \GLN_{n-i}(L)$ gives a canonical injection $\homo_{\sigma}(\GLN_{n-i}(L),E)\hookrightarrow \homo_{\sigma}(\bL_{\widehat{i}}(L),E)$.\;A choice of $\log_p(p)$ define a section $\homo_{\sigma}(\cO_L^{\times},E)\hookrightarrow \homo_{\sigma}(L^{\times},E)$ to the restriction map $\homo_{\sigma}(L^{\times},E)\hookrightarrow \homo_{\sigma}(\cO_L^{\times},E)$,\;and thus a section $\homo_{\sigma}(\GLN_{n-i}(\cO_L^{\times}),E)\hookrightarrow \homo_{\sigma}(\GLN_{n-i}(L),E)$ to the restriction map $\homo_{\sigma}(\GLN_{n-i}(L),E)\hookrightarrow \homo_{\sigma}(\GLN_{n-i}(\cO_L^{\times}),E)$,\;which induces an isomophism:
	\begin{equation*}
		\begin{aligned}
			\Big(\homo_{\sm}(\bL_{\widehat{i}}(L),E)\oplus_{\homo_{\sm}(L^{\times},E)}\homo(L^{\times},E)\Big)\oplus \homo_{\sigma}(\GLN_{n-i}(\cO_L^{\times}),E)\xrightarrow{\sim} \homo_{\sigma}(\bL_{\widehat{i}}(L),E).
		\end{aligned}
	\end{equation*}
	We thus get (similar to \cite[(59)]{BDcritical25}) isomorphsims $\homo_{\sigma}(\GLN_{n-i}(\cO_L^{\times}),E)\xrightarrow{\sim}\ext^1_{G,\sigma}\left(\pi_{\lalg},V_{i,u,\sigma}/\pi_{\lalg}\right)$ and 
	\begin{equation*}
		\begin{aligned}
			\Big(\homo_{\sm}(\bZ_{S^u_0}(L),E)\oplus_{\homo_{\sm}(L^{\times},E)}\homo(L^{\times},E)\Big)\oplus \homo_{\sigma}(\GLN_{n-i}(\cO_L^{\times}),E)\xrightarrow{\sim} \ext^1_{G,\sigma}\left(\pi_{\lalg},V_{i,u,\sigma}\right).
		\end{aligned}
	\end{equation*}
\end{rmk}




\subsubsection{Constructions of locally analytic representations}


For $u\in\sW_s$,\;consider the locally  $\bQ_p$-analytic  principal series:
\begin{equation}\label{locaparabolicind}
	\begin{aligned}
	&\mathrm{PS}_{u}(\underline{x},\bh):=\Big(\ind^G_{\op_{S^u_0}(L)}\pi_{\sm}(\underline{x}^u)\eta_{S^u_0}\otimes_E L_{S_0^u}(\lambda)\Big)^{\bQ_p-\ana}.\\
%	&\mathrm{PS}_{u,\sigma}(\underline{x},\bh):=\Big(\ind^G_{\op_{S^u_0}(L)}\pi_{\sm}(\underline{x}^u)\eta_{S^u_0}\otimes_E L_{S_0^u}(\lambda)\Big)^{\sigma-\ana}.\;
	\end{aligned}
\end{equation}
Note that the locally $\bQ_p$-algebraic vectors  $\mathrm{PS}^{\lalg}_{u}(\underline{x},\bh)$ in $\mathrm{PS}_{u}(\underline{x},\bh)$ is isomorphic to  $\soc_G\mathrm{PS}_{u}(\underline{x},\bh)=\pi_{\lalg}(\underline{x},\bh)$.\;Let $\mathrm{PS}_{u,1}(\underline{x},\bh)$ be the first two layers of the socle filtrartion of $\mathrm{PS}_{u}(\underline{x},\bh)$.\;By \cite[The main theorem]{orlik2015jordan}),\;
\[\mathrm{PS}_{u,1}(\underline{x},\bh)=\bigoplus_{\pi_{\lalg}(\underline{x},\bh)}^{{j\in \Delta_s,\sigma\in \Sigma_L}}V_{t^u_j,[u]_j,\sigma}.\;\]
Consider the amalgamated sum:
\[\pi^{\flat}_{1}(\underline{x},\bh):=\bigoplus_{\pi_{\lalg}(\underline{x},\bh)}^{\substack{j\in \Delta_s,u\in\sW_{s}^{\widehat{j},\emptyset}\\\sigma\in \Sigma_L}}V_{t^u_j,u,\sigma},\;\]
so that we have a short exact sequence:
\begin{equation}\label{firstwhole}
	0\rightarrow \pi_{\lalg}(\underline{x},\bh) \rightarrow \pi^{\flat}_{1}(\underline{x},\bh)\rightarrow \oplus_{\substack{j\in \Delta_s,u\in\sW_{s}^{\widehat{j},\emptyset}\\\sigma\in \Sigma_L}}C_{t^u_j,u,\sigma}\rightarrow 0.\;
\end{equation}
Moreover,\;$\pi^{\flat}_{1}(\underline{x},\bh)$ is the unique quotient of the amalgamated sum 
\[\bigoplus^{u\in \sW_{s}}_{\pi_{\lalg}(\underline{x},\bh)}\mathrm{PS}_{u,1}(\underline{x},\bh)\]
of socle $\pi_{\lalg}(\underline{x},\bh)$.\;We have injection $\mathrm{PS}_{u,1}(\underline{x},\bh)\hookrightarrow \pi^{\flat}_{1}(\underline{x},\bh)$ of $G$-representations for any $u\in \sW_{s}$.\;


For $u\in \sW_s$,\;we consider the natural map
\begin{equation}
	\begin{aligned}
		\homo\big(\bZ_{S^u_0}(L),E\big)&\rightarrow\ext^1_{G}\big(\mathrm{PS}_{u}(\underline{x},\bh),\mathrm{PS}_{u}(\underline{x},\bh)\big),\;\\
		\psi&\mapsto \Big(\ind^G_{\op_{S^u_0}(L)}\pi_{\sm}(\underline{x}^u)\eta_{S^u_0}\otimes_EL_{S^u_0}(\lambda)\otimes_E(1+(\psi\circ\mathrm{det}_{\bL_{S_0^u}})\epsilon)\Big)^{\bQ_p-\ana}.
	\end{aligned}
\end{equation}
In particular,\;by Schraen's spectral sequence \cite[(4.37)]{schraen2011GL3},\;we have an isomorphism
\begin{equation}
	\begin{aligned}
		\homo\big(\bZ_{S^u_0}(L),E\big)&\xrightarrow{\sim}\ext^1_{G}\big(\pi_{\lalg}(\underline{x},\bh),\mathrm{PS}_{u}(\underline{x},\bh)\big).
	\end{aligned}
\end{equation}
By \cite[Lemma 2.26]{2019DINGSimple},\;this map factors through (induced by the injection $\mathrm{PS}_{u,1}(\underline{x},\bh)\hookrightarrow \mathrm{PS}_{u}(\underline{x},\bh)$)
\[\ext^1_{G}\big(\pi_{\lalg}(\underline{x},\bh),\mathrm{PS}_{u,1}(\underline{x},\bh)\big)\rightarrow \ext^1_{G}\big(\pi_{\lalg}(\underline{x},\bh),\mathrm{PS}_{u}(\underline{x},\bh)\big),\]
then we obtain a map $\homo\big(\bZ_{S^u_0}(L),E\big)\rightarrow\ext^1_{G}\big( \pi_{\lalg}(\underline{x},\bh),\mathrm{PS}_{u,1}(\underline{x},\bh)\big)$.\;Composing with the push-forward map via the injection $\mathrm{PS}_{u,1}(\underline{x},\bh)\hookrightarrow \pi^{\flat}_1(\underline{x},\bh)$,\;we actually obtain a map:
\[\zeta_{u}:\homo\big(\bZ_{S^u_0}(L),E\big)\rightarrow\ext^1_{G}\big(\pi_{\lalg}(\underline{x},\bh),\pi^{\flat}_1(\underline{x},\bh)\big).\;\]




\begin{pro}\label{extpi1repn}
	\begin{itemize}
		\item[(1)] We have  $\dim_E\ext^1_{G}\big(\pi_{\lalg}(\underline{x},\bh),\mathrm{PS}_{u,1}(\underline{x},\bh)\big)=s(1+d_L)$ and an exact sequence:
		\begin{equation}\label{extenforPSw}
			\begin{aligned}
				0\rightarrow \ext^1_{G}\big(\pi_{\lalg}(\underline{x},\bh),\pi_{\lalg}(\underline{x},\bh)\big)\rightarrow &\;\ext^1_{G}\big(\pi_{\lalg}(\underline{x},\bh),\mathrm{PS}_{u,1}(\underline{x},\bh)\big)\\
				&\rightarrow \oplus_{\substack{j\in \Delta_s\\\sigma\in \Sigma_L}}  \ext^1_{G}\Big(\pi_{\lalg}(\underline{x},\bh),C_{t_j^u,[u]_j,\sigma}\Big)\rightarrow 0 .
			\end{aligned}
		\end{equation}
		Thus,\;we get an ismorphism $\zeta_u:\homo\big(\bZ_{S^u_0}(L),E\big)\xrightarrow{\sim}\ext^1_{u}\big(\pi_{\lalg}(\underline{x},\bh),\mathrm{PS}_{u,1}(\underline{x},\bh)\big)$.\;
		\item[(2)]We have $\ext^1_{G}\left(\pi_{\lalg}(\underline{x},\bh),\pi^{\flat}_1(\underline{x},\bh)\right)=s+(2^s-1)d_L$ and an short exact sequence:
		\begin{equation}
			\begin{aligned}
				0\rightarrow \ext^1_{G}\big(\pi_{\lalg}(\underline{x},\bh),\pi_{\lalg}(\underline{x},\bh)\big)\rightarrow &\;\ext^1_{G}\big(\pi_{\lalg}(\underline{x},\bh),\pi^{\flat}_1(\underline{x},\bh)\big)\\
				&\rightarrow \oplus_{\substack{j\in \Delta_s,u\in \sW_{s}^{\widehat{j},\emptyset}\\\sigma\in \Sigma_L}}  \ext^1_{G}\left(\pi_{\lalg}(\underline{x},\bh),C_{t^u_j,u,\sigma}\right)\rightarrow 0 .
			\end{aligned}
		\end{equation}
	\end{itemize}
\end{pro}
\begin{proof}
	For $(1)$,\;by Proposition \ref{reinterforextgp},\;the second last map in (\ref{extenforPSw}) is surjective.\;On the other hand,\;
	the natural morphism (by Remark \ref{rmkforextgroupsigma})
	\[\homo\big(\bZ_{S^u_0}(L),E\big)\rightarrow \Big(\homo_{\sm}(\bZ_{S^u_0}(L),E)\oplus_{\homo_{\sm}(L^{\times},E)}\homo(L^{\times},E)\Big)\oplus \Big(\oplus_{\sigma,i\in \Delta\backslash S_0^u}\homo_{\sigma}(\GLN_{n-i}(\cO_L^{\times}),E)\Big)\]
	is an isomorphism,\;therefore $\zeta_u$ is an ismorphism.\;When $u$ varying,\;we prove that the second last map in $(2)$ is also surjective since it is covered by the last term in (\ref{extenforPSw}) with $u$ varying.\;The dimension of $\ext^1_{G}\left(\pi_{\lalg}(\underline{x},\bh),\pi^{\flat}_1(\underline{x},\bh)\right)$ is equal to $s+d_L+d_L\sum_{j\in \Delta_s}\#\sW_{s}^{\widehat{j},\emptyset}=s+(2^s-1)d_L$.\;
\end{proof}

%Therefore,\;let $\ext^1_{u}\left(\pi_{\lalg}(\underline{x},\bh),\pi^{\flat}_1(\underline{x},\bh)\right)$ be the image of  $\ext^1_{G}\left(\pi_{\lalg}(\underline{x},\bh),\pi^{\flat}_{1}(\underline{x},\bh)_u\right)$ in the full space $\ext^1_{G}\left(\pi_{\lalg}(\underline{x},\bh),\pi^{\flat}_{1}(\underline{x},\bh)\right)$ via the injection $\pi^{\flat}_{1}(\underline{x},\bh)_u\hookrightarrow \pi^{\flat}_{1}(\underline{x},\bh)$.\;From Proposition \ref{extpi1repn},\;we see that $\ext^1_{u}\left(\pi_{\lalg}(\underline{x},\bh),\pi^{\flat}_1(\underline{x},\bh)\right)$ is also equal to $\mathrm{Im}(\zeta_u)$.\;Denote by
%\[\ext^1_{g}\left(\pi_{\lalg}(\underline{x},\bh),\pi^{\flat}_{1}(\underline{x},\bh)\right)\subset \ext^1_{g'}\left(\pi_{\lalg}(\underline{x},\bh),\pi^{\flat}_{1}(\underline{x},\bh)\right)\Big(\subset\ext^1_{u}\left(\pi_{\lalg}(\underline{x},\bh),\pi^{\flat}_{1}(\underline{x},\bh)\right)\Big)\]
%the respective image of  $\ext^{1,\lalg}_{G}\left(\pi_{\lalg}(\underline{x},\bh),\pi_{\lalg}(\underline{x},\bh)\right)$,\;$\ext^1_{G}\left(\pi_{\lalg}(\underline{x},\bh),\pi_{\lalg}(\underline{x},\bh)\right)$ via (\ref{extenforPSw}).\;

\begin{rmk}\label{rmkforsigmaextANA}
Fix $\sigma\in\Sigma_L$,\;replacing the locally $\bQ_p$-analytic parabolic induction in (\ref{locaparabolicind}) by locally $\sigma$-analytic parabolic induction,\;we obtain a $\sigma$-analytic version of all above discussion.\;We get a locally $\sigma$-analytic representation  $\pi^{\flat}_{1,\sigma}(\underline{x},\bh)$ which lies in the following exact sequence:
\begin{equation}\label{firstwholesigma}
	0\rightarrow \pi_{\lalg}(\underline{x},\bh) \rightarrow \pi^{\flat}_{1,\sigma}(\underline{x},\bh)\rightarrow \oplus_{\substack{j\in \Delta_s\\u\in \sW_{s}^{\widehat{j},\emptyset}}}C_{t_j^u,u,\sigma}\rightarrow 0.\;
\end{equation}
Note that $\pi^{\flat}_{1,\sigma}(\underline{x},\bh)=\bigoplus^{\pi_{\lalg}(\underline{x},\bh)}_{j\in \Delta_s,u\in\sW_{s}^{\widehat{j},\emptyset}}V_{t^u_j,u,\sigma}$ and $\pi^{\flat}_{1}(\underline{x},\bh)=\bigoplus^{\pi_{\lalg}(\underline{x},\bh)}_{\sigma\in\Sigma_L}\pi^{\flat}_{1,\sigma}(\underline{x},\bh)$.\;We have  a short exact sequence:
\begin{equation}
	\begin{aligned}
		0\rightarrow \ext^1_{G,\sigma}\big(\pi_{\lalg}(\underline{x},\bh),\pi_{\lalg}(\underline{x},\bh)\big)\rightarrow &\;\ext^1_{G,\sigma}\big(\pi_{\lalg}(\underline{x},\bh),\pi^{\flat}_{1,\sigma}(\underline{x},\bh)\big)\\
		&\rightarrow \oplus_{\substack{j\in \Delta_s\\u\in \sW_{s}^{\widehat{j},\emptyset}}}  \ext^1_{G,\sigma}\left(\pi_{\lalg}(\underline{x},\bh),C_{t^u_j,u,\sigma}\right)\rightarrow 0 .
	\end{aligned}
\end{equation}
Moreover,\;$\dim_E\ext^1_{G,\sigma}\left(\pi_{\lalg}(\underline{x},\bh),\pi^{\flat}_{1,\sigma}(\underline{x},\bh)\right)=s+(2^s-1)$.\;In a similar way,\;for $u\in\sW_s$,\;we get a map:
\[\zeta_{u,\sigma}:\homo_{\sigma}(\bZ_{S^u_0}(L),E)\rightarrow\ext^1_{G,\sigma}\big(\pi_{\lalg}(\underline{x},\bh),\pi^{\flat}_{1,\sigma}(\underline{x},\bh)\big).\;\]
\end{rmk}








\subsection{Main results}\label{reconstD}


This section follows the route in the recent work \cite[Section 2.4]{BDcritical25} of Breuil-Ding.\;Fix $\sigma\in \Sigma_L$,\;we will construct a map which only depends  on a choice of $\log_p(p)\in E$  (recall the locally $\sigma$-analytic representation  $\pi^{\flat}_{1,\sigma}(\underline{x},\bh)$  in Remark \ref{rmkforsigmaextANA}):
	\begin{equation}\label{constructionvarphi}
	\begin{aligned}
		t^{\flat}_{D_{\sigma}}:\ext^1_{G,\sigma}\left(\pi_{\lalg}(\underline{x},\bh),\pi^{\flat}_{1,\sigma}(\underline{x},\bh)\right)\rightarrow \ext^{1,\flat}_{\gr}({\Dpik},{\Dpik})\oplus\homo_{\fil}(D_{\sigma},D_{\sigma}).\;
	\end{aligned}
\end{equation}
(see Section \ref{reinterfor33} for the spaces $\ext^{1,\flat}_{\gr}({\Dpik},{\Dpik})$ and $\homo_{\fil}(D_{\sigma},D_{\sigma})$) and 
\begin{itemize}
	\item[(a)] describe the image of $t^{\flat}_{D_{\sigma}}$.\;
	\item[(b)] describe what information of  Hodge filtration  is determined by the kernel of  $t^{\flat}_{D_{\sigma}}$.\;
\end{itemize}
In the sequel,\;put
\[\Delta'=\{t_j^u:u\in \sW_s,j\in \Delta_s\}=\cup_{u\in \sW_s}\Delta\backslash S_0^u\subseteq\Delta.\]
For $i\in\Delta'$,\;put $\cI_{i}^{\flat}:=\{u\in \sW_s:\;i=t_{l_i(u)}^u \text{\;for some\;} 1\leq l_i(u)\leq s-1\}=\{u\in \sW_s:i\in \Delta\backslash S_0^u\}$.\;For $u\in \cI_{i}^{\flat}$,\;let 
\[I_i(u)=\{u^{-1}(1),\cdots,u^{-1}(l_i(u))\},\]
and $I_i(u)^c=\{1,\cdots,s\}\backslash I_i(u)$.\;In particular,\;if $S_0=\emptyset$,\;then $\Delta'=\Delta$ and $\cI_{i}^{\flat}\cong \{I\subseteq\{1,\cdots,n\}:\;|I|=i\}$.\;

%$\cI_{i}^{\flat}\cong \{(a_1,\cdots,a_i):\; a_j\in\{1,\cdots,n\},\;1\leq j\leq i\}$.\;

Recall that we have fixed a basis $\{e_{1,j,\sigma}\}_{1\leq j\leq r_1},\cdots,\{e_{s,j,\sigma}\}_{1\leq j\leq r_s}$ of $D_{\sigma}^1,\cdots,D_{\sigma}^s$ respectively,\;and we relabel them with $e_{1,\sigma},\cdots,e_{n,\sigma}$  by using lexicographical ordering (a basis of $D_{\sigma}=\oplus_{1\leq i\leq s}D_{\sigma}^i$).\;For $1\leq i\leq s$,\;let 
\[I_i=\{q:\;1\leq q\leq n,\;e_{q,\sigma}=e_{i,j,\sigma} \text{\;for some\;} 1\leq j\leq r_i\}=\{t_{i-1}^1+1,\cdots,t_i^1\}.\]
For $J\subseteq \{1,\cdots,n\}$,\;put
\[e_J:=\wedge_{j\in J}e_{j,\sigma}\in \bigwedge_E\nolimits^{\!|J|}\!D_{\sigma}.\]
For $1\leq i\leq s$,\;put 
\[\be_i:=\wedge_{j=1}^{r_i}e_{i,j,\sigma}=e_{I_i}\in\bigwedge\nolimits_E^{\!r_i}\!D^i_{\sigma}.\;\]
For $I\subseteq \{1,\cdots,s\}$,\;put $D_{\sigma}[I]:=\oplus_{i\in I}D_{\sigma}^i$ and 
\[\be_I:=\wedge_{i\in I}\be_i\in \bigwedge_E\nolimits^{\!\sum_{i\in I}r_i}\!D_{\sigma}[I]\subset \bigwedge_E\nolimits^{\!\sum_{i\in I}r_i}\!D_{\sigma}.\;\]


\subsubsection{Preliminaries}






%\begin{rmk}
% (Subspace filtration) There exists the induced subspace filtration $\fil_{\mathrm{sub}}^{\bullet}(D_{\sigma,(i)}^{\flat})$ of $\fil_H^{\bullet}(D_{\sigma})$ on $D_{\sigma,(i)}^{\flat}$.\;In an adapted basis of $D_{\sigma}=  D_{\sigma}^{\star,(i)} \oplus D^{\flat}_{\sigma,(i)}$,\;by the non-critical assumption,\;we see that $\gr_{\mathrm{sub}}^{-\bh_j}(D^{\flat}_{\sigma,(i)})$ if and only if $1\leq j\leq n-d_{(i)}$ and $	\fil_{\mathrm{sub}}^{-\bh_j}(D^{\flat}_{\sigma,(i)})\cong \fil_H^{-\bh_{j}}(D_{\sigma})/\fil_H^{-\bh_{n-d_{(i)}+1}}(D_{\sigma})$ when $1\leq j\leq n-d_{(i)}$.\;
%
%\end{rmk}





For each $u\in \cI_{i}^{\flat}$,\;we fix an
isomorphism of $1$-dimensional $E$-vector spaces (as in \cite[(43)]{BDcritical25}):
\begin{equation}
	\epsilon_{u,i}:\ext^1_{G,\sigma}\big(C_{i,u,\sigma},\pi_{\lalg}(\underline{x},\bh)\big)\xrightarrow{\sim}E\be_{I_i(u)^c}\in \bigwedge\nolimits_E^{\!n-i}\!D_{\sigma}.
\end{equation}
Then we get for each $i\in\Delta'$ a map:
\begin{equation}\label{dfnforeplisonistar}
	\epsilon_{i}:=\bigoplus_{u\in \cI_{i}^{\flat}}\epsilon_{u,i}:\ext^1_{G,\sigma}\Big(\bigoplus_{u\in \cI_{i}^{\flat}}C_{i,u,\sigma},\pi_{\lalg}(\underline{x},\bh)\Big)\rightarrow\bigwedge\nolimits_E^{\!n-i}\!D_{\sigma}.\;
\end{equation}
Let $\big(\bigwedge_{E}^{\!n-i}\!D_{\sigma}\big)^{\flat}$ be the image of $\epsilon_{i}$.\;Let
\[\big(\bigwedge\nolimits_{E}^{\!n-i}\!D_{\sigma}\big)^{\flat,c}=E\big\langle e_J:\;J\subseteq \{1,\cdots,n\},|J|=n-i,e_J\neq \be_{I_i(u)^c},\;\forall\;u\in \cI_{i}^{\flat}\big\rangle.\]
Then $\bigwedge_E^{n-i}D_{\sigma}=\big(\bigwedge\nolimits_{E}^{\!n-i}\!D_{\sigma}\big)^{\flat}\oplus\big(\bigwedge\nolimits_{E}^{\!n-i}\!D_{\sigma}\big)^{\flat,c}$ and  $\dim_E\big(\bigwedge_{E}^{\!n-i}\!D_{\sigma}\big)^{\flat}=|\cI_{i}^{\flat}|$.\;In the sequel,\;we identity $\big(\bigwedge_{E}^{\!n-i}\!D\big)^{\flat}$ with the quotient $\bigwedge_{E}^{\!n-i}\!D/\big(\bigwedge_{E}^{\!n-i}\!D\big)^{\flat,c}$.\;



For $1\leq i\leq n-1$,\;we define the following $E$-line of $\bigwedge_E^{n-i}D_{\sigma}$:
\begin{equation}
	\begin{aligned}
		\fil_i^{\max}(D_{\sigma})&:= \fil_H^{-\bh_{n,\sigma}}(D_{\sigma})\wedge \fil_H^{-\bh_{n-1,\sigma}}(D_{\sigma})\wedge\cdots\wedge \fil_H^{-\bh_{i+1,\sigma}}(D_{\sigma})\\
		&\xrightarrow{\sim}\bigwedge\nolimits_E^{\!n-i}\!\fil_H^{-\bh_{i+1,\sigma}}(D_{\sigma})\subseteq\bigwedge\nolimits_E^{\!n-i}\!D_{\sigma} .\;
	\end{aligned}
\end{equation}
For each $u\in \cI_{i}^{\flat} $,\;the non-critical assumption implies that the coefficient of $\be_{I_i(u)^c}$  in 	$\fil_i^{\max}(D_{\sigma})$ is non-zero.\;

The composition of the surjections in \cite[(68) \& (69)]{BDcritical25} and the inverse of the isomorphism \cite[(69)]{BDcritical25} gives:
\begin{equation}\label{mapforgisigma}
	\begin{aligned}
		g_{i,\sigma}:\homo_E\Big(\bigwedge\nolimits_E^{\!n-i}\!&D_{\sigma},\fil_i^{\max}(D_{\sigma})\Big)=\homo_E\Big(\bigwedge\nolimits_E^{\!n-i}\!D_{\sigma},\bigwedge\nolimits_E^{\!n-i}\!\fil_H^{-\bh_{i+1,\sigma}}(D_{\sigma})\Big)\\&\twoheadrightarrow \homo_E\Big(\bigwedge\nolimits_E^{\!n-i-1}\!\fil_H^{-\bh_{i+1,\sigma}}(D_{\sigma})\wedge D_{\sigma} ,\bigwedge\nolimits_E^{\!n-i}\!\fil_H^{-\bh_{i+1,\sigma}}(D_{\sigma})\Big)\\
		&\xrightarrow{\sim}\left\{f\in\homo_E\big(D_{\sigma},\fil_H^{-\bh_{i+1,\sigma}}(D_{\sigma})\big):f|_{\fil_H^{-\bh_{i+1,\sigma}}(D_{\sigma})} \text{scalar}\right\}\hookrightarrow \homo_{\fil}\big(D_{\sigma},D_{\sigma}\big),
	\end{aligned}
\end{equation}
the third inclusion is given by \cite[Lemma 2.2.5 (ii)]{BDcritical25}.\;By \cite[Lemma 2.2.5 (ii)]{BDcritical25},\;we have a surjection
\[\bigoplus_{i=1}^n\left\{f\in\homo_E\big(D_{\sigma},\fil_H^{-\bh_{i+1,\sigma}}(D_{\sigma})\big):f|_{\fil_H^{-\bh_{i+1,\sigma}}(D_{\sigma})} \text{scalar}\right\}\twoheadrightarrow \homo_{\fil}(D_{\sigma},D_{\sigma}).\]
Let $\Delta':=\{i_1<i_2<\cdots<i_{|\Delta'|}\}$ and $S:=\{i_{l_1}<i_{l_2}<\cdots<i_{l_{|S|}}\}\subseteq \Delta'$ for $1\leq l_1<\cdots<l_{|S|}\leq |\Delta'|$.\;The partial filtration $(\fil_H^{\bh_{j+1},\sigma}(D_{\sigma}),j\in S)$  on $D$ induces a natural decreasing filtration on $\bigwedge\nolimits_{E}^{\!n-i}\!D$.\;In  \cite[(97)-(99)]{BDcritical25},\;the author consider the \textit{one but last filtration} of this induced filtration:
\[\fil_{S,i}^{2^{\mathrm{nd}}-\max}(D_{\sigma})\]
Note that $\fil_i^{\max}(D_{\sigma})$ is just the last step of this induced filtration by \cite[Remark 2.3.6,\;(ii)]{BDcritical25}.\;See \cite[(97)-(99)]{BDcritical25} for the following definitions of $\fil_{S,i}^{2^{\mathrm{nd}}-\max}(D_{\sigma})$:
\begin{equation}
	\begin{aligned}
		&\fil_{S,i_{l_{|S|}}}^{2^{\mathrm{nd}}-\max}(D_{\sigma})=\Big(\bigwedge\nolimits_E^{\!n-i_{l_{|S|}}-1}\!\fil_H^{-\bh_{i_{l_{|S|}}+1,\sigma}}(D_{\sigma})\Big)\wedge \fil_H^{-\bh_{i_{l_{|S|-1}}+1,\sigma}}(D_{\sigma})\\
		&\fil_{S,i_{l_{j}}}^{2^{\mathrm{nd}}-\max}(D_{\sigma})=\Big(\bigwedge\nolimits_E^{\!n-i_{l_{j+1}}}\!\fil_H^{-\bh_{i_{l_{j+1}}+1,\sigma}}(D_{\sigma})\Big)\wedge\Big(\bigwedge\nolimits_E^{\!i_{l_{j+1}}-i_{l_{j}}-1}\!\fil_H^{-\bh_{i_{l_{j}}+1,\sigma}}(D_{\sigma})\Big)\wedge \fil_H^{-\bh_{i_{l_{j-1}}+1,\sigma}}(D_{\sigma})
	\end{aligned}
\end{equation}
 for $j\in\{1,\cdots,|S|-1\}$,\;In particular,\;if $S=\{i\}$,\;then $\fil_{\{i\},i}^{2^{\mathrm{nd}}-\max}(D_{\sigma})=\bigwedge\nolimits_E^{\!n-i-1}\!\fil_H^{-\bh_{i+1,\sigma}}(D_{\sigma})\wedge D_{\sigma}$ (see the second row of (\ref{mapforgisigma})).\;



For $i\in \Delta'$,\;define 
\begin{equation}
	\begin{aligned}
		D_{\sigma,(i)}^{\flat}:=\oplus_{u\in \cI_{i}^{\flat}}\oplus_{l_i(u)+1\leq q\leq s}D_{\sigma}^{u^{-1}(q)}=\oplus_{u\in \cI_{i}^{\flat}}\oplus_{j\in I_i(u)^c}D_{\sigma}^{j}=D_{\sigma}[\cS_i]
	\end{aligned}
\end{equation}
with $\cS_i:=\cup_{u\in \cI_{i}^{\flat}}I_i(u)^c\subseteq\{1,\cdots,s\}$.\;Put 
\[D_{\sigma}^{\flat,(i)}:=D_{\sigma}[\{1,\cdots,s\}\backslash\cS_i]=\oplus_{q\in \{1,\cdots,s\}\backslash\cS_i}D_{\sigma}^q.\]
Thus $D_{\sigma}=D_{\sigma}^{\flat,(i)}\oplus D^{\flat}_{\sigma,(i)}$.\;By definition,\;we have inclusions
\[\big(\bigwedge\nolimits_{E}^{\!n-i}\!D_{\sigma}\big)^{\flat}\hookrightarrow\bigwedge\nolimits_E^{\!n-i}\!D_{\sigma,(i)}^{\flat}\hookrightarrow \bigwedge\nolimits_E^{\!n-i}\!D_{\sigma}.\]

%Let $g_{i,\sigma}^{\flat}$ be the restriction of $g_{i,\sigma}$ on $\homo_E\Big(\big(\bigwedge\nolimits_{E}^{\!n-i}\!D_{\sigma}\big)^{\flat},\fil_i^{\max}(D_{\sigma})^{\flat}\Big)$.\;

In the sequel,\;let $\fil_{S,i}^{2^{\mathrm{nd}}-\max}(D)^{\flat}$  (resp.,\;$\fil_i^{\max}(D_{\sigma})^{\flat}$)
be the image of $\fil_{S,i}^{2^{\mathrm{nd}}-\max}(D)$ (resp.,\;$\fil_i^{\max}(D_{\sigma})$) in $ \bigwedge_E^{n-i}D_{\sigma}\twoheadrightarrow \bigwedge_E^{n-i}D_{\sigma}/\big(\bigwedge_{E}^{\!n-i}\!D_{\sigma}\big)^{\flat,c}\cong \big(\bigwedge_{E}^{\!n-i}\!D_{\sigma}\big)^{\flat}$.\;

\begin{lem}\label{lemforwedge} 
	The image under $g_{i,\sigma}$ of $\homo_E\Big(\big(\bigwedge_{E}^{\!n-i}\!D_{\sigma}\big)^{\flat},\bigwedge\nolimits_E^{\!n-i}\!\fil_H^{-\bh_{i+1,\sigma}}(D_{\sigma})\Big)$ is 
	\begin{equation}\label{imageofgflat}
		\begin{aligned}
			&\;\left\{f\in\homo_E\big(D_{\sigma}/D_{\sigma}^{\flat,(i)},\fil_H^{-\bh_{i+1}}(D_{\sigma})\big):f|_{\fil_H^{-\bh_{i+1}}(D_{\sigma})} \text{\;scalar\;}\right\} \Big(\hookrightarrow \homo_{\fil}\big(D_{\sigma,(i)}^{\flat},D_{\sigma}\big)\Big).
		%	\xrightarrow[\sim]{(\ref{comparefilwithflat})}&\;\left\{f\in\homo_E\big(D_{\sigma,(i)}^{\flat},\fil_{\mathrm{quot}}^{-\bh_{i+1}}(D_{\sigma})\big):f|_{\fil_{\mathrm{quot}}^{-\bh_{i+1}}(D_{\sigma,(i)}^{\flat})} \text{\;scalar\;}\right\}\hookrightarrow \homo_{\fil}\big(D_{\sigma,(i)}^{\flat},D_{\sigma,(i)}^{\flat}\big).\;
		\end{aligned}
	\end{equation}
Moreover,\;the kernel under $g_{i,\sigma}$ of $\homo_E\Big(\big(\bigwedge_{E}^{\!n-i}\!D_{\sigma}\big)^{\flat},\bigwedge\nolimits_E^{\!n-i}\!\fil_H^{-\bh_{i+1,\sigma}}(D_{\sigma})\Big)$ is 
\begin{equation}
	\begin{aligned}
		&\;\homo_E\Big(\big(\bigwedge\nolimits_{E}^{\!n-i}\!D_{\sigma}\big)^{\flat}/\fil_{\{i\},i}^{2^{\mathrm{nd}}-\max}(D_{\sigma})^{\flat},\fil_i^{\max}(D_{\sigma})^{\flat}\Big)
	%	\xrightarrow[\sim]{(\ref{comparefilwithflat})}&\;\homo_E\Big(\big(\bigwedge\nolimits_{E}^{\!n-i}\!D_{\sigma,(i)}^{\flat}\big)^{\flat}/\fil_{\{i\},i}^{2^{\mathrm{nd}}-\max}(D_{\sigma,(i)}^{\flat})^{\flat},\fil_i^{\max}(D_{\sigma,(i)}^{\flat})^{\flat}\Big).\;
	\end{aligned}
\end{equation}
\end{lem}
\begin{proof}For each $e_t\in D_{\sigma}^{\flat,(i)}$,\;$x\wedge e_t\notin \big(\bigwedge\nolimits_{E}^{\!n-i}\!D_{\sigma}\big)^{\flat}$ for any $x\in \bigwedge\nolimits_E^{\!n-i-1}\!\fil_H^{-\bh_{i+1}}(D_{\sigma})$,\;then $f$ vanishes on $D_{\sigma}^{\flat,(i)}$ by definition.\;When $f$ belongs to the space in (\ref{imageofgflat}),\;there exists $g\in \homo_E\big(\bigwedge\nolimits_E^{\!n-i}\!D_{\sigma},\bigwedge_E^{n-i}\fil_H^{-\bh_{i+1}}(D_{\sigma})\big)$ such that $g_{i,\sigma}(g)=f$,\;then   $g_{i,\sigma}\big(g|_{\big(\bigwedge\nolimits_{E}^{\!n-i}\!D_{\sigma}\big)^{\flat}}\big)=f$.\;This proves the surjectivity.\;By \cite[Remark 2.3.6,\;(ii)]{BDcritical25},\;we have
\[\ker(g_{i,\sigma})= \homo_E\Big(\big(\bigwedge\nolimits_{E}^{\!n-i}\!D_{\sigma}\big)/\fil_{\{i\},i}^{2^{\mathrm{nd}}-\max}(D_{\sigma}),\fil_i^{\max}(D_{\sigma})\Big).\]
This deduces the second assertion.\;
\end{proof}	

\begin{lem}\label{hodgegapforquofil}
	Consider the induced subspace filtration $\fil_H^{\bullet}(D_{\sigma}^{\flat,(i)})$ of $\fil_H^{\bullet}(D_{\sigma})$ on $D_{\sigma}^{\flat,(i)}$ and the quotient filtration $\fil_{\mathrm{quot}}^{\bullet}(D^{\flat}_{\sigma,(i)})$ on $D^{\flat}_{\sigma,(i)}\cong D_{\sigma}/D_{\sigma}^{\flat,(i)}$.\;Let  $\bd_i:=\dim_ED_{\sigma}^{\flat,(i)}\leq i$.\;Then
	\begin{equation}
	\begin{aligned}
		\gr_H^{-\bh_{j,\sigma}}(D_{\sigma}^{\flat,(i)})\neq 0,\quad \big(\text{resp.,\;}\gr_{\mathrm{quot}}^{-\bh_{j,\sigma}}(D^{\flat}_{\sigma,(i)})\big)
	\end{aligned}
	\end{equation}
	if and only if $1\leq j\leq \bd_i$ (resp.,\;$\bd_i+1\leq j\leq n$).\;Moreover,\;we have
		\begin{equation}\label{comparefilwithflat}
		\begin{aligned}
			\fil_H^{-\bh_{j,\sigma}}(D_{\sigma}^{\flat,(i)})/\cong \fil_H^{-\bh_{j,\sigma}}(D_{\sigma})/\fil_H^{-\bh_{\bd_i+1,\sigma}}(D_{\sigma})\quad (\text{resp.,\;}\fil_{\mathrm{quot}}^{-\bh_{j,\sigma}}(D^{\flat}_{\sigma,(i)})\cong\fil_H^{-\bh_{j,\sigma}}(D_{\sigma})).
		\end{aligned}
	\end{equation}
%	when	$1\leq j\leq d_{(i)}$ (resp.,\;$d_{(i)}+1\leq j\leq n$).\;
%	In particular,\;we have a strict exact sequence of filtered vector spaces:
%	\[0\rightarrow \fil_H^{\bullet}(D_{\sigma}^{\flat,(i)})\rightarrow \fil_H^{\bullet}(D_{\sigma})\rightarrow \fil_{\mathrm{quot}}^{\bullet}(D^{\flat}_{\sigma,(i)})\rightarrow 0,\]
%	such that 
\end{lem}
\begin{proof}
	These follow from the non-critical assumption in an adapted basis of $D_{\sigma}=D^{\flat}_{\sigma,(i)}\oplus D_{\sigma}^{\flat,(i)} $.\;
\end{proof}


\begin{lem}\label{filforhomospace}
For $i\in \Delta'$,\;we have 
	\begin{equation*}
	\begin{aligned}
&\;\left\{f\in\homo_E\big(D_{\sigma}/D_{\sigma}^{\flat,(i)},\fil_H^{-\bh_{i+1,\sigma}}(D_{\sigma})\big):f|_{\fil_H^{-\bh_{i+1,\sigma}}(D_{\sigma})} \text{\;scalar\;}\right\}\\\cong&\; E\oplus \homo_E\Big(\fil_H^{-\bh_{\bd_i+1,\sigma}}(D_{\sigma})/\fil_H^{-\bh_{i+1,\sigma}}(D_{\sigma}),\fil_H^{-\bh_{i+1,\sigma}}(D_{\sigma})\Big)\\\cong&\;E\oplus\bigoplus_{\substack{\bd_i+1\leq a\leq i\\i+1\leq b\leq n}}\homo_E(\gr_H^{-\bh_{a,\sigma}}(D_{\sigma}),\gr_H^{-\bh_{b,\sigma}}(D_{\sigma})).\;
		%	\xrightarrow[\sim]{(\ref{comparefilwithflat})}&\;\left\{f\in\homo_E\big(D_{\sigma,(i)}^{\flat},\fil_{\mathrm{quot}}^{-\bh_{i+1}}(D_{\sigma})\big):f|_{\fil_{\mathrm{quot}}^{-\bh_{i+1}}(D_{\sigma,(i)}^{\flat})} \text{\;scalar\;}\right\}\hookrightarrow \homo_{\fil}\big(D_{\sigma,(i)}^{\flat},D_{\sigma,(i)}^{\flat}\big).\;
	\end{aligned}
\end{equation*}
\end{lem}
\begin{proof}We have $\fil_H^{-\bh_{i+1,\sigma}}(D_{\sigma})=\oplus_{i+1\leq b\leq n}\gr_H^{-\bh_{b,\sigma}}(D_{\sigma})$ and 
\begin{equation*}
\begin{aligned}
	&\;D_{\sigma}/(D_{\sigma}^{\flat,(i)}+\fil_H^{-\bh_{i+1,\sigma}}(D_{\sigma}))\\
	=&\;\bigoplus_{1\leq a\leq i}\frac{\fil_H^{-\bh_{a,\sigma}}(D_{\sigma})+D_{\sigma}^{\flat,(i)}}{\fil_H^{-\bh_{a+1,\sigma}}(D_{\sigma})+D_{\sigma}^{\flat,(i)}}=\bigoplus_{1\leq a\leq i}\frac{\fil_H^{-\bh_{a,\sigma}}(D_{\sigma})}{\fil_H^{-\bh_{a+1,\sigma}}(D_{\sigma})+D_{\sigma}^{\flat,(i)}\cap \fil_H^{-\bh_{a,\sigma}}(D_{\sigma})},\\
	=&\;\bigoplus_{\bd_i+1\leq a\leq i}\gr_H^{-\bh_{a,\sigma}}(D_{\sigma})=\fil_H^{-\bh_{\bd_i+1,\sigma}}(D_{\sigma})/\fil_H^{-\bh_{i+1,\sigma}}(D_{\sigma}).\;
\end{aligned}
\end{equation*}
where the third identity follows from ${\fil_H^{-\bh_{a,\sigma}}(D_{\sigma})=\fil_H^{-\bh_{a+1,\sigma}}(D_{\sigma})+D_{\sigma}^{\flat,(i)}\cap \fil_H^{-\bh_{a,\sigma}}}(D_{\sigma})$ when $1\leq a\leq \bd_i$ (by Lemma \ref{hodgegapforquofil} and the non-critical assumption).\;
\end{proof}
 For any $1\leq a,b\leq n$,\;put
 \begin{equation}
 	\begin{aligned}
 		&\cS_{a,b}(S):=\{i\in S:\bd_i+1\leq a\leq i\leq b-1\},\\
 	%	&\cS^0_{a,b}(S):=\{i\in \cS_{a,b}(S):S\cap \{a,a+1,\cdots,b-1\}=\{i\},d_{(i)}+1\leq a\}
 	\end{aligned}
 \end{equation}
 Let
 \begin{equation}
	\begin{aligned}
		\cJ_i:=&\;\{(a,b):\;1\leq a,b\leq n,\;\cS_{a,b}(S)=\{i\}\},\\
		\cJ^1_i:=&\;\{(a,b):\;1\leq a,b\leq n,\;\cS_{a,b}(S)=\{i\}=S\cap \{a,a+1,\cdots,b-1\}\}.\;
	\end{aligned}
\end{equation}

\begin{lem}\label{lemforcardone} Let 
\begin{equation*}
	\cJ^{2}_i:=\cJ_i\backslash \cJ^1_i=\{(a,b):\;1\leq a,b\leq n,\;\cS_{a,b}(S)=\{i\},\;|S\cap \{a,a+1,\cdots,b-1\}|>1\}.\;
\end{equation*}
Then $\cJ^{2}_i=\emptyset$ and hence
	\[\cJ_{l_j}=\cJ^1_{l_j}=\{(a,b):\;\max\{\bd_i,l_{j-1}\}+1\leq a\leq l_j, l_j+1\leq b\leq l_{j+1}\}\]
	for $1\leq j\leq |S|$.\;
\end{lem}
\begin{proof}
	Suppose that $S\cap \{a,a+1,\cdots,b-1\}=\{i_{l_q},\cdots,i_{l_r}\}$ such that $q\leq j<r$.\;Therefore,\;there exists $j'$ such that  $q\leq j'\leq r$ and $j'\neq j$,\;and then $a\leq i_{l_{j'}}\leq \bd_i<\bd_i+1\leq i_{l_j}$,\;this is impossible.\;
\end{proof}








\subsubsection{Construction of map $t^{\flat}_{D_{\sigma}}$}



%For $1\leq i\leq s$,\;we put $I_i:=\{q:e_q=e_{i,j,\sigma},1\leq j\leq r_i\}$.\;

Similar to \cite[(46)]{BDcritical25},\;we consider the morphisms of $E$-vector spaces:
\begin{equation}\label{compositiondfnpisi}
	\begin{aligned}
		&\ext^1_{G,\sigma}\Big(\bigoplus_{u\in \cI_{i}^{\flat}}C_{i,u,\sigma},\pi_{\lalg}(\underline{x},\bh)\Big)\otimes \ext^1_{G,\sigma}\Big(\bigoplus_{u\in \cI_{i}^{\flat}}C_{i,u,\sigma},\pi_{\lalg}(\underline{x},\bh)\Big)^{\vee}\\
		\xrightarrow{\sim}&\;\ext^1_{G,\sigma}\Big(\Big(\bigoplus_{u\in \cI_{i}^{\flat}}C_{i,u,\sigma}\Big)\otimes\ext^1_{G,\sigma}\Big(\bigoplus_{u\in \cI_{i}^{\flat}}C_{i,u,\sigma},\pi_{\lalg}(\underline{x},\bh)\Big) ,\pi_{\lalg}(\underline{x},\bh)\Big)\\
		\xrightarrow{\sim}&\;\ext^1_{G,\sigma}\Big(\Big(\bigoplus_{u\in \cI_{i}^{\flat}}C_{i,u,\sigma}\Big)\otimes_E\fil_i^{\max}(D_{\sigma})^{\flat},\pi_{\lalg}(\underline{x},\bh)\Big),
	\end{aligned}
\end{equation}
where the second morphism is the push-forward induced by the composition
\[\fil_i^{\max}(D_{\sigma})^{\flat}\hookrightarrow \big(\bigwedge\nolimits_{E}^{\!n-i}\!D_{\sigma}\big)^{\flat}\xrightarrow{\epsilon_i^{-1}}\ext^1_{G,\sigma}\Big(\bigoplus_{u\in \cI_{i}^{\flat}}C_{i,u,\sigma},\pi_{\lalg}(\underline{x},\bh)\Big).\] 
Let $\pi^{\flat}_{s_i,\sigma}(\underline{x},\bh)$ (resp.,\;$\pi^{\flat}_{s_i,u,\sigma}(\underline{x},\bh)$) be a representative of the image for the canonical vector of the first term of (\ref{compositiondfnpisi}) by the composition (\ref{compositiondfnpisi}),\;which lies in the extension group 
\[\ext^1_{G,\sigma}\Big(\Big(\bigoplus_{u\in \cI_{i}^{\flat}}C_{i,u,\sigma}\Big)\otimes_E\fil_i^{\max}(D_{\sigma})^{\flat},\pi_{\lalg}(\underline{x},\bh)\Big),\;\text{resp.,\;}\ext^1_{G,\sigma}\Big(C_{i,u,\sigma}\otimes_E\fil_i^{\max}(D_{\sigma})^{\flat},\pi_{\lalg}(\underline{x},\bh)\Big).\]
We have injections $\pi_{\lalg}(\underline{x},\bh)\hookrightarrow \pi^{\flat}_{s_i,u,\sigma}(\underline{x},\bh) \hookrightarrow \pi^{\flat}_{s_i,\sigma}(\underline{x},\bh)$ of $G$-representations and a canonical isomorphism
\[\bigoplus^{u\in \cI_{i}^{\flat}}_{\pi_{\lalg}(\underline{x},\bh)}\pi^{\flat}_{s_i,u,\sigma}(\underline{x},\bh)\xrightarrow{\sim }\pi^{\flat}_{s_i,\sigma}(\underline{x},\bh).\;\]

\begin{pro}\label{construext1}Denote
	\begin{equation}
		\homo^{\flat,\dagger}_{\fil}(D_{\sigma},D_{\sigma}):=\sum_{i\in \Delta'}g_{i,\sigma}\Big(\homo_E\Big(\big(\bigwedge\nolimits_{E}^{\!n-i}\!D_{\sigma}\big)^{\flat},\fil_i^{\max}(D_{\sigma})^{\flat}\Big)\Big)\hookrightarrow \homo_{\fil}(D_{\sigma},D_{\sigma}).
	\end{equation}
There is a surjection of finite dimensional $E$-vector spaces which only depends on the $(\epsilon_{u,i})_{i\in \Delta',u\in \cI_{i}^{\flat}}$ and  on a choice of $\log_p(p)\in E$:
\[t^{\flat}_{D_{\sigma}}:\ext^1_{G,\sigma}\big(\pi_{\lalg}(\underline{x},\bh),\pi^{\flat}_{1,\sigma}(\underline{x},\bh)\big)\rightarrow\ext^{1,\flat}_{\gr}({\Dpik},{\Dpik})\oplus\homo^{\flat,\dagger}_{\fil}(D_{\sigma},D_{\sigma}).\]
Taking sum over $\sigma\in \Sigma_L$ leads to a surjection:
\[t^{\flat}_{D}:\ext^1_{G}\big(\pi_{\lalg}(\underline{x},\bh),\pi^{\flat}_{1}(\underline{x},\bh)\big)\rightarrow\ext^{1,\flat}_{\gr}({\Dpik},{\Dpik})\oplus\homo^{\flat,\dagger}_{\fil}(D,D),\]
where $\homo^{\flat,\dagger}_{\fil}(D,D):=\oplus_{\sigma\in \Sigma_L}\homo^{\flat,\dagger}_{\fil}(D_{\sigma},D_{\sigma})$.\;
\end{pro}
\begin{proof}We define $t^{\flat}_{D_{\sigma}}$ in restriction to $\ext^1_{\sigma}\big(\pi_{\lalg}(\underline{x},\bh),\pi_{\lalg}(\underline{x},\bh)\big)$ firstly.\;This is essentially the Step 2 in the proof of \cite[Proposition 2.2.4]{BDcritical25}.\;It suffices to use \cite[(63)]{BDcritical25} and replace \cite[(64) \& (65)]{BDcritical25} with 
\begin{equation}
	\begin{aligned}
		t^{\flat}_{D_{\sigma}}|_{\homo_{\mathrm{sm}}(\bZ_{S_0}(L),E)}&\xrightarrow{\sim }\ext^{1,\flat}_{\gr}({\Dpik},{\Dpik})= \prod_{i=1}^s\ext^{1}_g(E_i[1/t],E_i[1/t])\\
		(\psi_1,\cdots,\psi_s)&\mapsto \Big(E_{i}[1/t]\otimes_{\cR_{E,L}}\cR_{E[\epsilon]/\epsilon^2,L}(1+\psi_{i}\epsilon)\Big)_{1\leq i\leq s}.
	\end{aligned}
\end{equation}
The Step 3 in the proof of \cite[Proposition 2.2.4]{BDcritical25} need more modification.\;Choose $i\in \Delta'$,\;wed have isomorphisms 
	\begin{equation}
		\begin{aligned}
			\ext^1_{G,\sigma}\Big(\pi_{\lalg}(\underline{x},\bh),\pi^{\flat}_{s_i,\sigma}(\Dpik)/\pi_{\lalg}(\underline{x},\bh)\Big)&\xrightarrow{\sim}\ext^1_{G,\sigma}\Big(\pi_{\lalg}(\underline{x},\bh),\Big(\bigoplus_{u\in \cI_{i}^{\flat}}C_{i,u,\sigma}\Big)\otimes_E\fil_i^{\max}(D_{\sigma})^{\flat}\Big)\\
			&\xleftarrow{\sim}\ext^1_{G,\sigma}\Big(\pi_{\lalg}(\underline{x},\bh),\bigoplus_{u\in\cI_{i}^{\flat}}C_{i,u,\sigma}\Big)\otimes_E\fil_i^{\max}(D_{\sigma})^{\flat}\\
			&\xrightarrow[(\ref{dualpairing})]{\sim}\ext^1_{G,\sigma}\Big(\bigoplus_{u\in \cI_{i}^{\flat}}C_{i,u,\sigma},\pi_{\lalg}(\underline{x},\bh)\Big)^{\vee}\otimes_E\fil_i^{\max}(D_{\sigma})^{\flat}\\
			&\xleftarrow{\sim}\homo_E\Big(\big(\bigwedge\nolimits_{E}^{\!n-i}\!D_{\sigma}\big)^{\flat},\fil_i^{\max}(D_{\sigma})^{\flat}\Big).
		\end{aligned}
	\end{equation}
We get the desired surjection $t^{\flat}_{D_{\sigma}}$ by taking sum over  $i\in \Delta'$.\;
\end{proof}
\begin{rmk}
Similar to \cite[Proposition 2.2.7]{BDcritical25},\;we can show that $t^{\flat}_{D_{\sigma}}$ does not depends on the choices of $(\epsilon_{u,i})_{i\in \Delta',u\in \cI_{i}^{\flat}}$.\;
\end{rmk}

\subsubsection{Image of $t^{\flat}_{D_{\sigma}}$}

For  $u\in \sW_s$,\;let $u^{\sharp}\in \sW_n$ be the unique element sending $e_{i,j,\sigma}$ for $1\leq j\leq r_i$ to $e_{u^{-1}(i),j,\sigma}$ for $1\leq j\leq r_i$,\;which thus induces a permutation on $e_{1,\sigma},\cdots,e_{n,\sigma}$.\;


For $i\in \Delta'$ and $u\in \cI_{i}^{\flat}$,\;we  define $\homo_{\fil,\bF_u}^i(D_{\sigma},D_{\sigma})\subseteq \homo_{\fil,\bF_u}(D_{\sigma},D_{\sigma})$ by
\begin{equation}\label{dfnforhomofilmaxilevi}
	\begin{aligned}
		\homo_{\fil,\bF_u}^i(D_{\sigma},D_{\sigma}):=\left\{f\in \homo_{\fil,\bF_u}(D_{\sigma},D_{\sigma}):\exists\;a,b\in E\; \text{such that\;} f(e_{(u^{\sharp})^{-1}(j),\sigma})=ae_{(u^{\sharp})^{-1}(j),\sigma},\right.\\ \left. \forall\;1\leq j\leq i,f(e_{(u^{\sharp})^{-1}(j),\sigma})-be_{(u^{\sharp})^{-1}(j),\sigma}\in \oplus_{l=1}^{i}Ee_{(u^{\sharp})^{-1}(l),\sigma},\forall\;i+1\leq j\leq n\right\},
	\end{aligned}
\end{equation}
Keep the notation in Section \ref{reinterfor33}.\;Then we have (note that $\tau_{\widehat{i},\sigma}\subseteq \tau_{S_0^u,\sigma}\subseteq \fp_{S^u_0,\sigma}$,\;see Remark \ref{explainHomasLiealg})
\begin{equation}
	\begin{aligned}
		&\homo^i_{\fil,\bF_u}(D_{\sigma},D_{\sigma})\cong \mathrm{Ad}_{(u^{\sharp})^{-1}}(\tau_{\widehat{i},\sigma})\cap \mathrm{Ad}_{g_{\sigma}}(\fb_{\sigma}).\;
	\end{aligned}
\end{equation}
As in \cite[(134) \& (137)]{BDcritical25},\;sending $f$ to $(a\log_{p,\sigma},b\log_{p,\sigma})$ ($a,b$ are given in the definition (\ref{dfnforhomofilmaxilevi})) induces a canonical map 
\[f_{i,\sigma}:\homo^i_{\fil,\bF_u}(D_{\sigma},D_{\sigma})\rightarrow \homo_{\sigma}(\bL_{\widehat{i}}(\cO_L),E),\]
which is an isomorphism by comparing dimensions.\;Indeed,\;we have (see Remark \ref{explainHomasLiealg})
\begin{equation}
	\begin{aligned}
		\mathrm{Ad}_{(u^{\sharp})^{-1}}(\tau_{\widehat{i},\sigma})\cap \mathrm{Ad}_{g_{\sigma}}(\fb_{\sigma})=\mathrm{Ad}_{(u^{\sharp})^{-1}b_{\sigma}(u)}(\fz_{\widehat{i},\sigma}),\;
	\end{aligned}
\end{equation}
and thus $\dim_E\homo^i_{\fil,\bF_u}(D_{\sigma},D_{\sigma})=2$.\;Put
\begin{equation}\label{dfnforhomoflat}
	\homo^{\flat}_{\fil}(D_{\sigma},D_{\sigma}):=\sum_{i\in\Delta',u\in \cI_{i}^{\flat}}\homo^i_{\fil,\bF_u}(D_{\sigma},D_{\sigma})\subseteq\homo_{\fil}(D_{\sigma},D_{\sigma}).\;
\end{equation}

\begin{pro}\label{dfnisoforflatfilhomo}
	$\homo^{\flat}_{\fil}(D_{\sigma},D_{\sigma})\cong \mathrm{Ad}_{g_{\sigma}}(\fb_{\sigma})_{S_0}^{\flat}=\sum_{u\in \sW_s}\mathrm{Ad}_{(u^{\sharp})^{-1}}(\tau_{S^u_0,\sigma})\cap \mathrm{Ad}_{g_{\sigma}}(\fb_{\sigma})$.\;
\end{pro}
\begin{proof}
	It suffices to study the envelope $\sum_{i\in\Delta',u\in \cI_{i}^{\flat}}\mathrm{Ad}_{(u^{\sharp})^{-1}}(\tau_{\widehat{i},\sigma})\cap \mathrm{Ad}_{g_{\sigma}}(\fb_{\sigma})$ in $\mathrm{Ad}_{g_{\sigma}}(\fb_{\sigma})$.\;The left hand side equals to 
	$\sum_{u\in \sW_s}\sum_{1\leq j\leq s}\mathrm{Ad}_{(u^{\sharp})^{-1}}(\tau_{\widehat{t^u_j},\sigma})\cap \mathrm{Ad}_{g_{\sigma}}(\fb_{\sigma})$.\;By (\ref{noncriticalintersection}),\;we have $\mathrm{Ad}_{(u^{\sharp})^{-1}}(\tau_{\widehat{{t^u_j}},\sigma})\cap \mathrm{Ad}_{g_{\sigma}}(\fb_{\sigma})=\mathrm{Ad}_{(u^{\sharp})^{-1}b_{\sigma}(u)}(\fz_{\widehat{t^u_j},\sigma})$.\;Therefore,\;we deduce that 
	\begin{equation}
		\begin{aligned}
			\sum_{1\leq j\leq s}\mathrm{Ad}_{(u^{\sharp})^{-1}}(\tau_{\widehat{t^u_j},\sigma})\cap \mathrm{Ad}_{g_{\sigma}}(\fb_{\sigma})=\sum_{1\leq j\leq s}\mathrm{Ad}_{(u^{\sharp})^{-1}b_{\sigma}(u)}(\fz_{\widehat{t^u_j},\sigma})\\
			=\mathrm{Ad}_{(u^{\sharp})^{-1}b_{\sigma}(u)}(\fz_{S^u_0,\sigma})=\mathrm{Ad}_{(u^{\sharp})^{-1}}(\tau_{S_0,\sigma})\cap \mathrm{Ad}_{g_{\sigma}}(\fb_{\sigma}),
		\end{aligned}
	\end{equation}
	so that $\homo^{\flat}_{\fil}(D_{\sigma},D_{\sigma})=\mathrm{Ad}_{g_{\sigma}}(\fb_{\sigma})_{S_0}^{\flat}$.\;
\end{proof}
Consider the following composition of isomorphisms (similar to \cite[(142)-(144)]{BDcritical25}):
\begin{equation}\label{tDmapforiusigma}
	\begin{aligned}
		\ext^1_{G,\sigma}\left(\pi_{\lalg}(\underline{x},\bh),V
		_{i,u,\sigma}\right)&\xleftarrow[\text{Prop.\;}\ref{reinterforextgp}]{\sim}\homo_{\sm}(\bZ_{S^u_0}(L),E)\oplus_{\homo_{\sm}(\bL_{\widehat{i}}(L),E)}\homo_{\sigma}(\bL_{\widehat{i}}(L),E)\\&\xleftarrow{\sim}\homo_{\sm}(\bZ_{S^u_0}(L),E)\oplus\homo_{\sigma}(\bL_{\widehat{i}}(\cO_L),E)\\
		&\xrightarrow[\oplus_{\sigma}f_{i,\sigma}^{-1}]{\sim}\ext^{1,\flat}_{\gr}({\Dpik},{\Dpik})\oplus\homo^i_{\fil,\bF_u}(D_{\sigma},D_{\sigma})\\
		&\hookrightarrow \ext^{1,\flat}_{\gr}({\Dpik},{\Dpik})\oplus\homo_{\fil}(D_{\sigma},D_{\sigma}).\;
	\end{aligned}
\end{equation}
Taking sum of  (\ref{tDmapforiusigma}) over all the $i\in\Delta'$ and $u\in \cI_{i}^{\flat}$ also leads to a surjection:
\[\ext^1_{G,\sigma}\left(\pi_{\lalg}(\underline{x},\bh),\pi^{\flat}_{1,\sigma}(\underline{x},\bh)\right)\twoheadrightarrow\ext^{1,\flat}_{\gr}({\Dpik},{\Dpik})\oplus\homo^{\flat}_{\fil}(D_{\sigma},D_{\sigma}).\]
The remainder of this section will compare such surjection with $t^{\flat}_{D_{\sigma}}$ and   show that 
\[\homo^{\flat,\dagger}_{\fil}(D_{\sigma},D_{\sigma})=\homo^{\flat}_{\fil}(D_{\sigma},D_{\sigma}).\]






%Consider the intersection
%\[\Big(\bigwedge\nolimits_E^{\!n-i-1}\!\fil_H^{-\bh_{i,\sigma}}(D_{\sigma})\wedge D_{\sigma}\Big)\cap \big(\bigwedge\nolimits_{E}^{\!n-i}\!W_{\sigma}\big)^{\flat}\]	
%in $\bigwedge\nolimits_E^{\!n-i}\!D_{\sigma}$.\;It is easy to see that it is contained in the $\bigwedge\nolimits_E^{\!n-i-1}\!\fil_H^{-\bh_{i,\sigma}}(D_{\sigma})\wedge (D_{\sigma}/D_{\sigma}^{(i)})$.\;

\begin{pro}\label{comparesharpflat}
Let $i\in \Delta'$ and $u\in \cI_{i}^{\flat}$.\;The composition (given in (\ref{tDmapforiusigma}))
\begin{equation}
	\begin{aligned}
		\homo_{\sm}(\bZ_{S^u_0}(L),E)\oplus\homo_{\sigma}(\bL_{\widehat{i}}(\cO_L),E)\xrightarrow[f_{i,\sigma}^{-1}]{\sim}\ext^{1,\flat}_{\gr}({\Dpik},{\Dpik})\oplus\homo^i_{\fil,\bF_u}(D_{\sigma},D_{\sigma})
	\end{aligned}
\end{equation}
coincides with the composition
\begin{equation}
	\begin{aligned}
		\homo_{\sm}(\bZ_{S^u_0}(L),E)\oplus\homo_{\sigma}(\bL_{\widehat{i}}(\cO_L),E)&\xrightarrow{\sim}\ext^1_{G,\sigma}\left(\pi_{\lalg}(\underline{x},\bh),V_{i,u,\sigma}\right)\\&\xrightarrow{t^{\flat}_{D_{\sigma}}}\ext^{1,\flat}_{\gr}({\Dpik},{\Dpik})\oplus\homo^{\flat,\dagger}_{\fil}(D_{\sigma},D_{\sigma}).\;
	\end{aligned}
\end{equation}In particular,\;we have $\homo^{\flat,\dagger}_{\fil}(D_{\sigma},D_{\sigma})\cong \homo^{\flat}_{\fil}(D_{\sigma},D_{\sigma})$.\;
\end{pro}
\begin{proof}
This is essentially \cite[Proposition 2.5.4]{BDcritical25}.\;Only need to show that the key statement of  \cite[Lemma 2.2.9 (iii)]{BDcritical25} in the last paragraph of the proof of \cite[Proposition 2.5.4]{BDcritical25} is also suitable for our setting.\;Keep the notation in the proof of \cite[Proposition 2.5.4]{BDcritical25}.\;For $c_u\in \ext^1_{G,\sigma}\left(\pi_{\lalg}(\underline{x},\bh),C_{i,u,\sigma}\right)$,\;we claim that: 
\begin{itemize}
	\item[(1)] For $t\in I_j$ and $j\in I_i(u)$,\;$t^{\flat}_{D_{\sigma}}(c_u)(e_t)=0$.\;
	\item[(2)] For $t\in I_j$ and $j\not\in I_i(u)$,\;there exists a unique $\lambda_u\in E$ such that $c_{u,i}=\lambda_{u,i}\log$ where $\log_{p,\sigma}\in \ext^1_{G,\sigma}\left(\pi_{\lalg}(\underline{x},\bh),V_{i,u,\sigma}\right)$ is the image of $\log_{p,\sigma}\in \homo_{\sigma}(\GLN_{n-i}(\cO_L),E)$ (see Remark \ref{rmkforextgroupsigma}) and 
	\[t^{\flat}_{D_{\sigma}}(c_u)(e_t)-\lambda_{u,i}e_t\in\bigoplus_{t'\in I_{q},q\in I_i(u)}Ee_{t'}\]
\end{itemize}
We drop $\sigma$ in the subscript of $e_{i,\sigma}$ for simplicity.\;The argument in \cite[Page 33]{BDcritical25} and the diagram \cite[(78)]{BDcritical25} are still true for our setting,\;when we replace $\bigwedge\nolimits_E^{\!n-i}\!D_{\sigma}$ (resp.,\;$e_{I^c}$,\;resp,\;$\fil_i^{\max}(D_{\sigma})$, resp.,\;$C(I)$) with $\big(\bigwedge\nolimits_{E}^{\!n-i}\!D_{\sigma}\big)^{\flat}$ (resp.,\;$\be_{I_i(u)^c}$,\;resp.,\;$\fil_i^{\max}(D_{\sigma})^{\flat}$,\;resp.,\;$C_{i,u,\sigma}$).\;For any $t\in I_j$ with $j\in I_i(u)^c$,\;define 
\[\fil_H^{-\bh_{i+1,\sigma}}(D_{\sigma})^{(t)}=\fil_H^{-\bh_{i+1,\sigma}}(D_{\sigma})\cap\big(\oplus_{j\neq t}e_{j}\big).\]
We also have $\fil_H^{-\bh_{i+1,\sigma}}(D_{\sigma})^{(t)}=n-i-1$ and for any non-zero $x\in \bigwedge\nolimits_E^{\!n-i-1}\!\fil_H^{-\bh_{i+1,\sigma}}(D_{\sigma})^{(t)}$,\;we have 
\[x\wedge e_t\notin \Big(\bigoplus_{J\neq I_i(u)^c}E\be_J\Big)\oplus \big(\bigwedge\nolimits_{E}^{\!n-i}\!D_{\sigma}\big)^{\flat,c},\]
which implies that $x\wedge e_t\notin \bigoplus_{J\neq I_i(u)^c}E\be_J$.\;Note that 
\[e^{\star}_{I_i(u)^c}(x\wedge e_t)-x\wedge e_t\in \Big(\bigoplus_{J\neq I_i(u)^c}E\be_J\Big)\oplus \big(\bigwedge\nolimits_{E}^{\!n-i}\!D_{\sigma}\big)^{\flat,c}.\]
By definition of $t^{\flat}_{D_{\sigma}}(\log_{p,\sigma})$,\;this deduces that $x\wedge (t^{\flat}_{D_{\sigma}}(\log_{p,\sigma})(e_t)-e_t)\in\bigoplus_{J\neq I_i(u)^c}E\be_J$,\;and thus 
\begin{equation}\label{proofforzerocoeff}
	t^{\flat}_{D_{\sigma}}(c_u)(e_t)-e_t\in\bigoplus_{t'\in I_{q},q\neq j}Ee_{t'}
\end{equation}
It remains to show that the coefficient of $e_{j'}$ in the right hand side is zero when $t'\in I_{q}$ and $q\not\in I_i(u)$.\;Indeed,\;choose the same elements $\{f_{t''}:{t''}\in I_j,j\in I_i(u)^c\}$ as in  \cite[(84)]{BDcritical25},\;which form a basis of $\fil_H^{-\bh_{i+1,\sigma}}(D_{\sigma})$.\;Suppose that the coefficient of $\be_{t'}$ in the right hand side is non-zero.\;For any $t'\neq t$,\;we also see that the coefficient of $\be_{I_{i}(u)^c}$ is zero in 
\[\big(\wedge_{t''\neq t'}f_{t''}\big)\wedge e_t\in \bigwedge\nolimits_E^{\!n-i-1}\!\fil_H^{-\bh_{i+1,\sigma}}(D_{\sigma})\wedge D_{\sigma}.\]
But $\big(\wedge_{t''\neq t'}f_{t''}\big)\wedge t^{\flat}_{D_{\sigma}}(\log_{p,\sigma})(e_t)=0$ in $\bigwedge\nolimits_E^{\!n-i}\!D_{\sigma}$ and $e_{t'}$ cannot appears in $f_{t''}$ if $t''\neq t'$,\;then by (\ref{proofforzerocoeff}),\;the coefficient of $\be_{I_{i}(u)^c}$ in $\big(\wedge_{t''\neq t'}f_{t''}\big)\wedge e_t$ must be non-zero,\;this is impossible.\;This completes the proof.\;
\end{proof} 

For $\star\in\{\emptyset,\sigma\}$,\;we  have the following composition 
\begin{equation}\label{dfnforgammaD}
	\begin{aligned}
		\gamma^{\flat}_{\Dpik,\star}:\bigoplus_{u\in\sW_s}\overline{\ext}^{1,\flat}_{\star,u}(\Dpik,\Dpik)
		&\xrightarrow[\sim]{\kappa^{\flat}_{\Dpik,\star}} \bigoplus_{u\in\sW_s}\homo_{\star}\big(\bZ_{S^u_0}(L),E\big)
		\xrightarrow{\zeta_{\Dpik,\star}} \ext^1_{G}\left(\pi_{\lalg}(\underline{x},\bh),\pi^{\flat}_{1,\star}(\underline{x},\bh)\right).\;
	\end{aligned}
\end{equation}
where $\kappa^{\flat}_{\Dpik,\star}:=\oplus_{u\in\sW_s}\kappa_{u}$ is the first isomorphism and $\zeta_{\Dpik,\star}:=\oplus_{u\in\sW_s}\zeta_{u,\star}$ (thus $\gamma^{\flat}_{\Dpik,\star}=\zeta_{\Dpik,\star}\circ \kappa^{\flat}_{\Dpik,\star}$).\;Consider the natural map:
\[\overline{g}^{\flat}_{\Dpik,\star}:\bigoplus_{u\in\sW_s}\overline{\ext}^{1,\flat}_{\star,u}(\Dpik,\Dpik)\rightarrow \overline{\ext}^{1}(\Dpik,\Dpik).\]
The same strategy as in   \cite[Proposition 2.5.5,\;Corollary 2.5.6]{BDcritical25} and Proposition \ref{comparesharpflat} gives:
\begin{pro}\label{mainthmforfullrefine} $f^{\flat}_{\Dpik,\sigma}\circ\overline{g}^{\flat}_{\Dpik,\sigma}=t^{\flat}_{D_{\sigma}}\circ\gamma^{\flat}_{\Dpik,\sigma}$.\;After taking the sum over $\sigma\in \Sigma_L$,\;$f^{\flat}_{\Dpik}\circ\overline{g}^{\flat}_{\Dpik}=t^{\flat}_{D}\circ\gamma^{\flat}_{\Dpik}$.\;
\end{pro}

Let $\pi^{\flat}_{\min}(\Dpik)_{\star}$ be  the tautological extension of $\ker(t^{\flat}_{D_{\star}})\otimes_E\pi_{\lalg}(\underline{x},\bh)$ by $\pi^{\flat}_{1,\star}(\underline{x},\bh)$.\;Then 
\[\ker(t^{\flat}_{D})=\oplus_{\sigma\in \Sigma_L}\ker(t^{\flat}_{D_{\sigma}}),\quad\pi^{\flat}_{\min}(\Dpik)=\bigoplus_{\pi_{\lalg}(\underline{x},\bh)}^{\sigma\in\Sigma_L}\pi^{\flat}_{\min}(\Dpik)_{\sigma}.\]
Similar to \cite[Corollary 2.2.13]{BDcritical25},\;we can show that  $\pi^{\flat}_{\min}(\Dpik)_{\star}$ does not depends on any choices of $(\epsilon_{u,i})_{i\in \Delta',u\in \cI_{i}^{\flat}}$ and $\log_p(p)$.\;




\subsubsection{The $p$-adic Hodge parameters determined by the kernel of  $t^{\flat}_{D_{\sigma}}$}

This section describes what information of $p$-adic Hodge parameters is determined by the kernel of  $t^{\flat}_{D_{\sigma}}$ or   $\pi^{\flat}_{\min}(\Dpik)_{\sigma}$.\;For $S\subseteq \Delta'$,\;put
\[\pi^{\flat}_{S,\sigma}(\underline{x},\bh):=\bigoplus^{i\in S}_{\pi_{\lalg}(\underline{x},\bh)}\pi^{\flat}_{s_i,\sigma}(\underline{x},\bh).\]
Let $\ext^1_{G,\mathrm{inf}}\big(\pi_{\lalg}(\underline{x},\bh),\pi^{\flat}_{S,\sigma}(\underline{x},\bh)\big)$ be the subspace of $\ext^1_{G}\big(\pi_{\lalg}(\underline{x},\bh),\pi^{\flat}_{S,\sigma}(\underline{x},\bh)\big)$ of locally $\bQ_p$-analytic extensions with an infinitesimal characters.\;Put
\[\ker(t^{\flat}_{D_{\sigma}})_S:=\ker(t^{\flat}_{D_{\sigma}}|_{\ext^1_{G,\mathrm{inf}}(\pi_{\lalg}(\underline{x},\bh),\pi^{\flat}_{S,\sigma}(\underline{x},\bh))}),\quad \ker(t^{\flat}_{D_{\sigma}}):=\ker(t^{\flat}_{D_{\sigma}})_{\Delta'}\]
for simplicity.\;Note that $\ker(t^{\flat}_{D_{\sigma}})_{S_1}\subseteq \ker(t^{\flat}_{D_{\sigma}})_{S_2}\subseteq \ker(t^{\flat}_{D_{\sigma}})$ if $S_1\subseteq S_2\subseteq \Delta'$.



Let $\Delta':=\{i_1<i_2<\cdots<i_{|\Delta'|}\}$ and $S:=\{i_{l_1}<i_{l_2}<\cdots<i_{l_{|S|}}\}$ for $1\leq l_1<\cdots<l_{|S|}\leq |\Delta'|$.\;Consider the following compositions:
\begin{equation}\label{compositionfinalSver}
	\begin{aligned}
	\beta_S:	\ext^1_{G,\sigma}\Big(\pi_{\lalg}(\underline{x},\bh),\pi^{\flat}_{S,\sigma}(\underline{x},\bh)\Big)&\twoheadrightarrow \ext^1_{G,\sigma}\Big(\pi_{\lalg}(\underline{x},\bh),\pi^{\flat}_{S,\sigma}(\underline{x},\bh)/\pi_{\lalg}(\underline{x},\bh)\Big)\\
		&\xrightarrow{\sim}\bigoplus_{i\in S}\ext^1_{G,\sigma}\Big(\pi_{\lalg}(\underline{x},\bh),\pi_{s_i}(D_{\sigma})/\pi_{\lalg}(\underline{x},\bh)\Big)\\
		&\xrightarrow{\sim}\bigoplus_{i\in S}\homo_E\Big(\big(\bigwedge\nolimits_{E}^{\!n-i}\!D_{\sigma}\big)^{\flat},\fil_i^{\max}(D_{\sigma})^{\flat}\Big).\;
	\end{aligned}
\end{equation}
%and its projection to the $i$-th component of the latter space:
%\begin{equation}\label{compositionfinalSver}
%	\begin{aligned}
%		\beta_{S,i}:\ext^1_{G,\sigma}\Big(\pi_{\lalg}(\underline{x},\bh),\pi^{\flat}_{S,\sigma}(\underline{x},\bh)\Big)\twoheadrightarrow\homo_E\Big(\big(\bigwedge\nolimits_{E}^{\!n-i}\!D_{\sigma}\big)^{\flat},\fil_i^{\max}(D_{\sigma})^{\flat}\Big).\; 
%	\end{aligned}
%\end{equation}
By Lemma \ref{lemforwedge},\;we obtain the following  surjection
\begin{equation}
	\begin{aligned}
		g_{\sigma}:\bigoplus_{i\in S}\homo_E\Big(\big(\bigwedge\nolimits_{E}^{\!n-i}\!D_{\sigma}\big)^{\flat},&\fil_i^{\max}(D_{\sigma})^{\flat}\Big)\\&\twoheadrightarrow \sum_{i\in S}\left\{f\in\homo_E\big(D_{\sigma,(i)}^{\flat},\fil_H^{-\bh_{i+1,\sigma}}(D_{\sigma})\big):f|_{\fil_H^{-\bh_{i+1,\sigma}}(D_{\sigma})} \text{scalar}\right\}.
	\end{aligned}
\end{equation}








\begin{pro}\label{propfor2ndfil}
	\begin{itemize}
		\item[(1)] The  image under (\ref{compositionfinalSver}) of the subspace $\ext^1_{G,\mathrm{inf}}\Big(\pi_{\lalg}(\underline{x},\bh),\pi^{\flat}_{S,\sigma}(\underline{x},\bh)\Big)$  is 
		\[\homo_E\Big(\big(\bigwedge\nolimits_{E}^{\!n-i}\!D_{\sigma}\big)^{\flat}/\fil_i^{\max}(D_{\sigma})^{\flat},\fil_i^{\max}(D_{\sigma})^{\flat}\Big).\]
		\item[(2)] Let $i=i_{l_j}\in S$.\;Then the  image under (\ref{compositionfinalSver}) of the subspace $\ker(t^{\flat}_{D_{\sigma}})_S$ is	the  inverse image of 
		\[\Big\{f\in\homo_E\Big(\fil_H^{-\bh_{\bd_i+1,\sigma}}/\fil_H^{-\bh_{i+1,\sigma}},\fil_H^{-\bh_{i+1,\sigma}}\Big):f\Big(\fil_H^{-\bh_{{\max\{i_{l_{j-1}},\bd_i\}}+1,\sigma}}/\fil_H^{-\bh_{i+1,\sigma}}\Big)\subseteq\fil_H^{-\bh_{i_{l_{j+1}}+1,\sigma}}\Big\}\]
		via $g_{i,\sigma}$ (we omit $(D_{\sigma})$ in $\fil_H^{-\bh_{\bullet,\sigma}}(D_{\sigma})$ for simplicity),\;which is a subspace of
		\[\homo_E\Big(\big(\bigwedge\nolimits_{E}^{\!n-i}\!D_{\sigma}\big)^{\flat}/\fil_{S,i}^{2^{\mathrm{nd}}-\max}(D_{\sigma})^{\flat},\fil_i^{\max}(D_{\sigma})^{\flat}\Big)\]
		(compare with \cite[(99) \& (102)]{BDcritical25}).\;
		%The lands in
	%	
		%	If $D_{\sigma,(i_{l_{|S|}})}^{\flat}\subseteq D_{\sigma,(i_{l_{|S|-1}})}^{\flat}\subseteq\cdots \subseteq D_{\sigma,(i_{l_{1}})}^{\flat}$,\;it is surjective.\;\bigoplus_{\substack{1\leq a<b\leq n\\i\in \cS_{a,b}(S),|\cS_{a,b}(S))|>1}}\homo_E(\gr_H^{-\bh_{a}},\gr_H^{-\bh_{b}}
	
	\end{itemize}
\end{pro}
\begin{proof}
We examine the proof of \cite[Propositions 2.3.4 \& 2.3.7]{BDcritical25}.\;We only need to take the sum for $\{I_i(u)\}_{u\in \cI_i^{\flat}}$ in  \cite[(96)]{BDcritical25},\;so that we need to replace $\fil_i^{\max}(D_{\sigma})$ with $\fil_i^{\max}(D_{\sigma})^{\flat}$,\;$(1)$ follows.\;It remains to study the image of $\ker(g_{\sigma})$ via $g_{i,\sigma}$ when $i\in S$.\;When $S_0=\emptyset$,\;then $\Delta'=\Delta$ and $D_{\sigma,(i)}^{\flat}=D_{\sigma}$ for all $i$,\;\cite[Proposition 2.3.7]{BDcritical25} prove that the subspace $\homo_E\Big(\big(\bigwedge\nolimits_{E}^{\!n-i}\!D_{\sigma}\big)^{\flat}/\fil_{S,i}^{2^{\mathrm{nd}}-\max}(D_{\sigma})^{\flat},\fil_i^{\max}(D_{\sigma})^{\flat}\Big)$ in $(2)$ gives such image.\;We draw the matrix of each $f_j\in \homo_E(D_{\sigma},D_{\sigma})$ (note that $f_j\in \homo_E(D_{\sigma,(j)}^{\flat},D_{\sigma})$ in our situation) in an adapted basis of $D_{\sigma}$ for the filtration $\fil^{-\bh_{i,\sigma}}_H(D_{\sigma})$,\;then we get the matrix form of $f_{i_{l_1}}+\cdots+f_{i_{l_{|S|}}}$.\;If $f_{i_{l_1}}+\cdots+f_{i_{l_{|S|}}}=0$,\;then $f_{j}$ also satisfies that  (see \cite[(102)]{BDcritical25})
\begin{equation*}
	\begin{aligned}
		&\Big\{f_{|S|}\in \homo_E\Big(D_{\sigma}/\big(D_{\sigma}^{\flat,(i_{l_{|S|}})}+\fil_H^{-\bh_{i_{l_{|S|}}+1,\sigma}}\big),\fil_H^{-\bh_{i_{l_{|S|}}+1,\sigma}}\Big):\;f_{|S|}\Big(\fil_H^{-\bh_{i_{l_{|S|-1}}+1,\sigma}}/\fil_H^{-\bh_{i_{l_{|S|}}+1,\sigma}}\Big)=0\Big\},\\
		&\Big\{f_j\in\homo_E\Big(D_{\sigma}/\big(D_{\sigma}^{\flat,(i_{l_{j}})}+\fil_H^{-\bh_{i_{l_{j}}+1,\sigma}}\big),\fil_H^{-\bh_{i_{l_{j}}+1,\sigma}}\Big):f_j\Big(\fil_H^{-\bh_{i_{l_{j-1}}+1,\sigma}}/\fil_H^{-\bh_{i_{l_{j}}+1,\sigma}}\Big)\subseteq\fil_H^{-\bh_{i_{l_{j+1}}+1,\sigma}}\Big\}
	\end{aligned}
\end{equation*}
when $j<|S|$.\;Then the same proof of \cite[(102)]{BDcritical25} shows that the image of $\ker(g_{\sigma})$ must lands in the space $\homo_E\Big(\big(\bigwedge\nolimits_{E}^{\!n-i}\!D_{\sigma}\big)^{\flat}/\fil_{S,i}^{2^{\mathrm{nd}}-\max}(D_{\sigma})^{\flat},\fil_i^{\max}(D_{\sigma})^{\flat}\Big)$.\;In general,\;the irredularity of the shape $\bL_{S_0}$ makes $\ker(g_{\sigma})$ smaller.\;Note that
\[\homo_E(D_{\sigma},D_{\sigma})=\bigoplus_{1\leq a,b\leq n}\homo_E\big(\gr_H^{-\bh_{a,\sigma}},\gr_H^{-\bh_{b,\sigma}}\big)\]
We use Lemma \ref{filforhomospace} and study the kernel of 
\[ \bigoplus_{i\in S}\bigoplus_{\substack{\bd_i+1\leq a\leq i\\i+1\leq b\leq n}}\homo_E\big(\gr_H^{-\bh_{a,\sigma}},\gr_H^{-\bh_{b,\sigma}}\big)\rightarrow \bigoplus_{1\leq a,b\leq n}\homo_E\big(\gr_H^{-\bh_{a,\sigma}},\gr_H^{-\bh_{b,\sigma}}\big).\]
We see that the kernel of the map 
\[g_{a,b}:\bigoplus_{i\in \cS_{a,b}(S)}\homo_E\big(\gr_H^{-\bh_{a,\sigma}},\gr_H^{-\bh_{b,\sigma}}\big)\rightarrow \bigoplus_{1\leq a,b\leq n}\homo_E\big(\gr_H^{-\bh_{a,\sigma}},\gr_H^{-\bh_{b,\sigma}}\big).\]
has $E$-dimension $|\cS_{a,b}(S)|-1$.\;Moreover,\;for any $i\in \cS_{a,b}(S)$,\;the projection of  $\ker(g_{a,b})$ to the  $i$-component $\homo_E(\gr_H^{-\bh_{a,\sigma}},\gr_H^{-\bh_{b,\sigma}})$ is surjective if only if $|\cS_{a,b}(S)|>1$.\;Therefore,\;for $i=i_{l_j}$,\;we see that $g_{i,\sigma}(\ker (g_{\sigma}))$ is equals to
\begin{equation*}
	\begin{aligned}
		&\;\bigoplus_{\substack{1\leq a<b\leq n\\i\in \cS_{a,b}(S),|\cS_{a,b}(S)|>1}}\homo_E\big(\gr_H^{-\bh_{a,\sigma}},\gr_H^{-\bh_{b,\sigma}}\big)=\bigoplus_{b=i+1}^n\bigoplus_{\substack{\bd_i+1\leq a\leq i\\ |\cS_{a,b}(S)|>1}}\homo_E\big(\gr_H^{-\bh_{a,\sigma}},\gr_H^{-\bh_{b,\sigma}}\big)\\
       		=&\;\homo_E\Big(\fil_H^{-\bh_{\bd_i+1,\sigma}}/\fil_H^{-\bh_{i+1,\sigma}},\fil_H^{-\bh_{i+1,\sigma}}\Big)\big/\bigoplus_{(a,b)\in\cJ^1_i}\homo_E\big(\gr_H^{-\bh_{a,\sigma}},\gr_H^{-\bh_{b,\sigma}}\big)\\
       	=&\;\Big\{f_i\in\homo_E\Big(\fil_H^{-\bh_{\bd_i+1,\sigma}}/\fil_H^{-\bh_{i+1,\sigma}},\fil_H^{-\bh_{i+1,\sigma}}\Big):f_i\Big(\fil_H^{-\bh_{{\max\{i_{l_{j-1}},\bd_i\}}+1,\sigma}}/\fil_H^{-\bh_{i+1,\sigma}}\Big)\subseteq\fil_H^{-\bh_{i_{l_{j+1}}+1,\sigma}}\Big\}.
	\end{aligned}
\end{equation*}
by Lemma \ref{lemforcardone}.\;We completes the proof.\;
\end{proof}	





\begin{thm}\label{Hodgeparadeter}$\ker(t^{\flat}_{D_{\sigma}})_S$ determines (recall that $S:=\{i_{l_1}<i_{l_2}<\cdots<i_{l_{|S|}}\}$)
	\begin{equation}
	\begin{aligned}
		\Big\{\fil_H^{-\bh_{{\max\{i_{l_{j}},\bd_{i_{l_t}}\}}+1,\sigma}}(D_{\sigma})\oplus D_{\sigma}/\fil_H^{-\bh_{\bd_{i_{l_t}},\sigma}}(D_{\sigma})\Big\}_{1\leq j\leq |S|}
	\end{aligned}
\end{equation}	
In particular,\;if $\bd_{i_{l_t}}=0$,\;then $\ker(t^{\flat}_{D_{\sigma}})_S$ determines $\{\fil_H^{-\bh_{i_{l_{j}}+1,\sigma}}(D_{\sigma})\}_{1\leq j\leq |S|}$.\;
\end{thm}
\begin{proof}
	Let $t=|S|$.\;For $0\leq j\leq t-1$,\;put $S_j:=\{i_{l_1},i_{l_2},\cdots,i_{l_j},i_{l_t}\}$.\;Note that $S_0=\{i_{l_t}\}$ and $S_{t-1}=S$.\;Similar to the Step 2 of the proof of \cite[2.3.10]{BDcritical25} and Proposition \ref{propfor2ndfil},\;the following data is determined by  $\ker(t^{\flat}_{D_{\sigma}})_S$:
	\begin{equation}
		\begin{aligned}
			&\Big\{\homo_E\Big(\fil_H^{-\bh_{\bd_{i_{l_t}}+1,\sigma}}/\fil_H^{-\bh_{{\max\{i_{l_{j}},\bd_{i_{l_t}}\}}+1,\sigma}},\fil_H^{-\bh_{i_{l_t}+1,\sigma}}\Big)\Big\}_{0\leq j\leq t-1}
		\end{aligned}
	\end{equation}
(note that $i_0=0$) and 
	\begin{equation}
	\begin{aligned}
		&\homo_E\Big(\fil_H^{-\bh_{\bd_{i_{l_t}}+1,\sigma}}/\fil_H^{-\bh_{i_{l_{t}}+1,\sigma}},\fil_H^{-\bh_{i_{l_t}+1,\sigma}}\Big).
	\end{aligned}
\end{equation}
By the Lemma \ref{lem:recover-subspace-relative},\;$\ker(t^{\flat}_{D_{\sigma}})_S$ determines 
	\begin{equation}
	\begin{aligned}
		\Big\{\fil_H^{-\bh_{{\max\{i_{l_{j}},\bd_{i_{l_t}}\}}+1,\sigma}}\oplus D_{\sigma}/\fil_H^{-\bh_{\bd_{i_{l_t}},\sigma}}\Big\}_{1\leq j\leq t}.
	\end{aligned}
\end{equation}
We complete the proof.\;
\end{proof}

\begin{cor}\label{Hodgeparadetercor} For $1\leq t\leq |\Delta'|$,\;$\ker(t^{\flat}_{D})$ determines  (recall that $\Delta':=\{i_1<i_2<\cdots<i_{|\Delta'|}\}$)
	\begin{equation}
		\begin{aligned}
			\Big\{\fil_H^{-\bh_{{\max\{i_{{j}},\bd_{i_t}\}}+1,\sigma}}(D_{\sigma})\oplus D_{\sigma}/\fil_H^{-\bh_{\bd_{i_t},\sigma}}(D_{\sigma})\Big\}_{1\leq j\leq t}
		\end{aligned}
	\end{equation}	
	In particular,\;if $\bd_{i_{t}}=0$,\;then $\ker(t^{\flat}_{D_{\sigma}})$ determines $\{\fil_H^{-\bh_{i_{j}+1,\sigma}}(D_{\sigma})\}_{1\leq j\leq t}$.\;
\end{cor}



\begin{lem}\label{lem:recover-subspace-relative}
	Let $H'\subseteq H\subseteq D$ be subspaces of a finite-dimensional $E$-vector space $D$, and let $F$ be a
	nonzero $E$-vector space.\;Fix a quotient map $D\twoheadrightarrow H$ and consider the following composition 
	\[q:D\twoheadrightarrow H\twoheadrightarrow H/H'.\]
	Let $W:=\{\,f\circ q:f\in\Hom_E(H/H',F)\,\}
	\subseteq\Hom_E(D,F)$.\;Then
	$	\bigcap_{\varphi\in W}\ker(\varphi)=q^{-1}(0)=H'+\ker(D\twoheadrightarrow H)$ 
	and $H'=H\cap\bigcap_{\varphi\in W}\ker(\varphi).$

\end{lem}

\begin{proof}
	Since every element of $W$ factors through $q$, we have $q^{-1}(0)\subseteq\ker(\varphi)$ for every $\varphi\in W$.\;Hence $q^{-1}(0)\subseteq
	\bigcap_{\varphi\in W}\ker(\varphi).$\;Conversely,\;let $x\notin q^{-1}(0)$.\;Then $q(x)\neq0$ in $H/H'$.\;Since $F\neq0$,\; choose a nonzero vector $y\in F$ and a linear functional $	\lambda\in(H/H')^\vee$	such that $	\lambda(q(x))\neq0$.\;Define
	\[
	f:H/H'\longrightarrow F,\qquad
	\bar h\longmapsto\lambda(\bar h)y.
	\]
	Then $\varphi:=f\circ q\in W$ and $	\varphi(x)=\lambda(q(x))y\neq0$.\;Hence $x\notin
	\bigcap_{\varphi\in W}\ker(\varphi).$.\;We completes the proof.\;
\end{proof}




%\begin{cor}\label{Hodgeparadetercor}
%	$\{\ker(t^{\flat}_{D_{\sigma}})_{\{1,\cdots,j\}}\}_{1\leq j\leq n-1}$ and thus $\ker(t^{\flat}_{D_{\sigma}})$ determines $\{\fil_{H}^{-\bh_{j+1}}(D_{\sigma})^{\flat}_{(j)}\}_{1\leq j\leq n-1}$.\;
%\end{cor}


%For $i,j\in\Delta'$ and $i\leq j$,\;let $\fil_{H}^{-\bh_{j+1}}(D_{\sigma})^{\flat}_{(j)}$
%be  the image of $\fil_{H}^{-\bh_{i+1}}(D_{\sigma})$ in $D_{\sigma}/D_{\sigma}^{\flat,(j)}=D_{\sigma,(j)}^{\flat}$.\;Thus,\;we deduce that $\ker(t^{\flat}_{D_{\sigma}})_S$ determines $\{\fil_{H}^{-\bh_{i_j+1}}(D_{\sigma})^{\flat}_{(i_t)}\}_{1\leq j\leq t}$.\;The results follow.\;


\begin{exmp}
A typical example is $\bL_{S_0}\cong \GL_1^{\oplus (n-m)}\times \GLN_m$.\;As $E$-vector spaces,\;$D_{\sigma}=D_{\sigma}^1\oplus\cdots D_{\sigma}^{n-m}\oplus D_{\sigma}^{n-m+1}$ with $\dim_ED_{\sigma}^j=1$ for $1\leq j\leq n-m$ and $\dim_ED
_{\sigma}^{n-m+1}=m$.\;For $i\in \Delta'$,\;we have
\begin{equation*}
	D_{\sigma,(i)}=\left\{
	\begin{array}{ll}
		D_{\sigma}/D_{\sigma}^{n-m+1},\;&n-i<m\\
		D_{\sigma} ,\;&n-i\geq m
	\end{array}
	\right.,\;\bd_i=\left\{
	\begin{array}{ll}
		m,\;&n-i<m\\
		0,\;&n-i\geq m
	\end{array}
	\right.,\;
\end{equation*}
We divide into two cases.\;
\begin{itemize}
	\item[(1)] $m\leq n-m$.\;In this case,\;$\Delta'=\Delta$.\;For $n-i\geq m$,\;$\ker(t^{\flat}_{D_{\sigma}})$ determines $\{\fil_H^{-\bh_{j+1,\sigma}}(D_{\sigma})\}_{1\leq j\leq i}$.\;If $n-i<m$,\;$\ker(t^{\flat}_{D_{\sigma}})$ determines 
	\begin{equation*}
		\begin{aligned}
			\Big\{\fil_H^{-\bh_{j+1,\sigma}}(D_{\sigma})\oplus D_{\sigma}/\fil_H^{-\bh_{m,\sigma}}(D_{\sigma})\Big\}_{m \leq j\leq i}\cup \Big\{\fil_H^{-\bh_{m+1,\sigma}}(D_{\sigma})\oplus D_{\sigma}/\fil_H^{-\bh_{m,\sigma}}(D_{\sigma})\Big\}.
		\end{aligned}
	\end{equation*}	
	\item[(2)] $m>n-m$.\;In this case,\;$\Delta'=\Delta\backslash \{n-m+1,\cdots,m-1\}$.\;For $n-i\geq m$ (so that $i\leq n-m$ and $i\in \Delta'$),\;$\ker(t^{\flat}_{D_{\sigma}})$ determines $\{\fil_H^{-\bh_{j+1,\sigma}}(D_{\sigma})\}_{1\leq j\leq i}$.\;If $i>m-1$,\;$\ker(t^{\flat}_{D_{\sigma}})$ determines 
	\begin{equation*}
		\begin{aligned}
			\Big\{\fil_H^{-\bh_{j+1,\sigma}}(D_{\sigma})\oplus D_{\sigma}/\fil_H^{-\bh_{m,\sigma}}(D_{\sigma})\Big\}_{m \leq j\leq i}\cup \Big\{\fil_H^{-\bh_{m+1,\sigma}}(D_{\sigma})\oplus D_{\sigma}/\fil_H^{-\bh_{m,\sigma}}(D_{\sigma})\Big\}.
		\end{aligned}
	\end{equation*}	
\end{itemize}
\end{exmp}





\subsection{Universal extensions}\label{sectionforunverext}

Let $R_{\Dpik}$ be the universal deformation ring of deformations of $\Dpik$ over $\Art_E$.\;For $?\in\{\flat,0\}$,\;let $R^{?}_{\Dpik,u}$ be the universal deformation ring that pro-represents the functor $\cF_{\Dpik,\cF_u}^{?}$ over $\Art_E$,\;and let $R_{\Dpik,g}$ be the universal deformation ring of de Rham deformations.\;Let $I^{?}_u$ (resp.,\;$I_{g}$) be the kernel of $R_{\Dpik}\rightarrow R^{?}_{\Dpik,u}$ (resp.,\;$R_{\Dpik}\rightarrow R_{\Dpik,g}$).\;Put $R^{?}_{\Dpik,u,g}=R_{\Dpik}/(I_g+I^{?}_u)$.\;We have natural surjections $R_{\Dpik}\twoheadrightarrow R^{?}_{\Dpik,u}\rightarrow  R^{?}_{\Dpik,u,g}$.\;For a continuous character $\delta$ of $\bZ_{S^u_0}(L)$,\;denote by the $R_{S_0^u,\delta}$ (resp.\;$R_{S_0^u,\delta,g}$) the universal deformations ring (resp.,\;the universal deformations ring of locally algebraic deformations) of $\delta$ over $\Art_E$.\;We have  a natural surjection $R_{S_0^u,\delta}\twoheadrightarrow R_{S_0^u,\delta,g}$.\;

For a complete local Noetherian $E$-algebra $R$,\;we use $\fm_R$ (or use $\fm$ for simplicity when it does not cause confusion) to denote its maximal ideal.\;Then for $\ast\in\{\emptyset,g,u\}$ and $?\in\{\flat,0,\emptyset\}$,\;we have  $(\fm_{R^{?}_{\Dpik,\ast}}/\fm^2_{R^{?}_{\Dpik,\ast}})^{\vee}\cong \ext^{1,?}_{\ast}(\Dpik,\Dpik)$.\;Moreover,\;we have
\[(\fm_{R_{S_0^u,\delta}}/\fm^2_{R_{S_0^u,\delta}})^{\vee}\cong \homo({\bZ_{S^u_0}(L)},E),(\fm_{R_{S_0^u,\delta,g}}/\fm^2_{R_{S_0^u,\delta,g}})^{\vee}\cong \homo_{\mathrm{sm}}({\bZ_{S^u_0}(L)},E).\]
For $u\in \sW_s$,\;we have the following commutative Cartesian diagram over $\Art_E$ (for $?\in\{\flat,0\}$):
\begin{equation}
	\xymatrix{
		R_{S_0^u,1}/\fm^2 \ar@{->>}[r] \ar[d]^{\kappa^{\flat}_u} & R_{S_0^u,1,g}/\fm^2 \ar[d]^{\kappa^{\flat}_u} 
		\\
		R^{\flat}_{\Dpik,u}/\fm^2 \ar@{->>}[r]	&  R^{\flat}_{\Dpik,u,g}/\fm^2},\quad \xymatrix{
		R_{S_0^u,1}/\fm^2 \ar@{->>}[r] \ar[d]^{\kappa^{0}_u} & R_{S_0^u,1,g}/\fm^2 \ar[d]^{\kappa^{0}_u} 
		\\
		R^{0}_{\Dpik,u}/\fm^2 \ar@{->>}[r]	&  R_{\Dpik,g}/\fm^2}.
\end{equation}
Let $\FZ_{\Omega}$ be the Bernstein centre over $E$ associated to $\mathrm{PS}^{\infty}_{1}(\underline{x})$ of $G$.\;Recall that $\FZ_{\Omega_{S^u_0}}$ is the Bernstein centre over $E$ associated to $\pi_{\sm}(\underline{x}^u)$ of $\bL_{S_0^u}(L)$.\;Let $\widehat{\FZ}_{\Omega,\underline{x}}$ (resp.,\;$\widehat{\FZ}_{\Omega_{S_0^u},\underline{x}^u}$) be the completion of $\FZ_{\Omega}$ (resp.,\;$\FZ_{\Omega_{S^u_0}}$) at $\mathrm{PS}^{\infty}_{u}(\underline{x})$ (resp.,\;$\pi_{\sm}(\underline{x}^u)$).\;We have natural isomorphisms $\widehat{\FZ}_{\Omega,\underline{x}}\xrightarrow{\sim} \widehat{\FZ}_{\Omega_{S_0^u},\underline{x}^u}\xrightarrow{\sim} R_{S_0^u,1,g}$ (note that $\widehat{\FZ}_{\Omega,\underline{x}}\cong R_{S_0^u,1,g}\cong R_{S_0,1,g}$ by sending a smooth deformation $\chi$ of $1$ to $\big(\ind^G_{\op_{S^u_0}(L)}\pi_{\sm}(\underline{x}^u)\eta_{S^u_0}\chi\big)^{\infty}$,\;see Lemma \ref{smhomoindepent}).\;Moreover,\;by Lemma \ref{smhomoindepent},\;we see that the composition 
\[A_0:=\widehat{\FZ}_{\Omega,\underline{x}}/\fm^2\xrightarrow{\sim} R_{S_0^u,1,g}\hookrightarrow R^{\flat}_{\Dpik,u,g}/\fm^2\]
is independent of the choice of $u$.\;For $?\in\{\emptyset,\flat,0\}$,\;let 
\[A^{?}_{\Dpik}:=R^{?}_{\Dpik}/\fm^2\times_{R_{\Dpik,g}} A_0,\text{\;and\;}A^{?}_{\Dpik,u}:=R^{?}_{\Dpik,u}/\fm^2\times_{R^{\flat}_{\Dpik,u,g}} A_{0}.\;\]
Then the tangent space of $A^{?}_{\Dpik}$ (resp.,\;$A^{?}_{\Dpik,u}$) is naturally isomorphism to $\overline{\ext}^{1,?}(\Dpik,\Dpik)$ (resp.,\;$\overline{\ext}^{1,?}_u(\Dpik,\Dpik)$), and the tangent space of $A^{?}_{\Dpik,u}$ is naturally isomorphic to $\homo\big(\bZ_{S^u_0}(L),E\big)$.\;Let $\cI^{?}_{u}$ be the kernel of the natural morphism $A_{\Dpik}\rightarrow A^{?}_{\Dpik,u}$.\;Let $\cI^{?}_{\Dpik}$ be the kernel of the natural morphism $A_{\Dpik}\rightarrow \prod_{u\in \sW_s}A^{?}_{\Dpik,u}$,\;then we get an injection $A^{?}_{\Dpik}:=A_{\Dpik}/\cI^{?}_{\Dpik}\hookrightarrow\prod_{u\in \sW_s}A^{?}_{\Dpik,u}$ (In particular,\;the infinite fern in Theorem \ref{infinitePoten}   deduces that $I_{\Dpik}=0$).\;

Let $\pi^{\flat}_{1}(\underline{x},\bh)^{\univ}$ (resp.,\;$\pi^{\flat}_{1}(\underline{x},\bh)_u^{\univ}$) be the universal tautological extension of
\[\ext^1_{G}\big(\pi_{\lalg}(\underline{x},\bh),\pi^{\flat}_1(\underline{x},\bh)\big)\otimes_E \pi_{\lalg}(\underline{x},\bh),\;\quad(\text{resp.,\;}\ext^1_{G}\left(\pi_{\lalg}(\underline{x},\bh),\mathrm{PS}_{u,1}(\underline{x},\bh)\right)\otimes_E \pi_{\lalg}(\underline{x},\bh))\]
 by $\pi^{\flat}_{1}(\underline{x},\bh)$.\;Let $\widetilde{\delta}_u^{\univ}$ be the universal extension of $\homo(\bZ_{S^u_0}(L),E)\otimes_E\delta_u$ by $\delta_u$.\;Similar to the proof of \cite[Lemma 3.35]{ParaDing2024},\;the induced representation $I^G_{\op_{S^u_0}(L)}\pi_{\sm}(\underline{x}^u)\eta_{S^u_0}\widetilde{\delta}_u^{\univ}$  is the universal extension of $\homo(\bZ_{S^u_0}(L),E)\otimes \pi_{\lalg}(\underline{x},\bh)$ by $\mathrm{PS}_{u,1}(\underline{x},\bh)$.\;We have hence  isomorphisms of $G$-representations
\begin{equation}
	\begin{aligned}
		&\left(I^G_{\op_{S_0}(L)}\pi_{\sm}(\underline{x}^u)\eta_{S^u_0}\widetilde{\delta}_u^{\univ}\right)\oplus_{\mathrm{PS}_{u,1}(\underline{x},\bh)}\pi^{\flat}_{1}(\underline{x},\bh)\xrightarrow{\sim }\pi^{\flat}_{1}(\underline{x},\bh)_u^{\univ}.\;
	\end{aligned}
\end{equation}
For $?\in\{\flat,0\}$,\;we have a natural action of $A^{?}_{\Dpik,u}\hookrightarrow R_{S_0^u,1}/\fm^2$ on $\widetilde{\delta}_u^{\univ}$ where $y\in \fm_{A^{?}_{\Dpik,u}}/\fm^2_{A^{?}_{\Dpik,u}}\cong \homo\big(\bZ_{S^u_0}(L),E\big)^{\vee}$ acts via
$y:\widetilde{\delta}_u^{\univ}\twoheadrightarrow \homo\big(\bZ_{S^u_0}(L),E\big)\otimes_E\delta_u\xrightarrow{y}\delta_u\hookrightarrow \widetilde{\delta}_u^{\univ}$.\;Similarity,\;$\pi^{\flat}_{1}(\underline{x},\bh)_u^{\univ}$ admits an $A^{?}_{\Dpik,u}$-action,\;which is given by 
\[y:\pi^{\flat}_{1}(\underline{x},\bh)_u^{\univ}\rightarrow \ext^1_{u}\left(\pi_{\lalg}(\underline{x},\bh),\pi^{\flat}_{1}(\underline{x},\bh)\right)\otimes \pi_{\lalg}(\underline{x},\bh)\xrightarrow{y}\pi_{\lalg}(\underline{x},\bh)\hookrightarrow \pi^{\flat}_{1}(\underline{x},\bh)_u^{\univ},\]
where $y\in \fm_{A^{?}_{\Dpik,u}}/\fm^2_{A^{?}_{\Dpik,u}}\cong \homo\big(\bZ_{S^u_0}(L),E\big)^{\vee}\xrightarrow{\zeta_u}\big(\zeta_u\big(\homo\big(\bZ_{S^u_0}(L),E\big)\big)^{\vee}$.\;The injection
\[I^G_{\op_{S^u_0}(L)}\pi_{\sm}(\underline{x}^u)\eta_{S^u_0}\widetilde{\delta}_u^{\univ}\hookrightarrow \pi^{\flat}_{1}(\underline{x},\bh)_u^{\univ}\] is $A^{?}_{\Dpik,u}$-equivalent.\;

\begin{thm}\label{actionADonuniversal}
For $?\in\{\emptyset,\flat,0\}$,\;there is a unique $A^{?}_{\Dpik}$-action (the $A_{\Dpik}$-action factors through the natural surjections $A_{\Dpik}\twoheadrightarrow  A^0_{\Dpik}\twoheadrightarrow  A^{\flat}_{\Dpik}$) on $\pi^{\flat}_{1}(\underline{x},\bh)^{\univ}$ such that for all $u\in \sW_s$,\;we have an $A^{?}_{\Dpik,u}\times G$-equivariant injection
	\[\pi^{\flat}_{1}(\underline{x},\bh)_u^{\univ}\hookrightarrow \pi^{\flat}_{1}(\underline{x},\bh)^{\univ}[\cI^{?
	}_u].\]
\end{thm}
\begin{proof}
	By Theorem \ref{mainthmforfullrefine},\;we define an $A_{\Dpik}$-action on $\pi^{\flat}_{1}(\underline{x},\bh)^{\univ}$ by letting $y\in \fm_{A_{\Dpik}}/\fm^2_{A_{\Dpik}}\cong \overline{\ext}^1(\Dpik,\Dpik)\rightarrow \ext^1_{G}\left(\pi_{\lalg}(\underline{x},\bh),\pi^{\flat}_1(\underline{x},\bh)\right)^{\vee}$ act via
	\[y:\pi^{\flat}_{1}(\underline{x},\bh)^{\univ}\rightarrow \ext^1_{G}\left(\pi_{\lalg}(\underline{x},\bh),\pi^{\flat}_1(\underline{x},\bh)\right)\otimes_E\pi_{\lalg}(\underline{x},\bh)\xrightarrow{y} \pi_{\lalg}(\underline{x},\bh)\hookrightarrow \pi^{\flat}_{1}(\underline{x},\bh)^{\univ}.\;\]
By definition,\;such $A_{\Dpik}$-action factors through quotients $A_{\Dpik}\twoheadrightarrow  A^0_{\Dpik}\twoheadrightarrow  A^{\flat}_{\Dpik}$.\;This action satisfies the property in the theorem.\;The uniqueness follows from the fact that $\pi^{\flat}_{1}(\underline{x},\bh)^{\univ}$ is generated by $\pi^{\flat}_{1}(\underline{x},\bh)_u^{\univ}$.\;
\end{proof}



\section{Local-global compatibility}




\subsection{Patched Bernstein eigenvarieties and Bernstein parabolic varieties}\label{BENVARPARAVAR}

In this section,\;we recall briefly  the patched Bernstein eigenvariety  and Bernstein paraboline variety  of Breuil-Ding (see \cite[Section 3.3,\;Section 4.2]{Ding2021}).\;

We use the following patched setting and follow the notation of \cite[Section 4.1.1]{2019DINGSimple} and \cite[Section 2]{PATCHING2016}.\;We have a patched Galois deformation ring $R_{\infty}=R_{\infty}^{\fp}\widehat{\otimes} \defvarring$,\;where $\fp$ is a $p$-adic place over a totally real field $F^+$,\;$L:=F^+_{\fp}$,\;and $\overline{r}: \gal_L \rightarrow \GLN_n(k_E)$ is a continuous representation and  $\defvarring$ is  the maximal reduced and $p$-torsion free quotient of the universal $\co_E$-lifting ring of $\overline{r}$.\;We have a $R_{\infty}$-admissible unitary representation $\Pi_{\infty}$ (i.e.,\;the so-called patched Banach representation) of $G$ over $E$.\;We refer to \textit{loc.cit} for detail.\;Let $\Pi_\infty^{R_\infty-\ana}$ be the subrepresentation of $G$ of locally $R_\infty$-analytic vectors of $\Pi_\infty$ (see \cite[Section 3.1]{breuil2017interpretation}).\;



Put $\FX_{\infty}:=(\Spf R_{\infty})^{\rig}$,\;$\FX_{\infty}^{{\fp}}:=(\Spf R^{\fp}_{\infty})^{\rig}$ and  $\mathfrak{X}_{\overline{r}}=(\spf\;R_{\overline{r}}^{\square})^{\rig}$.\;We have thus
$\FX_{\infty}:=\FX_{\infty}^{{\fp}}\times \mathfrak{X}_{\overline{r}}$.\;As in Section \ref{Omegafil},\;for $u\in\sW_s$,\;we fix a cuspidal Bernstein component $\Omega_{S^u_0}$ of the Levi subgroup $\bL_{S^u_0}$.\;Let $\bh:=(\hpi_{1,\sigma},\hpi_{2,\sigma},\cdots,\hpi_{n,\sigma})_{\sigma\in \Sigma_L}\in X_{\Delta}^+$.\;We put $\hpi_{i}=(\hpi_{i,\sigma})_{\sigma\in \Sigma_L}$ for $1\leq i\leq n$ and put ${\bm\lambda}_\bh=(\hpi_{i,\sigma}+i-1)_{\sigma\in \Sigma_L,1\leq i\leq n}=\bh-\theta$.\;Let 
\[\mathcal{E}_{\omepik,\fp,{\bm\lambda}_\bh}^\infty(\overline{\rho})\hookrightarrow\FX_{\infty}\times\sbanpiku\times \rigchu\] 
be the patched Bernstein eigenvariety associated to $\Pi_{\infty}$ and  $\Omega_{S_0^u}$ (see \cite[Section 4.1]{He20222},\;this is an easy variation of the patched eigenvariety introduced in  \cite{Ding2021}).\;There exists a natural coherent sheaf $\cM_{\omepik,{\bm\lambda}_\bh}^{\infty}$ over $\FX_{\infty}\times\sbanpiku\times \rigchu$ such that we have a $\bZ_{S_0^u}(L)\times R_{\infty}$-equivalent isomorphism
\[\Gamma\Big(\FX_{\infty}\times\sbanpiku\times \rigchu,\cM_{\omepik,{\bm\lambda}_\bh}^{\infty}\Big)\cong B_{\omepik,{\bm\lambda}_\bh}(\Pi_\infty^{R_\infty-\ana})^\vee,\]
and $\mathcal{E}_{\omepik,\fp,{\bm\lambda}_\bh}^\infty(\overline{\rho})$ is
defined as the Zariski-closed support of $\cM_{\omepik,{\bm\lambda}_\bh}^{\infty}$ in $\FX_{\infty}\times\sbanpiku\times \rigchu$.\;In particular,\;for $y=(\rho_y, \pi_{y},\chi_y)=(\rho_{y}^{\fp},\rho_{y,\fp},\pi_{y},\chi_y)\in \FX_{\infty}^{{\fp}}\times \mathfrak{X}_{\overline{r}}\times\sbanpiku\times \rigchu$,\;$y\in \mathcal{E}_{\omepik,\fp,{\bm\lambda}_\bh}^\infty(\overline{\rho})$ if and only if
\[\homo_{\bL_{S^u_0}(L)}\Big(\pi_{y}\otimes_E\big((\chi_y)_{\varpi_{L}}\circ\mathrm{ det}_{\bL_{S^u_0}(L)}\big)\otimes_E L_{S^u_0}({\bm\lambda}_\bh),J_{\bP_{S^u_0}(L)}\big(\Pi_\infty^{R_\infty-\ana}[\fm_y] \big)\Big)\neq0.\]
where $\fm_y=(\fm_y^{\fp},\fm_{y,\fp})$ is the corresponding maximal ideal of $R_{\infty}[1/p]$.\;



Let $\defvaru\hookrightarrow\mathfrak{X}_{\overline{r}}^\Box\times \sbanpiku\times\rigchu$ be the  Bernstein paraboline varieties \cite[Section 4.2]{Ding2021} of type {$(\omepik,\bh)$},\;which parametrizes local Galois representations that admit $\Omega_{S_0^u}$-filtrations.\;By \cite[Theorem 4.2.5,\;Corollary 4.2.5]{Ding2021} or \cite[Proposition 4.2]{He20222},\;the rigid space $\defvaru$ is equidimensional of dimension $n^2+\left(\frac{n(n-1)}{2}+s\right)d_L$.\;Note that ``$n^2$" comes from the framing.\;Let $x=(\rho_x,\underline{x},\undelram)\in \defvaru$,\;then $D_{\rig}(\rho_x)$ admits an $\omepik$-filtration $\cF=\{\fil_i^\cF D_{\rig}(\rho_x)\}$ such that,\;for all $1=1,\cdots,s$,\;
\[\gr_{i}^{\cF}D_{\rig}(\rho_x)\otimes_{\cR_{k(x),L}}\cR_{k(x),L}((\delta_i^0)^{-1}_{\varpi_L})\big[1/t\big]=\Delta_{x_i}\big[1/t\big].\]
The Bernstein paraboline variety $\defvaru$ can be viewed as a local analogue of the patched Bernstein eigenvariety $\mathcal{E}_{\omepik,\fp,{\bm\lambda}_\bh}^\infty(\overline{\rho})$.\;Consider the composition
\begin{equation}\label{composition}
	\begin{aligned}
		\mathcal{E}_{\omepik,\fp,{\bm\lambda}_\bh}^\infty(\overline{\rho})\hooklongrightarrow \FX_{\overline{\rho}^{\fp}}^\Box&\times \mathfrak{X}_{\overline{r}}^\Box \times \sbanpiku\times \rigchu\\
		&\xrightarrow{\iota_{\omepik}} \FX_{\overline{\rho}^{\fp}}^\Box\times \mathfrak{X}_{\overline{r}}^\Box \times \sbanpiku\times \rigchu.
	\end{aligned}
\end{equation}
where $\iota_{\omepik}:\sbanpiku\xrightarrow{\sim}\sbanpiku$ is an isomorphism given in the argument before \cite[(3.32)]{Ding2021}.\;By an easy variation of the proof of \cite[Theorem 3.3.9]{Ding2021},\;the composition (\ref{composition}) factors through $\FX_{\overline{\rho}^{\fp}}^\Box\times \defvarrho$,\;i.e.,\;
\begin{equation}\label{mapbersteineigenvarpikdefvarrho}
	\Lambda:\mathcal{E}_{\omepik,\fp,{\bm\lambda}_\bh}^\infty(\overline{\rho})\hooklongrightarrow \FX_{\overline{\rho}^{\fp}}^\Box\times \iota_{\omepik}^{-1}(\defvaru).
\end{equation}	
%It induces an isomorphism between $\mathcal{E}_{\omepik,\fp,{\bm\lambda}_\bh}^\infty(\overline{\rho})$ with a union of some irreducible components of $\FX_{\overline{\rho}^{\fp}}^\Box\times \iota_{\omepik}^{-1}(\defvaru)$ equipped with the reduced closed rigid subspace structure.\;

\subsection{Main results on local-global compatibility}

This section follows the route of \cite[Section 4.1]{ParaDing2024}.\;Let $\rho\in \FX_{\infty}$ be a continuous Galois representation such that  the  following Hypothesis holds.\;Let $\fm_{\rho}=(\fm^{\fp},\fm_{\fp})$ be  the corresponding maximal ideal of $R_{\infty}[1/p]$.\;
\begin{hypothesis}\label{hyongaloisrep1}(Keep the notation in  Section \ref{Omegafil} and Definition \ref{noncriticalassumption})
	\begin{itemize}
		\item[(a)] $\rho_L:=\rho_{\fp}:=\rho|_{\gal_{F^+_{\fp}}}$ is a potentially crystalline  $p$-adic Galois representation with distinct Hodge-Tate weights $\bh:=(\hpi_{1}>\hpi_{2}>\cdots>\hpi_{n} )$.\;
		\item[(b)] Let $\Dpik:=D_{\rig}(\rho_L)$ be the associated $(\varphi,\Gamma)$-module over $\cR_{E,L}$ of rank $n$.\;For  $u\in \sW_s$,\;$\Dpik$ admits a non-critical generic $\omepik$-filtration $\cF_u$ with parameters  $(\underline{x}^u_+,\underline{\delta}^{0,u})\in \sbanpiku\times \rigchu$.\;
		\item[(c)] For $u\in \sW_s$,\;put $y_{u}:=(\rho,\iota^{-1}_{\omepik}(\underline{x}^u_+,\underline{\delta}^{0,u}))\in \FX_{\infty}\times\sbanpiku\times \rigchu$.\;Then $y_u\in \mathcal{E}_{\omepik,\fp,{\bm\lambda}_\bh}^\infty(\overline{\rho})$ (and thus $x_u:=(\rho_L,\underline{x}^u_+,\underline{\delta}^{0,u})\in \defvaru$).\;
	\end{itemize}
\end{hypothesis}

By \cite[Propositions 4.10 $\&$ 4.11]{He20222},\;we have 
\begin{lem}
$x_u$ (resp.,\;$y_{u}$) is a smooth point of $\defvaru$ (resp.,\;$\mathcal{E}_{\omepik,\fp,{\bm\lambda}_\bh}^\infty(\overline{\rho})$).\;Moreover,\;$y_{u}$ does not admit companion points of non-dominant weight.\;
\end{lem}
\begin{rmk}
By the argument in the proof of \cite[Theorem 4.35]{PATCHING2016},\;the coherent sheaf $\cM_{\omepik,{\bm\lambda}_\bh}^{\infty}$ is locally free of rank $1$ at all $y_{u}$.\;
\end{rmk}



Recall that the completion of $R_{\overline{r}}^{\square}[1/p]$ at $\rho_L$ is natural to $R^{\Box}_{\rho_L}\cong R^{\Box}_{\Dpik}$,\;where $R^{\Box}_{\rho_L}$ is the framed universal deformation ring of $\rho_L$ over $\Art_E$.\;Let $R^{\Box}_{\Dpik}:=R_{\Dpik}\otimes_{R_{\rho_L}}R^{\Box}_{\rho_L}$.\;Put $R_{\Dpik,u}^{\Box,?}:=R_{\Dpik}^{\Box}\widehat{\otimes}_ER^{?}_{\Dpik,u}$ for $u\in \sW_{s}$ and $?\in\{\flat,0\}$.\;Let $(\defvaru)^{\widehat{}}_{x_u}$ be the completion of $\defvaru$ at $x_u$.\;As $x_u$ is non-critical,\;by \cite[Proposition 6.4.5]{Ding2021} and \cite[Proposition 4.9]{He20222},\;$(\defvaru)^{\widehat{}}_{x_u}$ is naturally isomorphic to $R_{\Dpik,u}^{\Box,0}$.\;

Keep the notation in Section \ref{sectionforunverext}.\;Let $\fa\supseteq\fm^2_{R^{\Box}_{\Dpik}}$ be an ideal of $R^{\Box}_{\Dpik}$ satisfies that $\fa/\fm^2_{R^{\Box}_{\Dpik}}\oplus \fm_{A_{\Dpik}}/\fm^2_{A_{\Dpik}}\xrightarrow{\sim } \fm_{R^{\Box}_{\Dpik}}/\fm^2_{R^{\Box}_{\Dpik}}$.\;Then the composition $A_{\Dpik}\hookrightarrow R^{\Box}_{\Dpik}/\fm^2_{R^{\Box}_{\Dpik}}\twoheadrightarrow R^{\Box}_{\Dpik}/\fa$ is an isomorphism (i.e.,\;the ideal $\fa$ deletes the information comes from framing).\;For $u\in \sW_s$ and $?\in\{\emptyset,\flat,0\}$,\;we use $\fa$ to denote its image in $R^{\Box,?}_{\Dpik,u}$.\;Thus we get $\fa/\fm^2_{R^{\Box,?}_{\Dpik,u}}\oplus \fm_{A^{?}_{\Dpik,u}}/\fm^2_{A^{?}_{\Dpik,u}}\xrightarrow{\sim } \fm_{R^{\Box,?}_{\Dpik,u}}/\fm^2_{R^{\Box,?}_{\Dpik,u}}$ and $\fa/\fm^2_{R^{\Box,?}_{\Dpik,u,g}}\oplus \fm_{A_{0,u}}/\fm^2_{A_{0,u}}\xrightarrow{\sim } \fm_{R^{\Box,?}_{\Dpik,g}}/\fm^2_{R^{\Box,?}_{\Dpik,g}}$,\;so that the natural maps
$A^{?}_{\Dpik,u}\hookrightarrow R^{\Box,?}_{\Dpik,u}/\fm^2_{R^{\Box,?}_{\Dpik,u}}\twoheadrightarrow R^{\Box,?}_{\Dpik,u}/\fa$ and $A_{0,u}\hookrightarrow R^{\Box,?}_{\Dpik,u,g}/\fm^2_{R^{\Box,?}_{\Dpik,u,g}}\twoheadrightarrow R^{\Box,?}_{\Dpik,u,g}/\fa$ are isomorphisms.\;We use $\fa\subseteq R_{\overline{r}}^{\square}[1/p]$ denote the preimage of $\fa\subseteq R_{\overline{r}}^{\square}$.\;Let $\fm_{\rho}:=(\fm_{\fp},\fm^{\fp})\subseteq\fa_{\rho}:=(\fa,\fm^{\fp})\subseteq R_{\infty}[1/p]$.\;

Our goal is to investigate $\Pi_\infty^{R_\infty-\ana}[\fa_{\rho}]$,\;which is equipped with a natural $R^{\Box}_{\Dpik}/\fa\cong A_{\Dpik}$-action.\;By \cite[Proposition 4.33]{PATCHING2016},\;$\Pi_\infty^{R_\infty-\ana}[\fm_{\rho}]^{\lalg}\cong \pi_{\lalg}(\underline{x},\bh)$.\;Similar to \cite[Lemma 4.2]{ParaDing2024},\; $\fa+\fm_{A_{\Dpik}}=\fm_{R^{\Box}_{\Dpik}}$ and hence $\fa_{\rho}+\fm_{A_{\Dpik}}=\fm_{\rho}$,\;we get that 
$\Pi_\infty^{R_\infty-\ana}[\fm_{\rho}]=\Pi_\infty^{R_\infty-\ana}[\fa_{\rho}][\fm_{A_{\Dpik}}]$.\;Similar to \cite[Lemma 4.2 (2)]{ParaDing2024},\;we can prove that
\[\homo_G(\Pi_\infty^{R_\infty-\ana}[\fm_{\rho}]^{\lalg},\Pi_\infty^{R_\infty-\ana}[\fa_{\rho}])\cong \homo_G(\Pi_\infty^{R_\infty-\ana}[\fm_{\rho}]^{\lalg},\Pi_\infty^{R_\infty-\ana}[\fm_{\rho}])\cong E.\]
Let $U_u=U_{\fp,u}\times U^{\fp}_u\subseteq  \iota^{-1}_{\omepik}\big(\defvaru\big)\times\FX_{\overline{\rho}^{\fp}}^{\Box}$ be a smooth affinoid neighourhood of $y_u$ such $y_{u'}\not\in U_{\fp,u}$ for $u'\neq u$.\;Let $\fm_{y_{u,\fp}}$ be the maximal ideal of $\cO(U_{\fp,u})$ at $y_{u,\fp}$ and $\fa\supset \fm^2_{y_{u,\fp}}$ be the closed ideal generated by the above ideal $\fa\subseteq R_{\overline{r}}^{\square}[1/p]$.\;In the sequel,\;we put
\[\cM_{y_u}:=\cM_{\omepik,{\bm\lambda}_\bh}^{\infty}(U_u)/(\fm_{y_{u,\fp}}+\fm^{\fp}),\widetilde{\cM}_{y_u}:=\cM_{\omepik,{\bm\lambda}_\bh}^{\infty}(U_u)/(\fa+\fm^{\fp})\]
for simplicity.\;There are natural $\bZ_{S^u_0}(L)\times R_{\infty}$-equivariant injections:
\[\cM_{y_u}\hookrightarrow \widetilde{\cM}_{y_u}\hookrightarrow J_{\bP^u_{S_0}}(\Pi_\infty^{R_\infty-\ana})[\fa_{\rho}].\]
Since $\cM_{\infty}$ is locally free of rank $1$ at $y_{u}$,\;we obtain that $\widetilde{\cM}_{y_u}\cong R^{\Box,0}_{\Dpik,u}/\fa_{\rho}$,\;so that $\dim_E\widetilde{\cM}_{y_u}=1+(s+sd_L)$.\;In this case,\;similar to the argument around \cite[(4.5)]{ParaDing2024},\;we have a $\bZ_{S^u_0}(L)\times A^{0}_{\Dpik,u}$-equivariant isomorphism $\widetilde{\cM}^{\vee}_{y_u}\cong \pi_{\sm}(\underline{x}^u)\eta^{-1}_{S^u_0}\otimes_E\widetilde{\delta}_u^{\univ}$ (note that the $A^{0}_{\Dpik,u}$-action factors through $A^{0}_{\Dpik,u}\twoheadrightarrow A^{\flat}_{\Dpik,u}$).\;Similar to the proof of \cite[Lemma 4.19]{He20222},\;by using the non-critical assumption of $y_u$,\;the maps $\cM_{y_u}\hookrightarrow  \Pi_\infty^{R_\infty-\ana}[\fm_{\rho}]$ and $ \widetilde{\cM}_{y_u}\hookrightarrow \Pi_\infty^{R_\infty-\ana}[\fa_{\rho}]$ are balanced,\;hence induces a $G\times R_{\infty}$-equivariant injection:
\begin{equation}
	\begin{aligned}
		&\iota_u:I^G_{\op_{S^u_0}(L)}\pi_{\sm}(\underline{x}^u)\eta_{S^u_0}\widetilde{\delta}_u^{\univ}\hookrightarrow \Pi_\infty^{R_\infty-\ana}[\fa_{\rho}].\;
	\end{aligned}
\end{equation}
where the $R_{\infty}$-action on $I^G_{\op_{S^u_0}(L)}\pi_{\sm}(\underline{x}^u)\eta_{S^u_0}\widetilde{\delta}_u^{\univ}$ is induced by $R_{\infty}\rightarrow (R^{\Box,0}_{\Dpik,u}/\fa)\otimes_E (R^{\fp}_{\infty}[1/p]/\fm^{\fp})\xrightarrow{\sim}A_{\Dpik,u}^{0}\twoheadrightarrow (R^{\Box,\flat}_{\Dpik,u}/\fa)\otimes_E (R^{\fp}_{\infty}[1/p]/\fm^{\fp})\xrightarrow{\sim}A_{\Dpik,u}^{\flat}$.\;Moreover,\;we have an injection of locally analytic representations $\mathrm{PS}_{u}(\underline{x},\bh)\hookrightarrow I^G_{\op_{S^u_0}(L)}\pi_{\sm}(\underline{x}^u)\eta_{S^u_0}\widetilde{\delta}_u^{\univ}$.\;


Let $\widetilde{\pi}$ be the subrepresentation of $\Pi_\infty^{R_\infty-\ana}[\fa_{\rho}]$ generated by $\mathrm{Im}(\iota_u)$  for all $u\in \sW_s
$.\;Note that $\widetilde{\pi}$ has an $A_{\Dpik}$-action via $A_{\Dpik}\xrightarrow{\sim}R^{\Box}_{\Dpik}/\fa \xrightarrow{\sim}R_{\infty}[1/p]/\fa_{\rho}$.\;Similar to \cite[Lemma 4.4]{ParaDing2024},\;we have
\begin{pro}\label{imiotawmx}
 $\iota_u(\mathrm{PS}_{u,1}(\underline{x},\bh))=(\mathrm{Im}(\iota_u))[\fm_{\rho}]$.\;
 \end{pro}
\begin{proof}
Note that the composition $\mathrm{PS}_{u}(\underline{x},\bh)\hookrightarrow I^G_{\op_{S^u_0}(L)}\pi_{\sm}(\underline{x}^u)\eta_{S^u_0}\widetilde{\delta}_u^{\univ}\hookrightarrow \Pi_\infty^{R_\infty-\ana}[\fa_{\rho}]$ sends $\mathrm{PS}^{\lalg}_{u}(\underline{x},\bh)$ to $\Pi_\infty^{R_\infty-\ana}[\fm_{\rho}]$.\;Since  $\homo_{G}(\mathrm{PS}_{u,1}(\underline{x},\bh),\Pi_\infty^{R_\infty-\ana}[\fm_{\rho}])\xrightarrow{\sim}\homo_{G}(\mathrm{PS}^{\lalg}_{u}(\underline{x},\bh),\Pi_\infty^{R_\infty-\ana}[\fm_{\rho}])$, then the same strategy as in the first paragraph of the proof of \cite[Lemma 4.6]{ParaDing2024} shows that the previous composition has image in $\Pi_\infty^{R_\infty-\ana}[\fm_{\rho}]$.\;Using the same strategy as in \cite[Lemma 4.6]{ParaDing2024},\;the injection $\iota_u(\mathrm{PS}_{u,1}(\underline{x},\bh))\hookrightarrow(\mathrm{Im}(\iota_u))[\fm_{\rho}]$ is surjective too.\;
\end{proof}


\begin{thm}\label{mainthmglobal}We have an $A_{\Dpik}\times G$-equivariant isomorphism $\widetilde{\pi}\cong \pi^{\flat}_{1}(\underline{x},\bh)^{\univ}$.\;In particular,\;we have a $\GLN_n(L)$-equivariant injection
	\[\pi^{\flat}_{\min}(\Dpik)\cong \pi^{\flat}_{1}(\underline{x},\bh)^{\univ}[\fm_{A_{\Dpik}}]\cong \widetilde{\pi}[\fm_{A_{\Dpik}}]\hookrightarrow  \Pi_\infty^{R_\infty-\ana}[\fa_{\rho}][\fm_{A_{\Dpik}}]\cong \Pi_\infty^{R_\infty-\ana}[\fm_{\rho}].\]
\end{thm}
\begin{proof}
	Similar to the argument in the proof \cite[Theorem 4.5,\; Corollary 4.6]{ParaDing2024},\;we can also see that $\widetilde{\pi}$ is a subrepresentation of $\pi^{\flat}_{1}(\underline{x},\bh)^{\univ}$ and contains all $I^G_{\op_{S^u_0}(L)}\pi_{\sm}(\underline{x}^u)\eta_{S^u_0}\widetilde{\delta}_u^{\univ}$ so that $\widetilde{\pi}\cong \pi^{\flat}_{1}(\underline{x},\bh)^{\univ}$.\;
\end{proof}


\section{Appendix:\;envelope of intersections of Borel and parabolic subalgebras  in Lie algebras}\label{reproofforinfinite}




\begin{lem}\label{Inifinitefernpoten} Keep the non-critical assumption in Remark \ref{explainHomasLiealg}.\;Then $\sum_{u\in \sW_s}\mathrm{Ad}_{(u^{\sharp})^{-1}}(\fp_{S^u_0,\sigma})\cap \mathrm{Ad}_{g_{\sigma}}(\fb_{\sigma})=\mathrm{Ad}_{g_{\sigma}}(\fb_{\sigma})$.\;
\end{lem}
\begin{proof}Keep the notation in Remark \ref{explainHomasLiealg}.\;Note that $\mathrm{Ad}_{g_{\sigma}}(\fb_{\sigma})=\mathrm{Ad}_{b_{\sigma}w_{0,\sigma}}(\fb_{\sigma})$ and 
\[\sum_{u\in \sW_s}\mathrm{Ad}_{(u^{\sharp})^{-1}}(\fp_{S^u_0,\sigma})\cap \mathrm{Ad}_{b_{\sigma}w_{0,\sigma}}(\fb_{\sigma})=\sum_{u\in \sW_s}\mathrm{Ad}_{b_{\sigma}}(\mathrm{Ad}_{b_{\sigma}^{-1}(u^{\sharp})^{-1}}(\fp_{S^u_0,\sigma})\cap \mathrm{Ad}_{w_{0,\sigma}}(\fb_{\sigma})).\;\]
Therefore,\;it suffices to show that 
 \[\mathrm{Ad}_{w_{0,\sigma}}(\fb_{\sigma})=\sum_{u\in \sW_s}\mathrm{Ad}_{b_{\sigma}^{-1}(u^{\sharp})^{-1}}(\fp_{S^u_0,\sigma})\cap \mathrm{Ad}_{w_{0,\sigma}}(\fb_{\sigma}).\]
 Let $e^{i,j}$ be the elementary $(n\times n)$-matrix whose $(l,m)$-entry is given by $\delta_{i,l}\delta_{j,m}$.\;For any $i>j$,\;we claim that there is some permutation
$u_{i,j}\in\sW_s$ and $i-j$ scalars $\{x_{i,l}\}_{j+1\leq l\leq i}$ in $E$ such that the matrix $a^{i,j}:=e^{i,j}+\sum_{l=j+1}^ix_{i,l}e^{i,l}$ lies in $\mathrm{Ad}_{b_{\sigma}^{-1}(u^{\sharp})^{-1}}(\fp_{S^u_0,\sigma})$.\;Recall the previous basis 
$\{e_{1,j,\sigma}\}_{1\leq j\leq r_1},\cdots,\{e_{s,j,\sigma}\}_{1\leq j\leq r_s}$ (resp.,\;$e_{1,\sigma},\cdots,e_{n,\sigma}$).\;Recall that this basis gives a $\bP_{S^u_0}$-parabolic filtration 
\[V_{\bullet}=(0\subseteq V_0\subset V_1\subset\cdots\subset V_s=E^n),\]
 where $V_i=\langle \{e_{u^{-1}(1),j,\sigma}\}_{1\leq j\leq r_{u^{-1}(1)}},\cdots,\{e_{u^{-1}(i),j,\sigma}\}_{1\leq j\leq r_{u^{-1}(i)}}\rangle$.\;Let $s_{i,j}=(i,j)\in \sW_s$.\;Suppose that $t^u_{j'}+1\leq j\leq t^u_{j'+1}$ and $t^u_{i'}+1\leq j\leq t^u_{i'+1}$ for $\leq j'\leq i'\leq s$.\;Consider the following basis 
\[\cB=\{b^{-1}_{\sigma}(e_{1,\sigma}),\cdots,b^{-1}_{\sigma}(e_{j-1,\sigma}),e_{j,\sigma},b^{-1}_{\sigma}(e_{j+1,\sigma}),\cdots,b^{-1}_{\sigma}(e_{i,\sigma}),e_{i+1,\sigma},\cdots,e_{n,\sigma}\},\]
of $E^n$ (this is also a basis since $b^{-1}_{\sigma}$ is upper triangular).\;Define $\pi(x)=0$ for $x\in \cB\backslash\{e_{j,\sigma}\}$ and $\pi(e_{j,\sigma})=e_{i,\sigma}$.\;We also put $V'_l=\langle e_{1,\sigma},\cdots,e_{l,\sigma}\rangle$ for $1\leq l\leq n$.\;Note that
\[b_{\sigma}\pi b^{-1}_{\sigma}(s_{i,j}V_{l})=\left\{
\begin{array}{ll}
	0,\;&1\leq l \leq i-1,\\
	E\langle b_{\sigma}e_{i,\sigma}\rangle,\;&i\leq l \leq n.
\end{array}
\right.\]
Since $\langle b_{\sigma}e_{i,\sigma}\rangle\subseteq V'_i\subseteq s_{i,j}V_{i'}$ and $V_i\subseteq V_{i'}$,\;we find a morphism $\pi:E^n\rightarrow E^n$ such that the endomorphism $b_{\sigma}\pi b^{-1}_{\sigma}$ stabilizes the $\bP_{S^u_0}$-parabolic filtration $s_{i,j}V_{\bullet}$.\;Then the desired $a^{i,j}$ is the matrix in the standard basis.\;Similar to the last paragraph in \cite[Lemma 2.1]{DensityHMS},\;such $a^{i,j}$ also form a basis of the $\mathrm{Ad}_{w_{0,\sigma}}(\fb_{\sigma})$ which lie in the envelope $\sum_{u\in \sW_s}\mathrm{Ad}_{b_{\sigma}^{-1}(u^{\sharp})^{-1}}(\fp_{S^u_0,\sigma})\cap \mathrm{Ad}_{w_{0,\sigma}}(\fb_{\sigma})$.\;
\end{proof}


%\begin{pro}\label{envelooefortau}
%We have $\dim_E\sum_{u\in \sW_s}\mathrm{Ad}_{(u^{\sharp})^{-1}}(\tau_{S_0,\sigma})\cap \mathrm{Ad}_{g_{\sigma}}(\fb_{\sigma})$.\;
%\end{pro}
%\begin{proof}By (\ref{noncriticalintersection}),\;we have 
%\[\sum_{u\in \sW_s}\mathrm{Ad}_{(u^{\sharp})^{-1}}(\tau_{S_0,\sigma})\cap \mathrm{Ad}_{g_{\sigma}}(\fb_{\sigma})=\sum_{u\in \sW_s}\mathrm{Ad}_{(u^{\sharp})^{-1}}(\fz_{S_0,\sigma})\cap \mathrm{Ad}_{g_{\sigma}}(\fb_{\sigma}).\]
%Moreover,\;note that $\mathrm{Ad}_{(u^{\sharp})^{-1}}(\fz_{S_0,\sigma})\cap \mathrm{Ad}_{b_{\sigma}w_{0,\sigma}}(\fb_{\sigma})=\mathrm{Ad}_{b_{\sigma}}(\mathrm{Ad}_{b_{\sigma}^{-1}(u^{\sharp})^{-1}}(\fz_{S_0,\sigma})\cap \mathrm{Ad}_{w_{0,\sigma}}(\fb_{\sigma}))$.\;Note that
%\[\sum_{u\in \sW_s}\mathrm{Ad}_{b_{\sigma}^{-1}(u^{\sharp})^{-1}}(\fz
%_{S_0,\sigma})\cap \overline{\fb}_{\sigma}\cong \sum_{u\in \sW_s}\mathrm{Ad}_{b_{s,\sigma}^{-1}u}(\ft
%_{s,\sigma})\cap \overline{\fb}_{s,\sigma}\] for some $b_{s,\sigma}\in \bB_s(E)$.\;We thus return to the case in \cite[Theorem 2.3]{DensityHMS}.
%\end{proof}













\bibliographystyle{plain}
\begin{thebibliography}{10}
	\bibitem{ALLENPARK}
	P. B. Allen, 
	\newblock Deformations of polarized automorphic Galois representations and adjoint Selmer groups. (English summary)
	\newblock {\em Duke Math. J.} 165 (2016), no. 13, 2407–2460.
	11F80 (11F70 11R34)
	
\bibitem{AST2009324R10}
	J. Bellaïche,  G. Chenevie
	\newblock Families of Galois representations and Selmer groups. 
	\newblock {\em Astérisque} No. 324 (2009), xii+314 pp. ISBN: 978-2-85629-264-8


	
	
	\bibitem{berger2008equations}
	L. Berge
	\newblock {Équations différentielles $p$-adiques et $(\varphi,N)$-modules filtrés}.
	\newblock {\em Astérisque} No. 319 (2008), 13–38. ISBN: 978-2-85629-256-3
	


	\bibitem{bernstein1984centre}
	J. N. Bernstein,  and P. Deligne.
	\newblock {Le ``centre'' de Bernstein. (French) [The Bernstein ``center'']}.
	\newblock Edited by P. Deligne. {\em Travaux en Cours, Representations of reductive groups over a local field,} 1–32, Hermann, Paris, 1984.

	 \bibitem{breuil2019ext1}
	C. Breuil
	\newblock {${\ext}^1$ localement analytique et compatibilit{\'e}
		local-global}.
	\newblock {\em Amer. J. Math.} 141 (2019), no. 3, 611–703.
	11S37 (14D24)



	
	\bibitem{Ding2021}
	C. Breuil,  Y. Ding, 
	\newblock Bernstein eigenvarieties.
	\newblock {\em  Peking Mathematical Journal } (2024) 7:471–642.
	
		\bibitem{HigherLinvariantsGL3(Qp)}
	C. Breuil,  Y. Ding, 
	\newblock {Higher $L$-invariants for $\GLN_3(\bQ_p)$ and local-global
		compatibility}.
	\newblock {\em Camb. J. Math.} 8 (2020), no. 4, 775–951.
	
	\bibitem{BDcritical25}
	C. Breuil,  Y. Ding, 
	\newblock {Hodge filtration and crystalline representations of $\GLN_n$}.
	\newblock {\em arXiv:2512.12153, preprint}.
	
		\bibitem{Dingext1cusp}
	Y. Ding,
	\newblock {Locally analytic $\ext^1$ for $\GLN_2(\bQ_p)$ in de Rham non-trianguline case}
	\newblock {\em Represent. Theory } 26 (2022), 122-133.

	\bibitem{2015Ding}
Y. Ding,
\newblock {$\mathcal{L}$-invariants, partially de Rham families, and local-global compatibility}
\newblock {\em Ann. Inst. Fourier (Grenoble)} 67 (2017), no. 4, 1457–1519.
	
	\bibitem{breuil2020probleme}
	C. Breuil, Y. Ding,
	\newblock {SUR UN PROBL{\`E}ME DE COMPATIBILIT{\'E} LOCAL-GLOBAL LOCALEMENT ANALYTIQUE}.
	\newblock {\em Memoirs of the American Mathematical Society}, 2020.
	
	\bibitem{BQ24}
	C. Breuil, Z. Qian,
	\newblock {Splitting and expliciting the de Rham complex of the Drinfeld space.}
	\newblock {\em arXiv:2407.06651, preprint}


\bibitem{breuil2017interpretation}
C. Breuil, E. Hellmann, and B. Schraen,
\newblock Une interpr{\'e}tation modulaire de la vari{\'e}t{\'e} trianguline.
\newblock {\em Math. Ann.} 367 (2017), no. 3-4, 1587–1645.




	
	\bibitem{breuil2007first}
	C. Breuil, P. Schneide
	\newblock First steps towards $p$-adic langlands functoriality. 
	\newblock {\em J. Reine Angew. Math.} 610 (2007), 149–180.

	
	\bibitem{bushnell2016admissible}
	C. J. Bushnell,  P. C. Kutzko,
	\newblock {\em The Admissible Dual of $\GLN_n$ via Compact Open
		Subgroups}, 
	\newblock {\em Annals of Mathematics Studies, 129.} Princeton University Press, Princeton, NJ, 1993. xii+313 pp. ISBN: 0-691-03256-4; 0-691-02114-7

%	\bibitem{Gvalued}
%	R. Bellovin, T. Gee, 
%	\newblock {{$G$}-valued local deformation rings and global lifts}.
%	\newblock {\em	Algebra Number Theory} 13(2019), no.2, 333–378.
	
	\bibitem{PATCHING2016}
	A. Caraiani, M. Emerton, T. Gee, D. Geraghty, V.
	Pa\v{s}k\={u}nas, and S. W. Shin.
	\newblock Patching and the p-adic local Langlands correspondence. 
	\newblock {\em Camb. J. Math.} 4 (2016), no. 2, 197–287.

	
	\bibitem{chenevier2011infinite}
	G. Chenevier.
	\newblock On the infinite fern of Galois representations of unitary type. 
	\newblock {\em Ann. Sci. Éc. Norm. Supér.} (4) 44 (2011), no. 6, 963–1019.

	

   
     \bibitem{ParaDing2024}
    Y. Ding,
    \newblock {$p$-adic Hodge parameters in the crystalline representations of $\GLN_n$}
    \newblock {\em arXiv}.
	
	\bibitem{2019DINGSimple}
	Y. Ding, 
	\newblock {Simple $\cL$-invariants for $\GL_{n}$}.
	\newblock {\em Trans. Amer. Math. Soc.} 372 (2019), no. 11, 7993–8042.

   \bibitem{Dingsocle}
Y. Ding, 
\newblock {Companion points and locally analytic socle for $\GL_2(L)$}.
\newblock {\em Israel Journal of Mathematics} Vol. 231, Issue 1, (2019), p. 47-122.
	
	\bibitem{emerton2006jacquet}
	M. Emerton.
	\newblock {Jacquet modules of locally analytic representations of p-adic reductive groups. I. Construction and first properties.}
	\newblock {\em Ann. Sci. École Norm. Sup.} (4) 39 (2006), no. 5, 775–839.

	


	
	\bibitem{Emerton2007summary}
	M. Emerton, 
	\newblock Locally analytic representation theory of p-adic reductive groups: a summary of some recent developments.
	\newblock {\em L-functions and Galois representations}, 407–437,
	London Math. Soc. Lecture Note Ser., 320, Cambridge Univ. Press, Cambridge, 2007.
	
		\bibitem{emerton2017locally}
	M. Emerton, 
	\newblock {\em Locally analytic vectors in representations of locally $p$-adic analytic groups}.
	\newblock {\em Mem. Amer. Math. Soc.} 248 (2017), no. 1175, iv+158 pp. ISBN: 978-0-8218-7562-9; 978-1-4704-4052-7

%	\bibitem{emerton2014geometric}
%	M. Emerton and Toby Gee.
%	\newblock A geometric perspective on the breuil--m{\'e}zard conjecture.  
%	\newblock {\em J. Inst. Math. Jussieu} 13 (2014), no. 1, 183–223.

%    \bibitem{emerton2023introduction}
%    M. Emerton,\;T. Gee and E. Hellmann, 
%    \newblock {An introduction to the categorical $p$-adic Langlands program}.
%    \newblock {\em arXiv:} 2210.01404 (2023).
    
	
%	\bibitem{1987Catehighetweight}
%	T. J. Enright,  B. Shelton, 
%	\newblock {Categories of highest weight modules: applications to classical Hermitian symmetric pairs}.
%	\newblock {\em Mem. Amer. Math. Soc.} 67 (1987), no. 367, iv+94 pp.
	
	
	\bibitem{2022ext1hyq}
	Y. He.
	\newblock {Extensions of locally analytic generalized parabolic Steinberg
		representations}.
	\newblock {\em arXiv:\;2211.00476, preprint}.
	
	\bibitem{He20222}
	Y. He,
	\newblock Parabolic Simple $\mathcal{L}$-Invariants.
	\newblock {\em arXiv:\;2211.10847, preprint}.
	
		\bibitem{HEparaforsemitable}
	Y. He,
	\newblock Towards the $p$-adic Hodge parameters in the semistable representation of $\GLN_n(\bQ_p)$.
	\newblock {\em In preparation}.
	

	
	
	\bibitem{humphreysBGG}
	J. E. Humphreys, 
	\newblock {\em Representations of Semisimple Lie Algebras in the BGG Category $\cO$}, volume~94.
	\newblock Graduate Studies in Mathematics, 94. American Mathematical Society, Providence, RI, 2008. xvi+289 pp. ISBN: 978-0-8218-4678-0
	
	\bibitem{DensityHMS}
	E. Hellmann, C. M. Margerin, B. Schraen 
	\newblock {Density of automorphic points in deformation
		rings of polarized global Galois representations.}
	\newblock {\em Duke Mathematical Journal}, 171(13) 2699-2752 15
	September 2022.
	

	\bibitem{liu2007cohomology}
	R. Liu.
	\newblock Cohomology and duality for ($\varphi$, $\Gamma$)-modules over the
	robba ring.
	\newblock {\em Int. Math. Res. Not. IMRN 2008}, no. 3, Art. ID rnm150, 32 pp. 
	
%	 \bibitem{matrixsch}
%	E. Mille B. Sturmfels,
%	\newblock {Combinatorial commutative algebra}.
%	\newblock {\em Graduate Texts in Mathe
%		matics}  vol. 227, Springer-Verlag, New York, 2005.
%	
	
	\bibitem{nakamura2009classification}
	K. Nakamura,
	\newblock {Classification of two-dimensional split trianguline representations
	of p-adic fields.}
	\newblock {\em Compos. Math.} 145 (2009), no. 4, 865–914.

%	\bibitem{XLROBBA}
%	\textit{K.~S.~Kedlaya}, \textit{J.~Pottharst} and  \textit{L.~Xiao},
%	\newblock {Cohomology of arithmetic families of {$(\varphi,\Gamma)$}-modules}.
%		\newblock {\em J. Amer. Math. Soc.} \textbf{27} (2014), no. 4, 1043-1115.
	
	
	
	
	\bibitem{orlik2015jordan}
	S. Orlik,
	\newblock On Jordan-Hölder series of some locally analytic representations.
	\newblock {\em J. Amer. Math. Soc.} 28 (2015), no. 1, 99–157.
%
%	
%	\bibitem{orlik2014jordan}
%	S. Orlik, B. Schraen.
%	\newblock The {Jordan-H{\"o}lder} series of the locally analytic steinberg representation. 
%	\newblock {\em Doc. Math.} 19 (2014), 647–671.
%	
%	\bibitem{BSCONJ}
%	A. Pyvovarov,
%	\newblock On the Breuil-Schneider conjecture generic case. 
%	\newblock {\em Algebra Number Theory} 15 (2021), no. 2, 309–339.
%
%	
	\bibitem{Dilogarithm}
	Z. Qian, 
	\newblock {Dilogarithm and higher {$\mathcal{L}$}-invariants for {${\rm GL}_3({\BQ}_p)$}}.
	\newblock {\em	Represent. Theory} 25 (2021), 344–411.
%	
%	\bibitem{wholeLINV}
%	Z. Qian,
%	\newblock {On generalization of Breuil-Schraen's $\sL$-invariants to $\GLN_{n}$.}
%	\newblock {\em arXiv:\;2210.01381, preprint}
	
	\bibitem{schraen2011GL3}
	B.  Schraen, 
	\newblock {Repr{\'e}sentations localement analytiques de
		$\mathrm{GL}_3(\mathbb{Q}_p)$}.
	\newblock {\em Ann. Sci. {\'E}c. Norm. Sup{\'e}r.} (4) 44 (2011), no. 1, 43–145. 
	
	
		\bibitem{Coadmissible}
	P. Schneider and J. Teitelbaum, 
	\newblock {Algebrasof $p$-adic distributions and admissible representations}.
	\newblock {\em Invent. Math.} 153	(2003), 145–196. 
		
	
%	
%	\bibitem{schneider2002banach}
%	P. Schneide J. Teitelbaum,  
%	\newblock Banach space representations and Iwasawa theory. 
%	\newblock {\em Israel J. Math.} 127 (2002), 359–380.

%   \bibitem{Stack}
%   The Stacks project,
%   \newblock {https://stacks.math.columbia.edu/.}
	
   \bibitem{scholze2013local}
   P. Scholze,
  \newblock The local Langlands correspondence for $\GLN_{n}$ over $p$-adic fields, and the cohomology of compact unitary Shimura varieties.  Shimura varieties, 251–265,
  \newblock London Math. Soc. Lecture Note Ser., 457, Cambridge Univ. Press, Cambridge, [2020], ©2020. 
  
%  \bibitem{MSW}
%  W. S. Weirich, 
%  \newblock {The trianguline variety at semi-stable points}.
%  \newblock {\em PhD.,\;thesis}.
  
  
%  \bibitem{wu2021local}
%  Z. Wu, 
%  \newblock {Local models for the trianguline variety and partially classical families}.
%  \newblock {\em arXiv:\;2103.03823}  (2021).
%  
%  \bibitem{wu2023localcompan2}
%  Z. Wu,
%  \newblock {Companion points on the eigenvariety with non-regular weights.}
%  \newblock {\em Int. Math. Res. Not.} (2023), no.9, 8008–8032.
%  
%  \bibitem{wu2024translation}
%  Z. Wu,
%  \newblock {Geometric translations of $(\varphi,\Gamma)$-modules for $\GLN_{2}(L)$}
%  \newblock {\em arXiv:\;2405.16637} (2024).
  
  \bibitem{av1980induced2}
  A. V. Zelevinsky,
  \newblock {Induced representations of reductive $p$-adic groups. II. On irreducible representations of $\GLN(n)$}.
  \newblock {\em Ann. Sci. École Norm. Sup.} (4) 13 (1980), no. 2, 165–210.
  
  
 
  
 
  
 
\end{thebibliography}
\printindex
\end{document}

\subsubsection{Preliminaries on Hodge parameters}
This section follow the route in [DingParameter,\;Section 2.2-2.4].\;In this section,\;we assume that $d\geq 2$ and write $\Dpik$ with the $\cF_{I,w_0}$ paraboline filtration,\;i.e.,\;
\[\Dpik=E_{1}-E_{2}-\cdots-E_{d}:=C_1-E_{d}.\]
We consider the following refinement of $\Dpik$
\[\Dpik=E'_{d}-E'_{1}-\cdots-E'_{d-1}:=D'_{d}-C_1',\]
Then we obtain a natural morphism $\iota_{\Dpik}:C_1\rightarrow C_1'$,\;which is induced by the composition $C_1\hookrightarrow \Dpik\twoheadrightarrow  C_1'$.

%We also have two exact sequences:
%\begin{equation}
%	\begin{aligned}
	%		&0\rightarrow D_1^{E_1}\rightarrow\Dpik(E_1)\rightarrow E_1\rightarrow 0,\\
	%		&0\rightarrow F_1'\rightarrow\Dpik(E_1)\rightarrow  F_1^{E_1}\rightarrow 0.\;
	%	\end{aligned}
%\end{equation}


As in [Proposition 2.1]{DingParameter},\;we have
\begin{pro} We have $\dim_E\homo(C_1,C_1')=\left\{
	\begin{array}{ll}
		1,\;&d=2\\
		2,\;&d\geq 3
	\end{array}
	\right.$
\end{pro}

Under the following cup-product
\begin{equation}
	\ext^1(E_d,C_1)\times \homo(C_1,C_1')\rightarrow \ext^1({E_d},C_1').
\end{equation}
Similar to the argument in [DingParameter,Proposition 2.3],\;we see that $E[\Dpik]=[\iota_{\Dpik}]^{\perp}$,\;and in particular,\;$\Dpik$ is determined by $C_1$,\;$C_1'$,\;and $\iota_{\Dpik}$.\;

This section follow the route in [DingParameter,\;Section 2.2-2.4].\;In this section,\;we assume that $d\geq 2$ and write $\Dpik$ with the $\cF_{I,w_0}$ paraboline filtration,\;i.e.,\;
\[\Dpik=D_{I^d}-D_{I^{d-1}}-\cdots-D_{I^1}.\]

For any $(\varphi,\Gamma)$-submodule $E_1\subseteq D_{I^1}$ of rank $r\leq n_1$,\;we set $E_1^c:=D_{I^1}/E_1$.\;Put $D_1(E_1):=D_{I^d}-D_{I^{d_1}}-\cdots-D_{I^{2}}-E_1$ and $\Dpik(E_1):=\ker (\Dpik\twoheadrightarrow  E_1^c)$.\;

Secondly,\;we consider 
\[\Dpik=F_1'-D'_{I^d}-D'_{I^{d-1}}-\cdots-D'_{I^{2}}-F_1^c,\]
and put $C_1(F_1):=D'_{I^d}-D'_{I^{d-1}}-\cdots-D'_{I^{2}}-F_1^c$.\;We have two exact sequences:
\begin{equation}
	\begin{aligned}
		&0\rightarrow D_1(E_1)\rightarrow\Dpik\rightarrow E_1^c\rightarrow 0,\\
		&0\rightarrow F_1'\rightarrow\Dpik\rightarrow  C_1(F_1)\rightarrow 0.\;
	\end{aligned}
\end{equation}
We also write $D_1^{E_1}:=D_{I^d}-D_{I^{d_1}}-\cdots-D_{I^{2}}\hookrightarrow D_1(E_1)$ and $C_1^{F_1}:=D'_{I^d}-D'_{I^{d-1}}-\cdots-D'_{I^{2}}\hookrightarrow C_1(F_1)$.\;
Then we obtain a natural morphism $\iota_{\Dpik}(E_1,F_1):D_1^{E_1}\rightarrow C_1^{F_1}$,\;which is induced by the composition $D_1^{E_1}\hookrightarrow D_1(E_1)\hookrightarrow\Dpik\twoheadrightarrow  C_1(F_1)$,\;but this composition will factors through $C_1^{E_1}\hookrightarrow C_1(F_1) $ since $\homo(D_1^{E_1},F_1^c)=0$.\;When $E_1=F_1=D_{I^1}$,\;we omit the notation $E_1$ and $F_1$.\;


%We also have two exact sequences:
%\begin{equation}
%	\begin{aligned}
	%		&0\rightarrow D_1^{E_1}\rightarrow\Dpik(E_1)\rightarrow E_1\rightarrow 0,\\
	%		&0\rightarrow F_1'\rightarrow\Dpik(E_1)\rightarrow  F_1^{E_1}\rightarrow 0.\;
	%	\end{aligned}
%\end{equation}


As in [Proposition 2.1]{DingParameter},\;we have
\begin{pro} We have $\dim_E\homo(D_1^{E_1},C_1^{F_1})=\left\{
	\begin{array}{ll}
		1,\;&d=2\\
		2,\;&d\geq 3
	\end{array}
	\right.$
\end{pro}

Under the following cup-product
\begin{equation}
	\ext^1({E_1},D_1^{E_1})\times \homo(D_1^{E_1},C_1^{F_1})\rightarrow \ext^1({E_1},C_1^{F_1}).
\end{equation}
Similar to the argument in [Proposition 2.3]{DP},\;we see that $E[\Dpik({E_1})]=[\iota_{\Dpik}({E_1})]^{\perp}$,\;and in particular,\;$\Dpik({E_1})$ is determined by $D_1^{E_1}$,\;$C_1^{E_1}$,\;${E_1}$ and $\iota_{\Dpik}({E_1})$.\;

Recall the two exact sequences:
\begin{equation}
	\begin{aligned}
		0\rightarrow D_1\rightarrow\Dpik\rightarrow D_{d}\rightarrow 0,\\
		0\rightarrow D'_{d}\rightarrow\Dpik\rightarrow  C_1\rightarrow 0.\;
	\end{aligned}
\end{equation}

We then discuss the some parabolic deformations.\;For any $w\in\sW_d$,\;denote by $\ext^1_{\bL_I^w}(\Dpik,\Dpik)\subseteq \ext^1(\Dpik,\Dpik)$ the subspace of parabolic deformations with respect to the parabolic filtration $D_{I^{w(1)}}-D_{I^{w(2)}}-\cdots-D_{I^{w(d)}}$,\;i.e.,\;$\widetilde{D}\in \ext^1(\Dpik,\Dpik)$ belongs to $\ext^1_{\bL_I^w}(\Dpik,\Dpik)$ iff 
\[\widetilde{D}=\widetilde{D}_{I^{w(1)}}-\widetilde{D}_{I^{w(2)}}-\cdots-\widetilde{D}_{I^{w(d)}},\;\]
where $\widetilde{D}_{I^{w(i)}}$ is a deformation of ${D}_{I^{w(i)}}$ over $E[\epsilon]/\epsilon^2$ for $1\leq i\leq d$.\;We also consider another subspace $\ext^{1,\ast}_{\bL_I^w}(\Dpik,\Dpik)\subseteq \ext^1_{\bL_I^w}(\Dpik,\Dpik)$ such that 
\[\widetilde{D}=\widetilde{D}_{I^{w(1)}}-\widetilde{D}_{I^{w(2)}}-\cdots-\widetilde{D}_{I^{w(d)}},\;\]
with $\widetilde{D}_{I^{w(i)}}\cong {D}_{I^{w(i)}}\otimes_{\cR_{E,L}}\cR_{E[\epsilon]/\epsilon^2}(\psi_i)$ for $\psi_i\in \homo(L^{\times},E)$.\;Sending $\widetilde{D}$ to the parameters $(\psi_i)$  induces a map
\begin{equation}
	\kappa_w:\ext^{1,\ast}_{\bL_I^w}(\Dpik,\Dpik)\rightarrow \homo(\bZ_I^w(L),E).\;
\end{equation}
with $\bZ_I^w\cong \BG_m^d$.\;


We define a subspace $\homo_w(\bZ_I^w(L),E)$ of $\homo(\bZ_I^w(L),E)$ as:
\begin{equation}
	\homo_w(\bZ_I^w(L),E):=\left\{(\psi_i)\in \homo(\bZ_I^w(L),E)|\psi_j-\psi_i\in E\val_p,\text{if\;} i<j,w(j)=w(i)+1\right\}.
\end{equation}
By definition,\;the subspace $\homo_{\sm}(\bZ_I^w(L),E)\subseteq \homo(\bZ_I^w(L),E)$ of smooth characters is contained in $\homo_w(\bZ_I^w(L),E)$.\;It is easy to see that $\dim_E\homo_w(\bZ_I^w(L),E)=2d-|w(R_d^+)\cap \Delta_d|$.\;



Recall that $\ext^1_g(\Dpik,\Dpik)\subseteq \ext^1(\Dpik,\Dpik)$ the subspace of de rham deformations.\;

\begin{pro}
	\begin{itemize}
		\item[(1)] $\dim_E\ext^1(\Dpik,\Dpik)=1+n^2$ and $\dim_E\ext^1_{\bL_I^w}(\Dpik,\Dpik)=1+\sum_{i\leq j}n_{w(i)}n_{w(j)}=1+\dim \bP_I^w=1+\dim \bP_I$ for all $w\in \sW_d$.\;
		\item[(2)] $\dim_E\ext^{1,\flat}_{\bL_I^w}(\Dpik,\Dpik)=1+\sum_{i< j}n_{w(i)}n_{w(j)}+d=1+\dim \bN_I+d$ for all $w\in \sW_d$.\;
		\item[(3)] For $w\in\sW_n$,\;$\kappa_w$ is a surjective onto $\homo_w(\bZ_I^w(L),E)$.\;
	\end{itemize}
\end{pro}
\begin{proof}
	We first prove Part $(1)$.\;$\dim_E\ext^1(\Dpik,\Dpik)=1+n^2$ is well known.\;We next compute $\dim_E\ext^1_{\bL_I^w}(\Dpik,\Dpik)$. We have a natural exact sequence of $(\varphi,\Gamma)$-module over $\cR_{E,L}$,\;
	\[0\rightarrow\homo_{\cR_{E,L}}(D_{I^{w(d)}},\Dpik)\rightarrow \EndO_{\bL_I^w}(\Dpik,\Dpik)\rightarrow \EndO_{\bL_I^w}(\Dpik_1^{n-1},\Dpik_1^{n-1})\rightarrow 0,\]
	where $\EndO_{\cF_w}(\Dpik,\Dpik):=\{f\in \EndO(\Dpik,\Dpik)| f(\fil^i_{\cF_w}(\Dpik))\subseteq \fil^i_{\cF_w}(\Dpik)\}$.\;Therefore,\;we see that $\EndO_{\cF_w}(\Dpik,\Dpik)$ is a $(\varphi,\Gamma)$-module over $\cR_{E,L}$ of rank $n+(n-1)+\cdots+1=\frac{n(n+1)}{2}$,\;and it is isomorphic to an successive extension of $(\varphi,\Gamma)$-modules $\Dpik_1^{i}\otimes_{\cR_{E,L}}\cR_{E,L}(\psi_{w(i)}z^{\bh_i})$ with $i=1,\cdots,n$.\;\end{proof}



For $w\in\sW_d$,\;Let $\ext^{1,\flat}_g(\Dpik,\Dpik)\subseteq\ext^{1,\flat}_{\bL_I^w}(\Dpik,\Dpik)$ be the preimage of the subspace $\homo_{\sm}(\bZ_I^w(L),E)\subseteq \homo_w(\bZ_I^w(L),E)$ of smooth characters via $\kappa_w$.\;Similar to the proof of [Lemma 2.12]
,\;we see that $\ext^{1,\flat}_g(\Dpik,\Dpik)$ is independent on the choice of $w$.\;In particular,\;we see that $\dim_E\ext^{1,\flat}_{g}(\Dpik,\Dpik)=1+\dim \bN_I$ for all $w\in \sW_d$.\;



Let $\ext^{1,\flat}_0(\Dpik,\Dpik):=\ker\kappa_w|_{\ext^{1,\flat}_g(\Dpik,\Dpik)}\subseteq\ext^{1,\flat}_g(\Dpik,\Dpik)$,\;which is not dependent on the choice of $w$.\;The above proposition also deduce that
\[\dim_E\ext^{1,\flat}
_0(\Dpik,\Dpik)=1+\dim \bN_I-d.\]

Furthermore,\;let $\homo_{g'}(\bZ_I^w(L),E)$ to be the subspace of characters $\psi=\psi_1+\psi_0\flat\det$ for $\psi_0\in\homo(\bQ_p,E)$ and $\psi_1\in \homo_{\sm}(\bZ_I^w(L),E)$.\;By definition,\;$\homo_{g'}(\bZ_I^w(L),E)\subseteq \homo_w(\bZ_I^w(L),E)$ for all $w$.\;Let $\ext^{1,\flat}_{g'}(\Dpik,\Dpik)$ be the preimage of $\homo_{g'}(\bZ_I^w(L),E)$ under $\kappa_w$.\;In particular,\;we have 
\[\dim_E\ext^{1,\flat}_{g'}(\Dpik,\Dpik)=2+\dim \bN_I.\;\]




For $\ext^{\ast}_{?}(\Dpik,\Dpik)\subseteq \ext^1(\Dpik,\Dpik)$ with $\ext^1_0(\Dpik,\Dpik)\subseteq\ext^1_{?}(\Dpik,\Dpik)$ and $\ast\in\{1,(1,\flat)\}$,\;we put
\[\overline{\ext}^{\ast}_{?}(\Dpik,\Dpik)=\ext^{\ast}_{?}(\Dpik,\Dpik)/\ext^1_0(\Dpik,\Dpik).\;\]
So we have 
\[\overline{\ext}^1_{g}(\Dpik,\Dpik)\xrightarrow{\sim }\homo_{\sm}(\bZ_I^w(L),E),\;\overline{\ext}^1_{w}(\Dpik,\Dpik)\twoheadrightarrow\homo_w(\bZ_I^w(L),E).\]
%\begin{rmk}
%	Note that $\dim_E\overline{\ext}^1_{w}(\Dpik,\Dpik)=2n$.\;If $w(R^+)\cap \Delta=\emptyset$,\;we still have an isomorphism $\overline{\ext}^1_{w}(\Dpik,\Dpik)\xrightarrow{\sim}\homo_w(\bZ_I^w(L),E)$.\;
%\end{rmk}





\subsubsection{Infinite fern for special semistable case}

Recall the two exact sequences:
\begin{equation}
	\begin{aligned}
		0\rightarrow D_1\rightarrow\Dpik\rightarrow D_{I^1}\rightarrow 0,\\
		0\rightarrow D'_{I^1}\rightarrow\Dpik\rightarrow  C_1\rightarrow 0.\;
	\end{aligned}
\end{equation}
Let $\cF$ (resp.,\;$\cG$) be the parabolic filtration $D_1\subseteq\Dpik$ (resp.,\;$D'_{I^1}\subseteq\Dpik$).\;


In this section,\;we assume that $d\geq 3$.\;We have natural surjections:
\begin{equation}
	\begin{aligned}
		\kappa_{\cF}:\ext^{1}_{\cF}(\Dpik,\Dpik)\twoheadrightarrow\ext^{1}(D_1,D_1)\times\ext^{1}(D_{I^1},D_{I^1}),\\
		\kappa_{\cG}:\ext^{1}_{\cG}(\Dpik,\Dpik)\twoheadrightarrow\ext^{1}(C_1,C_1)\times\ext^{1}(D'_{I^1},D'_{I^1}).\;
	\end{aligned}
\end{equation}
\begin{rmk}
	Using $\cF$ and $\cG$,\;there is no obstruction to the parabolic deformations $\cF$ and $\cG$.\;
\end{rmk}
For any $\iota\in\homo(D_1,C_1)$,\;we have the following pull-back (resp.,\;pull-forward) maps:
\begin{equation}
	\begin{aligned}
		\iota^{-}:\ext^1(C_1,D_1)\rightarrow \ext^1(D_1,D_1),\\
		\iota^{+}:\ext^1(C_1,D_1)\rightarrow \ext^1(C_1,C_1).\;
	\end{aligned}
\end{equation}
Put
\[\ext_\iota^1(D_1,D_1)=\iota^{-}(\ext^1(C_1,D_1)),\ext_\iota^1(C_1,C_1)=\iota^{+}(\ext^1(C_1,D_1)).\]
The higher intertwining concerns the intersection $ \ext^{1}_{\cF}(\Dpik,\Dpik)\cap \ext^{1}_{\cG}(\Dpik,\Dpik)$ in $\ext^{1}(\Dpik,\Dpik)$.\;



By the same argument as [Theorem 2.32]{DingParameter},\;we get the higher  for the special semistable case:
\begin{thm}
	For any $\widetilde{D}\in \ext^{1,\flat}_{\cF}(\Dpik,\Dpik)$ with $\kappa_{\cF}(\widetilde{D})=(\widetilde{D_1},\psi_n)$.\;The followings are equivalent:
	\begin{itemize}
		\item[(a)] $\widetilde{D}\in \ext^{1}_{\cF}(\Dpik,\Dpik)\cap \ext^{1}_{\cG}(\Dpik,\Dpik)$,\;
		\item[(b)] $\widetilde{D_1}\otimes_{\cR_{E[\epsilon]/\epsilon^2}}(\cR_{E[\epsilon]/\epsilon^2}(1-\psi_n\epsilon))\in \ext_{\iota_{\Dpik}}^1(D_1,D_1)$.\;
	\end{itemize}
	When $\widetilde{D}\in \ext^{1,\flat}_{\cF}(\Dpik,\Dpik)\cap \ext^{1,\flat}_{\cG}(\Dpik,\Dpik)$,\;then there exists $M\in\ext^1(C_1,D_1)$ such that
	\begin{equation}
		\begin{aligned}
			&\widetilde{D_1}=\iota_{\Dpik}^{-}(M)\otimes_{\cR_{E[\epsilon]/\epsilon^2}}(\cR_{E[\epsilon]/\epsilon^2}(1+\psi_n\epsilon)),\;\\
			&\widetilde{C_1}:=\kappa_{\cG,1}(\widetilde{D})=\iota_{\Dpik}^{+}(M)\otimes_{\cR_{E[\epsilon]/\epsilon^2}}(\cR_{E[\epsilon]/\epsilon^2}(1+\psi_n\epsilon)).\;
		\end{aligned}
	\end{equation}
	Moreover,\;we have $\dim_E \ext^{1}_{\cF}(\Dpik,\Dpik)\cap \ext^{1}_{\cG}(\Dpik,\Dpik)=1+(n^2-2n+2)$ and thus the following map is surjective:
	\[{\ext}^{1}_{\cF}(\Dpik,\Dpik)\oplus {\ext}^{1}_{\cG}(\Dpik,\Dpik)\twoheadrightarrow {\ext}^1(\Dpik,\Dpik).\]
\end{thm}

\begin{pro}
	We have a natural surjection
	\[\bigoplus_{w\in \sW_n}\overline{\ext}^1_{w}(\Dpik,\Dpik)\twoheadrightarrow \overline{\ext}^1(\Dpik,\Dpik).\]
\end{pro}



\subsubsection{Higher interwining for special semistable case}
Recall the two exact sequences:
\begin{equation}
	\begin{aligned}
		0\rightarrow D_1\rightarrow\Dpik\rightarrow D_{I^1}\rightarrow 0,\\
		0\rightarrow D'_{I^1}\rightarrow\Dpik\rightarrow  C_1\rightarrow 0.\;
	\end{aligned}
\end{equation}
Let $\cF$ (resp.,\;$\cG$) be the parabolic filtration $D_1\subseteq\Dpik$ (resp.,\;$D'_{I^1}\subseteq\Dpik$).\;


In this section,\;we assume that $d\geq 3$.\;We have natural surjections:
\begin{equation}
	\begin{aligned}
		\kappa_{\cF}:\ext^{1,\flat}_{\cF}(\Dpik,\Dpik)\twoheadrightarrow\ext^{1}(D_1,D_1)\times\ext^{1,\flat}(D_{I^1},D_{I^1})\cong \ext^{1}(D_1,D_1)\times \homo(L^{\times},E),\\
		\kappa_{\cG}:\ext^{1,\flat}_{\cG}(\Dpik,\Dpik)\twoheadrightarrow\ext^{1}(C_1,C_1)\times\ext^{1,\flat}(D'_{I^1},D'_{I^1})\cong \ext^{1}(C_1,C_1)\times \homo(L^{\times},E).\;
	\end{aligned}
\end{equation}
\begin{rmk}
	Using such two triangulations,\;there is no obstruction to the parabolic deformations $\cF$ and $\cG$.\;
\end{rmk}

We have $\dim_E\ext^{1,\flat}_{\ast}(\Dpik,\Dpik)=1+(n-n_1)^2+n_1^2+n_1(n-n_1)=1+n^2+n_1^2-nn_1$ and $\dim_E\ext^{1,\flat}_{\ast}(\Dpik,\Dpik)=1+n^2+n_1^2-nn_1-(n_1^2-1)=2+n^2-nn_1$ for $\ast\in\{\cF,\cG\}$.\;Similar to the discussion in [Corollary 2.15],\;we have
\begin{equation}
	\begin{aligned}
		\kappa_{\cF}:\overline{\ext}^{1}_{\cF}(\Dpik,\Dpik)\xrightarrow{\sim}\overline{\ext}^1(D_1,D_1)\times\homo(L^{\times},E),\\
		\kappa_{\cG}:\overline{\ext}^{1}_{\cG}(\Dpik,\Dpik)\xrightarrow{\sim}\overline{\ext}^1(C_1,C_1)\times\homo(L^{\times},E),\;
	\end{aligned}
\end{equation}
by comparing the dimensions.\;



For any $\iota\in\homo(D_1,C_1)$,\;we have the following pull-back (resp.,\;pull-forward) maps:
\begin{equation}
	\begin{aligned}
		\iota^{-}:\ext^1(C_1,D_1)\rightarrow \ext^1(D_1,D_1),\\
		\iota^{+}:\ext^1(C_1,D_1)\rightarrow \ext^1(C_1,C_1).\;
	\end{aligned}
\end{equation}
Put
\[\ext_\iota^1(D_1,D_1)=\iota^{-}(\ext^1(C_1,D_1)),\ext_\iota^1(C_1,C_1)=\iota^{+}(\ext^1(C_1,D_1)).\]
By the same argument as [Theorem 2.32]{DingParameter},\;we get the higher intertwining for the special crystalline case:
\begin{thm}
	For any $\widetilde{D}\in \ext^{1,\flat}_{\cF}(\Dpik,\Dpik)$ with $\kappa_{\cF}(\widetilde{D})=(\widetilde{D_1},\psi_1)$.\;The followings are equivalent:
	\begin{itemize}
		\item[(a)] $\widetilde{D}\in \ext^{1}_{\cF}(\Dpik,\Dpik)\cap \ext^{1}_{\cG}(\Dpik,\Dpik)$,\;
		\item[(b)] $\widetilde{D_1}\otimes_{\cR_{E[\epsilon]/\epsilon^2}}(\cR_{E[\epsilon]/\epsilon^2}(1-\psi_1\epsilon))\in \ext_{\iota_{\Dpik}}^1(D_1,D_1)$.\;
	\end{itemize}
	When $\widetilde{D}\in \ext^{1,\flat}_{\cF}(\Dpik,\Dpik)\cap \ext^{1,\flat}_{\cG}(\Dpik,\Dpik)$,\;then there exists $M\in\ext^1(C_1,D_1)$ such that
	\begin{equation}
		\begin{aligned}
			&\widetilde{D_1}=\iota_{\Dpik}^{-}(M)\otimes_{\cR_{E[\epsilon]/\epsilon^2}}(\cR_{E[\epsilon]/\epsilon^2}(1+\psi_1\epsilon)),\;\\
			&\widetilde{C_1}:=\kappa_{\cG,1}(\widetilde{D})=\iota_{\Dpik}^{+}(M)\otimes_{\cR_{E[\epsilon]/\epsilon^2}}(\cR_{E[\epsilon]/\epsilon^2}(1+\psi_1\epsilon)).\;
		\end{aligned}
	\end{equation}
	Moreover,\;we have $\dim_E \ext^{1}_{\cF}(\Dpik,\Dpik)\cap \ext^{1}_{\cG}(\Dpik,\Dpik)=?$.\;
\end{thm}
Define
\begin{equation}
	\begin{aligned}
		V^{\flat}(D_1,C_1):=&\left(\overline{\ext}^{1,\flat}(D_1,D_1)\times\homo(L^{\times},E)\right)\oplus \left(\overline{\ext}^{1,\flat}(C_1,C_1)\times\homo(L^{\times},E)\right),\\
		&\xleftarrow{\sim,(\kappa_{\cF},\kappa_{\cG})}\overline{\ext}^{1,\flat}_{\cF}(\Dpik,\Dpik)\oplus \overline{\ext}^{1,\flat}_{\cG}(\Dpik,\Dpik).\;
	\end{aligned}
\end{equation}
Therefore,\;we get the following exact sequence:
\begin{equation}
	0\rightarrow \cL(\Dpik,D_1,C_1)\rightarrow  V(D_1,C_1)\rightarrow \overline{\ext}^1(\Dpik,\Dpik),\;
\end{equation}
where $\cL(\Dpik,D_1,C_1)$ is the kernel of the surjection $V(D_1,C_1)\twoheadrightarrow\overline{\ext}^1(\Dpik,\Dpik)$,\;by higher intertwining,\;it consists of $((x,\psi),(y,\psi))\in V(D_1,C_1)$ such that $((x,\psi),(y,\psi))$ comes from some $M\in\ext^1(C_1,D_1)$.\;

\subsection{Deformations of $\cF_{u}$ or $\cF_{J_0(u)}$ for $u\in \sW_s$}


Recall for any $u\in \sW_s$,\;we define a
subset $J_0(u)$ of $\Delta$ which is characterized uniquely in the following way:\;$l\in \Delta\backslash J_0(u)$ if and only $j(u(l))\neq j(u(l+1))$.\;Suppose that $\bL_{J_0(u)}=\GLN_{h_1}\times\GLN_{h_2}\times\cdots\times \GLN_{h_t}$.\;By the definition of $J_0(u)$,\;$\Dpik^{u}_j:=D_{h_1+\cdots+h_{j-1}+1}^{h_1+\cdots+h_{j}}$ is a Steinberg $(\varphi,\Gamma)$-module.\;This induces an $\bP_{J_0(u)}$ parabolic filtration $\cF_{J_0(u)}$ on $\Dpik$.\;In particular,\;we have $\cF_u\subseteq \cF_{S_0,u}$.\;

We get a map
\begin{equation}
	\kappa_{S_0,u}:=\kappa_{\cF_{S_0,u}}:\ext^{1,\ast}_{{S_0,u}}(\Dpik,\Dpik):=\ext^{1,\ast}_{\cF_{S_0,u}}(\Dpik,\Dpik)\rightarrow \homo(\bZ_{J_0(u)}(L),E).\;
\end{equation}
with $\bZ_{J_0(u)}\cong \BG_m^t$.\;We define a subspace $\homo_u(\bZ_{J_0(u)}(L),E)$ of $\homo(\bZ_{J_0(u)}(L),E)$ as:
\begin{equation}
	\homo_u(\bZ_{J_0(u)}(L),E):=\left\{(\psi_i)\in \homo(\bZ_{J_0(u)}(L),E)|\psi_j-\psi_i\in E\val_p,\text{if\;} i<j,\Dpik^u_{i}\succ \Dpik^u_{j}\right\}.
\end{equation}
By definition,\;the subspace $\homo_{\sm}(\bZ_{J_0(u)}(L),E)\subseteq \homo(\bZ_{J_0(u)}(L),E)$ of smooth characters is contained in $\homo_u(\bZ_{J_0(u)}(L),E)$.\;By \cite{2019DINGSimple},\;$\kappa_{S_0,u}$ is a surjective onto $\homo_u(\bZ_{J_0(u)}(L),E)$.\;


Let $\ext^{1,\ast}_{g,{S_0,u}}(\Dpik,\Dpik)\subseteq\ext^{1,\ast}_{{S_0,u}}(\Dpik,\Dpik)$ is the preimage of $\homo_{\sm}(\bZ_{J_0(u)}(L),E)$ of via $\kappa_w$,\;and let $\ext^{1,\ast}_{g',S_0,u}(\Dpik,\Dpik)$ be the preimage of $\homo_{g'}(\bZ_{J_0(u)}(L),E)$ under $\kappa_{S_0,u}$.\;

Similar to the proof of Proposition \ref{dimesionforParafilgeneric},\;we have 
\begin{lem}Suppose $u\in  \sW_s$ satisfies that $\Dpik^u_{i}\succ \Dpik^u_{j}$ for any $ i<j$.\;Then $\ext^{1}_{g}(\Dpik,\Dpik)\subseteq \ext^{1}_{{S_0,u}}(\Dpik,\Dpik)$ and  $\ext^{1,\ast}_{g,{S_0,u}}(\Dpik,\Dpik)=\ext^{1,\ast}_{{S_0,u}}(\Dpik,\Dpik)\cap \ext^{1}_{g}(\Dpik,\Dpik)$.\;
\end{lem}
Similar to lemma \ref{smhomoindepent},\;we obtain
\begin{lem}
	The following diagram commutes
	\begin{equation}
		\xymatrix{ \ext^{1,\ast}_{g,{S_0,u_1}}(\Dpik,\Dpik)\cap \ext^{1,\ast}_{g,{S_0,u_2}}(\Dpik,\Dpik)\ar@{=}[d] \ar[r]_{\hspace{30pt}\kappa_{S_0,u_1}} & \homo_{\sm}(\bZ_{J_0(u_1)}(L),E) \ar[d]^{\sim}_{u_2^{-1}u_1}\\
			\ext^{1,\ast}_{g,{S_0,u_2}}(\Dpik,\Dpik)\cap\ext^{1,\ast}_{g,{S_0,u_1}}(\Dpik,\Dpik) \ar[r]_{\hspace{30pt}\kappa_{S_0,u_2}} & \homo_{\sm}(\bZ_{J_0(u_2)}(L),E) },
	\end{equation}
	for $u_1,u_2\in\sW_{s}$.\;
\end{lem}Put $\ext^{1,\ast}_{0,{S_0,u}}(\Dpik,\Dpik):=\ker\kappa_{S_0,u}$.\;Clearly,\;we have inclusions
\[\ext^{1,\ast}_{0,{S_0,u}}(\Dpik,\Dpik)\subseteq\ext^{1,\ast}_{g,{S_0,u}}(\Dpik,\Dpik)\subseteq \ext^{1,\ast}_{g',S_0,u}(\Dpik,\Dpik).\]



In particular,\;if $u\in \sW_{s}$,\;we have $J_0(u)=u(J_0)$ and $t=s$.\;In this case,\;we get the map
\begin{equation}
	\kappa_{S_0,u}:=\kappa_{\cF_{S_0,u}}:\ext^{1,\ast}_{{S_0,u}}(\Dpik,\Dpik):=\ext^{1,\ast}_{\cF_{S_0,u}}(\Dpik,\Dpik)\rightarrow \homo(\bZ_{J_0}^u(L),E).\;
\end{equation}
with $\bZ_{u(J_0)}\cong \bZ_{J_0}^u\cong \BG_m^s$.\;The subspace $\homo_u(\bZ_{J_0}^u(L),E)$ can be written as another way:
\begin{equation}
	\homo_u(\bZ_{J_0}^u(L),E):=\left\{(\psi_i)\in \homo(\bZ_{J_0}^u(L),E)|\psi_j-\psi_i\in E\val_p,\text{if\;} i<j,E_{u(i)}\succ E_{u(j)}\right\}.
\end{equation}
By definition,\;the subspace $\homo_{\sm}(\bZ_{J_0}^u(L),E)\subseteq \homo(\bZ_{J_0}^u(L),E)$ of smooth characters is contained in $\homo_u(\bZ_{J_0}^u(L),E)$.\;In particular,\;we have 
\begin{equation}
	\begin{aligned}
		&\dim_E\homo_1(\bZ_{J_0}^{1}(L),E)=2s-|I_0|,\;\\
		&\dim_E\homo_{w_{0,s,n}}(\bZ_{J_0}^{w_{0,s,n}}(L),E)=2s.\;
	\end{aligned}
\end{equation}



\begin{pro}We have the following facts.\;
	\begin{itemize}
		\item[(1)] We have $\dim_E\ext^{1}_{u}(\Dpik,\Dpik)=1+\frac{n(n+1)}{2}$,\;$\dim_E\ext^{1}_{S_0,u}(\Dpik,\Dpik)=1+\dim \bP_{J_0}$ and  $\dim_E\ext^{1,\ast}_{S_0,u}(\Dpik,\Dpik)=1+s+\dim \bN_{J_0}$.\;
		\item[(2)]$\kappa_{S_0,u}$ is a surjection onto $\homo_u(\bZ_{J_0}^u(L),E)$.\;
	\end{itemize}
\end{pro}
\begin{proof}
	
	We compute $\dim_E\ext^{1}_{u}(\Dpik,\Dpik)$ firstly.\;We follow the route in the proof of
	\cite[Proposition 3.6]{2019DINGSimple}.\;Let $\EndO_{\cF_u}(\Dpik)$ be the saturated $(\varphi,\Gamma)$-submodule of $\EndO(\Dpik)$ that preserves the triangulation $\cF_u$.\;Then we have a natural exact sequence of 
	$(\varphi,\Gamma)$-modules over $\cR_{E,L}$:
	\[0\rightarrow \homo_{\cR_{E,L}}(\cR_{E,L}(\delta_n),\Dpik
	)\rightarrow \EndO_{\cF_u}(\Dpik)\rightarrow \EndO_{\cF_u|_{\Dpik_1^{n-1}}}(\Dpik_1^{n-1})\rightarrow 0.\]
	Note that $\hH^0_{(\varphi,\Gamma)}(\homo_{\cR_{E,L}}(\cR_{E,L}(\delta_n),\Dpik))=0$ (since thata $\delta_n$ has Hodge-Tate weight $\bh_n$ and we only have $\cR_{E,L}(\delta_nz^{\bh_1-\bh_n})\hookrightarrow\Dpik$).\;By induction on $n$,\;we show that $\hH^0_{(\varphi,\Gamma)}(\EndO_{\cF_u}(\Dpik))\xrightarrow{\sim}\hH^0_{(\varphi,\Gamma)}(\EndO_{\cF_u|_{\Dpik_1^{n-1}}}(\Dpik_1^{n-1}))\cong E$.\;We are going to show that $\hH^2_{(\varphi,\Gamma)}(\EndO_{\cF_u}(\Dpik))=0$.\;By induction on $n$,\;we only need to show that $\hH^2_{\cR_{E,L}}(\homo_{\cR_{E,L}}(\cR_{E,L}(\delta_n),\Dpik
	))=0$.\;By Tate duality,\;we have isomorphism
	\[\hH^2_{(\varphi,\Gamma)}(\homo_{\cR_{E,L}}(\cR_{E,L}(\delta_n),\Dpik
	))\cong \hH^0_{(\varphi,\Gamma)}(\homo_{\cR_{E,L}}(\Dpik,
	\cR_{E,L}(\delta_n))).\]
	Note that $\homo_{\cR_{E,L}}(\cR_{E,L}(\delta_n),\Dpik
	)$ only depends on $\cR_{E,L}(\delta_n)$ and is independent on any triangulation on $\Dpik$.\;If for any subquotient $\cR_{E,L}(\delta')$,\;
	$\hH^2_{(\varphi,\Gamma)}(\homo_{\cR_{E,L}}(\cR_{E,L}(\delta_n),\cR_{E,L}(\delta')))=0$,\;we get the result.\;Otherwise,\;we assume that $\Dpik_1^{n-1}\twoheadrightarrow $ (with triangulation $\cF_u|_{\Dpik_1^{n-1}}\cR_{E,L}(\delta')$) such that $\hH^2_{(\varphi,\Gamma)}(\homo_{\cR_{E,L}}(\cR_{E,L}(\delta_n),\cR_{E,L}(\delta')))\neq0$,\;now we return to the case in \cite[Proposition 3.6]{2019DINGSimple},\;we also get the result.\;
	
	For $u\in \sW_{s}$,\;replacing the rank-$1$ object $\cR_{E,L}(\delta_n)$ in above discussion by the Steinberg block $E_{u(s)}$ (and changing $\cF_u$ to $\cF_{S_0,u}$,\;$\Dpik_1^{n-1}$ to $E_{u(1)}-\cdots-E_{u(s-1)}$),\;we can still prove $\hH^0_{(\varphi,\Gamma)}(\EndO_{\cF_{S_0,u}}(\Dpik))\cong E$ and $\hH^2_{(\varphi,\Gamma)}(\EndO_{\cF_{S_0,u}}(\Dpik))=0$ (note that the $\hH^2$-group between two Steinberg blocks $E_{u(i)}\succ E_{u(j)}$ can be finally reduced into the case in \cite[Proposition 3.6]{2019DINGSimple}).\;In particular,\;we obtain $\dim_E\ext^{1}_{S_0,u}(\Dpik,\Dpik)=1+\dim \bP_{J_0}$ and  $\dim_E\ext^{1,\ast}_{S_0,u}(\Dpik,\Dpik)=1+s+\dim \bN_{J_0}$.\;
	
	
	
	For Part $(2)$,\;write $\Dpik=D_1^u-\cdots-D_{h_u}^u$.\;We prove it by induction.\;Assume that $N_1:=D_1^u-\cdots-D_{h_u-1}^u$ satisfies the conclusion.\;There exists at most one $1\leq q\leq h_u-1$ such that $\Dpik^u_{p}\succ \Dpik^u_{h_u}$.\;Let $E_{u(q)}$ be the cosocle of $\Dpik^u_{p}$ and $\cR_{E,L}(\delta)$ be the socle of $\Dpik^u_{h_u}$.\;Let $N_2=[D_{1}^u-\cdots-D_{q}^u]\hookrightarrow N_1$ and let $N_1\rightarrow N_2'=[E_{u(q)}-\cdots-D_{h_u-1}^u]\twoheadrightarrow N_3':=N_2'/E_{u(q)}=[D_{q+1}^u-\cdots-D_{h_u-1}^u]$.\;W
	e consider the following commutative diagram:
	\begin{equation}
		\xymatrix{
			\ext^1(\cR_{E,L}(\delta),N'_2) \ar@{=}[d]  &\times & \ext^1(N'_2,N'_3)\ar[r]^{\cup}  &  \ext^2(\cR_{E,L}(\delta),N'_3)=0\\
			\ext^1(\cR_{E,L}(\delta),N_2') \ar@{=}[d] &\times & \ext^1(N_2',N_2')\ar[r]^{\cup} \ar[u]^{\kappa_1}  &  \ext^2(\cR_{E,L}(\delta),N_2') \ar@{=}[d] \ar[u]\\
			\ext^1(\cR_{E,L}(\delta),N'_2) \ar[d]^{g_1}    &\times & \ext^1(N'_2,E_{u(q)})\ar[r]^{\cup} \ar[u]^{\kappa_1'} &  \ext^2(\cR_{E,L}(\delta),E_{u(q)}) \ar@{=}[d] \\
			\ext^1(\cR_{E,L}(\delta),N'_3)   &\times & \ext^1(N'_3,E_{u(q)})\ar[r]^{\cup} \ar[u]^{\kappa_1} &  \ext^2(\cR_{E,L}(\delta),E_{u(q)}). }
	\end{equation}
	Let $D''=\Dpik/\ker(N_1\rightarrow N_2')$.\;Let $\cL_{\mathrm{FM}}(D'':N_2'):=(E[D''])^{\perp}$ (through the top cup-product),\;and let $l_{\mathrm{FM}}(D'':N_2'):=(E[D''])^{\perp}\subseteq \ext^1(N_2',E_{u(q)})$ (via the second row).\;We claim that 
	$\cL_{\mathrm{FM}}(D'':N_2'):=(E[D''])^{\perp}$ (resp.,\;$l_{\mathrm{FM}}(D'':N_2')$) has co-dimension $1$ in $\ext^1(N_2',N_2')$ (resp.,\;$\ext^1(N_2',E_{u(q)})$).\;
	
	We describe these subspaces explicitly.\;It easy to see that $\ext^2(N'_2,E_{u(q)})=0$ so that $\kappa_1$ is surjective.\;On the other hand,\;$\kappa_1'$ is injective since $\homo(N_2',N_2')\xrightarrow{\sim }\homo(N_2',N_3')$ is an ismorohism.\;Therefore,\;we have a short exact sequence of $E$-vector spaces:
	\[0\rightarrow l_{\mathrm{FM}}(D'':N_2')\rightarrow \cL_{\mathrm{FM}}(D'':N_2')\rightarrow \ext^1(N'_2,N'_3)\rightarrow 0.\]
	Note that $\homo(N'_2,E_{u(q)})=0$,\;we get that $\ker(\kappa_1)\cong \homo(E_{u(q)},E_{u(q)})\cong E$,\;so that $\dim_E\mathrm{Im}(\kappa_1)=\rk(N_3')-1$.\;Consider the following commutative diagram:
	\begin{equation}\label{N2N3ANDRANK1}
		\xymatrix{
			\ext^1(\cR_{E,L}(\delta),E_{u(q)}) \ar[d]^{g_1'} &\times & \ext^1(E_{u(q)},E_{u(q)})\ar[r]^{\cup}   &  \ext^2(\cR_{E,L}(\delta),E_{u(q)}) \ar@{=}[d] \\
			\ext^1(\cR_{E,L}(\delta),N'_2) \ar[d]^{g_1}    &\times & \ext^1(N'_2,E_{u(q)})\ar[r]^{\cup} \ar[u]^{\kappa_1''} &  \ext^2(\cR_{E,L}(\delta),E_{u(q)}) \ar@{=}[d] \\
			\ext^1(\cR_{E,L}(\delta),N'_3)   &\times & \ext^1(N'_3,E_{u(q)})\ar[r]^{\cup} \ar[u]^{\kappa_1} &  \ext^2(\cR_{E,L}(\delta),E_{u(q)}). }
	\end{equation}
	Since $\ext^2(N'_3,E_{u(q)})=0$,\;$\kappa_1''$ is surjective.\;Since $\homo(\cR_{E,L}(\delta),N'_3)=0$,\;$g_1''$ is injective.\;Observe the following commutative diagram:
	\begin{equation}
		\xymatrix{
			\ext^1(\cR_{E,L}(\delta),N'_3) \ar@{=}[d]   &\times & \ext^1(N'_3,\cR_{E,L}(\delta\epsilon))\ar[r]^{\cup}  &  \ext^2(\cR_{E,L}(\delta),\cR_{E,L}(\delta\epsilon))  \\
			\ext^1(\cR_{E,L}(\delta),N'_3)   &\times & \ext^1(N'_3,E_{u(q)})\ar[r]^{\cup} \ar[u]^{\eta} &  \ext^2(\cR_{E,L}(\delta),E_{u(q)})\ar[u]^{\eta'}_{\sim}, }
	\end{equation}
	where $\eta$ and $\eta'$ are induced by the injection  $E_{u(q)}\hookrightarrow \cR_{E,L}(\delta\epsilon)$ of $(\varphi,\Gamma)$-modules.\;The top cup-product is perfect by the Tate-duality.\;On the other hand,\;by \cite[Lemma 5.1.1]{breuil2020probleme},\;we have
	\begin{equation}
		\begin{aligned}
			&\;\homo((N'_3)^{\vee}\otimes_{\cR_{E,L}}E_{u(q)},(N'_3)^{\vee}_1\otimes_{\cR_{E,L}}\cR_{E,L}(\delta\epsilon))\\
			\cong&\; \hH^0(\gal_p,W_{\dr}^+((N'_3)^{\vee}\otimes_{\cR_{E,L}}\cR_{E,L}(\delta\epsilon))/W_{\dr}^+((N'_3)^{\vee}\otimes_{\cR_{E,L}}E_{u(q)})).\;
		\end{aligned}
	\end{equation}
	Then we oabtain $\dim_E\ker(\eta)=\rk(N'_3)$ and thus $\eta=0$,\;which implies that $\ext^1(N'_3,E_{u(q)})\subseteq (Eg_1([D'']))^{\perp}$ and $\kappa_1''(\ext^1(N'_3,E_{u(q)}))\subseteq l_{\mathrm{FM}}(D'':N_2') $.\;Observe the following commutative diagram:
	\begin{equation}
		\xymatrix{
			\ext^1(\cR_{E,L}(\delta),N'_2) \ar@{=}[d]   &\times & \ext^1(N'_2,\cR_{E,L}(\delta\epsilon))\ar[r]^{\cup}  &  \ext^2(\cR_{E,L}(\delta),\cR_{E,L}(\delta\epsilon))  \\
			\ext^1(\cR_{E,L}(\delta),N'_2)   &\times & \ext^1(N'_2,E_{u(q)})\ar[r]^{\cup} \ar[u]^{\eta_1} &  \ext^2(\cR_{E,L}(\delta),E_{u(q)})\ar[u]^{\eta'_1}_{\sim}, }
	\end{equation}
	where $\eta$ and $\eta'$ are induced by the injection  $E_{u(q)}\hookrightarrow \cR_{E,L}(\delta\epsilon)$ of $(\varphi,\Gamma)$-modules.\;The top cup-product is perfect by the Tate-duality.\;On the other hand,\;by \cite[Lemma 5.1.1]{breuil2020probleme},\;we have
	\begin{equation}
		\begin{aligned}
			&\;\homo((N'_3)^{\vee}\otimes_{\cR_{E,L}}E_{u(q)},(N'_3)^{\vee}_1\otimes_{\cR_{E,L}}\cR_{E,L}(\delta\epsilon))\\
			\cong&\; \hH^0(\gal_p,W_{\dr}^+((N'_2)^{\vee}\otimes_{\cR_{E,L}}\cR_{E,L}(\delta\epsilon))/W_{\dr}^+((N'_2)^{\vee}\otimes_{\cR_{E,L}}E_{u(q)}))\cong E^{\rk(N'_2)-1}.\;
		\end{aligned}
	\end{equation}
	Then we oabtain $\dim_E\ker(\eta_1)=\rk(N'_2)-2$.\;Therefore,\;$\ker(\eta_1)\subseteq\ext^1(\cR_{E,L}(\delta),N'_2)^{\perp}$.\;
	
	Let $e_{\mathrm{cris}}$ be the basis of $\ext^1_e(\cR_{E,L}(\delta),E_{u(q)}) $ (note that it is $1$-demensional) and $e_{\mathrm{cris}},e_{g}$ be the basis of $\ext^1(\cR_{E,L}(\delta),E_{u(q)})$.\;Recall that $\ext^1(E_{u(q)},E_{u(q)})=E\langle \val_p,\log_p\rangle$.\;Extend $\val_p,\log_p$ to a basis $\val_p,\log_p,e_1,\cdots,e_{\rk(N_2')-2}$ of $\ext^1(\cR_{E,L}(\delta),N'_2)$.\;Since the top bottom cup-product in 
	(\ref{N2N3ANDRANK1}) is prefect,\;we see that $E\langle \val_p,\log_p\rangle\subseteq \ext^1(\cR_{E,L}(\delta),N'_3)^{\perp}$ and $\val_p,\log_p$ are not contained in $\ext^1(\cR_{E,L}(\delta),N'_2)^{\perp}$ (note that the injection $g_1'$ in (\ref{N2N3ANDRANK1})),\;this implies $\ker(\eta_1)=E\langle e_1,\cdots,e_{\rk(N_2')-2}\rangle=\ext^1(N'_3,E_{u(q)})$.\;Note that $E[D'']=E(e_{\mathrm{cris}}+L)$ for some vector $L\in \ext^1(\cR_{E,L}(\delta),N'_3)$.\;So $\kappa_1''(l_{\mathrm{FM}}(D'':N_2'))$ equals to $E\val_p$ by using the fact that $\langle\val_p,L\rangle=\langle\log_p,L\rangle=0$ and $E\val_p=(Ee_{\mathrm{cris}})^\perp$ via the top bottom cup-product in 
	(\ref{N2N3ANDRANK1}).\;In particular,\;$\dim_El_{\mathrm{FM}}(D'':N_2')=(\rk N_3'-1)+1=1\rk N_2'-1$ and $\kappa_1''(l_{\mathrm{FM}}(D'':N_2'))=E\val_p\subseteq \ext^1(E_{u(q)},E_{u(q)})$.\;
	
	Finally,\;note that the so-called Colmez-Greenberg-Stevens (see \cite[Theorem 2.7]{HigherLinvariantsGL3(Qp)}) formula implies:\;for the tuple $(\widetilde{N_2'},\psi)\in \ext^{1}(N'_2,N_2')\times \homo(L^{\times},E)$,\;there exists a deformation  $\widetilde{N'_2}-\cR_{ E[\epsilon]/\epsilon^2}(\delta(1+\psi\epsilon))$ of $D''$ over $E[\epsilon]/\epsilon^2$  iff
	$\widetilde{N_2'}\otimes_{\cR_{E[\epsilon]/\epsilon^2}}(\cR_{E[\epsilon]/\epsilon^2}(1-\psi\epsilon))\in\cL_{\mathrm{FM}}(D'':N_2')$,\;thus $\psi_q-\psi_{h_u}\in E\val_p$.\;This completes the proof of $(2)$.\;
\end{proof}




Let $\ext^{1,\ast}_{g,{S_0,u}}(\Dpik,\Dpik)\subseteq\ext^{1,\ast}_{{S_0,u}}(\Dpik,\Dpik)$ is the preimage of the space $\homo_{\sm}(\bZ^u_{J_0}(L),E)$ of via $\kappa_w$,\;and let $\ext^{1,\ast}_{g',S_0,u}(\Dpik,\Dpik)$ be the preimage of $\homo_{g'}(\bZ^u_{J_0}(L),E)$ under $\kappa_{S_0,u}$.\;Put $\ext^{1,\ast}_{0,{S_0,u}}(\Dpik,\Dpik):=\ker\kappa_{S_0,u}$.\;We set $\overline{\ext}^{1,0}_{{S_0,u}}(\Dpik,\Dpik):=\ext^{1,\ast}_{{S_0,u}}(\Dpik,\Dpik)/\ext^{1,\ast}_{0,{S_0,u}}(\Dpik,\Dpik)$,\;thus we get an isomorphism:
\[\kappa_{S_0,u}:\overline{\ext}^{1,0}_{{S_0,u}}(\Dpik,\Dpik)\xrightarrow{\sim}\homo_u(\bZ^u_{J_0}(L),E).\]
In the sequel,\;we set
\begin{equation}
	\ext^{1,\ast}_{g}(\Dpik,\Dpik):=\ext^{1,\ast}_{g,{J_0,1}}(\Dpik,\Dpik),\;\ext^{1,\ast}_{0,J_0}(\Dpik,\Dpik)=\ext^{1}_{0,{J_0,1}}(\Dpik,\Dpik).\;
\end{equation}
Note that
\begin{equation}
	\begin{aligned}
		&\dim_E\ext^{1,\ast}_{g}(\Dpik,\Dpik)=1+\dim \bN_{J_0}+|I_0|,\\
		&\dim_E\ext^{1,\ast}_{0}(\Dpik,\Dpik)=1+\dim \bN_{J_0}-s+|I_0|.
	\end{aligned}
\end{equation}



\begin{lem}
	We have $\ext^{1,\ast}_{g,{S_0,u}}(\Dpik,\Dpik)=\ext^{1,\ast}_{{S_0,u}}(\Dpik,\Dpik)\cap \ext^{1,\ast}_{g}(\Dpik,\Dpik)$ and $\ext^{1,\ast}_{0,{S_0,u}}(\Dpik,\Dpik)=\ext^{1,\ast}_{{S_0,u}}(\Dpik,\Dpik)\cap \ext^{1,\ast}_{0}(\Dpik,\Dpik)$.\;
\end{lem}
\begin{proof}
	By definition,\;we have $\ext^{1,\ast}_{{S_0,u}}(\Dpik,\Dpik)\cap \ext^{1,\ast}_{g}(\Dpik,\Dpik)\subseteq \ext^{1,\ast}_{g,{S_0,u}}(\Dpik,\Dpik)$.\;Conversely,\;for any $(\psi_i)\in \homo_{\sm}(\bZ_{J_0}^u(L),E)$,\;we will show that \[\widetilde{D}=\widetilde{E}_{u(1)}-\widetilde{E}_{u(2)}-\cdots-\widetilde{E}_{u(t)},\;\]
	with $\widetilde{E}_{u(i)}\cong E_{u(i)}\otimes_{\cR_{E,L}}\cR_{E[\epsilon]/\epsilon^2}(1+\psi_i\epsilon)$ is de Rham.\;For $i<j$,\;we first show that any element in
	\[\hH^1\left(\gal_{p},W(E_{u(j)}^{\vee}\otimes_{\cR_{E,L}}{E_{u(i)}})\otimes_{\cR_{E,L}}\cR_{E[\epsilon]/\epsilon^2}(1+(\psi_j-\psi_i)\epsilon)\right)\]
	are de Rham.\;If $E_{u(i)}\succ E_{u(j)}$ and $\psi_j=\psi_i$ for $i<j$,\;the de Rhamness follows from  the same argument as in the proof of \cite[Lemma 3.11 (2)]{He20222} (replaced by the block version).\;If the above condition not hold,\;the de Rhamness has been proved in Proposition \ref{dimesionforParafilgeneric}.\;Now using \cite[Appendix,\;Proposition A.3]{Dingsocle},\;we can show that the de Rhamness of  $\widetilde{E}_{u(1)}-\widetilde{E}_{u(2)}-\cdots-\widetilde{E}_{u(j)}$ (with $j$ increasing) step by step.\;Indeed,\;we can show that $\dim_E\ext^1_{g}(E_{u(j+1)}\otimes_{\cR_{E,L}}\cR_{E[\epsilon]/\epsilon^2}(1+\psi_{j+1}\epsilon),\widetilde{E}_{u(1)}-\widetilde{E}_{u(2)}-\cdots-\widetilde{E}_{u(j)})=\dim_E\ext^1(E_{u(j+1)}\otimes_{\cR_{E,L}}\cR_{E[\epsilon]/\epsilon^2}(1+\psi_{j+1}\epsilon),\widetilde{E}_{u(1)}-\widetilde{E}_{u(2)}-\cdots-\widetilde{E}_{u(j)})$.\;We obtain the first assertion.\;The second statement follows from the first result directly.\;
\end{proof}



For any subspace $V\subseteq \ext^1(\Dpik,\Dpik)$,\;we put $\overline{V}^{\flat}:=V/V\cap \ext^{1,\ast}_{0}(\Dpik,\Dpik)$.\;In this case,\;we obtain a 
natural morphism 
\[\bigoplus_{u\in \sW_{s}}\overline{\ext}^{1,0}_{S_0,u}(\Dpik,\Dpik)\rightarrow \overline{\ext}^1(\Dpik,\Dpik)={\ext}^1(\Dpik,\Dpik)/\ext^{1,\ast}_{0}(\Dpik,\Dpik),\]
which is  a statement on block-version infinite fern.\;Denote by $\overline{\ext}_{J_0,\Sigma}^{1,0}(\Dpik,\Dpik)$ its image.\;


Finally,\;we discuss $\cF_{u}$.\;Put $\ext^{1}_{{u}}(\Dpik,\Dpik):=\ext^{1,\ast}_{\cF_{u}}(\Dpik,\Dpik)$.\;By definition,\;$\ext^{1,\ast}_{{S_0,u}}(\Dpik,\Dpik)\subseteq \ext^{1}_{{u}}(\Dpik,\Dpik)\subseteq \ext^{1}_{{S_0,u}}(\Dpik,\Dpik)$.\ It is easy to see that $\homo_{\sm}(\bZ_{J_0(u)}(L),E)\subseteq \homo(\bZ_{S_0}(L),E)$.\;Similarity,\;let $\ext^{1}_{g,u}(\Dpik,\Dpik)$ (resp.,\;$\ext^{1}_{g',u}(\Dpik,\Dpik)$) be the preimage of $\homo_{\sm}(\bZ_{J_0(u)}(L),E)$ (resp.,\;$\homo_{g'}(\bZ_{J_0(u)}(L),E)$) of via $\kappa_u$.\;Put $\ext^{1}_{0,u}(\Dpik,\Dpik):=\ker\kappa_{u}$.\;It is easy to see that $\ext^{1}_{0,u}(\Dpik,\Dpik)\subseteq\ext^{1}_{g,u}(\Dpik,\Dpik)\subseteq \ext^{1}_{g',u}(\Dpik,\Dpik)$.\;Similarity,\;we have 

\begin{lem}
	We have $\ext^{1}_{g,u}(\Dpik,\Dpik)=\ext^{1}_{u}(\Dpik,\Dpik)\cap \ext^{1}_{g}(\Dpik,\Dpik)$ and $\ext^{1}_{0,u}(\Dpik,\Dpik)=\ext^{1}_{u}(\Dpik,\Dpik)\cap \ext^{1}_{0}(\Dpik,\Dpik)$.\;Moreover,\;we have  $\ext^{1,\ast}_{g,{S_0,u}}(\Dpik,\Dpik)=\ext^{1,\ast}_{{S_0,u}}(\Dpik,\Dpik)\cap \ext^{1}_{g,u}(\Dpik,\Dpik)$ and $\ext^{1,\ast}_{0,{S_0,u}}(\Dpik,\Dpik)=\ext^{1,\ast}_{{S_0,u}}(\Dpik,\Dpik)\cap \ext^{1}_{0,u}(\Dpik,\Dpik)$.\;
\end{lem}
Put $\ext^{1}_{0}(\Dpik,\Dpik)=\ext^{1}_{0,1}(\Dpik,\Dpik)$.\;For any subspace $V\subseteq \ext^1(\Dpik,\Dpik)$, we put $\overline{V}:=V/V\cap \ext^{1}_{0}(\Dpik,\Dpik)$.\;In this case,\;we also obtain a 
natural morphism 
\begin{equation}\label{infintefernfortriang}
	\bigoplus_{u\in \sW_s}\overline{\ext}^{1}_{u}(\Dpik,\Dpik)\twoheadrightarrow \overline{\ext}^1(\Dpik,\Dpik).
\end{equation}
Dented by $\overline{\ext}^1_{\Sigma}(\Dpik,\Dpik)$ its image.\;Moreover,\;we see that $\ext^{1,\ast}_{0,{S_0,u}}(\Dpik,\Dpik)$ is a subspace of $\ext^{1}_{0,u}(\Dpik,\Dpik)$ and $\ext^{1,\ast}_{0,J_0}(\Dpik,\Dpik)\subseteq \ext^{1}_{0}(\Dpik,\Dpik)$.\;It is clear that  $\overline{\ext}_{J_0,\Sigma}^{1,0}(\Dpik,\Dpik)\hookrightarrow \overline{\ext}^1_{\Sigma}(\Dpik,\Dpik)$.\;


\subsection{Hodge parameters and Steinberg blocks/crystalline higher intertwining}\label{bolcksHI}

Recall the two exact sequences:
\begin{equation}
	\begin{aligned}
		0\rightarrow E_1^{s-1}\rightarrow\Dpik\rightarrow E_s\rightarrow 0,\\
		0\rightarrow E'_s\rightarrow\Dpik\rightarrow  (E_1^{s-1})'\rightarrow 0.\;
	\end{aligned}
\end{equation}
In this section,\;we write $M:=M$ and $M':=N$ for simplicity.\;

Let $\cF$ (resp.,\;$\cG$) be the parabolic filtration $M\subseteq\Dpik$ (resp.,\;$E'_s\subseteq\Dpik$).\;By our manage of $\cF$,\;the natural morphism
\[\kappa_{\cF}:\ext^{1}_{\cF}(\Dpik,\Dpik)\rightarrow\ext^{1}(M,M)\times\ext^{1}(E_s,E_s)\]
is surjective,\;since for each $1\leq i\leq s-1$,\;$E_i\succ E_s$ not holds.\;In particular,\;we have natural surjection:
\begin{equation}
	\begin{aligned}
		\kappa_{\cF}=(\kappa_{\cF,1},\kappa_{\cF,2}):\ext^{1,\flat,+}_{\cF}(\Dpik,\Dpik)\twoheadrightarrow\ext^{1}(M,M)\times\ext^{1,\ast}(E_s,E_s)\cong \ext^{1}(M,M)\times \homo(L^{\times},E).\;
	\end{aligned}
\end{equation}
On the other hand,\;consider the natural morphism
\[\kappa_{\cG}:\ext^{1}_{\cG}(\Dpik,\Dpik)\rightarrow\ext^{1}(N,N)\times\ext^{1}(E_s,E_s).\]
If for each $1\leq i\leq s-1$,\;$E_s\succ E_i$ not holds.\;Then $\kappa_{\cG}$ is also surjective.\;In general,\;this morphism may not a surjection,\;it image was described by the so-called Fontaine-Mazur $\cL$-invariants.\;Without loss generality,\;we assume that $E'_{s}\succ E'_1\hookrightarrow N$.\;


\begin{lem}We have a short exact sequence:
	\begin{equation}
		\begin{aligned}
			0\rightarrow \ext^1(N,\Dpik)\xrightarrow{\iota}\ext^1(\Dpik,\Dpik)\xrightarrow{\kappa }\ext^{1}(E'_s,\Dpik)\rightarrow0.\;
		\end{aligned}
	\end{equation}
	The subspace $\ext^{1}_{\cF}(\Dpik,\Dpik)$ is the kernel of the following composition of the maps:
	\[\ext^{1}(\Dpik,\Dpik)\xrightarrow{\kappa} \ext^{1}(E'_s,\Dpik)\xrightarrow{\kappa_2}\ext^{1}(E'_s,N).\]
\end{lem}
\begin{proof}
	By the hypothesis on $\Dpik$,\;we have a long exact sequence:
	\begin{equation}
		\begin{aligned}
			0\rightarrow
			\homo_{(\varphi,\Gamma)}(N,\Dpik)&\rightarrow\homo_{(\varphi,\Gamma)}(\Dpik,\Dpik)\rightarrow\homo_{(\varphi,\Gamma)}(E'_s,\Dpik)\\&\rightarrow	\ext^{1}(N,\Dpik)\xrightarrow{\iota}\ext^1(\Dpik,\Dpik)\xrightarrow{\kappa }\ext^{1}(E'_s,\Dpik)\rightarrow \ext^{2}(N,\Dpik).\;
		\end{aligned}
	\end{equation}
	We see that $\iota$ is injective since the third arrow is an isomorphism (using the fact that $\Dpik$ is not split and noting that both source and target are $1$-dimensional $E$-vector spaces).\;We show that $\kappa$ is surjective,\;i.e.,\;$\ext^{2}(N,\Dpik)=0$.\;By d\'{e}vissage it is enough to show that $\ext^2(E'_i,\Dpik)=0$ for $1\leq i\leq s-1$.\;By Tate duality,\;we have isomorphisms $\ext^2(E'_i,\Dpik)=\hH^2(\Dpik\otimes_{\cR_{E,L}}(E'_i)^{\vee})=\hH^0(\Dpik^{\vee}\otimes_{\cR_{E,L}}E'_i\otimes_{\cR_{E,L}}\cR_{E,L}(\epsilon))$,\;its vanishing follows from the non-split assumption on $\Dpik$.\;On the other hand,\;the kernel of $\kappa_2$ is $\ext^{1}(E'_s,E'_s)$ since $\homo(E'_s,N)=0$.\;Therefore,\;$\ext^{1}_{\cF}(\Dpik,\Dpik)=\kappa^{-1}(\ext^{1}(E'_s,E'_s))$.\;
\end{proof}

Let $E_s'\twoheadrightarrow\cR_{E,L}(\delta) $ be the unique rank-$1$ quotient of $(\varphi,\Gamma)$-modules.\;Consider the cup-products in  the following commutative diagram:
\begin{equation}
	\xymatrix{\ext^1(N,E_s') \ar[d]^{f_1} &\times & \ext^1(N,N)\ar[r]^{\cup} \ar@{=}[d] &  \ext^2(N,E_s') \ar[d]^{\sim}\\
		\ext^1(N,\cR_{E,L}(\delta))  &\times & \ext^1(N,N)\ar[r]^{\cup} &  \ext^2(N,\cR_{E,L}(\delta))}.
\end{equation}
Let $\cL_{\mathrm{FM}}(\Dpik:N):=(E[\Dpik])^{\perp}$ and $\cL^1_{\mathrm{FM}}(\Dpik:N):=(E[\Dpik/(E_s')_{1}^{r_s-1}])^{\perp}$.\;Since $\ext^2(N,(E_s')_{1}^{r_s-1})=0$,\;we get that $f_1$ is surjective.\;Note that $\cL^1_{\mathrm{FM}}(\Dpik:N)$ is of co-dimension $1$ in $\ext^1(N,N)$.\;By definition,\;$\cL_{\mathrm{FM}}(\Dpik:N)\subseteq \cL^1_{\mathrm{FM}}(\Dpik:N)$.\;Then we have  the following so-called Colmez-Greenberg-Stevens formula.\;

\begin{pro}
	For $(\widetilde{N},\psi)\in \ext^{1}(N,N)\times \homo(L^{\times},E)$,\;there exists a $\widetilde{D}\in \ext^{1,\flat,+}_{\cG}(\Dpik,\Dpik)$ such that $\kappa_{\cG}(\widetilde{D})=(\widetilde{N},\psi)$ iff
	$\widetilde{N}\otimes_{\cR_{E[\epsilon]/\epsilon^2}}(\cR_{E[\epsilon]/\epsilon^2}(1-\psi\epsilon))\in\cL^1_{\mathrm{FM}}(\Dpik:N)$.\;
\end{pro}




\begin{pro}There is a natural exact sequence:
	\begin{equation}
		0\rightarrow \ext^1(N,E_s')/[\Dpik]\rightarrow \ext^{1,\flat,+}_{\cG}(\Dpik,\Dpik)\rightarrow \cL^1_{\mathrm{FM}}(\Dpik:N)\oplus \homo(L^{\times},E)\rightarrow 0.\;
	\end{equation}
	
	%,\;and more precisely,\;it lies in a short exact sequence\[0\rightarrow \cL^1_{\mathrm{FM}}(\Dpik:N)\rightarrow \sL^1_{\mathrm{FM}}(\Dpik:N)\rightarrow \homo(L^{\times},E)\rightarrow 0.\]
\end{pro}
\begin{proof}
	By the above proof,\;we get an exact sequence
	\begin{equation}
		\begin{aligned}
			0
			\xrightarrow{\iota}\ext^{1}(N,\Dpik)\rightarrow	\ext^{1}_{\cG}(\Dpik,\Dpik)\xrightarrow{\kappa}\ext^{1}(E'_s,E'_s)\rightarrow 0.\;
		\end{aligned}
	\end{equation}
	In particular,\;$\ext^{1,\flat,+}_{\cG}(\Dpik,\Dpik)$ is the inverse image of $\homo(L^{\times},E)\subseteq \ext^{1}(E'_s,E'_s)$ via $\kappa$.\;So for $\widetilde{D}\in \ext^{1,\flat,+}_{\cG}(\Dpik,\Dpik)$ with $\kappa_{\cG}(\widetilde{D})=(\widetilde{N},\psi)\in \ext^{1}(N,N)\times \homo(L^{\times},E)$,\;$\widetilde{D}\in \ker \kappa$ iff $\psi=0$.\;On the other hand,\;the composition 
	\[\ker(\kappa)\xrightarrow{\iota^{-1}}\ext^{1}(N,\Dpik)\xrightarrow{\kappa_1}\ext^{1}(N,N),\]
	has image $\cL_{\mathrm{FM}}(\Dpik:N)$.\;On the other hand,\;consider the exact sequence:
	\begin{equation}
		0\rightarrow \homo(N,N)\rightarrow	\ext^1(N,E_s')\rightarrow \ext^1(N,\Dpik).\;
	\end{equation}
	The image of the first map is $E[\Dpik]$.\;Thus  we get $\ext^1(N,E_s')/[\Dpik]\hookrightarrow \ker(\kappa)$.\;Since $\cL_{\mathrm{FM}}(\Dpik:N)$ has co-dimension $1$ in $\ext^{1}(N,N)$.\;We deduce a short exact sequence:
	\[0\rightarrow \ext^1(N,E_s')/[\Dpik]\rightarrow \ker(\kappa)\rightarrow \cL^1_{\mathrm{FM}}(\Dpik:N)\rightarrow 0\]
	by comparing dimensions.\;On the other hand,\;we have a short exact sequence:
	\[0\rightarrow \homo(E_s',E_s')\rightarrow \EndO_{\cG}(\Dpik,\Dpik)\rightarrow \homo(N,\Dpik)\rightarrow 0.\]
	The long exact sequence associated it deduces that $\homo(L^{\times},E)$ is a direct summand of $\ext^{1,\flat,+}_{\cG}(\Dpik,\Dpik)$.\;The proposition follows.\;
\end{proof}
\begin{rmk}
	By definition,\;${\ext}^{1,0}_{g}(\Dpik,\Dpik)\subseteq \ker(\kappa)$,\;it image in $\cL_{\mathrm{FM}}(\Dpik:N)$ is ${\ext}^{1,0}_{g}(N,N)$,\;so we deduce that $\dim_E(\ext^1(N,E_s')/[\Dpik])\cap {\ext}^{1,0}_{g}(\Dpik,\Dpik)=(1+\dim \bN_{J_0}-s)-(1+\dim \bN_{J_0(M)}-(s-1))=r_s(n-r_s)-1$(note that $\dim \bN_{J_0}=\dim \bN_{J_0(M)}+r_s(n-r_s)$).\;We thus obtain an exact sequence:
	\[0\rightarrow (\ext^1(N,E_s')/[\Dpik])\cap {\ext}^{1,0}_{g}(\Dpik,\Dpik)\rightarrow  {\ext}^{1,0}_{g}(\Dpik,\Dpik)\rightarrow {\ext}^{1,0}_{g}(N,N)\rightarrow 0.\]
\end{rmk}
We have $\dim_E\ext^{1,\ast}_{\ast}(\Dpik,\Dpik)=1+(n-r_s)^2+r_s^2+r_s(n-r_s)=1+n^2+r_s^2-nr_s$ and $\dim_E\ext^{1,\ast}_{\ast}(\Dpik,\Dpik)=1+n^2+r_s^2-nr_s-(r_s^2-1)=2+n^2-nr_s$ for $\ast\in\{\cF,\cG\}$.\;Similar to the discussion in \cite[Corollary 2.15]{ParaDing2024},\;we have
\begin{equation}
	\begin{aligned}
		\kappa_{\cF}:\overline{\ext}^{1,0}_{\cF}(\Dpik,\Dpik)\xrightarrow{\sim}\overline{\ext}^1(M,M)\times\homo(L^{\times},E).\;
	\end{aligned}
\end{equation}
by comparing the dimensions.\;Indeed,\;the dimension of $\overline{\ext}^{1}_{\cF}(\Dpik,\Dpik)$ is equal to $2+n^2-nr_s-(1+\dim \bN_{J_0}-s)=1+n^2-nr_s+s-\dim \bN_{J_0}$ and the dimension of right side is equal to $3+(n-r_s)^2-(1+\dim \bN_{J_0(M)}-(s-1))$.\;Note that $\dim \bN_{J_0}=\dim \bN_{J_0(M)}+r_s(n-r_s)$,\;we see that  $3+(n-r_s)^2-(1+\dim \bN_{J_0(M)}-(s-1))=1+(n-r_s)^2+s-\dim \bN_{J_0}+r_s(n-r_s)=1+n^2-nr_s+s-\dim \bN_{J_0}$.\;


For general case,\;we have an isomorphism:
\begin{equation}
	\kappa_{\cG}=(\kappa_{\cG,1},\kappa_{\cG_2}):\overline{\ext}^{1,0}_{\cG}(\Dpik,\Dpik)\xrightarrow{\sim}\cL^1_{\mathrm{FM}}(\Dpik:N)/\overline{\ext}^{1,0}_{g}(N,N)\oplus \homo(L^{\times},E),\;
\end{equation}
%and an exact sequence $0\rightarrow \cL^1_{\mathrm{FM}}(\Dpik:N)/\overline{\ext}^{1,0}_{g}(N,N)\rightarrow \sL^1_{\mathrm{FM}}(\Dpik:N)/\overline{\ext}^{1,0}_{g}(N,N)\rightarrow \homo(L^{\times},E)\rightarrow 0.$
If $\kappa_{\cG}$ is still surjective,\;we also have the isomorphism:
\begin{equation}
	\kappa_{\cG}:\overline{\ext}^{1,0}_{\cG}(\Dpik,\Dpik)\xrightarrow{\sim}\overline{\ext}^1(N,N)\times\homo(L^{\times},E).\;
\end{equation}


For any $\iota\in\homo(M,N)$,\;we have the following pull-back (resp.,\;pull-forward) maps:
\begin{equation}
	\begin{aligned}
		\iota^{-}:\ext^1(N,M)\rightarrow \ext^1(M,M),\\
		\iota^{+}:\ext^1(N,M)\rightarrow \ext^1(N,N).\;
	\end{aligned}
\end{equation}
Put
\[\ext_\iota^1(M,M)=\iota^{-}(\ext^1(N,M)),\ext_\iota^1(N,N)=\iota^{+}(\ext^1(N,M)).\]


\begin{pro}Let $\iota\in\homo(M,N)$ be an injection.\;Put $t'_s:=|\{(i,j):1\leq j\leq i\leq \min\{t_{s-1},j+r_s-1\}\}|$ (note that $t_s'=t_{s-1}$ if $r_s=1$).\;
	\begin{itemize}
		\item[(1)] $\dim_E\ext_\iota^1(M,M)=\dim_E\ext_\iota^1(N,N)=1+(n-r_s)^2-t'_s$.\;
		\item[(2)] $\ext_g^1(M,M)\subseteq \ext_\iota^1(M,M)$ and $\ext_g^1(N,N)\subseteq \ext_\iota^1(N,N)$.\;
	\end{itemize}
\end{pro}
\begin{proof}We just prove the result for $M$.\;The same strategy is also suitable for $N$.\;By d\'{e}vissage,\;the kernel of $\iota^{-}$ is $\homo((N)^{\vee}\otimes_{\cR_{E,L}}M,D^{\vee}_1\otimes_{\cR_{E,L}}M)$.\;By \cite[Lemma 5.1.1]{breuil2020probleme},\;we have
	\begin{equation}
		\begin{aligned}
			&\;\homo((N)^{\vee}\otimes_{\cR_{E,L}}M,D^{\vee}_1\otimes_{\cR_{E,L}}M)\\
			\cong&\; \hH^0(\gal_p,W_{\dr}^+(D^{\vee}_1\otimes_{\cR_{E,L}}M)/W_{\dr}^+((N)^{\vee}\otimes_{\cR_{E,L}}M))\\
			\cong&\;\hH^0(\gal_p,\oplus_{1\leq i\leq t_{s-1}}\oplus_{1\leq j\leq t_{s-1} }t^{\bh_{i}-\bh_j}B_{\dr}^+/t^{\bh_i-\bh_{j+r_s}}B_{\dr}^+).\;
		\end{aligned}
	\end{equation}
	This implies $\dim_E\ker\iota^{-}=|\{(i,j):1\leq j\leq i\leq \min\{t_{s-1},j+r_s-1\}\}|-1$ and thus $\dim_E\mathrm{Im}\iota^{-}=(n-r_s)^2-|\{(i,j):1\leq j\leq i\leq \min\{t_{s-1},j+r_s-1\}\}|+1$.\;
\end{proof}


By the same argument as \cite[Theorem 2.32,\;Corollary 2.33]{ParaDing2024},\;we get the result for the block's higher intertwining :
\begin{thm}
	For any $\widetilde{D}\in \ext^{1,\flat,+}_{\cF}(\Dpik,\Dpik)$ with $\kappa_{\cF}(\widetilde{D})=(\widetilde{M},\psi_1)$.\;The followings are equivalent:
	\begin{itemize}
		\item[(a)] $\widetilde{D}\in \ext^{1}_{\cF}(\Dpik,\Dpik)\cap \ext^{1}_{\cG}(\Dpik,\Dpik)$,\;
		\item[(b)] $\widetilde{M}\otimes_{\cR_{E[\epsilon]/\epsilon^2}}(\cR_{E[\epsilon]/\epsilon^2}(1-\psi_1\epsilon))\in \ext_{\iota_{\Dpik}}^1(M,M)$.\;
	\end{itemize}
	When $\widetilde{D}\in \ext^{1,\flat,+}_{\cF}(\Dpik,\Dpik)\cap \ext^{1,\flat,+}_{\cG}(\Dpik,\Dpik)$,\;then there exists $M\in\ext^1(N,M)\cap (\iota_D^{+})^{-1}(\cL^1_{\mathrm{FM}}(\Dpik:N)\otimes_{\cR_{E[\epsilon]/\epsilon^2}}(\cR_{E[\epsilon]/\epsilon^2}(1-\psi_1\epsilon))$ such that
	\begin{equation}
		\begin{aligned}
			&\widetilde{M}=\iota_{\Dpik}^{-}(M)\otimes_{\cR_{E[\epsilon]/\epsilon^2}}(\cR_{E[\epsilon]/\epsilon^2}(1+\psi_1\epsilon)),\;\\
			&\widetilde{N}:=\kappa_{\cG,1}(\widetilde{D})=\iota_{\Dpik}^{+}(M)\otimes_{\cR_{E[\epsilon]/\epsilon^2}}(\cR_{E[\epsilon]/\epsilon^2}(1+\psi_1\epsilon)).\;
		\end{aligned}
	\end{equation}
	Moreover,\;we have $\dim_E \ext^{1,\flat,+}_{\cF}(\Dpik,\Dpik)\cap \ext^{1,\flat,+}_{\cG}(\Dpik,\Dpik)=2+n^2-nr_s-t_s'$,\;thus the imagem of ${\ext}^{1,0}_{\cF}(\Dpik,\Dpik)\oplus {\ext}^{1,0}_{\cG}(\Dpik,\Dpik)$ in ${\ext}^{1}(\Dpik,\Dpik)$,\;denoted by $\ext^{1,\ast}_{\cF+\cG}(\Dpik,\Dpik)$,\;has dimension $2+n^2-nr_s+t_s'$.\;
\end{thm}
\begin{proof}
	We have $\dim_E \ext^{1,\flat,+}_{\cF}(\Dpik,\Dpik)\cap \ext^{1,\flat,+}_{\cG}(\Dpik,\Dpik)=2+n^2-nr_s-t_s'$.\;Thus the image of ${\ext}^{1,0}_{\cF}(\Dpik,\Dpik)\oplus {\ext}^{1,0}_{\cG}(\Dpik,\Dpik)$ in ${\ext}^{1}(\Dpik,\Dpik)$ has dimension $2+n^2-nr_s+t_s'$.\;
\end{proof}

Finally,\;we define
\begin{equation}
	\begin{aligned}
		V^{\flat}(M,N):=&\left(\overline{\ext}^{1,0}(M,M)\times\homo(L^{\times},E)\right)\oplus \left(\overline{\ext}^{1,0}(N,N)\times\homo(L^{\times},E)\right),\\
		&\xleftarrow{\sim,(\kappa_{\cF},\kappa_{\cG})}\overline{\ext}^{1,0}_{\cF}(\Dpik,\Dpik)\oplus \overline{\ext}^{1,0}_{\cG}(\Dpik,\Dpik).\;
	\end{aligned}
\end{equation}
Therefore,\;we get the following exact sequence:
\begin{equation}
	0\rightarrow \cL(\Dpik,M,N)\rightarrow  V(M,N)\rightarrow \ext^{1,\ast}_{\cF+\cG}(\Dpik,\Dpik)\rightarrow 0,\;
\end{equation}
where $\cL(\Dpik,M,N)$ is the kernel of the surjection $V(M,N)\twoheadrightarrow\ext^{1,\ast}_{\cF+\cG}(\Dpik,\Dpik)$.\;By higher intertwining,\;it consists of $((x,\psi),(y,\psi))\in V(M,N)$ such that $((x,\psi),(y,\psi))$ comes from some $M\in\ext^1(N,M)$.\;

\begin{conjecture}Denote by $\overline{\ext}_{J_0}^{1,0}(\Dpik,\Dpik)$ the image of the natural morphism 
	\[\bigoplus_{u\in \sW_{s}}\overline{\ext}^{1,0}_{S_0,u}(\Dpik,\Dpik)\rightarrow \overline{\ext}^1(\Dpik,\Dpik).\]
	Then $\dim_E\overline{\ext}_{J_0}^{1,0}(\Dpik,\Dpik)=1+2s+\dim\bN_{J_0}$.\;
\end{conjecture}

Indeed,\;we define
\begin{equation}
	\pi_{\lalg}(\underline{x},\bh)-C_{t^u_j,u,\sigma}:=\pi_{\lalg}(\underline{x},\bh)\oplus_{Q_{\Delta}^{\diamond}(\emptyset,I_u^-,\lambda),\iota}[Q_{\Delta}^{\diamond}(\emptyset,I_u^-,\lambda)-C_{t^u_j,u,\sigma}]
\end{equation}
via the natural embedding $\iota:Q_{\Delta}^{\diamond}(\emptyset,I_u^-,\lambda)\hookrightarrow \pi_{\lalg}(\underline{x},\bh)$ and 
\begin{equation}
	C_{t^u_j,u,\sigma}-\pi_{\lalg}(\underline{x},\bh):=[C_{t^u_j,u,\sigma}-Q_{\Delta}^{\diamond}(I_u^{+},\emptyset,\lambda)]\times_{Q_{\Delta}^{\diamond}(I_u^{+},\emptyset,\lambda),\beta}\pi_{\lalg}(\underline{x},\bh)
\end{equation}
via the natural quotient map $\beta:\pi_{\lalg}(\underline{x},\bh)\twoheadrightarrow Q_{\Delta}^{\diamond}(I_u^{+},\emptyset,\lambda)$.\;Indeed,\;the map $\ext^1_G\left(C_{t^u_j,u,\sigma},Q_{\Delta}^{\diamond}(\emptyset,I_u^{-},\lambda)\right)\rightarrow \ext^1_G\left(C_{t^u_j,u,\sigma},\pi_{\lalg}(\underline{x},\bh)\right)$ induced by injection $Q_{\Delta}^{\diamond}(\emptyset,I_u^{-},\lambda)\hookrightarrow \pi_{\lalg}(\underline{x},\bh)$ is an isomorphism.\;Indeed,\;
The extension of $\pi_{\lalg}(\underline{x},\bh)\otimes_E\underline{\sL}_{J_0}^{\diamond}(\Dpik)$ by $\pi^{\diamond}_{J_0,1}(\underline{x},\bh)$ gives a locally analytic representation $\pi^{\diamond}_{I_0,\min}(\Dpik)$, which encodes the ``crystalline" Hodge parameters (more precisely,\;the information $\{\iota^{l,q,\diamond}_{\Dpik}\}_{1\leq r\leq r+1<q\leq s}$) among Steinberg blocks.\;Indeed,\;we have an easier locally analytic representation $\pi^{\diamond,-}_{I_0,\min}(\Dpik)\hookrightarrow \pi^{\diamond}_{I_0,\min}(\Dpik)$ that encodes necessary information,\;which is spirited by  Theorem \ref{higherdetercry}.\;For $1\leq l<l+1<q\leq s$,\;put $w_{l,q}=s_ls_{l+1}\cdots s_q\in\sW_s$.\;Then consider the locally analytic representation
$\pi_{l,q}(\cL_{l,q})$ with the following form (see (\ref{basiclinespecial})) :
\begin{equation}
	\pi_{l,q}(\cL_{l,q}):=Q_{\Delta}^{\diamond}(\emptyset,I_{j_s(w_{l,q})}^-,\lambda)-C_{j_s(s_l),j_s(w_{l,q})}-Q_{\Delta}^{\diamond}(I_{j_s(w_{l,q})}^+,\lambda).
\end{equation}
Note that $I_{j_s(w_{l,q})}^-=t_{q}$.\;Taking amalgamated sum,\;we put
\[\pi^{\diamond,-}_{I_0,\min}(\Dpik)=\bigoplus^{3\leq q\leq s}_{\pi_{\lalg}}\bigoplus^{1\leq l<l+1<s}_{Q_{\Delta}^{\diamond}(\emptyset,\{t_{q}\},\lambda)}\pi_{l,q}(\cL_{l,q}).\]
Then $\pi^{\diamond,-}_{I_0,\min}(\Dpik)$ encodes $\{\iota^{l,q,\diamond}_{\Dpik}\}_{1\leq r\leq r+1<q\leq s}$ by choosing $E$-lines $\cL_{l,q}$ carefully.\;In $\pi^{\diamond,-}_{I_0,\min}(\Dpik)$,\;we have $\frac{(s-1)(s-2)}{2}$'s extra locally algebraic consistents.\;


Similarity,\;we have an easier locally analytic representation $\pi^{-}_{\min}(B)\hookrightarrow \pi_{\min}(B)$ that determine $B$,\;which is spirited by  Theorem \ref{higherdetercry}.\;For $b+1\leq l<l+1<q\leq b+d$,\;put $w_{l,q}=s_ls_{l+1}\cdots s_q$.\;Then consider the locally analytic representation
$\pi_{l,q}(\cL_{l,q})$ with the following form (see (\ref{basiclinespecial})) :
\begin{equation}
	\pi_{l,q}(\cL_{l,q}):=Q_{\Delta_d}^{\diamond}(\emptyset,\{q\},\bh)-C_{s_l,w_{l,q}}-Q_{\Delta_d}^{\diamond}(\{l+1,q+1\},\bh).
\end{equation}
Taking amalgamated sum,\;we put
\[\pi^{-}_{\min}(B)=\bigoplus^{b+3\leq q\leq b+d}_{Q_{\Delta_d}^{\diamond}(\emptyset,\bh)}\bigoplus^{b+1\leq l<l+1<q}_{Q_{\Delta_d}^{\diamond}(\emptyset,\{q\},\bh)}\pi_{l,q}(\cL_{l,q}).\]
By Theorem \ref{higherdetercry},\;$\pi^{-}_{\min}(B)$ determines $B$ by choosing $E$-lines $\cL_{l,q}$ carefully.\;In $\pi^{-}_{\min}(B)$,\;we have $\frac{(d-1)(d-2)}{2}$'s extra locally algebraic consistents,\;which is much smaller than the $2^d-\frac{d(d+1)}{2}-1+\# I_0\backslash J_0(B)$.\;

///////////////////////////////////////////deformations

For $u\in\sW_s$,\;recall the  $\omepik$-filtration $\cF_{u}$:
\begin{equation}
	\Dpik:=E'_{u(1)}-E'_{u(2)}-\cdots-E'_{u(m)},
\end{equation}
We applying the discussion in Section\ref{genedeformations} to $\cG=\cF_{u}$ for $u\in \sW_{s}$.\;Write $\ext^{1}_{u}(\Dpik,\Dpik):=\ext^{1}_{\cF_{u}}(\Dpik,\Dpik)$ for simplicity.\;We have the following inclusion of extension groups:
\[\ext^{1,\diamond}_{u}(\Dpik,\Dpik)\subseteq\ext^{1,\flat}_{u}(\Dpik,\Dpik)\subseteq\ext^{1,0}_{u}(\Dpik,\Dpik)\subseteq \ext^{1}_{u}(\Dpik,\Dpik).\]
We obtain the natural maps
\begin{equation}
	\begin{aligned}
		\kappa^{\ast}_{u}:\ext^{1,\ast}_{u}(\Dpik,\Dpik):=\ext^{1,\ast}_{\cF_{u}}(\Dpik,\Dpik)\twoheadrightarrow\homo(\bZ_{S^u_0}(L),E)
	\end{aligned}
\end{equation}
with $\bZ_{S^u_0}\cong \BG_m^s$ and $\ast\in\{\diamond,\flat,0\}$.\;For $u\in \sW_{s}$,\;let  $\ext^{1,\ast}_{g,u}(\Dpik,\Dpik)\subseteq\ext^{1,\ast}_{u}(\Dpik,\Dpik)$ be the preimage of $\homo_{\sm}(\bZ_{S^u_0}(L),E)$ of via $\kappa_w$.\;Let $\ext^1_g(\Dpik,\Dpik)$ be the subspace of de Rham deformations of $\Dpik$.\;



\begin{pro}\label{dimesionforParafilgeneric}We have the following facts.\;
	\begin{itemize}
		\item[(1)] We have $\dim_E\ext^1(\Dpik,\Dpik)=1+n^2$ and $\dim_E\ext^1_{{u}}(\Dpik,\Dpik)=1+\dim \bP_{S^u_0}=1+\dim \bP_{S_0}$.\;
		\item[(2)] We have $\dim_E\ext^{1,\flat}_{u}(\Dpik,\Dpik)=1+\dim \bN_{S_0}+m$,\;$\dim_E\ext^{1,0}_{u}(\Dpik,\Dpik)=1+\dim \bN_{S_0}+m$ and $\dim_E\ext^{1,\diamond}_{\cF_{u}}(\Dpik,\Dpik)=1+\frac{m(m+1)}{2}$.\;
		\item[(3)] For $u\in \sW_{s}$,\;we have $\ext^1_g(\Dpik,\Dpik)\subseteq \ext^{1}_{u}(\Dpik,\Dpik)$.\;Moreover,\;we have 
		\[\ext^{1,\ast}_{g,u}(\Dpik,\Dpik)=\ext^{1,\ast}_{u}(\Dpik,\Dpik)\cap \ext^1_g(\Dpik,\Dpik).\;\]
		
	\end{itemize}
\end{pro}	
\begin{proof}
	Part $(1)$ follows from [Proposition 3.6].\;By the formally smoothness in [Proposition 3.7].\;Part $(2)$ follows from the codimension of $\prod_{i=1}^m\ext^{1,\flat}(E_{u(i)},E_{u(i)})\cong \homo(\bQ^{\times}_p,E)$ (resp.,\;$\prod_{i=1}^m\ext^{1,\flat,0}(E_{u(i)},E_{u(i)})$) in $\prod_{i=1}^m\ext^1(E_{u(i)},E_{u(i)})$ is $m+\dim \bL_{S^u_0}-2m=\dim \bL_{S^u_0}-m$ (resp.,\;).\;For $\dim_E\ext^{1,\diamond}_{u}(\Dpik,\Dpik)$,\;we prove the result by using induction on $m$.\;Put $\Dpik_0^{m-1}:=E'_{u(1)}-E'_{u(2)}-\cdots-E'_{u(m-1)}$ and suppose that $\dim_E\ext^{1,\diamond}_{{u}}(\Dpik_0^{m-1},\Dpik_0^{m-1})=1+\frac{(m-1)m}{2}$.\;But the extensions dimension of the desired forms  in $(\ref{pictureforblockext})$ is $m-1$,\;so $\dim_E\ext^{1,\diamond}_{{u}}(\Dpik,\Dpik)=1+\frac{(m-1)m}{2}+1+(m-1)=1+\frac{m(m+1)}{2}$.\;For $(3)$,\;let $\widehat{D}=D-D\in \ext^1_g(\Dpik,\Dpik)$.\;By \cite[Appendix,\;Proposition A.3]{Dingsocle},\;for $i<j$,\;we have $\dim_E\hH^1_g(\gal_L,W(E_{u(i)}^{\vee}\otimes_{\cR_{E,L}}{E_{u(j)}}))=0$,\;where $W(-)$ means the associated $E$-$B$ pair.\;The d\'{e}vissage argument shows the inclusion $\ext^1_g(\Dpik,\Dpik)\subseteq \ext^{1}_{u}(\Dpik,\Dpik)$.\;We now show $(3)$,\;since $\cR_{E[\epsilon]/\epsilon^2}(1+\psi_i\epsilon)$ iff $\psi\in \homo_{\sm}(\bQ^{\times}_p,E)$,\;we deduce that $\ext^{1,\ast}_{u}(\Dpik,\Dpik)\cap \ext^1_g(\Dpik,\Dpik)\subseteq \ext^{1,\ast}_{g,u}(\Dpik,\Dpik)$.\;On the other hand,\;for $(\psi_i)\in \homo_{\sm}(\bZ_{S^u_0}(L),E)$,\;we will show that \[\widetilde{D}=\widetilde{E}_{u(1)}-\widetilde{E}_{u(2)}-\cdots-\widetilde{E}_{u(m)},\;\]
	with $\widetilde{E}_{u(i)}\cong E_{u(i)}\otimes_{\cR_{E,L}}\cR_{E[\epsilon]/\epsilon^2}(1+\psi_i\epsilon)$ is de Rham.\;For $i<j$,\;we first show that any element in
	\[\hH^1\left(\gal_{p},W(E_{u(j)}^{\vee}\otimes_{\cR_{E,L}}{E_{u(i)}})\otimes_{\cR_{E,L}}\cR_{E[\epsilon]/\epsilon^2}(1+(\psi_j-\psi_i)\epsilon)\right)\]
	are de Rham.\;However,\;since the Hodge-Tate weights of $E_{u(j)}^{\vee}\otimes_{\cR_{E,L}}{E_{u(i)}}$ are positive,\;we deduce from \cite[Appendix,\;Proposition A.3]{Dingsocle} that $\hH_g^1(\gal,-)=\hH^1(\gal,-)$.\;Now using \cite[Appendix,\;Proposition A.5]{Dingsocle},\;we can show that the de Rhamness of  $\widetilde{E}_{1}-\widetilde{E}_{2}-\cdots-\widetilde{E}_{j}$ (with $j$ increasing) step by step.\;This completes the proof of $(3)$.\;
\end{proof}

An block variation of \cite[Lemma 2.12]{ParaDing2024} gives  the following lemma.\;
\begin{lem}\label{smhomoindepent}For $\ast\in\{\diamond,\flat\}$,\;$\ext^{1,\ast}_{{u}}(\Dpik,\Dpik)\cap \ext^1(\Dpik,\Dpik)$ is independent on the choice of $u$,\;and the following diagram commutes
	\begin{equation}
		\xymatrix{ \ext^{1,\ast}_{g,u_1}(\Dpik,\Dpik)\ar@{=}[d] \ar[r]_{\kappa_{u_1}} & \homo_{\sm}(\bZ^{u_1}_{S_0}(L),E) \ar[d]^{\sim}_{u_2^{-1}u_1}\\
			\ext^{1,\ast}_{g,u_2}(\Dpik,\Dpik) \ar[r]_{\kappa_{u_2}} & \homo_{\sm}(\bZ^{u_2}_{S_0}(L),E) },
	\end{equation}
	for $u_1,u_2\in\sW_{s}$.\;
\end{lem}
\begin{proof}
	It suffices to prove the statement for the case that $\ast=\flat$ ($\ast=\diamond$ is just a restriction) anad $u_2^{-1}u_1$ is a simple reflection $j_m(s_k)\in\sW_{s}$.\;Choose $\widetilde{D}\in  \ext^{1,\ast}_{g,,u_1}(\Dpik,\Dpik)$.\;Suppose that \[\widetilde{D}=\widetilde{E}_{u_1(1)}-\widetilde{E}_{u_1(2)}-\cdots-\widetilde{E}_{u_1(m)},\;\]
	with $\widetilde{E}_{u_1(i)}\cong E_{u_1(i)}\otimes_{\cR_{E,L}}\cR_{E[\epsilon]/\epsilon^2}(1+\psi_i\epsilon)$.\;Now by assumption $u_1(j)=u_2(j)$ for $j\neq k,k+1$.\;Thus,\;for $j<k$ or $j>k+1$,\;we have $\fil^j_{\cF_{u_1}}\widetilde{D}=\fil^j_{\cF_{u_2}}\widetilde{D}$.\;Since $\homo(\widetilde{E}_{u_1(1)},\widetilde{D})\cong \homo(\widetilde{E}_{u_1(1)},\widetilde{E}_{u_1(1)})\cong E[\epsilon]/\epsilon^2$,\;using d\'{e}vissage for $\fil^j_{\cF_{u_2}}\widetilde{D}$,\;we deduce that $\homo(\widetilde{E}_{u_1(1)},\widetilde{E}_{u_2(1)})\cong E[\epsilon]/\epsilon^2$,\;hence $\widetilde{E}_{u_1(1)}=\widetilde{E}_{u_2(1)}$.\;Now consider $\widetilde{D}/\widetilde{E}_{u_1(1)}$,\;repeat the same argument,\;we will get that $\widetilde{E}_{u_1(j)}=\widetilde{E}_{u_2(j)}$ for $j<k$.\;For $j=k$,\;we see that $\homo(\widetilde{E}_{u_1(k)},\widetilde{E}_{u_2(k+1)})\cong E[\epsilon]/\epsilon^2$,\;so that $\widetilde{E}_{u_1(k)}=\widetilde{E}_{u_2(k+1)}$.\;Exchanging $\cF_{u_1}$ and $\cF_{u_2}$,\;we obtain $\widetilde{E}_{u_1(k+1)}=\widetilde{E}_{u_2(k)}$.\;For $j>k+1$,\;we consider $\widetilde{D}/\fil^{k+1}_{\cF_{u_1}}$,\;and we see that $\widetilde{E}_{u_1(j)}=\widetilde{E}_{u_2(j)}$ for $j>k+1$.\;In conclusion,\;we see that $\widetilde{D}\in  \ext^{1,\ast}_{g,u_2}(\Dpik,\Dpik)$ and satisfies the above commutative diagram.\;We complete the proof.\;
\end{proof}

Put $ \ext^{1,\ast}_{g}(\Dpik,\Dpik):=\ext^{1,\ast}_{{u}}(\Dpik,\Dpik)\cap \ext^1_g(\Dpik,\Dpik)$.\;Let $\ext^{1,\ast}_0(\Dpik,\Dpik):=\ker\kappa^{\ast}_{u}\subseteq\ext^{1,\ast}_g(\Dpik,\Dpik)$ (for some $u$,\;but by above lemma,\;it is independent on the choice of $u$).\;Moreover,\;we deduce that 
\[\dim_E\ext^{1,\flat}_0(\Dpik,\Dpik)=1+\dim \bN_{S_0}-m.\]
We also let $\ext^{1,\flat}_{g'}(\Dpik,\Dpik)$ be the preimage of $\homo_{g'}(\bZ_{S^u_0}(L),E)$ under $\kappa^{\flat}_{u}$.\;We have
\[\dim_E\ext^{1,\ast}_{g'}(\Dpik,\Dpik)=2+\dim \bN_{S_0}.\]
We also have
\[\dim_E\ext^{1,\diamond}_0(\Dpik,\Dpik)=1+\frac{m(m-1)}{2}-m.\]
Let $\ext^{1,\diamond}_{g'}(\Dpik,\Dpik)$ be the preimage of $\homo_{g'}(\bZ_{S^u_0}(L),E)$ under $\kappa^{\diamond}_{u}$.\;We have
\[\dim_E\ext^{1,\diamond}_{g'}(\Dpik,\Dpik)=2+\frac{m(m+1)}{2}.\]
For any subspace $V\subseteq \ext^1(\Dpik,\Dpik)$,\;we put
\[\overline{V}=V/V\cap\ext^1_0(\Dpik,\Dpik).\;\]
So we also have ismorphisms $\overline{\ext}^{1,\flat}_{g}(\Dpik,\Dpik)\xrightarrow{\sim }\homo_{\sm}(\bZ_{S^u_0}(L),E),\;\overline{\ext}^{1,\flat}_{u}(\Dpik,\Dpik)\xrightarrow{\sim}\homo(\bZ_{S^u_0}(L),E)$,\;and
$\overline{\ext}^{1,\diamond}_{g}(\Dpik,\Dpik)\xrightarrow{\sim }\homo_{\sm}(\bZ_{S^u_0}(L),E),\;\overline{\ext}^{1,\diamond}_{u}(\Dpik,\Dpik)\xrightarrow{\sim}\homo(\bZ_{S^u_0}(L),E)$.\;Finally,\;we obtain a commutative diagram:
\begin{equation}
	\xymatrix{ \overline{\ext}^{1,\diamond}_{u}(\Dpik,\Dpik) \ar@{->>}[d]_{\kappa^{\diamond}_{u}} \ar@{^(->}[r] & \overline{\ext}^{1,0}_{u}(\Dpik,\Dpik)\ar@{->>}[dl]^{\kappa^{\flat}_{u}} \\
		\homo(\bZ_{S^u_0}(L),E)  &  },
\end{equation}



We end this section with certain statement on block-version infinite fern.\;
\begin{conjecture}Denote by ${\ext}_{S_0}^{1,\flat}(\Dpik,\Dpik)$ (resp.,\;$\overline{\ext}_{S_0}^{1,\flat}(\Dpik,\Dpik)$) the image of the natural morphism 
	\[\bigoplus_{u\in \sW_{s}}{\ext}^{1,\flat}_{u}(\Dpik,\Dpik)\rightarrow {\ext}^1(\Dpik,\Dpik),\;\text{resp.,\;}\bigoplus_{u\in \sW_{s}}\overline{\ext}^{1,\flat}_{u}(\Dpik,\Dpik)\rightarrow \overline{\ext}^1(\Dpik,\Dpik).\]
	Then $\dim_E{\ext}_{S_0}^{1,\flat}(\Dpik,\Dpik)=1+m+2\dim\bN_{S_0}=1+n^2-\dim\bL_{S_0}+m$ and  $\dim_E\overline{\ext}_{S_0}^{1,\flat}(\Dpik,\Dpik)=2m+\dim\bN_{S_0}$.\;
\end{conjecture}


\begin{conjecture}
	Denote by ${\ext}_{S_0}^{1,\diamond}(\Dpik,\Dpik)$ (resp.,\;$\overline{\ext}_{S_0}^{1,\diamond}(\Dpik,\Dpik)$) the image of the natural morphism 
	\[\bigoplus_{u\in \sW_{s}}{\ext}^{1,\diamond}_{u}(\Dpik,\Dpik)\rightarrow \overline{\ext}^1(\Dpik,\Dpik),\;\text{resp.,\;}\bigoplus_{u\in \sW_{s}}\overline{\ext}^{1,\diamond}_{u}(\Dpik,\Dpik)\rightarrow {\ext}^1(\Dpik,\Dpik).\]
	Then $\dim_E{\ext}_{S_0}^{1,0}(\Dpik,\Dpik)=1+m^2$ and $\dim_E\overline{\ext}_{S_0}^{1,0}(\Dpik,\Dpik)=\frac{m(m+1)}{2}+m$.\;
\end{conjecture}












\subsection{Hodge parameters and blocks higher intertwining}\label{bolcksHI}


Let $\cF$ (resp.,\;$\cG$) be the parabolic filtration $M\subseteq\Dpik$ (resp.,\;$E'_s\subseteq\Dpik$).\;By our manage of $\cF$,\;the natural morphism
\begin{equation}
	\begin{aligned}
		&\kappa_{\cF}:\ext^{1}_{\cF}(\Dpik,\Dpik)\rightarrow\ext^{1}(M,M)\times\ext^{1}(E_s,E_s),\\
		&\kappa_{\cG}:\ext^{1}_{\cG}(\Dpik,\Dpik)\rightarrow\ext^{1}(M',M')\times\ext^{1}(E_s,E_s).
	\end{aligned}
\end{equation}
is surjective by the generic assumption.\;For $\ast\in\{\flat,0\}$,\;we set 
\begin{equation}
	\begin{aligned}
		&\ext^{1,\ast,+}_{\cF}(\Dpik,\Dpik):=\ext^{1}_{\cF}(\Dpik,\Dpik)\times_{\ext^{1}(E_s,E_s)}\ext^{1,\ast}(E_s,E_s),\\
		&\ext^{1,\ast,+}_{\cG}(\Dpik,\Dpik):=\ext^{1}_{\cG}(\Dpik,\Dpik)\times_{\ext^{1}(E_s',E_s')}\ext^{1,\ast}(E_s',E_s'),\\
	\end{aligned}
\end{equation}
In particular,\;when $\ast=\flat$,\;we have natural surjection:
\begin{equation}
	\begin{aligned}
		&\kappa_{\cF}=(\kappa_{\cF,1},\kappa_{\cF,2}):\ext^{1,\flat,+}_{\cF}(\Dpik,\Dpik)\twoheadrightarrow\ext^{1}(M,M)\times \homo(L^{\times},E),\;\\
		&\kappa_{\cG}=(\kappa_{\cG,1},\kappa_{\cG,2}):\ext^{1,\flat,+}_{\cG}(\Dpik,\Dpik)\rightarrow\ext^{1}(M',M')\times \homo(L^{\times},E).
	\end{aligned}
\end{equation}
For any $\iota\in\homo(M,M')$,\;we have the following pull-back (resp.,\;pull-forward) maps:
\begin{equation}
	\begin{aligned}
		&\iota^{-}:\ext^1(M',M)\rightarrow \ext^1(M,M),\\
		&\iota^{+}:\ext^1(M',M)\rightarrow \ext^1(M',M').\;
	\end{aligned}
\end{equation}
Put
\[\ext_\iota^1(M,M)=\iota^{-}(\ext^1(M',M)),\ext_\iota^1(M',M')=\iota^{+}(\ext^1(M',M)).\]


\begin{pro}Let $\iota\in\homo(M,M')$ be an injection.\;Put $t'_s:=\#\{(i,j):1\leq j\leq i\leq \min\{t_{s-1},j+n_s-1\}\}$ (note that $t_s'=t_{s-1}$ if $n_s=1$).\;
	\begin{itemize}
		\item[(1)] $\dim_E\ext_\iota^1(M,M)=\dim_E\ext_\iota^1(M',M')=1+(n-n_s)^2-t'_s$.\;
		\item[(2)] $\ext_g^1(M,M)\subseteq \ext_\iota^1(M,M)$ and $\ext_g^1(M',M')\subseteq \ext_\iota^1(M',M')$.\;
	\end{itemize}
\end{pro}
\begin{proof}We just prove the result for $M$.\;The same strategy is also suitable for $M'$.\;By d\'{e}vissage,\;the kernel of $\iota^{-}$ is $\homo(M'^{\vee}\otimes_{\cR_{E,L}}M,D^{\vee}_1\otimes_{\cR_{E,L}}M)$.\;By \cite[Lemma 5.1.1]{breuil2020probleme},\;we have
	\begin{equation}
		\begin{aligned}
			&\;\homo(M'^{\vee}\otimes_{\cR_{E,L}}M,D^{\vee}_1\otimes_{\cR_{E,L}}M)\\
			\cong&\; \hH^0(\gal_p,W_{\dr}^+(D^{\vee}_1\otimes_{\cR_{E,L}}M)/W_{\dr}^+(M'^{\vee}\otimes_{\cR_{E,L}}M))\\
			\cong&\;\hH^0(\gal_p,\oplus_{1\leq i\leq t_{s-1}}\oplus_{1\leq j\leq t_{s-1} }t^{\bh_{i}-\bh_j}B_{\dr}^+/t^{\bh_i-\bh_{j+n_s}}B_{\dr}^+).\;
		\end{aligned}
	\end{equation}
	This implies $\dim_E\ker\iota^{-}=\#\{(i,j):1\leq j\leq i\leq \min\{t_{s-1},j+n_s-1\}\}-1$ and thus $\dim_E\mathrm{Im}\iota^{-}=(n-n_s)^2-\#\{(i,j):1\leq j\leq i\leq \min\{s_{m-1},j+n_s-1\}\}+1$.\;
\end{proof}


By the same argument as \cite[Theorem 2.32,\;Corollary 2.33]{ParaDing2024},\;we get the result for the block's higher intertwining :
\begin{thm}\label{mainthmonhigherintw}
	For any $\widetilde{D}\in \ext^{1,\flat,+}_{\cF}(\Dpik,\Dpik)$ with $\kappa_{\cF}(\widetilde{D})=(\widetilde{M},\psi_1)$.\;The followings are equivalent:
	\begin{itemize}
		\item[(a)] $\widetilde{D}\in \ext^{1}_{\cF}(\Dpik,\Dpik)\cap \ext^{1}_{\cG}(\Dpik,\Dpik)$,\;
		\item[(b)] $\widetilde{M}\otimes_{\cR_{E[\epsilon]/\epsilon^2}}(\cR_{E[\epsilon]/\epsilon^2}(1-\psi_1\epsilon))\in \ext_{\iota_{\Dpik}}^1(M,M)$.\;
	\end{itemize}
	When $\widetilde{D}\in \ext^{1,\flat,+}_{\cF}(\Dpik,\Dpik)\cap \ext^{1,\flat,+}_{\cG}(\Dpik,\Dpik)$,\;then there exists $M\in\ext^1(M',M)$ such that
	\begin{equation}
		\begin{aligned}
			&\widetilde{M}=\iota_{\Dpik}^{-}(M)\otimes_{\cR_{E[\epsilon]/\epsilon^2}}(\cR_{E[\epsilon]/\epsilon^2}(1+\psi_1\epsilon)),\;\\
			&\widetilde{M'}:=\kappa_{\cG,1}(\widetilde{D})=\iota_{\Dpik}^{+}(M)\otimes_{\cR_{E[\epsilon]/\epsilon^2}}(\cR_{E[\epsilon]/\epsilon^2}(1+\psi_1\epsilon)).\;
		\end{aligned}
	\end{equation}
	Moreover,\;we have $\dim_E \ext^{1,\flat,+}_{\cF}(\Dpik,\Dpik)\cap \ext^{1,\flat,+}_{\cG}(\Dpik,\Dpik)=2+n^2-nn_s-t_s'$,\;thus the imagem of ${\ext}^{1,0}_{\cF}(\Dpik,\Dpik)\oplus {\ext}^{1,0}_{\cG}(\Dpik,\Dpik)$ in ${\ext}^{1}(\Dpik,\Dpik)$,\;denoted by $\ext^{1,\ast}_{\cF+\cG}(\Dpik,\Dpik)$,\;has dimension $2+n^2-nn_s+t_s'$.\;
\end{thm}
\begin{proof}
	We have $\dim_E \ext^{1,\flat,+}_{\cF}(\Dpik,\Dpik)\cap \ext^{1,\flat,+}_{\cG}(\Dpik,\Dpik)=2+n^2-nn_s-t_s'$.\;Since $\dim_E{\ext}^{1,0}_{\cF}(\Dpik,\Dpik)=1+1+(n-n_s)^2+n_s(n-n_s)=2+n(n-n_s)$,\;we deduce that the image of ${\ext}^{1,0}_{\cF}(\Dpik,\Dpik)\oplus {\ext}^{1,0}_{\cG}(\Dpik,\Dpik)$ in ${\ext}^{1}(\Dpik,\Dpik)$ has dimension $4+2n^2-2nn_s-(2+n^2-nn_s-t_s')=2+n^2-nn_s+t_s'$.\;
\end{proof}
For convenient,\;we manage the special semistable blocks $D_{i}$ such that $f_1\geq f_2\geq\cdots\geq n_s$.\;Thus,\;we have 
$s'_m:=\#\{(i,j):1\leq j\leq i\leq \min\{t_{s-1},j+n_s-1\}\}=\sum_{j\leq t_{s-1}-n_s+1}n_s+\sum_{t_{s-1}-n_s+1<j\leq t_{s-1}}(t_{s-1}-j+1)=(t_{s-1}-n_s+1)n_s+\frac{n_s(n_s-1)}{2}=t_{s-1}n_s-\frac{n_s(n_s-1)}{2}=nn_s-n_s^2-\frac{n_s(n_s-1)}{2}$.\;Then $\dim_E \ext^{1,\flat,+}_{\cF}(\Dpik,\Dpik)\cap \ext^{1,\flat,+}_{\cG}(\Dpik,\Dpik)=2+(n-n_s)^2+\frac{n_s(n_s-1)}{2}$ and $\dim_E\ext^{1,\ast}_{\cF+\cG}(\Dpik,\Dpik)=2+n^2-n_s^2-\frac{n_s(n_s-1)}{2}$.\;

Define
\begin{equation}
	\begin{aligned}
		V^{\flat}(M,M'):=&\left(\overline{\ext}^{1,0}(M,M)\times\homo(L^{\times},E)\right)\oplus \left(\overline{\ext}^{1,0}(M',M')\times\homo(L^{\times},E)\right),\\
		&\xleftarrow{\sim,(\kappa_{\cF},\kappa_{\cG})}\overline{\ext}^{1,\flat,+}_{\cF}(\Dpik,\Dpik)\oplus \overline{\ext}^{1,\flat,+}_{\cG}(\Dpik,\Dpik).\;
	\end{aligned}
\end{equation}
Therefore,\;we get the following exact sequence:
\begin{equation}
	0\rightarrow \cL(\Dpik,M,M')\rightarrow  	V^{\flat}(M,M')\rightarrow \overline{\ext}^{1,0}_{\cF+\cG}(\Dpik,\Dpik)\rightarrow 0,\;
\end{equation}
where $\cL(\Dpik,M,M')$ is the kernel of the surjection $V(M,M')\twoheadrightarrow\overline{\ext}^{1,0}_{\cF+\cG}(\Dpik,\Dpik)$.\;By higher intertwining,\;it consists of $((x,\psi),(y,\psi))\in V(M,M')$ such that $((x,\psi),(y,\psi))$ comes from some $M\in\ext^1(M',M)$.\;Note that  $\dim_E\overline{\ext}^{1,0}_{\cF+\cG}(\Dpik,\Dpik)=1+n^2-n_s^2-\frac{n_s(n_s-1)}{2}-\dim \bN_{S_0}+m$.\;Note that $\dim_EV^{\flat}(M,M')=\dim_E\overline{\ext}^{1,0}_{\cF}(\Dpik,\Dpik)+\dim_E\overline{\ext}^{1,0}_{\cG}(\Dpik,\Dpik)=2(2+n(n-n_s)-(1+\dim\bN_{S_0}-m))=2(1+n(n-n_s)-\dim\bN_{S_0}+m)$.\;This shows that $\dim_E\cL(\Dpik,M,M')=1+2n(n-n_s)-\dim\bN_{S_0}+m-(n^2-n_s^2-\frac{n_s(n_s-1)}{2})=1+(n-n_s)^2-\dim\bN_{S_0}+m+\frac{n_s(n_s-1)}{2})$.\;

We construct a certain subspace ${\ext}_{S_0}^{1,\diamond,-}(\Dpik,\Dpik)$ of ${\ext}_{S_0}^{1,\diamond}(\Dpik,\Dpik)$ by induction on $m$.\;Recall that $\cF$ (resp.,\;$\cG$) is the parabolic filtration $M\subseteq\Dpik$ (resp.,\;$E'_s\subseteq\Dpik$).\;We define a subspace $\ext^{1,\diamond,+}_{\cF}(\Dpik,\Dpik)\subseteq \ext^{1}_{\cF}(\Dpik,\Dpik)$,\;$\widetilde{D}=[\widetilde{M}-\widetilde{D}_m]\in \ext^{1,\diamond,+}_{\cF}(\Dpik,\Dpik)$ iff $\widetilde{M}\in {\ext}_{I_0(M)}^{1,\diamond,-}(M,M)$ (by inductive hypothesis),\;$\widetilde{D}_m\in \ext^{1,\ast}(E_s,E_s)$,\;and $\widetilde{D}$ has the form
\begin{equation}\label{pictureforblockext}
	\xymatrix{ M\ar@{-}[dr]_{[\widetilde{M}]} \ar@{-}[drr]_{v} \ar@{-}[r]^{E[\Dpik]} & E_s \ar@{-}[dr]^{[\widetilde{D}_m]} &\\
		&	M \ar@{-}[r]_{E[\Dpik]} & E_s }.
\end{equation}
for $v\in V_{m-1}$.\;Similar for $\cG$.\;Let $\ext^{1,\diamond,-}_{S_0}(\Dpik,\Dpik)$ be the image of 
\[\ext^{1,\diamond,+}_{\cF}(\Dpik,\Dpik)\oplus \ext^{1,\diamond,+}_{\cG}(\Dpik,\Dpik)\rightarrow  \ext^{1}_{\cF}(\Dpik,\Dpik).\;\]
This completes the inductive construction of $\ext^{1,\diamond,-}_{S_0}(\Dpik,\Dpik)$.\;

//////////////////////////////////////////////////////////

\section{Appendix}

For $?\in\{\flat,0\}$,\;sending $\widetilde{D}$ to the parameters $(\psi_i)$  induces a map
\begin{equation}
	\kappa^{?}_{\cG}:\ext^{1,\flat}_{\cG}(\Dpik,\Dpik)\rightarrow \homo(\bZ_{\cG}(L),E).\;
\end{equation}
with $\bZ_{\cG}\cong \BG_m^h$.\;Let $\ext^{1,?}_{g,\cG}(\Dpik,\Dpik)\subseteq\ext^{1,\flat}_{\cG}(\Dpik,\Dpik)$ be the preimage of the subspace $\homo_{\sm}(\bZ_{\cG}(L),E)\subseteq \homo(\bZ_{\cG}(L),E)$ of smooth characters via $\kappa_w$.\;


Choose $(\psi_i)\in \homo(\bZ_{\cG}(L),E)$,\;we construct a special deformation
\begin{equation}\label{blockwisedeformations}
	\widetilde{D}=\widetilde{F}_{1}-\widetilde{F}_{2}-\cdots-\widetilde{F}_{h}\in \ext^{1,\flat}_{\cG}(\Dpik,\Dpik)
\end{equation}
with $\widetilde{F}_{i}\cong 
{F}_{i}\otimes_{\cR_{E,L}}\cR_{E[\epsilon]/\epsilon^2}(1+\psi_i\epsilon)$ inductively.\;Suppose that $\widetilde{D}_1^{j-1}=\widetilde{F}_{1}-\widetilde{F}_{2}-\cdots-\widetilde{F}_1^{j-1}$ is constructed and  $D_{\cG,1}^{j-1}=\widetilde{D}_1^{j-1}\;\mathrm{mod}\;\epsilon$.\;Note that 
\begin{equation}
	\ext^1_{(\varphi,\Gamma)}(\widetilde{F}_{j},\widetilde{D}_1^{j-1})\xrightarrow{\sim}\hH^1_{(\varphi,\Gamma)}({F}^{\vee}_{j}\otimes_{\cR_{E,L}}\widetilde{D}_1^{j-1}\otimes_{\cR_{E,L}}\cR_{E[\epsilon]/\epsilon^2}(1-\psi_j\epsilon))
\end{equation}
We have an exact sequence:
\[0\rightarrow\hH^1_{(\varphi,\Gamma)}({F}^{\vee}_{j}\otimes_{\cR_{E,L}}D_{\cG,1}^{j-1}) \rightarrow \hH^1_{(\varphi,\Gamma)}({F}^{\vee}_{j}\otimes_{\cR_{E,L}}\widetilde{D}_1^{j-1}\otimes_{\cR_{E,L}}\cR_{E[\epsilon]/\epsilon^2}(1-\psi_j)\epsilon)) \xrightarrow{g} \hH^1_{(\varphi,\Gamma)}({F}^{\vee}_{j}\otimes_{\cR_{E,L}}D_{\cG,1}^{j-1}).\]
For any $1\leq l\leq j-1$,\;we have an exact sequence:
\[0\rightarrow\hH^1_{(\varphi,\Gamma)}({F}^{\vee}_{j}\otimes_{\cR_{E,L}}D_{\cG,1}^{l-1}) \rightarrow \hH^1_{(\varphi,\Gamma)}({F}^{\vee}_{j}\otimes_{\cR_{E,L}}D_{\cG,1}^{j-1}) \xrightarrow{g_l} \hH^1_{(\varphi,\Gamma)}({F}^{\vee}_{j}\otimes_{\cR_{E,L}}D_{\cG,l}^{j-1}).\]
Note that $g_l([D_{\cG,1}^{j}])=[D_{\cG,l}^{j}]$.\;Define $V^{\cG}_{j-1}:=E\langle [D_{\cG,l}^{j}]:\;1\leq l\leq j-1\rangle$ of $\hH^1_{(\varphi,\Gamma)}({F}^{\vee}_{j}\otimes_{\cR_{E,L}}D_{\cG,1}^{j-1})$.\;Note that $\dim_EV^{\cG}_{j-1}=j-1$.\;

Suppose that the element $L=[D_{\cG,1}^{j}]\in\hH^1_{(\varphi,\Gamma)}({F}^{\vee}_{j}\otimes_{\cR_{E,L}}D_{\cG,1}^{j-1})$ lies in the image of $g$,\;consider subspace $V^{\cG}_{j-1}=j-1\subseteq \hH^1_{(\varphi,\Gamma)}({F}^{\vee}_{j}\otimes_{\cR_{E,L}}D_{\cG,1}^{j-1})$ containing $L$,\;we have a $j$-dimensional subspace $\ext^1({F}^{\vee}_{j}\otimes_{\cR_{E,L}}\widetilde{D}_1^{j-1},L)$ which fixes into a short exact sequence of $E$-vector spaces:
\[0\rightarrow V^{\cG}_{j-1}\rightarrow \ext^1({F}^{\vee}_{j}\otimes_{\cR_{E,L}}\widetilde{D}_1^{j-1},L)\rightarrow E[L]\rightarrow 0.\;\]
Furthermore,\;$\ext^1({F}^{\vee}_{j}\otimes_{\cR_{E,L}}\widetilde{D}_1^{j-1},L)$ contains all the (isomorphic classes of  extensions of the following forms:
\begin{equation}\label{pictureforblockext}
	\xymatrix{ D_{j-1}\ar@{-}[dr]_{[\widetilde{D}_{j-1}]} \ar@{-}[drr]_{v_2} \ar@{-}[r]^{aL} & F_j \ar@{-}[dr]^{\psi_j} &\\
		&	D_{j-1} \ar@{-}[r]_{bL} & F_j }.
\end{equation}
for $v\in V$ and $a,b\in E$.\;In the sequel,\;we let $\ext^{1,\diamond}_{\cG}(\Dpik,\Dpik)$ be the subspace of deformations constructed in (\ref{blockwisedeformations}).\;


We construct a certain subspace ${\ext}_{S_0}^{1,\diamond,-}(\Dpik,\Dpik)$ of ${\ext}_{S_0}^{1,\diamond}(\Dpik,\Dpik)$ by induction on $m$.\;Recall that $\cF$ (resp.,\;$\cG$) is the parabolic filtration $M\subseteq\Dpik$ (resp.,\;$E'_s\subseteq\Dpik$).\;We define a subspace $\ext^{1,\diamond,+}_{\cF}(\Dpik,\Dpik)\subseteq \ext^{1}_{\cF}(\Dpik,\Dpik)$,\;$\widetilde{D}=[\widetilde{M}-\widetilde{D}_m]\in \ext^{1,\diamond,+}_{\cF}(\Dpik,\Dpik)$ iff $\widetilde{M}\in {\ext}_{I_0(M)}^{1,\diamond,-}(M,M)$ (by inductive hypothesis),\;$\widetilde{D}_m\in \ext^{1,\ast}(E_s,E_s)$,\;and $\widetilde{D}$ has the form
\begin{equation}\label{pictureforblockext}
	\xymatrix{ M\ar@{-}[dr]_{[\widetilde{M}]} \ar@{-}[drr]_{v} \ar@{-}[r]^{E[\Dpik]} & E_s \ar@{-}[dr]^{[\widetilde{D}_m]} &\\
		&	M \ar@{-}[r]_{E[\Dpik]} & E_s }.
\end{equation}
for $v\in V_{m-1}$.\;Similar for $\cG$.\;Let $\ext^{1,\diamond,-}_{S_0}(\Dpik,\Dpik)$ be the image of 
\[\ext^{1,\diamond,+}_{\cF}(\Dpik,\Dpik)\oplus \ext^{1,\diamond,+}_{\cG}(\Dpik,\Dpik)\rightarrow  \ext^{1}_{\cF}(\Dpik,\Dpik).\;\]
This completes the inductive construction of $\ext^{1,\diamond,-}_{S_0}(\Dpik,\Dpik)$.\;

To characterize  $\overline{\ext}^{1,\diamond}_{\Sigma}(\Dpik,\Dpik)$ more precisely,\;we make the following consideration.\;


Recall that $\cF_1$ (resp.,\;$\cG_1$) is the parabolic filtration $M\subseteq\Dpik$ (resp.,\;$D'_s\subseteq\Dpik$) of $\Dpik$.\;We set 
\begin{equation}
	\begin{aligned}
		&\overline{\ext}^{1,\diamond,+}_{\cF_1}(\Dpik,\Dpik):=\overline{\ext}^{1,\flat,+}_{\cF_1}(\Dpik,\Dpik)\times_{\overline{\ext}^{1}(M,M)}\ext^{1,\diamond}_{\Sigma}(M,M),\\
		&\overline{\ext}^{1,\diamond,+}_{\cG_1}(\Dpik,\Dpik):=\overline{\ext}^{1,\flat,+}_{\cG_1}(\Dpik,\Dpik)\times_{\overline{\ext}^{1}(M',M')}\overline{\ext}^{1,\diamond}_{\Sigma}(M',M'),\\
	\end{aligned}
\end{equation}
From the proof of Theorem \ref{mainthmforspecialblock},\;we will get the following property of the image $\ext^{1,\diamond}_{\Sigma}(\Dpik,\Dpik)$.\;
\begin{cor}\label{diamondsurjhypoforI0}
	We have surjections $\overline{\ext}^{1,\diamond,+}_{\cF_1}(\Dpik,\Dpik)\oplus \overline{\ext}^{1,\diamond,+}_{\cG_1}(\Dpik,\Dpik)\twoheadrightarrow \overline{\ext}^{1,\diamond}_{\Sigma}(\Dpik,\Dpik)$.\;
\end{cor}
\begin{rmk}
	This is an analogue of Theorem \ref{mainthmonhigherintw} in \cite{ParaDing2024}.\;Furthermore,\;we suspect that $\overline{\ext}^{1,\diamond}_{\Sigma}(\Dpik,\Dpik)$  comes from 
	a subspace $\ext^{1,\diamond}_{\Sigma}(\Dpik,\Dpik)$ (of dimension $1=\frac{m(m+1)}{2}$) of ${\ext}^{1}(\Dpik,\Dpik)$.\;In Appendix,\;we discuss the pure Galois description of $\ext^{1,\diamond}_{\Sigma}(\Dpik,\Dpik)$.\;
\end{rmk}



\begin{conjecture}
	Denote by ${\ext}_{S_0}^{1,\diamond}(\Dpik,\Dpik)$ (resp.,\;$\overline{\ext}_{S_0}^{1,\diamond}(\Dpik,\Dpik)$) the image of the natural morphism 
	\[\bigoplus_{u\in \sW_{s}}{\ext}^{1,\diamond}_{u}(\Dpik,\Dpik)\rightarrow \overline{\ext}^1(\Dpik,\Dpik),\;\text{resp.,\;}\bigoplus_{u\in \sW_{s}}\overline{\ext}^{1,\diamond}_{u}(\Dpik,\Dpik)\rightarrow {\ext}^1(\Dpik,\Dpik).\]
	Then $\dim_E{\ext}_{S_0}^{1,0}(\Dpik,\Dpik)=1+m^2$ and $\dim_E\overline{\ext}_{S_0}^{1,0}(\Dpik,\Dpik)=\frac{m(m+1)}{2}+m$.\;
\end{conjecture}

We also have
\[\dim_E\ext^{1,\diamond}_0(\Dpik,\Dpik)=1+\frac{m(m-1)}{2}-m.\]
Let $\ext^{1,\diamond}_{g'}(\Dpik,\Dpik)$ be the preimage of $\homo_{g'}(\bZ_{S^u_0}(L),E)$ under $\kappa^{\diamond}_{u}$.\;We have
\[\dim_E\ext^{1,\diamond}_{g'}(\Dpik,\Dpik)=2+\frac{m(m+1)}{2}.\]
For any subspace $V\subseteq \ext^1(\Dpik,\Dpik)$,\;we put
\[\overline{V}=V/V\cap\ext^1_0(\Dpik,\Dpik).\;\]
So we also have ismorphisms $\overline{\ext}^{1,\flat}_{g}(\Dpik,\Dpik)\xrightarrow{\sim }\homo_{\sm}(\bZ_{S^u_0}(L),E),\;\overline{\ext}^{1,\flat}_{u}(\Dpik,\Dpik)\xrightarrow{\sim}\homo(\bZ_{S^u_0}(L),E)$,\;and
$\overline{\ext}^{1,\diamond}_{g}(\Dpik,\Dpik)\xrightarrow{\sim }\homo_{\sm}(\bZ_{S^u_0}(L),E),\;\overline{\ext}^{1,\diamond}_{u}(\Dpik,\Dpik)\xrightarrow{\sim}\homo(\bZ_{S^u_0}(L),E)$.\;Finally,\;we obtain a commutative diagram:
\begin{equation}
	\xymatrix{ \overline{\ext}^{1,\diamond}_{u}(\Dpik,\Dpik) \ar@{->>}[d]_{\kappa^{\diamond}_{u}} \ar@{^(->}[r] & \overline{\ext}^{1,0}_{u}(\Dpik,\Dpik)\ar@{->>}[dl]^{\kappa^{\flat}_{u}} \\
		\homo(\bZ_{S^u_0}(L),E)  &  },
\end{equation}



We follow the notation of \cite[Section 4.1.1]{2019DINGSimple} and \cite[Section 2]{PATCHING2016}.\;Suppose that $p\nmid 2n$, and let $\overline{r}: \gal_L \rightarrow \GLN_n(k_E)$ be a continuous representation such that $\overline{r}$ admits a potentially crystalline lift $r_{\mathrm{pot.diag}}: \gal_L \rightarrow \GLN_n(E)$  of regular weight $\xi$ which is potentially diagonalisable.\;We can find a triple $(F,F^+, \overline{\rho})$,\;where $F$ is an imaginary CM field with maximal totally real subfield $F^+$,\;and $\overline{\rho}:\gal_{F^+} \rightarrow \cG_n(k_E)$ is a \emph{suitable globalisation} (cf. \cite[Section 2.1]{PATCHING2016}) of $\overline{r}$.\;In particular,\;for any place $v|p$ of $F^+$,\;$v$ splits in $F$, and has $F^+_v\cong L$.\;There is a place $\widetilde{v}$ of $F$ lying over $v$ with $\overline{\rho}|_{\gal_{F_{\widetilde{v}}}}\cong \overline{r}$.\;

Let $S_p$ be the set of places of $F^+$ above $p$.\;By \cite[Section 2.3]{PATCHING2016},\;we can find the following objects
\begin{equation*}
	\{\widetilde{G}, v_1, \fp\in S_p, \{U_m\}_{m\in \BZ}\},
\end{equation*}
where $\widetilde{G}$ is a certain definite unitary group over $F^+$,\;$v_1$ is a certain finite place of $F^+$ prime to $p$,\;and $\{U_m=\prod_{v} U_{m,v}\}_{m\in \BZ_{\geq 0}}$ is a tower of certain compact open subgroups of $\widetilde{G}(\BA_{F^+}^{\infty})$ (see \cite[Section 4.1.1,\;Page 8028]{2019DINGSimple} or \cite[Section 2,\;Page 214]{PATCHING2016} for a precise description).\;We mention that $U_m$ has full level $\fp^m$ (resp.,\;level $\widetilde{G}(\cO_{F^+_v})$) at $\fp$ (resp.,\;$v\in S_p\setminus\{\fp\}$).\;

Let $\tau$ be the inertial type of $r_{\mathrm{pot.diag}}$.\;As in \cite[Section 4.1.1,\;Page 8029]{2019DINGSimple} or  \cite[Section 2,\;Page 215]{PATCHING2016},\;we consider the space of $p$-adic automorphic forms 	$\widehat{S}_{\xi,\tau}(U^{\fp}, \co_E)$ and $\widehat{S}_{\xi,\tau}(U^{\fp},E)$ (roughly speaking,\;the space of $p$-adic algebraic automorphic forms of fixed type $\sigma(\tau)$ (see \cite[Theorem 3.7]{PATCHING2016},\;the ``inertial local Langlands correspondence") at the place $S_p\setminus\{\fp\}$,\;full level at $\fp$,\;and whose weight is $0$ at places above $\fp$,\;and given by the regular weight $\xi$ at each of the places in $S_p\setminus\{\fp\}$).\;Note that $\widehat{S}_{\xi,\tau}(U^{\fp},E)$ is a Banach space for the supermum norm and is equipped with a continuous (unitary) action of $\GLN_n(L)$ (by right translation on functions).\;The space $\widehat{S}_{\xi,\tau}(U^{\fp},E)$  is also equipped with a faithful  action of a certain commutative global Hecke algebra $\bT^{S_p,\univ}$ over $\cO_E$ which is generated by some prime-to-$p$ Hecke operators (see \cite[Section 4.1.1,\;Page 8029]{2019DINGSimple}).\;We can associate to $\overline{\rho}$ a maximal ideal $\fm_{\overline{\rho}}$ of $\bT^{S_p,\univ}$.\;
Let $\widehat{S}_{\xi,\tau}(U^{\fp},\ast)_{\fm_{\overline{\rho}}}$ be the localization of $\widehat{S}_{\xi,\tau}(U^{\fp},\ast)$ at $\fm_{\overline{\rho}}$ for $\ast\in\{\co_E,E\}$.\;Then the action of $\bT^{S_p,\univ}$ on the localization $\widehat{S}_{\xi,\tau}(U^{\fp}, \co_E)_{\fm_{\overline{\rho}}}$ factors through certain  Hecke algebra
$\bT_{\xi,\tau}^{S_p}(U^{\fp}, \co_E)_{\fm_{\overline{\rho}}}$ (see \cite[Page 8029]{2019DINGSimple}).\;We also see that $\widehat{S}_{\xi,\tau}(U^{\fp},E)_{*}$ with $*\in \{\fm_{\overline{\rho}}, \emptyset\}$ are admissible unitary Banach representation of with invariant lattice $\widehat{S}_{\xi,\tau}(U^{\fp}, \co_E)_*$.\;

We denote by $R_{\widetilde{v}}^{\square}$ the maximal reduced and $p$-torsion free quotient of the universal $\co_E$-lifting ring of $\overline{\rho}_{\widetilde{v}}:=\overline{\rho}|_{\gal_{F_{\widetilde{v}}}}$ ($\cong \overline{r}$,\;and therefore $R_{\widetilde{v}}^{\square}\cong \defvarring$). For $v\in S_p\backslash\{\fp\}$, we denote by $R_{\widetilde{v}}^{\square, \xi,\tau}$ for the reduced and $p$-torsion free quotient of $R_{\widetilde{v}}^{\square}$ corresponding to potentially crystalline lifts of weight $\xi$ and inertial type $\tau$.\;Consider the following global deformation problem (in the terminology of \cite{clozel2008automorphy},\;see also \cite[Section 2.4]{PATCHING2016})
\begin{equation*}
	\begin{aligned}
		\cS&=\bigg\{{F}/{F}^+,T^+,T,\cO_E,\overline{\rho},\chi_{\mathrm{cyc}}^{1-n}\delta_{{F}/{F}^+}^n,\{R_{\widetilde{v}_1}^{\square}\}\cup
		\{R_{\fp}^{\square}\}\cup \{R_{\widetilde{v}}^{\square, \xi,\tau}\}_{v\in S_p\backslash \{\fp\}}\bigg\}
	\end{aligned}
\end{equation*}
Then by  \cite[Proposition 2.2.9]{clozel2008automorphy} (see also \cite[Section 2.4]{PATCHING2016}),\;this deformation problem is represented by a universal deformation ring $R_{\cS}^{\univ}$.\;Note that we have a natural morphism $R_{\cS}^{\univ}\rightarrow \bT_{\xi,\tau}^{S_p}(U^{\fp}, \co_E)_{\fm_{\overline{\rho}}}$.\;

Following \cite[Section 4.1.1]{2019DINGSimple} (or \cite[Section 2.8]{PATCHING2016}) we put
\begin{eqnarray*}
	R^{\loc}:=R_{\widetilde{\fp}}^{\square} \widehat{\otimes} \Big(\widehat{\otimes}_{S_p\backslash\{\fp\}}R_{\widetilde{v}}^{\square, \xi,\tau}\Big)\widehat{\otimes} R_{\widetilde{v_1}}^{\square} ,
\end{eqnarray*}
where all completed tensor products are taken over $\cO_E$.\;We put $g:=q-[F^+:\BQ]\frac{n(n-1)}{2}$,\;where $q\geq [F^+:\BQ]\frac{n(n-1)}{2}$ is a certain integer as in  \cite[Section 2.6,\;Page 217]{PATCHING2016}).\;We now put 
\begin{eqnarray*}
	R_{\infty}&:=&R^{\loc}\llbracket x_1,\cdots, x_g\rrbracket,\\
	S_{\infty}&:=&\co_E\llbracket z_1,\cdots, z_{n^2(|S_p|+1)}, y_1,\cdots, y_q\rrbracket,
\end{eqnarray*}
where $x_i$, $y_i$, $z_i$ are formal variables.\;By the end of \cite[Section 4.1.1]{2019DINGSimple} (or \cite[Section 2.8]{PATCHING2016}),\;we get the following objects:
\begin{enumerate}
	\item a continuous $R_{\infty}$-admissible unitary representation $\Pi_{\infty}$ of $G=\GLN_n(L)$ over $E$ together with a $G$-stable and $R_{\infty}$-stable unit ball $\Pi_{\infty}^o\subset \Pi_{\infty}$;
	\item a morphism of local $\co_E$-algebras $S_{\infty}\ra R_{\infty}$ such that $M_{\infty}:= \homo_{\co_L}(\Pi_{\infty}^o, \co_E)$ is finite projective as $S_{\infty}\llbracket \GLN_n(\co_L)\rrbracket$-module;
	\item a closed ideal $\fa$ of $R_{\infty}$, a surjection $R_{\infty}/\fa R_{\infty}\twoheadrightarrow R_{\cS}^{\univ}$ and a  $G   \times R_{\infty}/\fa R_{\infty}$-invariant isomorphism $\Pi_{\infty}[\fa]\cong \widehat{S}_{\xi,\tau}(U^{\fp},E)_{\fm_{\overline{\rho}}}$, where $R_{\infty}$ acts on $\widehat{S}_{\xi,\tau}(U^{\fp},E)_{\fm_{\overline{\rho}}}$ via $R_{\infty}/\fa R_{\infty}\twoheadrightarrow  R_{\cS}^{\univ}$.
\end{enumerate}

Let $R^{\fp}=\Big(\widehat{\otimes}_{v\in S_p\backslash\{\fp\}}R_{\widetilde{v}}^{\square, \xi,\tau}\Big)\widehat{\otimes} R_{\widetilde{v_1}}^{\square}$ and $R_{\infty}^{\fp}:=R^{\fp}\llbracket x_1,\cdots, x_g\rrbracket$.\;Then we have  $R^{\loc}=R^{\fp}\widehat{\otimes} \defvarring$ (recall that $R_{\widetilde{v}}^{\square}\cong \defvarring$ by definition) and $R_{\infty}=R_{\infty}^{\fp}\widehat{\otimes} \defvarring$.\;Let $\BU$ be the open unit ball in $\BA^1$.\;We put $\FX_{\overline{\rho}^{\fp}}^\Box:=(\Spf\;R^{\fp})^{\rig}$ and  $\mathfrak{X}_{\overline{r}}^\Box=(\Spf\;R_{\overline{r}}^\Box)^{\rig}$.\;Then we have $(\Spf\;R^{\fp}_{\infty})^{\rig}=\FX_{\overline{\rho}^{\fp}}^\Box\times \BU^g$.\;We have thus
$\FX_{\infty}:=(\Spf\;R_{\infty})^{\rig}\cong (\Spf\;R^{\fp}_{\infty})^{\rig}\times \mathfrak{X}_{\overline{r}}^\Box\cong \FX_{\overline{\rho}^{\fp}}^\Box\times \BU^g\times \mathfrak{X}_{\overline{r}}^\Box$.\;

%The patched module $\Pi_\infty$ defined in Section 4.1 is a $R_\infty$-module and 

In the sequel,\;for $u\in\sW_n$,\;we fis the cuspidal Bernstein component $\Omega_{S_0^u}:=\prod_{i=1}^s\Omega_{u(i)}$ of Levi subgroup $\bL_{S^u_0}:=\prod_{i=1}^s\GLN_{n_{u(i)}}$. 



Let $\bh:=(\hpi_{\tau,1},\hpi_{\tau,2},\cdots,\hpi_{\tau,n} )_{\tau\in \Sigma_L}$ be a strictly $S_0$-dominant weight.\;We put $\hpi_{i}=(\hpi_{\tau,i})_{\tau\in \Sigma_L}$ for $1\leq i\leq n$ and put ${\bm\lambda}_\bh=(\hpi_{\tau,i}+i-1)_{\tau\in \Sigma_L,1\leq i\leq n}$.\;We denote by $\Pi_\infty^{R_\infty-\ana}$ the subrepresentation of $G$ of locally $R_\infty$-analytic vectors of $\Pi_\infty$ (see \cite[Section 3.1]{breuil2017interpretation}).\;Recall that $\bZ_{S_0}(L)\cong Z_{S_0,\varpi_L}\times \bZ_{S_0}(\cO_L)$,\;where 
$Z_{S_0,\varpi_L}$ is the image of $\oplus_{i=1}^k\BZ\hookrightarrow \bZ_{S_0}(L),\;(m_i)\mapsto (\varpi_L^{m_i})$.\;

We recall the construction of certain  $R_{\infty}\times\FZ_{\omepik}\times Z_{S_0,\varpi_L}\times \bZ_{S_0}(\cO_L)$-module $B_{\omepik,{\bm\lambda}_\bh}(\Pi_\infty^{R_\infty-\ana})$ in \cite[Sections  3.1.1,\;3.1.2,\;3.3]{Ding2021}.\;

By an easy variation of the proof of \cite[Lemma 3.1.3]{Ding2021} and the argument before \cite[(3.28)]{Ding2021},\;we see that $B_{\omepik,{\bm\lambda}_\bh}(\Pi_\infty^{R_\infty-\ana})^\vee$ is a coadmissible module over $\cO(\FX_{\infty}\times\zuni\times\rigch)$, which corresponds to a coherent sheaf $\cM_{\omepik,{\bm\lambda}_\bh}^{\infty,0}$ over $\FX_{\infty}\times \zuni\times \rigch$.\;Since the action of $Z_{S_0,\varpi_L}$ factors through $\sbanpik$,\;we see that $\cM_{\omepik,{\bm\lambda}_\bh}^{\infty,0}$ gives rise to a coherent sheaf $\cM_{\omepik,{\bm\lambda}_\bh}^{\infty}$ over $\FX_{\infty}\times\sbanpik\times \rigch$ such that
\[\Gamma\Big(\FX_{\infty}\times\sbanpik\times \rigch,\cM_{\omepik,{\bm\lambda}_\bh}^{\infty}\Big)\cong B_{\omepik,{\bm\lambda}_\bh}(\Pi_\infty^{R_\infty-\ana})^\vee.\]
Let $\mathcal{E}_{\omepik,\fp,{\bm\lambda}_\bh}^\infty(\overline{\rho})\hookrightarrow\FX_{\infty}\times\sbanpik\times \rigch$ be the Zariski-closed support of $\cM_{\omepik,{\bm\lambda}_\bh}^{\infty}$.\;We call $\mathcal{E}_{\omepik,\fp,{\bm\lambda}_\bh}^\infty(\overline{\rho})$  the patched Bernstein eigenvariety.\;

By an easy variation of the proof of Proposition  3.3.2,\;Corollary \;3.3.3,\;Proposition\;3.3.4, Theorem 3.3.5 and Proposition\;3.3.6 in\cite{Ding2021}, we have
\begin{pro}\label{basicpropertyeigen}\hspace{20pt}
	\begin{itemize}
		\item[(1)]  For $x=(\fm_x, \pi_{x},\chi_x)\in \FX_{\infty}\times\sbanpik\times \rigch$,\;$x\in \mathcal{E}_{\omepik,\fp,{\bm\lambda}_\bh}^\infty(\overline{\rho})$ iff
		\[\homo_{\bL_{S_0}(L)}\Big(\pi_{x}\otimes_E((\chi_x)_{\varpi_{L}}\flat\mathrm{ det}_{\bL_{S_0}(L)})\otimes_E L_{S_0}({\bm\lambda}_\bh),J_{\bP_{S_0}(L)}(\Pi_\infty^{R_\infty-\ana}[\fm_y] )\Big)\neq0\]
		\item[(2)] The rigid space $\mathcal{E}_{\omepik,\fp,{\bm\lambda}_\bh}^\infty(\overline{\rho})$ is reduced and equidimensional of dimension
		\[g+kd_L+n^2(|S_p|+1)+[F^+:\BQ]\frac{n(n-1)}{2}.\]
		\item[(3)] The coherent sheaf $\cM_{\omepik,{\bm\lambda}_\bh}^{\infty}$ is Cohen-Macaulay over $\mathcal{E}_{\omepik,\fp,{\bm\lambda}_\bh}^\infty(\overline{\rho})$.\;
		\item[(4)] The set of very classical non-critical generic points  (see \cite[Defintion 3.2.7,\;(3.22)]{Ding2021} for the definitions of classical,\;very classical,\;non-critical,\;and generic points respectively) is Zarisiki-dense in $\mathcal{E}_{\omepik,\fp,{\bm\lambda}_\bh}^\infty(\overline{\rho})$ and is an accumulation set.\;The set of very classical non-critical generic points  accumulates at point $x=(\fm_x, \pi_{x},\chi_x)$ with $\chi_x$ locally algebraic.\;
	\end{itemize}
\end{pro}
%For $x=(\fm_x, \pi_{x},\chi_x)\in \mathcal{E}_{\omepik,\fp,{\bm\lambda}_\bh}^\infty(\overline{\rho})$,\;let $\rho_x$ be the $\gal_L$-representation associated to $x$ via $\fm_x\in (\Spf R_{\infty})^{\rig}=\FX_{\overline{\rho}^{\fp}}^\Box\times \BU^g\times \mathfrak{X}_{\overline{r}}^\Box$ in $\mathfrak{X}_{\overline{r}}^\Box$.\;Let $\mathbf{WD}_x=(\mathbf{WD}_{x,i})_{1\leq i\leq k}$ be the Weil-Deligne representation associated to $x$,\;i.e.,\;$\mathbf{WD}_{x,i}:=\rec(\pi_{x,i})(\frac{1-r}{2}-(i-1)r)$.\;
%We call a generic classical point $x$ is non-critical if $D^{\dag}_{\rig}(\rho_x)$ is non-critical in the sense of Section 2.4,\;with respect to the $\Omega$-filtration defined as in  \cite[(3.30),(3.24))]{Ding2021}.\;%Fix $\Omega$ a cuspidal Bernstein component of $\bL_{\Omega}(L)$,\;and denote by $\mathfrak{Z}_{\Omega}$ the corresponding Bernstein centre (over $E$).\;Let $V$ be an admissible locally analytic representation of $G$ over $E$.\;Then we see that $J_{\bP_\Omega}(V)$ is an essentially admissible representation of  $\bL_{\Omega}(L)$ (\cite{emerton2006jacquet}).\;Let $\underline{\lambda}$ be a integral $\bP_\Omega$-dominant weight of $G$,\;which corresponds an algebraic representation $L(\underline{\lambda})_\Omega$  of $\bL_{\Omega}(L)$ with weight $\underline{\lambda}$.\;We set
%\[J_{\bP_\Omega(L)}(V)_{\underline{\lambda}}:=\homo_{\sbanpik\times\rigch_{\Omega}^D}(L(\underline{\lambda})_\Omega,J_{\bP_\Omega(L)}(V)),\]
%which is closed $\bL_{\Omega}(L)$-subrepresentation of $J_{\bP_\Omega(L)}(V)\otimes_EL(\underline{\lambda})_\Omega^\vee$.\;Let $\sigma$ be  the irreducible smooth representation of $\bL_{\Omega}^0$ as in Section 2.3.\;We now put
%\[B_{\sigma,\ul{\lambda}}(V)\cong \big(J_{\bP_\Omega(L)}(V)_{\underline{\lambda}}\tene \cC^{\BQ_p-\an}(Z_{\bL_{\Omega}}^0,E)\ten \sigma^\vee\big)^{\bL_{\Omega}^0}.\]We explain the $\mathfrak{Z}_{\Omega}\times \iota_0(Z_{\bL_{\Omega}}(L))\times \iota_2(Z_{\bL_{\Omega}}^0)$ on $B_{\sigma,\ul{\lambda}}(V)$ as follows.\; Let $\iota_{2}(Z_{\bL_{\Omega}}(L))$ (resp.,\;$\iota_1(Z_{\bL_{\Omega}}^0)$) denotes the action of $Z_{\bL_{\Omega}}$ (resp.,\;$Z_{\bL_{\Omega}}^0$) on the$J_{\bP_\Omega(L)}(V)_{\underline{\lambda}}\tene \cC^{\BQ_p-\an}(Z_{\bL_{\Omega}}^0,E)$ by only acting on the term $J_{\bP_\Omega(L)}(V)_{\underline{\lambda}}$ (resp.,\;$\cC^{\BQ_p-\an}(Z_{\bL_{\Omega}}^0,E)$).\;The action $\iota_{2}(Z_{\bL_{\Omega}})\times \iota_1(Z_{\bL_{\Omega}}^0)$ is called standard action.\;On the other hand,\;we have\[B_{\sigma,\ul{\lambda}}(V)\cong \homo_{\bL_{\Omega}(L)}(\mathbf{c-ind}^{\bL_{\Omega}(L)}_{\bL_{\Omega}^0}\sigma,J_{\bP_\Omega(L)}(V)_{\underline{\lambda}}\tene \cC^{\BQ_p-\an}(Z_{\bL_{\Omega}}^0,E)).\]Then we deduce that $B_{\Omega,\ul{\lambda}}(V)$ is equipped with a natural action of $\mathfrak{Z}_{\Omega}\cong \EndO_{\bL_{\Omega}(L)}(\mathbf{c-ind}^{\bL_{\Omega}(L)}_{\bL_{\Omega}^0}\sigma)$.\;Meanwhile,\;we also use $\iota_0(Z_{\bL_{\Omega}}(L))$ to denotes the diagonal action of $Z_{\bL_{\Omega}}(L)$ on $B_{\sigma,\ul{\lambda}}(V)$,\;i.e.,\;with the natural action $Z_{\bL_{\Omega}}(L)$ on $J_{\bP_\Omega(L)}(V)_{\underline{\lambda}}\tene \cC^{\BQ_p-\an}(Z_{\bL_{\Omega}}^0,E)$ by acting via
%\[Z_{\bL_{\Omega}}(L)\xrightarrow{(\id,\;\pr_{\varpi_L}\flat\det_{\bL_{\Omega}}^{-1})}\iota_{2}(Z_{\bL_{\Omega}}(L))\times \iota_1(Z_{\bL_{\Omega}}^0).\]
%Then by definition,\;we see that $\iota_0(Z_{\bL_{\Omega}}(L))|_{Z_{\bL_{\Omega}}^0}=\psi_\sigma:=\sigma|_{Z_{\bL_{\Omega}}^0}$.\;This shows that
%$B_{\sigma,\ul{\lambda}}(V)$ is a closed subrepresentation of $\big(J_{\bP_\Omega(L)}(V)_{\underline{\lambda}}\tene \cC^{\BQ_p-\an}(Z_{\bL_{\Omega}}^0,E)\big)[Z_{\bL_{\Omega}}^0=\psi_\sigma]$.\;Therefore,\;we conclude that the $\iota_0(Z_{\bL_{\Omega}}(L))\times \iota_1(Z_{\bL_{\Omega}}^0)$ on $B_{\sigma,\ul{\lambda}}(V)$-action factors through the$\iota_0({Z_{\underline{\Omega}}})\times \iota_1(Z_{\bL_{\Omega}}^0)$.\;Since $\pr_{\varpi_L}\flat\det_{\bL_{\Omega}}^{-1}$ is trivial on $Z_{\varpi_L}$,\;we see that the restriction of standard action $\iota_{2}(Z_{\bL_{\Omega}})\times \iota_1(Z_{\bL_{\Omega}}^0)$ on $\iota_{2}(Z_{\varpi_L})$ coincides with the action $\iota_0(Z_{\bL_{\Omega}}(L))|_{Z_{\varpi_L}}=\iota_0(Z_{\varpi_L})$.\;This shows that the $\iota_0(Z_{\varpi_L})\times \iota_1(Z_{\bL_{\Omega}}^0)$-action coincide with standard action $\iota_{2}(Z_{\varpi_L})\times \iota_1(Z_{\bL_{\Omega}}^0)$.\;We can check that the natural action of $\mathfrak{Z}_{\Omega}$ on $B_{\sigma,\ul{\lambda}}(V)$ is commutes with the standard action of $\iota_0(Z_{\varpi_L})\times \iota_1(Z_{\bL_{\Omega}}^0)$.\;

%Furthermore,\;we have a natural map $Z_{\bL_{\Omega}}(L)\rightarrow  \EndO_{\bL_{\Omega}(L)}(\mathbf{c-ind}^{\bL_{\Omega}(L)}_{\bL_{\Omega}^0}\sigma)\cong\mathfrak{Z}_{\Omega}$ by setting $z$ to the homomorphism $v\mapsto z\cdot v$,\;for $v\in\mathbf{c-ind}^{\bL_{\Omega}(L)}_{\bL_{\Omega}^0}\sigma$.\;This also induce a $Z_{\bL_{\Omega}}(L)$-action on $B_{\sigma,\ul{\lambda}}(V)$ via the $\mathfrak{Z}_{\Omega}$-action on $B_{\sigma,\ul{\lambda}}(V)$.\;One can prove that this $Z_{\bL_{\Omega}}(L)$-action coincides with the $\iota_0(Z_{\bL_{\Omega}}(L))$-action.\;This shows that $\iota_0(Z_{\varpi_L})$-action factors through the $\mathfrak{Z}_{\Omega}$-action.\;

%Let $V$ be an admissible locally $\BQ_p$-analytic representation of $G$ over $E$,\;by the construction and \cite[Lem.\;3.1.3]{Ding2021},\;$B_{\sigma,\ul{\lambda}}(V)^\vee$ is a coadmissible module over $\cO(\zuni,E) \tene\cO(\rigch,E)$,\;and hence $B_{\sigma,\ul{\lambda}}(V)$ is an essentially admissible locally $\BQ_p$-analytic representation of $\iota_0(Z_{\varpi_L})\times \iota_1(Z_{\bL_{\Omega}}^0)$.\;The \cite[Lem.\;3.1.7]{Ding2021} shows that the $\mathfrak{Z}_{\Omega}\times \iota_0(Z_{\varpi_L})\times \iota_1(Z_{\bL_{\Omega}}^0)$-module $B_{\sigma,\ul{\lambda}}(V)$ does not depend on the choice $\sigma$.\;Therefore,\;we denote $B_{\Omega,\ul{\lambda}}(V):=B_{\sigma,\ul{\lambda}}(V)$.\;

%By the argument in \cite[Section 3.1.3]{Ding2021},\;we see that there exists a coherent sheaf $\cM_{\Omega,\ul{\lambda}}^0(V)$ over $\zuni\times \rigch$ such that $B_{\Omega,\ul{\lambda}}(V)\vee$ is isomorphic to the global sections of $\cM_{\Omega,\ul{\lambda}}^0(V)$.\;Since the action of $Z_{\underline{\omega}}$ factors  through $\sban$,\;we see that $\cM_{\Omega,\ul{\lambda}}^0(V)$ gives rise to a coherent sheaf $\cM_{\Omega,\ul{\lambda}}(V)$ over $\sban\times \rigerch$ such that \[\Gamma(\sban\times \rigerch,\cM_{\Omega,\ul{\lambda}}(V))\cong B_{\Omega,\ul{\lambda}}(V)^\vee.\]Let $\mathcal{E}_{\Omega,\ul{\lambda}}(V)$ be the Zariski-closed support of $\cM_{\Omega,\ul{\lambda}}(V)$.\;

We now recall the definition of Bernstein paraboline varieties \cite[Section 4.2]{Ding2021}.\;The Bernstein paraboline variety $\defvar$ of type {$(\omepik,\bh)$} is a subspace of $\mathfrak{X}_{\overline{r}}^\Box\times \sbanpik\times\rigch$.\;It contains a subspace $U^\Box_{\omepik,\mathbf{{h}}}(\overline{r})$  consists of the point $(\rho,\underline{x},\undelram)$ such that
\begin{itemize}\label{dfnvardef}
	\item[(1)] $(\underline{x},\undelram)\in \big(\sbanpik\times\rigch\big)^{\gen}$ (the set of generic points in $\sbanpik\times\rigch$,\;see \cite[Section 4.2]{Ding2021}),
	\item[(2)] $D_{\rig}(\rho)$ admits an $\omepik$-filtration $\cF=\fil_{\bullet}^{\cF} D_{\rig}(\rho)$ such that
	\begin{equation}\label{paradefiinj}
		\begin{aligned}
			\gr_i^{\cF}D_{\rig}(\rho)\otimes_{\cR_{k(x),L}}\cR_{k(x),L}((\delta_i^0)^{-1}_{\varpi_L})\hookrightarrow \Delta_{x_i}\otimes_{\cR_{k(x),L}}\cR_{k(x),L}(z^{\bh_{ir}})
		\end{aligned}
	\end{equation}
	and the image has Hodge-Tate weights $(\bh_{(i-1)r+1},\cdots,\bh_{ir})$.\;By using Berger's equivalence of categories \cite[Theorem  A]{berger2008equations} and comparing the Hodge-Tate weights,\;we see that (\ref{paradefiinj}) is equivalent to 
	\begin{equation}\label{dfnvardef1}
		\Delta_{x_i}\otimes_{\cR_{k(x),L}}\cR_{k(x),L}(z^{\bh_{(i-1)r+1}})  \hookrightarrow\gr_i^{\cF}D_{\rig}(\rho)\otimes_{\cR_{k(x),L}}\cR_{k(x),L}((\delta_i^0)^{-1}_{\varpi_L}).
	\end{equation}
\end{itemize}
We define $\defvar$ to be the Zariski-closure of $U^\Box_{\omepik,\mathbf{{h}}}(\overline{r})$ in $\mathfrak{X}_{\overline{r}}^\Box\times \sbanpik\times\rigch$.\;By \cite[Theorem\;4.2.5,\;Corollary\;4.2.5]{Ding2021},\;we have:
\begin{pro}\label{propertyparavar}\hspace{20pt}
	\begin{itemize}
		\item[(1)] The rigid space $\defvar$ is equidimensional of dimension $n^2+\left(\frac{n(n-1)}{2}+k\right)d_L$.\;
		\item[(2)] The set $U^\Box_{\omepik,\mathbf{{h}}}(\overline{r})$ is Zariski-open and  Zariski-dense in $\defvar$.
		\item[(3)] The rigid space  $U^\Box_{\omepik,\mathbf{{h}}}(\overline{r})$ is smooth over $E$,\;and the morphism $${\omega}|_{U^\Box_{\omepik,\mathbf{{h}}}(\overline{r})}:U^\Box_{\omepik,\mathbf{{h}}}(\overline{r})\rightarrow \sbanpik\times\rigch$$ is smooth.\;
		\item[(4)] Let $x=(\rho_x,\underline{x},\undelram)\in \defvar$,\;then $D_{\rig}(\rho_x)$ admits an $\omepik$-filtration $\cF=\{\fil_i^\cF D_{\rig}(\rho_x)\}$ such that,\;for all $1=1,\cdots,s$,\;
		\[\gr_{i}^{\cF}D_{\rig}(\rho_x)\otimes_{\cR_{k(x),L}}\cR_{k(x),L}((\delta_i^0)^{-1}_{\varpi_L})\bigg[\frac{1}{t}\bigg]=\Delta_{x_i}\bigg[\frac{1}{t}\bigg].\]
	\end{itemize}
\end{pro}
\begin{rmk}In general,\;$(\underline{x},((\delta_i^0)_{\varpi_L}z^{\bh_{ir}}))$ is not a parameter (recall Definition \ref{dfnomegafil}) of the $\omepik$-filtration $\cF$ in (4).\;By definition,\;for $x=(\rho_x,\underline{x},\undelram)\in U^\Box_{\omepik,\mathbf{{h}}}(\overline{r})$,\;we see that $(\underline{x},((\delta_i^0)_{\varpi_L}z^{\bh_{ir}}))$ is a parameter of  $\omepik$-filtration  $\cF$.\;
\end{rmk}


The (purely local) Bernstein paraboline variety is  closely related to the (global) patched Bernstein eigenvariety.\;Consider the composition
\begin{equation}\label{composition}
	\begin{aligned}
		\bersteineigenvarpik\hooklongrightarrow \FX_{\overline{\rho}^{\fp}}^\Box&\times \mathfrak{X}_{\overline{r}}^\Box \times \sbanpik\times \rigch\\
		&\xrightarrow{\iota_{\omepik}} \FX_{\overline{\rho}^{\fp}}^\Box\times \mathfrak{X}_{\overline{r}}^\Box \times \sbanpik\times \rigch.
	\end{aligned}
\end{equation}
where $\iota_{\omepik}:\sbanpik\xrightarrow{\sim}\sbanpik$ is the isomorphism such that $\pi_{\iota_{\omepik}(x)_i}=\pi_{x_i}\otimes_E\unr(q_L^{(i-1)r+\frac{r-1}{2}})\flat\det$ for $x=(x_i)_{1\leq i\leq k}\in \sbanpik$.\;An easy variation of the proof of \cite[Theorem 3.3.9]{Ding2021} asserts that
\begin{pro}\label{linkingpatched}The composition in (\ref{composition}) factors through $\FX_{\overline{\rho}^{\fp}}^\Box\times \defvarrho$,\;i.e.,\;
	\begin{equation}\label{mapbersteineigenvarpikdefvarrho}
		\Lambda:\bersteineigenvarpik\hooklongrightarrow \FX_{\overline{\rho}^{\fp}}^\Box\times \iota_{\omepik}^{-1}(\defvar).
	\end{equation}	
	It induces an isomorphism between $\bersteineigenvarpik$ with a union of irreducible components of $\FX_{\overline{\rho}^{\fp}}^\Box\times \iota_{\omepik}^{-1}(\defvar)$ equipped with the reduced closed rigid subspace structure.\;
\end{pro}



%\begin{rmk}\label{underthmlinkingpatched}In some sense,\;we can view deformation variety as purely local version of the patched Bernstein eigenvariety,\;But the constructions of  two objects are quite different.\;The first one is the support of the coherent sheaf defined by applying the ``Jacquet-Emerton'' functor to the locally analytic representation $\Pi_{\infty}^{R_{\infty}-\ana}$.\;In general,\;there is no purely local analogue of the patched module.\;\end{rmk}



%By definition,\;we have the following natural morphisms:
%\begin{equation}
%\begin{aligned}
%& {\omega}:\defvar\rightarrow \sbanpik\times\rigch,\;\\
%& \zeta:\defvar\rightarrow \mathfrak{X}_{\overline{r}}^\Box,\\
%\item ${\omega}=\varsigma_{\Omega}\flat\widetilde{\omega}:\defvar\rightarrow \sbanpik\times\rigch \rightarrow \rigerchl$,
%& \omega^1:\defvar\rightarrow \sbanpik\times\rigch\xrightarrow{\pr_1} \sban,\\
%& \omega^2:\defvar\rightarrow \sbanpik\times\rigch\xrightarrow{\pr_2} \rigerch.
%\end{aligned}
%\end{equation}
%Let $y=(\rho,\underline{x},\undelram)\in \defvar$,\;the tangent maps of the above morphisms are described as follows.\;
%\begin{itemize}
%\item $d\omega_y=d\omega^1_y\oplus d\omega^2_y:T_{\defvar,y}\rightarrow T_{\sbanpik\times\rigch,\omega(y)}= T_{\sban,\underline{x}}\oplus T_{\rigerch,\undelram}$.
%\item $d{\zeta}:T_{\defvar,x}\rightarrow T_{\mathfrak{X}_{\overline{r}}^\Box,\rho}$.\;Let $\bv\in T_{\defvar,x}$,\;seeing $d{\zeta}(\bv)$ as a element in $\ext^1(\rho,\rho)\cong \ext^1_{(\varphi,\Gamma)}(D_{\rig}(\rho),D_{\rig}(\rho))$ as an $k(x)[\epsilon]/\epsilon^2$-valued representation of $\gal_L$.\;Then the set of the Sen's $\tau$-weight of $d{\zeta}(\bv)$ is the set $\{\wt_\tau(\delta_{\bv,i}^0)+h_{j_i,\tau}|i=1,\cdots,r,j_i=s_{i-1}+1,\cdots.s_i\}$.\;        
%\item $d\widetilde{\omega}:T_{\defvar,x}\rightarrow T_{\sbanpik\times\rigch,\widetilde{\omega}(x)}$,\;
%\item $d\omega^1_y:T_{\defvar,x}\rightarrow T_{\sban,\underline{x}}$.\;
%For $d\omega_1:T_{\defvar,x}\rightarrow T_{\sban,\underline{x}}$,\;
%\item $d\omega_2:T_{\defvar,x}\rightarrow T_{\rigerch,\undelram}$.\;
%For $\bv\in T_{\defvar,x}$,\;we can see $\bv$ as a $E[\epsilon]/\epsilon^2$-valued point of $X$,\;and the composition
%\[\spr \big(E[\epsilon]/\epsilon^2\big)\xrightarrow{\bv} \defvar \rightarrow \rigerch\]gives rise to continuous character $\delta_{\bv}^0=\boxtimes_{i=1}^k\delta_{\bv,i}^0:{Z}_{L_\Omega}(\cO_L)\longrightarrow \big(E[\epsilon]/\epsilon^2\big)^\times$ such that $\delta_{\bv}^0\equiv \delta^0\mod \epsilon$.\;Then the map $d\omega_2$ sets $\bv\in T_{\defvar,x}$ to $(\wt_\tau(\delta_{\bv,i}^0(\delta_{i}^{0})^{-1}))_{1\leq i \leq k,\tau\in\Sigma_L}$.\;In \S 2.7,\;we see that $\{\psi_{\sigma,L}\}_{\sigma\in \Sigma_{L}}$ gives a basis of $\homo(\cO_L^{\times},E)$,\;then we can obtain the isomorphism $$T_{\rigerch,\undelram}\cong k(x)^{d_Lr}\xrightarrow{\sim}\homo({Z}_{L_\Omega}(\cO_L),k(x))=\prod_{i=1}^k\homo(\cO_L^{\times},E)$$ by sending $(a_{i,\tau})_{1\leq i\leq k,\tau\in\Sigma_L}$ to $\prod_{i=1}^k\big(\sum_{\tau\in\Sigma_L} a_{i,\tau}\psi_{i,\tau,L}\big)$,\;where.\;Using this isomorphism,\;the elements $d\omega_2(\bv)$ can be rewritten by $((\delta_{\bv,i}^0(\delta_{i}^{0})^{-1}-1)/\epsilon)_{1\leq i \leq k}$.\;
%\end{itemize}Let $X$ be a subspace of $\defvar$.\;The above discussion can be applied to $T_{X,x}$ by restriction.\;

%Since $\varsigma_\Omega$ is finite,\;the tangent map $d\varsigma_\Omega:T_{\sbanpik\times\rigch,\omega(x)}\xrightarrow{\sim} T_{\rigerchl,\omega(x)}$ induces an isomorphism.\;Therefore,\;we can identity $d\omega$ and $d\widetilde{\omega}$ via this isomorphism.\;

Let $W_{u}(\underline{x},\bh)=\ker(\mathrm{PS}_{u}(\underline{x},\bh)\twoheadrightarrow \mathrm{PS}_{u}(\underline{x},\bh))$ and let $V_u$ be the intersection of $I^G_{\op_{S^u_0}(L)}\pi_{\sm}(\underline{x}^u)\eta_{S^u_0}\widetilde{\delta}_u^{\univ}$  with  $W_{u}(\underline{x},\bh)$. Note that $W^{\lalg}_{u}(\underline{x},\bh):=W_{u}(\underline{x},\bh)\cap \mathrm{PS}_{u}^{\lalg}(\underline{x},\bh)$ is independent on the choice of $u$.\;Similar to the argument in the proof of \cite[Proposition 4.7 (a)]{2019DINGSimple},\;we see that $\iota_u$  factors through the projection $I^G_{\op_{S^u_0}(L)}\pi_{\sm}(\underline{x}^u)\eta_{S^u_0}\widetilde{\delta}_u^{\univ}\twoheadrightarrow I^G_{\op_{S^u_0}(L)}\pi_{\sm}(\underline{x}^u)\eta_{S^u_0}\widetilde{\delta}_u^{\univ}/V_w$ (based on the fact $y_u$ is non-critical).\;




\begin{proof}
	We prove this lemma by induction on $s=n$.\;Assume that $N_1:=\cR_{E,L}(\delta_1^u)-\cdots-\cR_{E,L}(\delta_{n-1}^u)$ satisfies the conclusion.\;There exists at most one $1\leq q\leq n-1$ such that $u(n)=u(q)+1$.\;Let $N_2=[\cR_{E,L}(\delta_1^u)-\cdots-\cR_{E,L}(\delta_q^u)]\hookrightarrow N_1$ and let $N_1\rightarrow N_2'=[\cR_{E,L}(\delta_q^u)-\cdots-\cR_{E,L}(\delta_{n-1}^u)]\twoheadrightarrow N_3':=N_2'/\cR_{E,L}(\delta_q^u)=[\cR_{E,L}(\delta_{q+1}^u)-\cdots-\cR_{E,L}(\delta_{n-1}^u)]$.\;We consider the following commutative diagram:
	\begin{equation}
		\xymatrix{
			\ext^1(\cR_{E,L}(\delta_{n}^u),N'_2) \ar@{=}[d]  &\times & \ext^1(N'_2,N'_3)\ar[r]^{\cup}  &  \ext^2(\cR_{E,L}(\delta_{n}^u),N'_3)=0\\
			\ext^1(\cR_{E,L}(\delta_{n}^u),N_2') \ar@{=}[d] &\times & \ext^1(N_2',N_2')\ar[r]^{\cup} \ar[u]^{\kappa_1}  &  \ext^2(\cR_{E,L}(\delta_{n}^u),N_2') \ar@{=}[d] \ar[u]\\
			\ext^1(\cR_{E,L}(\delta_{n}^u),N'_2) \ar[d]^{g_1}    &\times & \ext^1(N'_2,\cR_{E,L}(\delta_{q}^u))\ar[r]^{\cup} \ar[u]^{\kappa_1'} &  \ext^2(\cR_{E,L}(\delta_{n}^u),\cR_{E,L}(\delta_{q}^u)) \ar@{=}[d] \\
			\ext^1(\cR_{E,L}(\delta_{n}^u),N'_3)   &\times & \ext^1(N'_3,\cR_{E,L}(\delta_{q}^u))\ar[r]^{\cup} \ar[u]^{\kappa_1} &  \ext^2(\cR_{E,L}(\delta_{n}^u),\cR_{E,L}(\delta_{q}^u)). }
	\end{equation}
	Let $D''=\Dpik/\ker(N_1\rightarrow N_2')$.\;Let $\cL_{\mathrm{FM}}(D'':N_2'):=(E[D''])^{\perp}$ (through the top cup-product),\;and let $l_{\mathrm{FM}}(D'':N_2'):=(E[D''])^{\perp}\subseteq \ext^1(N_2',\cR_{E,L}(\delta_{q}^u))$ (via the second row).\;We claim that 
	$\cL_{\mathrm{FM}}(D'':N_2'):=(E[D''])^{\perp}$ (resp.,\;$l_{\mathrm{FM}}(D'':N_2')$) has co-dimension $1$ in $\ext^1(N_2',N_2')$ (resp.,\;$\ext^1(N_2',\cR_{E,L}(\delta_{q}^u))$).\;
	
	We describe these subspaces explicitly.\;It easy to see that $\ext^2(N'_2,\cR_{E,L}(\delta_{q}^u))=0$ so that $\kappa_1$ is surjective.\;On the other hand,\;$\kappa_1'$ is injective since $\homo(N_2',N_2')\xrightarrow{\sim }\homo(N_2',N_3')$ is an ismorohism.\;Therefore,\;we have a short exact sequence of $E$-vector spaces:
	\[0\rightarrow l_{\mathrm{FM}}(D'':N_2')\rightarrow \cL_{\mathrm{FM}}(D'':N_2')\rightarrow \ext^1(N'_2,N'_3)\rightarrow 0.\]
	Note that $\homo(N'_2,\cR_{E,L}(\delta_{q}^u))=0$,\;we get that $\ker(\kappa_1)\cong \homo(\cR_{E,L}(\delta_{q}^u),\cR_{E,L}(\delta_{q}^u))\cong E$,\;so that $\dim_E\mathrm{Im}(\kappa_1)=\rk(N_3')-1$.\;Consider the following commutative diagram:
	\begin{equation}\label{N2N3ANDRANK1}
		\xymatrix{
			\ext^1(\cR_{E,L}(\delta_{n}^u),\cR_{E,L}(\delta_{q}^u)) \ar[d]^{g_1'} &\times & \ext^1(\cR_{E,L}(\delta_{q}^u),\cR_{E,L}(\delta_{q}^u))\ar[r]^{\cup}   &  \ext^2(\cR_{E,L}(\delta_{n}^u),\cR_{E,L}(\delta_{q}^u)) \ar@{=}[d] \\
			\ext^1(\cR_{E,L}(\delta_{n}^u),N'_2) \ar[d]^{g_1}    &\times & \ext^1(N'_2,\cR_{E,L}(\delta_{q}^u))\ar[r]^{\cup} \ar[u]^{\kappa_1''} &  \ext^2(\cR_{E,L}(\delta_{n}^u),\cR_{E,L}(\delta_{q}^u)) \ar@{=}[d] \\
			\ext^1(\cR_{E,L}(\delta_{n}^u),N'_3)   &\times & \ext^1(N'_3,\cR_{E,L}(\delta_{q}^u))\ar[r]^{\cup} \ar[u]^{\kappa_1} &  \ext^2(\cR_{E,L}(\delta_{n}^u),\cR_{E,L}(\delta_{q}^u)). }
	\end{equation}
	Since $\ext^2(N'_3,\cR_{E,L}(\delta_{q}^u))=0$,\;$\kappa_1''$ is surjective.\;Since $\homo(\cR_{E,L}(\delta_{n}^u),N'_3)=0$,\;$g_1''$ is injective.\;Observe the following commutative diagram:
	\begin{equation}
		\xymatrix{
			\ext^1(\cR_{E,L}(\delta_{n}^u),N'_3) \ar@{=}[d]   &\times & \ext^1(N'_3,\cR_{E,L}(\delta\epsilon))\ar[r]^{\cup}  &  \ext^2(\cR_{E,L}(\delta_{n}^u),\cR_{E,L}(\delta\epsilon))  \\
			\ext^1(\cR_{E,L}(\delta_{n}^u),N'_3)   &\times & \ext^1(N'_3,\cR_{E,L}(\delta_{q}^u))\ar[r]^{\cup} \ar[u]^{\eta} &  \ext^2(\cR_{E,L}(\delta_{n}^u),\cR_{E,L}(\delta_{q}^u))\ar[u]^{\eta'}_{\sim}, }
	\end{equation}
	where $\eta$ and $\eta'$ are induced by the injection  $\cR_{E,L}(\delta_{q}^u)\hookrightarrow \cR_{E,L}(\delta_n^u\epsilon)$ of $(\varphi,\Gamma)$-modules.\;The top cup-product is perfect by the Tate-duality.\;On the other hand,\;by \cite[Lemma 5.1.1]{breuil2020probleme},\;we have
	\begin{equation}
		\begin{aligned}
			&\;\homo((N'_3)^{\vee}\otimes_{\cR_{E,L}}\cR_{E,L}(\delta_{q}^u),(N'_3)^{\vee}_1\otimes_{\cR_{E,L}}\cR_{E,L}(\delta_{n}^u\epsilon))\\
			\cong&\; \hH^0(\gal_{L},W_{\dr}^+((N'_3)^{\vee}\otimes_{\cR_{E,L}}\cR_{E,L}(\delta\epsilon))/W_{\dr}^+((N'_3)^{\vee}\otimes_{\cR_{E,L}}\cR_{E,L}(\delta_{q}^u))).\;
		\end{aligned}
	\end{equation}
	Then we oabtain $\dim_E\ker(\eta)=\rk(N'_3)$ and thus $\eta=0$,\;which implies that $\ext^1(N'_3,\cR_{E,L}(\delta_{q}^u))\subseteq (Eg_1([D'']))^{\perp}$ and $\kappa_1''(\ext^1(N'_3,\cR_{E,L}(\delta_{q}^u)))\subseteq l_{\mathrm{FM}}(D'':N_2') $.\;Observe the following commutative diagram:
	\begin{equation}
		\xymatrix{
			\ext^1(\cR_{E,L}(\delta_{n}^u),N'_2) \ar@{=}[d]   &\times & \ext^1(N'_2,\cR_{E,L}(\delta\epsilon))\ar[r]^{\cup}  &  \ext^2(\cR_{E,L}(\delta_{n}^u),\cR_{E,L}(\delta\epsilon))  \\
			\ext^1(\cR_{E,L}(\delta_{n}^u),N'_2)   &\times & \ext^1(N'_2,\cR_{E,L}(\delta_{q}^u))\ar[r]^{\cup} \ar[u]^{\eta_1} &  \ext^2(\cR_{E,L}(\delta_{n}^u),\cR_{E,L}(\delta_{q}^u))\ar[u]^{\eta'_1}_{\sim}, }
	\end{equation}
	where $\eta$ and $\eta'$ are induced by the injection  $\cR_{E,L}(\delta_{q}^u)\hookrightarrow \cR_{E,L}(\delta_{n}^u\epsilon)$ of $(\varphi,\Gamma)$-modules.\;The top cup-product is perfect by the Tate-duality.\;On the other hand,\;by \cite[Lemma 5.1.1]{breuil2020probleme},\;we have
	\begin{equation}
		\begin{aligned}
			&\;\homo((N'_3)^{\vee}\otimes_{\cR_{E,L}}\cR_{E,L}(\delta_{q}^u),(N'_3)^{\vee}_1\otimes_{\cR_{E,L}}\cR_{E,L}(\delta\epsilon))\\
			\cong&\; \hH^0(\gal_{L},W_{\dr}^+((N'_2)^{\vee}\otimes_{\cR_{E,L}}\cR_{E,L}(\delta\epsilon))/W_{\dr}^+((N'_2)^{\vee}\otimes_{\cR_{E,L}}\cR_{E,L}(\delta_{q}^u)))\cong E^{\rk(N'_2)-1}.\;
		\end{aligned}
	\end{equation}
	Then we oabtain $\dim_E\ker(\eta_1)=\rk(N'_2)-2$.\;Therefore,\;$\ker(\eta_1)\subseteq\ext^1(\cR_{E,L}(\delta_{n}^u),N'_2)^{\perp}$.\;
	
	Let $e_{\mathrm{cris}}$ be the basis of $\ext^1_e(\cR_{E,L}(\delta_{n}^u),\cR_{E,L}(\delta_{q}^u)) $ (note that it is $1$-demensional) and $e_{\mathrm{cris}},e_{g}$ be the basis of $\ext^1(\cR_{E,L}(\delta_{n}^u),\cR_{E,L}(\delta_{q}^u))$.\;Recall that $\ext^1(\cR_{E,L}(\delta_{q}^u),\cR_{E,L}(\delta_{q}^u))=E\langle \val_p,\log_p\rangle$.\;Extend $\val_p,\log_p$ to a basis $\val_p,\log_p,e_1,\cdots,e_{\rk(N_2')-2}$ of $\ext^1(\cR_{E,L}(\delta_{n}^u),N'_2)$.\;Since the top bottom cup-product in 
	(\ref{N2N3ANDRANK1}) is prefect,\;we see that $E\langle \val_p,\log_p\rangle\subseteq \ext^1(\cR_{E,L}(\delta_{n}^u),N'_3)^{\perp}$ and $\val_p,\log_p$ are not contained in the subspace $\ext^1(\cR_{E,L}(\delta_{n}^u),N'_2)^{\perp}$ (note that the injection $g_1'$ in (\ref{N2N3ANDRANK1})),\;this implies $\ker(\eta_1)=E\langle e_1,\cdots,e_{\rk(N_2')-2}\rangle=\ext^1(N'_3,\cR_{E,L}(\delta_{q}^u))$.\;Note that $E[D'']=E(e_{\mathrm{cris}}+L)$ for some vector $L\in \ext^1(\cR_{E,L}(\delta_{n}^u),N'_3)$.\;So we see that $\kappa_1''(l_{\mathrm{FM}}(D'':N_2'))$ equals to $E\val_p$ by using the fact that $\langle\val_p,L\rangle=\langle\log_p,L\rangle=0$ and $E\val_p=(Ee_{\mathrm{cris}})^\perp$ via the top bottom cup-product in 
	(\ref{N2N3ANDRANK1}).\;In particular,\;$\dim_El_{\mathrm{FM}}(D'':N_2')=(\rk N_3'-1)+1=1\rk N_2'-1$ and $\kappa_1''(l_{\mathrm{FM}}(D'':N_2'))=E\val_p\subseteq \ext^1(\cR_{E,L}(\delta_{q}^u),\cR_{E,L}(\delta_{q}^u))$.\;Finally,\;note that the so-called Colmez-Greenberg-Stevens formula(see \cite[Theorem 2.7]{HigherLinvariantsGL3(Qp)})  implies:\;for the tuple $(\widetilde{N_2'},\psi)\in \ext^{1}(N'_2,N_2')\times \homo(L^{\times},E)$,\;there exists a deformation  $\widetilde{N'_2}-\cR_{ E[\epsilon]/\epsilon^2}(\delta(1+\psi\epsilon))$ of $D''$ over $E[\epsilon]/\epsilon^2$  iff
	$\widetilde{N_2'}\otimes_{\cR_{E[\epsilon]/\epsilon^2}}(\cR_{E[\epsilon]/\epsilon^2}(1-\psi\epsilon))\in\cL_{\mathrm{FM}}(D'':N_2')$,\;thus $\psi_q-\psi_{s}\in E\val_p$.\;This completes the proof of $(2)$.\;
\end{proof}





%For any $u\in \sW_s$,\;and $1\leq j\leq s$,\;we define the following $E$-vector subspace of $\bigwedge_E^{n-t_j^u}\df(\Dpik_{\sigma}^{\flat,u})$:
%\begin{equation}
%	\begin{aligned}
	%			\fil_j^{\max}(\Dpik_{\sigma}^{\flat,u})&:= \fil^{\bh_{n-r_{u(s)},\sigma}}(\Dpik_{\sigma}^{\flat,u})\wedge \fil^{\bh_{n-r_{u(s)}-r_{u(s-1)},\sigma}}(\Dpik_{\sigma}^{\flat,u})\wedge\cdots\wedge \fil^{\bh_{n-\sum_{l=j+1}^sr_{u(l)},\sigma}}(\Dpik_{\sigma}^{\flat,u})\\
	%		&\xrightarrow{\sim}\text{$\bigwedge_E^{n-t_j^u}\fil^{\bh_{n-\sum_{l=j+1}^sr_{u(l)},\sigma}}(\Dpik_{\sigma}^{\flat,u})\subseteq\bigwedge_E^{n-t_j^u}\df(\Dpik_{\sigma}^{\flat,u})$} .\;
	%		\end{aligned}
%\end{equation}
%By definition,\;for $i\in\Delta'$ and $u\in \cI_{i}^{\flat}$,\;$\fil_{l_i(u)}^{\max}(\Dpik_{\sigma}^{\flat,u})$ is independent on the choice of $u$.\;Moreover,\;we have 
%\[\fil_i^{\max}(D_{\sigma})\hookrightarrow \fil_{l_i(u)}^{\max}(\Dpik_{\sigma}^{\flat,u}).\]
%\begin{lem}
%For $i\in\Delta'$ and $u\in \cI_{i}^{\flat}$,\;the coefficient of $E\be_{I_i(u)^c}$ in $\fil_i^{\max}(D_{\sigma})$ is non-zero.\;
%\end{lem}


For $1\leq j \leq s$,\;set $t_j^{u,c}:=\sum_{i=j}^sr_{u^{-1}(j)}$.\;Define a subspace
\[\homo^{\flat}_{\fil^{\flat,u}}(D_{\sigma},D_{\sigma}):=\left\{f\in\homo_E(D_{\sigma},\fil_H^{-\bh_{i,\sigma},\flat,u}(D_{\sigma})):f|_{\gr_{H}^{-\bh_{i,\sigma},\flat,u}(D_{\sigma})} \text{\;scalar},\;i=t_j^{u,c},\forall\;1\leq j \leq s \right\}.\]
For $i=t_j^{u,c}$ and $1\leq j \leq s$,\;recall the surjections in \cite[(68) \& (69)]{BDcritical25} and the isomorphism \cite[(69)]{BDcritical25}:
\begin{equation}
	\begin{aligned}
		\homo_E\Big(\bigwedge\nolimits_E^{\!n-i}\!D_{\sigma},&\bigwedge_E^{n-i}\fil_H^{-\bh_{i,\sigma},\flat,u}(D_{\sigma})\Big)\twoheadrightarrow \homo_E\Big(\bigwedge\nolimits_E^{\!n-i-1}\!\fil_H^{-\bh_{i,\sigma},\flat,u}(D_{\sigma})\wedge D_{\sigma} ,\bigwedge\nolimits_E^{\!n-i}\!\fil_H^{-\bh_{i,\sigma},\flat,u}(D_{\sigma})\Big)\\
		&\xleftarrow{\sim}\left\{f\in\homo_E(D_{\sigma},\fil_H^{-\bh_{i,\sigma},\flat,u}(D_{\sigma})):f|_{\fil_H^{-\bh_{i,\sigma},\flat,u}(D_{\sigma})} \text{scalar}\right\}\hookrightarrow \homo_{\fil}(D_{\sigma},D_{\sigma}),
	\end{aligned}
\end{equation}
the third inclusion is given by \cite[Lemma 2.2.5 (ii)]{BDcritical25}.\;By \cite[Lemma 2.2.5 (ii)]{BDcritical25},\;we have a surjection
\[\bigoplus_{j=1}^s\left\{f\in\homo_E(D_{\sigma},\fil_H^{-\bh_{i,\sigma},\flat,u}(D_{\sigma})):f|_{\fil_H^{-\bh_{i,\sigma},\flat,u}(D_{\sigma})} \text{scalar}\right\}\twoheadrightarrow \homo^{\flat}_{\fil^{\flat,u}}(D_{\sigma},D_{\sigma}).\]


With respect to the isomorphism $B^{\ast}_{\dr}\otimes_{\bQ_p}E\cong \prod_{\sigma\in \Sigma_L}B^{\ast}_{\dr,\sigma}$ ($B^{\ast}_{\dr,\tau}:=B_{\dr}\otimes_{L,\tau}E$) for $\ast\in\{\emptyset,+\}$,\;we have $\bW^{\ast}_{\Dpik}\cong\prod_{\sigma\in \Sigma_L}\bW^{\ast}_{\Dpik,\sigma}$  for $\ast\in\{\emptyset,+\}$.\;


If we consider the composition $E_1^{s-1}\hookrightarrow \Dpik \twoheadrightarrow (E_1^{s-1})'\hookrightarrow (\Dpik^{\flat}_{w_{1,s}})_2^{s}$,\;then it factors through $E_1^{s-1}\hookrightarrow(\Dpik^{\flat}_{1})_1^{s-1}$,\;and thus  a map
\begin{equation}\label{iotaDE-flat}
	\iota^{\flat}_{\Dpik}:(\Dpik^{\flat}_{1})_1^{s-1}\rightarrow (\Dpik^{\flat}_{w_{1,s}})_2^{s}.\;
\end{equation}
Consider the following cup product:
\begin{equation}\label{cupEcrystalline}
	\begin{aligned}
		&\ext_{(\varphi,\Gamma)}^1(\Delta_s[n],(\Dpik^{\flat}_{1})_1^{s-1})\times \homo((\Dpik^{\flat}_{1})_1^{s-1},(\Dpik^{\flat}_{w_{1,s}})_2^{s})\rightarrow \ext_{(\varphi,\Gamma)}^1(\Delta_s[n],(\Dpik^{\flat}_{w_{1,s}})_2^{s}).
	\end{aligned}
\end{equation}
\begin{pro}\label{cupproductPresflat}Under such cup product,\;$E[\Dpik]\subseteq [\iota_{\Dpik}]^{\perp}$.\;If $L=\bQ_p$ and $r_s=r_{s-1}=1$,\;then $E[\Dpik]=[\iota_{\Dpik}]^{\perp}$,\;and $\Dpik$ is determined by $(\Dpik^{\flat}_{1})_1^{s-1}$,\;$(\Dpik^{\flat}_{w_{1,s}})_2^{s}$,\;$E_s$ and $\iota^{\flat}_{\Dpik}$.\;
\end{pro}

\begin{proof}
	We follow the route in \cite[Proposition 2.4]{ParaDing2024}.\;As $\iota_{\Dpik}(E_s)$ factors through $\Dpik$,\;the map induced by the paring $\langle-,\iota_{\Dpik}\rangle$ is equal to the following composition
	\[\ext^1(\Delta_s[n],(\Dpik^{\flat}_{1})_1^{s-1})\xleftarrow{\sim}\ext^1(\Delta_s[n],E_{1}^{s-1})\rightarrow \ext^1(\Delta_s[n],\Dpik)\rightarrow \ext^1(\Delta_s[n],(\Dpik^{\flat}_{w_{1,s}})_2^{s}).\;\]
	The second map sends $\Dpik$ to zero,\;so $\langle\Dpik,\iota^{\flat}_{\Dpik}\rangle=0$.\;By d\'{e}vissage,\;the kernel of such composition is $\homo(\Delta_s,(\Dpik^{\flat}_{w_{1,s}})_2^{s}/(\Dpik^{\flat}_{1})_1^{s-1})$.\;By \cite[Lemma 5.1.1]{breuil2020probleme},\;we have
	\begin{equation}
		\begin{aligned}
			&\;\homo(\Delta_s[n],(\Dpik^{\flat}_{w_{1,s}})_2^{s}/(\Dpik^{\flat}_{1})_1^{s-1})\cong \hH^0(\gal_L,W_{\dr}^+(\Dpik^{\flat}_{w_{1,s}})_2^{s}\otimes \Delta_s[n]^{\vee})/W_{\dr}^+(\Dpik^{\flat}_{1})_1^{s-1}\otimes \Delta_s[n]^{\vee}))\\
			\cong&\;\hH^0(\gal_L,\oplus_{\tau\in\Sigma_L}\oplus_{1\leq i\leq s-1}(t^{\bh_{t_i+r_s,\tau}-\bh_{n,\tau}}B_{\dr}^+/t^{\bh_{t_i,\tau}-\bh_{n,\tau}}B_{\dr}^+)^{\oplus r_sr_i})=\oplus_{\tau\in\Sigma_L}E^{\oplus r_{s-1}r_s}.\;
		\end{aligned}
	\end{equation}
\end{proof}




For any $w\in \sW_s$,\;there exists a unique $(\varphi,\Gamma)$-module of the following form
\begin{equation}
	\Dpik^{\flat}_w:=\Delta_{\nu^{-1}(1)}[t_1^w]-\Delta_{\nu^{-1}(2)}[t_2^w]-\cdots-\Delta_{\nu^{-1}(s)}[t_s^w].\;
\end{equation}
such that $\Dpik^{\flat}_w[1/t]=\cM_{\Dpik}$ and $\Dpik\hookrightarrow \Dpik^{\flat}_w$.\;All above discussion are also suitable for $\Dpik^{\flat}_w$.\;In this case,\;we obtain
\begin{equation}
	\begin{aligned}
		\nu_{\sigma}:&\;\ext^{1}(\Dpik^{\flat}_w,\Dpik^{\flat}_w)\rightarrow \homo_{\fil^{\flat,w}}(D_{\sigma},D_{\sigma}),\\
		\nu:=\oplus_{\sigma\in\Sigma_L}\nu_{\sigma}:	&\;\ext^{1}(\Dpik^{\flat}_w,\Dpik^{\flat}_w)\rightarrow \bigoplus_{\sigma\in\Sigma_L}\homo_{\fil^{\flat,w}}(D_{\sigma},D_{\sigma}).\;
	\end{aligned}
\end{equation}
By \cite[Lemma 2.4.1]{BDcritical25},\;this morphism has kernel $\ext^{1}_g(\Dpik^{\flat}_w,\Dpik^{\flat}_w)$.\;Therefore,\;for $u\in\sW_s$ and  $?\in\{\flat,0\}$,\;the above morphisms induce
\begin{equation}\label{sendtofilD}
	\begin{aligned}
		\nu^?_{\sigma,u}:&\;\ext^{1,?}_u(\Dpik^{\flat}_w,\Dpik^{\flat}_w)/\ext^{1,?}_{g,u}(\Dpik^{\flat}_w,\Dpik^{\flat}_w)\rightarrow\homo^{\flat}_{\fil^{\flat,w},\bF_u}(D_{\sigma},D_{\sigma}).\;\\
		\nu^?_u:=\oplus_{\sigma\in\Sigma_L}\nu^?_{\sigma,u}:&\;\ext^{1,?}_u(\Dpik^{\flat}_w,\Dpik^{\flat}_w)/\ext^{1,?}_{g,u}(\Dpik^{\flat}_w,\Dpik^{\flat}_w)\rightarrow\bigoplus_{\sigma\in\Sigma_L}\homo^{\flat}_{\fil^{\flat,w},\bF_u}(D_{\sigma},D_{\sigma}).\;
	\end{aligned}
\end{equation}
Under the above fixed basis,\;we identify the space $\homo_{\fil^{\flat,w}}(D_{\sigma},D_{\sigma})$ (resp.,\;$\homo_{\fil^{\flat,w},\bF_u}(D_{\sigma},D_{\sigma})$,\;resp., $\homo^{\flat}_{\fil^{\flat,w},\bF_u}(D_{\sigma},D_{\sigma})$) with $\mathrm{Ad}_{g_{\sigma}}(\fp^w_{S_0,\sigma})$ (resp.,\;$\mathrm{Ad}_{(u^{\sharp})^{-1}}(\fp_{S_0,\sigma})\cap \mathrm{Ad}_{g_{\sigma}}(\fp^w_{S_0,\sigma})$,\;resp.,\;$\mathrm{Ad}_{(u^{\sharp})^{-1}}(\tau_{S_0,\sigma})\cap \mathrm{Ad}_{g_{\sigma}}(\fp^w_{S_0,\sigma})$).\;


For $u\in \sW_s$,\;the same strategy applies to ${\ext}^{1}_u(\Dpik^{\flat}_w,\Dpik^{\flat}_w)$ gives an exact sequence of $E$-vector spaces: 
\[0\rightarrow \prod_{i=1}^s\ext^{1,\flat}_g(\Delta_i,\Delta_i)\rightarrow {\ext}^{1}_u(\Dpik^{\flat}_w,\Dpik^{\flat}_w)/{\ext}^{1,\flat}_{0,u}(\Dpik^{\flat}_w,\Dpik^{\flat}_w)\rightarrow  \mathrm{Im}(\nu^{\flat,w}_u)\rightarrow 0,\]
for $?\in\{\emptyset,\flat,0\}$.\;Moreover,\;the following diagram is commutative:
\begin{equation}
	\xymatrix{
		0   \ar[r]  & \prod_{i=1}^s\ext^{1,\flat}_g(\Delta_i,\Delta_i) \ar@{=}[d]\ar[r]  & {\ext}^{1,?}_u(\Dpik^{\flat}_w,\Dpik^{\flat}_w)/{\ext}^{1,\flat}_{0,u}(\Dpik^{\flat}_w,\Dpik^{\flat}_w) \ar@{^(->}[d]\ar[r]  &  \mathrm{Im}(\nu^{\flat,w,?}_u) \ar@{^(->}[d] \ar[r] & 0  \\
		0    \ar[r] & \prod_{i=1}^s\ext^{1,\flat}_g(\Delta_i,\Delta_i) \ar[r] & {\ext}^{1}(\Dpik^{\flat}_w,\Dpik^{\flat}_w)/{\ext}^{1,\flat}_{0}(\Dpik^{\flat}_w,\Dpik^{\flat}_w)\ar[r]   &  \mathrm{Im}(\nu^{\flat,w}) \ar[r] & 0  }
\end{equation}
where $\ext^{1,\flat}_g(\Delta_i,\Delta_i)\cong\homo_{\mathrm{sm}}(L^{\times},E)$.\;



The following proposition is similar to \cite[Proposition 2.4]{ParaDing2024}.\;
\begin{pro}\label{cupproductPres}We have the following facts.\;
	\begin{itemize}
		\item[(1)] Consider the following cup products:
		\begin{equation}\label{cupEcrystalline}
			\xymatrix{
				\ext_{(\varphi,\Gamma)}^1(E_s,E_{1}^{s-1})  &\times & \homo(E_{1}^{s-1},(E_{1}^{s-1})')\ar[r]^{\cup_{\Dpik}} \ar[d]^{g'} &  \ext_{(\varphi,\Gamma)}^1(E_s,(E_{1}^{s-1})')\\
				\ext_{(\varphi,\Gamma)}^1(\Delta_s[n],E_{1}^{s-1}) \ar[u]^{\sim}  &\times & \homo(E_{1}^{s-1},(E_{1}^{s-1})')\ar[r]^{\cup_{\Dpik}^{\flat}}  &  \ext_{(\varphi,\Gamma)}^1(\Delta_s[n],(E_{1}^{s-1})')\ar[u]
			}.\;
		\end{equation}
		Under such two cup products,\;$E[\Dpik]\subseteq [\iota_{\Dpik}]^{\perp}$.\;
		\item[(2)] Assume that $L=\bQ_p$ and $r_s=1$.\;Then $E[\Dpik]=[\iota_{\Dpik}]^{\perp}$ under two cup products,\;$\Dpik$ is determined by $E_{1}^{s-1}$,\;$(E_{1}^{s-1})'$,\;$E_s\cong \Delta_s[n]$  and $\iota_{\Dpik}$.
	\end{itemize}	
\end{pro}

\begin{proof}
	We follow the route in \cite[Proposition 2.4]{ParaDing2024}.\;As $\iota_{\Dpik}$ factors through $\Dpik$,\;the map induced by the paring $\langle-,\iota_{\Dpik}\rangle$ is equal to the following composition (for $\ast\in\{E_{1}^{s-1},\Delta_s[n]\}$)
	\[\ext^1(\ast,E_{1}^{s-1})\rightarrow \ext^1(\ast,\Dpik)\rightarrow \ext^1(\ast,(E_{1}^{s-1})').\;\]
	The first map sends $\Dpik$ to zero,\;so $\langle\Dpik,\iota_{\Dpik}\rangle=0$.\;This is $(1)$.\;By d\'{e}vissage,\;the kernel of such composition is $\homo(E_s,(E_{1}^{s-1})'/E_{1}^{s-1})$.\;By \cite[Lemma 5.1.1]{breuil2020probleme},\;we have
	\begin{equation}
		\begin{aligned}
			&\;\homo(E_s,(E_{1}^{s-1})'/E_{1}^{s-1})\\
			\cong&\; \hH^0(\gal_L,W_{\dr}^+((E_{1}^{s-1})'\otimes E_s^{\vee})/W_{\dr}^+(E_{1}^{s-1}\otimes E_s^{\vee}))\\
			\cong&\;\hH^0(\gal_L,\oplus_{\tau\in\Sigma_L}\oplus_{1\leq i\leq t_{s-1}}\oplus_{t_{s-1}+1\leq j\leq t_{s-1}+r_s }t^{\bh_{i+r_s,\tau}-\bh_{j,\tau}}B_{\dr}^+/t^{\bh_{i,\tau}-\bh_{j,\tau}}B_{\dr}^+)\\=&\;\oplus_{\tau\in\Sigma_L}\oplus_{t_{s-1}+1-r_s\leq i\leq t_{s-1}}\oplus_{t_{s-1}+1\leq j\leq i+r_s}E.\;
		\end{aligned}
	\end{equation}
	So $\homo(E_s,(E_{1}^{s-1})'/E_{1}^{s-1})$ has dimension $d_L\sum_{j=1}^{ r_s}j=d_L\frac{r_s(r_s+1)}{2}$.\;If $r_s=1$ and $L=\bQ_p$,\;it has dimension $1$,\;so that this kernel is exactly generated by $[\Dpik]$.\;On the other hand,\;by using \cite[Lemma 5.1.1]{breuil2020probleme} again,\;we have
	\begin{equation}
		\begin{aligned}
			&\;\homo(\Delta_s[n],(E_{1}^{s-1})'/E_{1}^{s-1})\\
			\cong&\; \hH^0(\gal_L,W_{\dr}^+((E_{1}^{s-1})'\otimes \Delta_s[n]^{\vee})/W_{\dr}^+(E_{1}^{s-1}\otimes \Delta_s[n]^{\vee}))\\
			\cong&\;\hH^0(\gal_L,\oplus_{\tau\in\Sigma_L}\oplus_{1\leq i\leq t_{s-1}}(t^{\bh_{i+r_s,\tau}-\bh_{n,\tau}}B_{\dr}^+/t^{\bh_{i,\tau}-\bh_{n,\tau}}B_{\dr}^+)^{\oplus r_s})\\=&\;\oplus_{\tau\in\Sigma_L}E.\;
		\end{aligned}
	\end{equation}
	If $L=\bQ_p$ and $r_s=1$,\;it has dimension $1$,\;so that this kernel is exactly generated by $[\Dpik]$.\;This completes the proof of $(2)-(ii)$.\;
\end{proof}
Fix $\sigma\in\Sigma_L$,\;we let $\bI_{\sigma}(\bh)$ be the weight such that $\bI_{\sigma}(\bh)_{\sigma,i}=\bh_{\sigma,i}$ and $\bI_{\sigma}(\bh)_{\tau,i}=\bh_{\tau,n}$ for $\tau\neq \sigma$.\;Similar to \cite[Proposition 2.5]{ParaDing2024} (i.e.,\;using the theory of $E$-$B$-pairs),\;there exists a unique $(\varphi,\Gamma)$-module (up to isomorphism) $\Dpik_{\sigma}$ over $\cR_{E,L}$ such that $\Dpik\hookrightarrow \Dpik_{\sigma}$,\;$\Dpik[1/t]= \Dpik_{\sigma}[1/t]$ and the Hodge-Tate weights of $\Dpik_{\sigma}$ are $\bI_{\sigma}(\bh)$.\;The results in Proposition \ref{homd1c1} are also true by replacing $E_1^{s-1}$,\;$(E_1^{s-1})'$ with $(E_1^{s-1})_{\sigma}$ and $(E_1^{s-1})'_{\sigma}$.\;In a similar way,\;we obtain $\iota_{\Dpik_{\sigma}}:(E_1^{s-1})_{\sigma}\rightarrow (E_1^{s-1})'_{\sigma}$.\;We get the cup-product
\begin{equation}\label{cupEcrystallinesigma}
	\begin{aligned}
		&\ext_{(\varphi,\Gamma)}^1(E_s,(E_1^{s-1})_{\sigma})\times \homo((E_1^{s-1})_{\sigma},(E_{1}^{s-1})_{\sigma}')\rightarrow \ext_{(\varphi,\Gamma)}^1(E_s,(E_{1}^{s-1})_{\sigma}').
	\end{aligned}
\end{equation}
Under the cup product,\;$E[\Dpik_{\sigma}]\subseteq [\iota_{\Dpik_{\sigma}}]^{\perp}$.\;If $r_s=1$,\;we have $E[D_{\sigma}]=[\iota^{-}_{\Dpik_{\sigma}}]^{\perp}$,\;and $\Dpik_{\sigma}$ is determined by $(E_1^{s-1})_{\sigma}$,\;$(E_{1}^{s-1})_{\sigma}'$ and $\iota_{\Dpik_{\sigma}}$.\;

For any $\sigma\in\Sigma_L$,\;we have natural map
\begin{equation}
	\begin{aligned}
		\homo_{\fil}(D_{N,\sigma},D_{M,\sigma})\times\homo(M,N)&\rightarrow \homo_{\fil}(D_{M,\sigma},D_{M,\sigma})\times\homo_{\fil}(D_{M,\sigma},D_{M,\sigma}),\\
		(f,\iota)&\mapsto (\iota_{\sigma}\flat f,f\flat\iota_{\sigma}).\;
	\end{aligned}
\end{equation}


The injection $\iota_{\Dpik}$ induces an injection of filtered $E$-vector spaces:\;$\fil^{\bullet}_{H,\sigma} (D_{M,\sigma})\hookrightarrow \fil^{\bullet}_{H,\sigma} (D_{N,\sigma})$.\;
\begin{lem}The relative position of $\fil^{-\hpi_{\sigma,n}}_{H,\sigma} (D_{N,\sigma})$ and $\fil^{-\hpi_{\sigma,n-1}}_{H,\sigma}(D_{M,\sigma})$ in $\fil^{-\hpi_{\sigma,n-1}}_{H,\sigma} (D_{N,\sigma})$ determines $\iota_{\Dpik}$ in $\homo(M,N)$.\;Moreover,\;$\fil^{-\hpi_{\sigma,n-1}}_{H,\sigma} (D_{N,\sigma})$ is generated by $\fil^{-\hpi_{\sigma,n}}_{H,\sigma} (D_{N,\sigma})$ and $\fil^{-\hpi_{\sigma,n-1}}_{H,\sigma}(D_{M,\sigma})$.\;
\end{lem}
\begin{proof}
	In precise,\;write $\fil^{-\hpi_{\sigma,n-1}}_{H,\sigma}(D_{M,\sigma})=E(e_{n-1,\sigma}+\sum_{1\leq j\leq n-2}\cL_{n-1,j}e_{j,\sigma})$, $\fil^{-\hpi_{\sigma,n}}_{H,\sigma} (D_{N,\sigma})=E(e_{1,\sigma}+\sum_{2\leq j\leq n-1}\cL'_{1,j}e_{j,\sigma})$ and $\fil^{-\hpi_{\sigma,n-1}}_{H,\sigma} (D_{N,\sigma})=\fil^{-\hpi_{\sigma,n}}_{H,\sigma} (D_{N,\sigma})\oplus E(e_{2,\sigma}+\sum_{3\leq j'\leq n-1}\cL'_{2,j}e_{j,\sigma})$.\;Suppose that $\dim_E\homo(M,N)=2$.\;By Proposition \ref{homd1c1},\;and for $i\in\{1,\cdots,n-1\}$ we have an isomorphism $M^{\{1\}}\cong N_{\{1\}}$.\;We get that an isomorphism of induced filtered $E$-vector spaces 
	\[\fil^{\bullet}_{H,\sigma} (D_{M,\sigma}/\fil^{-\hpi_{\sigma,n}}_{H,\sigma} (D_{N,\sigma}))\xrightarrow{\sim }\fil^{\bullet}_{H,\sigma}(D_{N,\sigma}/e_{1,\sigma}).\;\]
	This shows that
	$E(e_{n-1,\sigma}+\sum_{2\leq j\leq n-2}\cL_{n-1,j}e_{j,\sigma})=E(e_{2,\sigma}+\sum_{3\leq j'\leq n-1}\cL'_{2,j}e_{j,\sigma})$.\;Put $\sL:=e_{n-1,\sigma}+\sum_{2\leq j\leq n-2}\cL_{n-1,j}e_{j,\sigma}$.\;We get that 
	\begin{equation}
		\begin{aligned}
			&\fil^{-\hpi_{\sigma,n-1}}_{H,\sigma}(D_{M,\sigma})=E(\sL+\cL_{n-1,1}e_{1,\sigma}),\;\fil^{-\hpi_{\sigma,n}}_{H,\sigma} (D_{N,\sigma})=E(e_{1,\sigma}+\sum_{2\leq j\leq n-1}\cL'_{1,j}e_{j,\sigma}),\;\\
			& \fil^{-\hpi_{\sigma,n-1}}_{H,\sigma} (D_{N,\sigma})=\fil^{-\hpi_{\sigma,n}}_{H,\sigma} (D_{N,\sigma})\oplus E\sL.
		\end{aligned}
	\end{equation}
	But $E(\sL+\cL_{n-1,1}e_{1,\sigma})\neq E(e_{n-1,\sigma}+\sum_{1\leq j\leq n-2}\cL_{n-1,j}e_{j,\sigma})$ by non-critical assumption (otherwise we deduce that $ \fil^{-\hpi_{\sigma,n-1}}_{H,\sigma} (D_{N,\sigma})=\fil^{-\hpi_{\sigma,n}}_{H,\sigma} (D_{N,\sigma})\oplus E\sL=Ee_{1,\sigma}\oplus E\sL$,\;this is impossible).\;This completes the proof.\;
\end{proof}

\begin{itemize}
	\item[(1)] 
	\item[(2)] For $i,j\in \{1,\cdots,s-1\}$,\;we have (put $b=(n-r_s-\frac{r_i+r_j+r_s-1}{2})(r_i+r_j-r_s)$) \[\dim_E(\ext_{\alpha_i}^1(M,M)+\ext_{\alpha_j}^1(M,M))=\left\{
	\begin{array}{ll}
		1+\Big((n-r_s)(n-2r_s)+\frac{1}{2}(r_s^2-r_s)+b\Big)d_L,\;&r_i+r_j\geq r_s\\
		1+\Big((n-r_s)(n-2r_s)+\frac{1}{2}(r_s^2-r_s)\Big)d_L,\;&r_i+r_j\leq r_s
	\end{array}
	\right.\]
\end{itemize}
(note that if $r_i+r_j\geq r_s$,\;then $b>0$).\;

For $(2)$,\;we first compute $\ext_{\alpha_i}^1(M,M)\cap\ext_{\alpha_j}^1(M,M)$.\;By definition,\;such intersection lies in the kernel of composition
\[g_{i,\sigma,j}:\ext^1(M,M)\xrightarrow{\kappa_{i,j}} \ext^1(M_{\{i,j\}},M)\xrightarrow{f_{i,j}} \ext^1(E_i'',M)\oplus \ext^1(E_j'',M).\]
We have the following facts (set $h=r_i+r_j$ for simplicity).\;
\begin{itemize}
	\item[(1)] $\kappa_{i,j}$ is surjective and $\ker(\kappa_{i,j})\cong \ext^1(M^{\{i,j\}},M)$.\;
	\item[(2)] $\dim_E\ker(f_{i,j})=(t_{s-1}h-\frac{r_i(r_i-1)}{2}-\frac{r_j(r_j-1)}{2})d_L-(t_{s-1}h-\frac{h(h-1)}{2})d_L-1=r_ir_jd_L-1
	$.\;
	\item[(3)] We have a short exact sequence $0\rightarrow \ker(\kappa_{i,j})\rightarrow \ker(g_{i,\sigma,j})\rightarrow \ker(f_{i,j})\rightarrow 0$.\;
\end{itemize}
We study the intersection $\ext_{\alpha_i}^1(M,M)\cap\ext_{\alpha_j}^1(M,M)\cap \ker(\kappa_{i,j})$ first.\;Consider the following commutative diagram of short exact sequences:
\begin{equation}
	\xymatrix{
		0   \ar[r]  & \ext^{1}(N_{\{i,j\}},M) \ar[d]\ar[r]  & \ext^{1}(N_{\{i\}},M) \ar[d]\ar[r]  &  \ext^{1}((E_{j}^{(2)})',M) \ar[d]\ar[r] & 0  \\
		0    \ar[r] & \ext^{1}(M^{\{i,j\}},M) \ar[r] & \ext^{1}(M^{\{i\}},M) \ar[r]  & \ext^{1}(E_{j}^{(2)},M) \ar[r] & 0  },
\end{equation}
where the Hodge-Tate weights of $(E_{j}^{(2)})'$ is $\bh_{r_i+r_s+1},\cdots,\bh_{r_i+r_s+1+r_j}$.\; 



is the image  of $\ext^1(N_{{\{i,j\}}},M)\rightarrow\ext^1(M^{\{i,j\}},M)$, 
which has dimension $t_{s-1}(t_{s-1}-h)d_L-(\max\{0,r_s-h\}t_{s-1}-\frac{1}{2}\max\{0,r_s-h\}(r_s+h-1))d_L=t_{s-1}(t_{s-1}
-h-\max\{0,r_s-h\})d_L+\frac{1}{2}\max\{0,r_s-h\}(r_s+h-1)d_L$ by Lemma \ref{lemforTh}.\;We need to study $\ker(g_{i,\sigma,j})\cap \ext_{\alpha_i}^1(M,M)\cap\ext_{\alpha_j}^1(M,M)$.\;Consider the following commutative diagram of short exact sequences:
\begin{equation}
	\xymatrix{
		0   \ar[r]  & \ext^{1}(M^{\{i,j\}},M) \ar[d]\ar[r]  & \ext^{1}(M,M) \ar@{=}[d]\ar[r]  &  \ext^{1}(M_{\{i,j\}},M) \ar[d]^{f_{i,j}} \ar[r] & 0  \\
		0    \ar[r] & \ext^{1}(M^{{\widehat{\ast}}},M) \ar[r] & \ext^{1}(M,M) \ar[r]  & \ext^{1}(M_{\{\ast\}},M) \ar[r] & 0  }  
\end{equation}
for $\ast\in\{i,j\}$.\;Any element $\alpha\in \ker(f_{i,j})$ corresponds to an element $\alpha'$ that belongs to the cokernel of the first vertical map by snake lemma.\;On the other hand,\;we see that $\ker(f_{i,j})$ is the intersection of the image of 
$\ext^{1}(E_{j}^{(2)},M)\rightarrow \ext^{1}(M_{\{i,j\}},M)$ and $\ext^{1}(E_{i}^{(2)},M)\rightarrow \ext^{1}(M_{\{i,j\}},M)$,\;where the Hodge-Tate weights of $E_{j}^{(2)}$ (resp.,\;$E_{i}^{(2)}$) is $\bh_{r_i+1},\cdots,\bh_{r_i+1+r_j}$ (resp.,\;$\bh_{r_j+1},\cdots,\bh_{r_j+1+r_i}$).




For the third vertical map,\;note that $\homo(((E_{j}^{(2)})')^{\vee}\otimes_{\cR_{E,L}}M,(E_{j}^{(2)})^{\vee}\otimes_{\cR_{E,L}}M)$ is computed by 

\begin{equation}
	\begin{aligned}
		&\;\hH^0(\gal_{L},W_{\dr}^+((E_{j}^{(2)})^{\vee}\otimes_{\cR_{E,L}}M)/W_{\dr}^+(((E_{j}^{(2)})')^{\vee}\otimes_{\cR_{E,L}}M))\\
		\cong&\;\hH^0(\gal_{L},\oplus_{\tau\in\Sigma_L}\oplus_{1\leq i\leq t_{s-1}}\oplus_{1\leq l\leq r_j }t^{\bh_{i,\tau}-\bh_{l+r_i,\tau}}B_{\dr,\tau}^+/t^{\bh_{i,\tau}-\bh_{l+r_s+r_i,\tau}}B_{\dr,\tau}^+).\;
	\end{aligned}
\end{equation}
\[T:=\{(i,j):1\leq l+r_i\leq i\leq \min\{t_{s-1},l+r_s+r_i-1\},1\leq l\leq r_j\}.\;\]
If $r_j+r_s+r_i-1\leq t_{s-1}$,\;we see that $|T|=r_jr_s$.\;If $r_j+r_s+r_i-1\geq t_{s-1}$,\;we see that $|T|$ equals to 
\begin{equation*}
	\begin{aligned}
		&\sum_{l=1}^{\min\{t_{s-1}+1-r_s-r_i,r_j\}}r_s+\sum_{l=t_{s-1}+2-r_s-r_i}^{r_j}(t_{s-1}-r_i-l+1)=(t_{s-1}+1-r_s-r_i)r_s+\sum^{-1+r_s}_{l=t_{s-1}-r_i-r_j+1}l,\\
		=&r_jr_s+\frac{1}{2}(n-(r_i+r_j)-2r_s)(r_s+r_i+r_j-1-t_{s-1}).
	\end{aligned}
\end{equation*}



In a conclusion,\;we see that 
\[\dim_E\ext_{\alpha_i}^1(M,M)\cap\ext_{\alpha_j}^1(M,M)=\left\{
\begin{array}{ll}
	t_{s-1}(t_{s-1}
	-h)d_L+r_ir_jd_L-1,\;&h\geq r_s\\
	t_{s-1}(t_{s-1}
	-r_s)d_L+\frac{1}{2}(r_s-h)(r_s+h-1)d_L+r_ir_jd_L-1,\;&h\leq r_s
\end{array}
\right.\]
Thus,\;we obtain (recall $b=(n-r_s-\frac{r_i+r_j+r_s-1}{2})(r_i+r_j-r_s)$)
\[\dim_E(\ext_{\alpha_i}^1(M,M)+\ext_{\alpha_j}^1(M,M))=\left\{
\begin{array}{ll}
	1+t_{s-1}(t_{s-1}-r_s)d_L+\frac{1}{2}(r_s^2-r_s)d_L+bd_L,\;&h\geq r_s\\
	1+t_{s-1}(t_{s-1}
	-r_s)d_L+\frac{1}{2}(r_s^2-r_s)d_L,\;&h\leq r_s
\end{array}
\right.\]

\begin{equation}\label{hoforalphai}
	\begin{aligned}
		&\;\dim_E\homo(N_{{\widehat{i}}}^{\vee}\otimes_{\cR_{E,L}}M,(M^{{\widehat{i}}})^{\vee}\otimes_{\cR_{E,L}}M)
		\cong\hH^0(\gal_{L},W_{\dr}^+((M^{{\widehat{i}}})^{\vee}\otimes_{\cR_{E,L}}M)/W_{\dr}^+(N_{{\widehat{i}}}^{\vee}\otimes_{\cR_{E,L}}M))\\
		\cong&\;\hH^0(\gal_{L},\oplus_{\tau\in\Sigma_L}\oplus_{1\leq i\leq t_{s-1}}\oplus_{1\leq j\leq t_{s-1}-r_i }t^{\bh_{i,\tau}-\bh_{j+r_i,\tau}}B_{\dr,\tau}^+/t^{\bh_{i,\tau}-\bh_{j+r_s,\tau}}B_{\dr,\tau}^+).\;
	\end{aligned}
\end{equation}
\begin{lem}Under the following perfect paring of finite dimensional vector spaces (induced by the wedge product)
	\[\bigwedge\nolimits_{E}^{\! i}\!D\times \bigwedge\nolimits_{E}^{\!n-i}\!D\rightarrow \bigwedge\nolimits_{E}^{\!n}\!D\cong E,\]
	we have $\big(\bigwedge\nolimits_{E}^{\!n-i}\!D\big)^{\flat}\cong\Big(\bigwedge\nolimits_{E}^{\!n-i}\!D^{\flat}\Big)^{\vee}$.\;
\end{lem}


\begin{lem}The subspace $\fil^{-\hpi_{\sigma,n-r_s}}_{H,\sigma} (D_{N,\sigma})$ is generated by $\fil^{-\hpi_{\sigma,n-r_s+1}}_{H,\sigma} (D_{N,\sigma})$ and $\fil^{-\hpi_{\sigma,n-r_s}}_{H,\sigma}(D_{M,\sigma})$. Furthermore,\;the relative position of $\fil^{-\hpi_{\sigma,n-r_s+1}}_{H,\sigma} (D_{N,\sigma})$ and $\fil^{-\hpi_{\sigma,n}}_{H,\sigma}(D_{M,\sigma})$ in $\fil^{-\hpi_{\sigma,n-1}}_{H,\sigma} (D_{N,\sigma})$ determines $\iota_{\Dpik}$ in $\homo(M,N)$.\;
\end{lem}
\begin{proof}
	In precise,\;we write $\fil^{-\hpi_{\sigma,n-r_s}}_{H,\sigma}(D_{M,\sigma})=E\sL$  with $\sL=e_{n-r_s,\sigma}+\sum_{1\leq j\leq n-r_s-1}\cL_{n-r_s,j}e_{j,\sigma}$, $\fil^{-\hpi_{\sigma,n}}_{H,\sigma} (D_{N,\sigma})=E\sL'$ with $\sL'=e_{1,\sigma}+\sum_{2\leq j\leq n-r_s}\cL'_{1,j}e_{j,\sigma}$ and $\fil^{-\hpi_{\sigma,n-1}}_{H,\sigma} (D_{N,\sigma})=\fil^{-\hpi_{\sigma,n}}_{H,\sigma} (D_{N,\sigma})\oplus W $ for some $W\subseteq D_{M,\sigma}/e_{1,\sigma}$.\;Then $\sL=\cL_{n-r_s,1}\sL'+w$ for some $w\in W$ (by non-critical assumption,\;$\cL_{n-r_s,1}\neq 0$).\;This completes the proof of the first assertion.\;
\end{proof}

The following lemma gives some explicit $p$-adic Hodge parameters that are determined by $\iota_{\Dpik}$.\;Since $\iota_{\Dpik}$ induces an injection of filtered $E$-vector spaces:\;$\fil^{\bullet}_{H,\sigma} (D_{M,\sigma})\hookrightarrow \fil^{\bullet}_{H,\sigma} (D_{N,\sigma})$,\;we obtain an inclusion $\fil^{-\hpi_{\sigma,n-r_s}}_{H,\sigma}(D_{M,\sigma})\hookrightarrow \fil^{-\hpi_{\sigma,n-r_s}}_{H,\sigma} (D_{N,\sigma})$.\;


Indeed,\;we have an easier locally analytic representation $[\pi_1^{\diamond}(\underline{x},\bh)-(\pi_{\lalg})^{\oplus\frac{(s-1)(s-2)}{2}}]\hookrightarrow \pi_{\min}(\Dpik)$ that encodes necessary information of Hodge parameters by choosing appropriate ``disjoint" extension lines.\;





%For $x=(\fm_x, \pi_{x},\chi_x)\in \mathcal{E}_{\omepik,\fp,{\bm\lambda}_\bh}^\infty(\overline{\rho})$,\;let $\rho_x$ be the $\gal_L$-representation associated to $x$ via $\fm_x\in (\Spf R_{\infty})^{\rig}=\FX_{\overline{\rho}^{\fp}}^\Box\times \BU^g\times \mathfrak{X}_{\overline{r}}^\Box$ in $\mathfrak{X}_{\overline{r}}^\Box$.\;Let $\mathbf{WD}_x=(\mathbf{WD}_{x,i})_{1\leq i\leq k}$ be the Weil-Deligne representation associated to $x$,\;i.e.,\;$\mathbf{WD}_{x,i}:=\rec(\pi_{x,i})(\frac{1-r}{2}-(i-1)r)$.\;
%We call a generic classical point $x$ is non-critical if $D^{\dag}_{\rig}(\rho_x)$ is non-critical in the sense of Section 2.4,\;with respect to the $\Omega$-filtration defined as in  \cite[(3.30),(3.24))]{Ding2021}.\;%Fix $\Omega$ a cuspidal Bernstein component of $\bL_{\Omega}(L)$,\;and denote by $\mathfrak{Z}_{\Omega}$ the corresponding Bernstein centre (over $E$).\;Let $V$ be an admissible locally analytic representation of $G$ over $E$.\;Then we see that $J_{\bP_\Omega}(V)$ is an essentially admissible representation of  $\bL_{\Omega}(L)$ (\cite{emerton2006jacquet}).\;Let $\underline{\lambda}$ be a integral $\bP_\Omega$-dominant weight of $G$,\;which corresponds an algebraic representation $L(\underline{\lambda})_\Omega$  of $\bL_{\Omega}(L)$ with weight $\underline{\lambda}$.\;We set
%\[J_{\bP_\Omega(L)}(V)_{\underline{\lambda}}:=\homo_{\sbanpik\times\rigch_{\Omega}^D}(L(\underline{\lambda})_\Omega,J_{\bP_\Omega(L)}(V)),\]
%which is closed $\bL_{\Omega}(L)$-subrepresentation of $J_{\bP_\Omega(L)}(V)\otimes_EL(\underline{\lambda})_\Omega^\vee$.\;Let $\sigma$ be  the irreducible smooth representation of $\bL_{\Omega}^0$ as in Section 2.3.\;We now put
%\[B_{\sigma,\ul{\lambda}}(V)\cong \big(J_{\bP_\Omega(L)}(V)_{\underline{\lambda}}\tene \cC^{\BQ_p-\an}(Z_{\bL_{\Omega}}^0,E)\ten \sigma^\vee\big)^{\bL_{\Omega}^0}.\]We explain the $\mathfrak{Z}_{\Omega}\times \iota_0(Z_{\bL_{\Omega}}(L))\times \iota_2(Z_{\bL_{\Omega}}^0)$ on $B_{\sigma,\ul{\lambda}}(V)$ as follows.\; Let $\iota_{2}(Z_{\bL_{\Omega}}(L))$ (resp.,\;$\iota_1(Z_{\bL_{\Omega}}^0)$) denotes the action of $Z_{\bL_{\Omega}}$ (resp.,\;$Z_{\bL_{\Omega}}^0$) on the$J_{\bP_\Omega(L)}(V)_{\underline{\lambda}}\tene \cC^{\BQ_p-\an}(Z_{\bL_{\Omega}}^0,E)$ by only acting on the term $J_{\bP_\Omega(L)}(V)_{\underline{\lambda}}$ (resp.,\;$\cC^{\BQ_p-\an}(Z_{\bL_{\Omega}}^0,E)$).\;The action $\iota_{2}(Z_{\bL_{\Omega}})\times \iota_1(Z_{\bL_{\Omega}}^0)$ is called standard action.\;On the other hand,\;we have\[B_{\sigma,\ul{\lambda}}(V)\cong \homo_{\bL_{\Omega}(L)}(\mathbf{c-ind}^{\bL_{\Omega}(L)}_{\bL_{\Omega}^0}\sigma,J_{\bP_\Omega(L)}(V)_{\underline{\lambda}}\tene \cC^{\BQ_p-\an}(Z_{\bL_{\Omega}}^0,E)).\]Then we deduce that $B_{\Omega,\ul{\lambda}}(V)$ is equipped with a natural action of $\mathfrak{Z}_{\Omega}\cong \EndO_{\bL_{\Omega}(L)}(\mathbf{c-ind}^{\bL_{\Omega}(L)}_{\bL_{\Omega}^0}\sigma)$.\;Meanwhile,\;we also use $\iota_0(Z_{\bL_{\Omega}}(L))$ to denotes the diagonal action of $Z_{\bL_{\Omega}}(L)$ on $B_{\sigma,\ul{\lambda}}(V)$,\;i.e.,\;with the natural action $Z_{\bL_{\Omega}}(L)$ on $J_{\bP_\Omega(L)}(V)_{\underline{\lambda}}\tene \cC^{\BQ_p-\an}(Z_{\bL_{\Omega}}^0,E)$ by acting via
%\[Z_{\bL_{\Omega}}(L)\xrightarrow{(\id,\;\pr_{\varpi_L}\flat\det_{\bL_{\Omega}}^{-1})}\iota_{2}(Z_{\bL_{\Omega}}(L))\times \iota_1(Z_{\bL_{\Omega}}^0).\]
%Then by definition,\;we see that $\iota_0(Z_{\bL_{\Omega}}(L))|_{Z_{\bL_{\Omega}}^0}=\psi_\sigma:=\sigma|_{Z_{\bL_{\Omega}}^0}$.\;This shows that
%$B_{\sigma,\ul{\lambda}}(V)$ is a closed subrepresentation of $\big(J_{\bP_\Omega(L)}(V)_{\underline{\lambda}}\tene \cC^{\BQ_p-\an}(Z_{\bL_{\Omega}}^0,E)\big)[Z_{\bL_{\Omega}}^0=\psi_\sigma]$.\;Therefore,\;we conclude that the $\iota_0(Z_{\bL_{\Omega}}(L))\times \iota_1(Z_{\bL_{\Omega}}^0)$ on $B_{\sigma,\ul{\lambda}}(V)$-action factors through the$\iota_0({Z_{\underline{\Omega}}})\times \iota_1(Z_{\bL_{\Omega}}^0)$.\;Since $\pr_{\varpi_L}\flat\det_{\bL_{\Omega}}^{-1}$ is trivial on $Z_{\varpi_L}$,\;we see that the restriction of standard action $\iota_{2}(Z_{\bL_{\Omega}})\times \iota_1(Z_{\bL_{\Omega}}^0)$ on $\iota_{2}(Z_{\varpi_L})$ coincides with the action $\iota_0(Z_{\bL_{\Omega}}(L))|_{Z_{\varpi_L}}=\iota_0(Z_{\varpi_L})$.\;This shows that the $\iota_0(Z_{\varpi_L})\times \iota_1(Z_{\bL_{\Omega}}^0)$-action coincide with standard action $\iota_{2}(Z_{\varpi_L})\times \iota_1(Z_{\bL_{\Omega}}^0)$.\;We can check that the natural action of $\mathfrak{Z}_{\Omega}$ on $B_{\sigma,\ul{\lambda}}(V)$ is commutes with the standard action of $\iota_0(Z_{\varpi_L})\times \iota_1(Z_{\bL_{\Omega}}^0)$.\;

%Furthermore,\;we have a natural map $Z_{\bL_{\Omega}}(L)\rightarrow  \EndO_{\bL_{\Omega}(L)}(\mathbf{c-ind}^{\bL_{\Omega}(L)}_{\bL_{\Omega}^0}\sigma)\cong\mathfrak{Z}_{\Omega}$ by setting $z$ to the homomorphism $v\mapsto z\cdot v$,\;for $v\in\mathbf{c-ind}^{\bL_{\Omega}(L)}_{\bL_{\Omega}^0}\sigma$.\;This also induce a $Z_{\bL_{\Omega}}(L)$-action on $B_{\sigma,\ul{\lambda}}(V)$ via the $\mathfrak{Z}_{\Omega}$-action on $B_{\sigma,\ul{\lambda}}(V)$.\;One can prove that this $Z_{\bL_{\Omega}}(L)$-action coincides with the $\iota_0(Z_{\bL_{\Omega}}(L))$-action.\;This shows that $\iota_0(Z_{\varpi_L})$-action factors through the $\mathfrak{Z}_{\Omega}$-action.\;

%Let $V$ be an admissible locally $\BQ_p$-analytic representation of $G$ over $E$,\;by the construction and \cite[Lem.\;3.1.3]{Ding2021},\;$B_{\sigma,\ul{\lambda}}(V)^\vee$ is a coadmissible module over $\cO(\zuni,E) \tene\cO(\rigch,E)$,\;and hence $B_{\sigma,\ul{\lambda}}(V)$ is an essentially admissible locally $\BQ_p$-analytic representation of $\iota_0(Z_{\varpi_L})\times \iota_1(Z_{\bL_{\Omega}}^0)$.\;The \cite[Lem.\;3.1.7]{Ding2021} shows that the $\mathfrak{Z}_{\Omega}\times \iota_0(Z_{\varpi_L})\times \iota_1(Z_{\bL_{\Omega}}^0)$-module $B_{\sigma,\ul{\lambda}}(V)$ does not depend on the choice $\sigma$.\;Therefore,\;we denote $B_{\Omega,\ul{\lambda}}(V):=B_{\sigma,\ul{\lambda}}(V)$.\;

%By the argument in \cite[Section 3.1.3]{Ding2021},\;we see that there exists a coherent sheaf $\cM_{\Omega,\ul{\lambda}}^0(V)$ over $\zuni\times \rigch$ such that $B_{\Omega,\ul{\lambda}}(V)\vee$ is isomorphic to the global sections of $\cM_{\Omega,\ul{\lambda}}^0(V)$.\;Since the action of $Z_{\underline{\omega}}$ factors  through $\sban$,\;we see that $\cM_{\Omega,\ul{\lambda}}^0(V)$ gives rise to a coherent sheaf $\cM_{\Omega,\ul{\lambda}}(V)$ over $\sban\times \rigerch$ such that \[\Gamma(\sban\times \rigerch,\cM_{\Omega,\ul{\lambda}}(V))\cong B_{\Omega,\ul{\lambda}}(V)^\vee.\]Let $\mathcal{E}_{\Omega,\ul{\lambda}}(V)$ be the Zariski-closed support of $\cM_{\Omega,\ul{\lambda}}(V)$.\;

Moreover,\;let $\mathrm{PS}^{\lalg}_{u}(\underline{x},\bh)^{\univ}$ (resp.,\;$\mathrm{PS}^{\lalg}_{u}(\underline{x},\bh)^{\univ,\flat}$) be the universal locally algebraic extension of $\ext^1_{\lalg}(\mathrm{PS}^{\lalg}_{u}(\underline{x},\bh),\mathrm{PS}^{\lalg}_{u}(\underline{x},\bh))\otimes_E\mathrm{PS}^{\lalg}_{u}(\underline{x},\bh)$ (resp.,\;$\ext^1_{\lalg}(\mathrm{PS}^{\lalg}_{u}(\underline{x},\bh),\mathrm{PS}^{\lalg}_{u}(\underline{x},\bh))\otimes_E\pi_{\lalg}(\underline{x},\bh)$) by $\mathrm{PS}^{\lalg}_{u}(\underline{x},\bh)$.\;


\begin{rmk}\label{modulispaceforpara}In the sequel,\;for $1\leq i\leq s$,\;let $\{e_{i,j}\}_{1\leq j\leq r_i}$ be a fixed basis of $\df(E_{i})$.\;With respect to the basis $\{e_{1,j}\}_{1\leq j\leq r_1},\cdots,\{e_{s,j}\}_{1\leq j\leq r_s}$ of $\df_{\Dpik}$ (i.e.,\;the $\Omega_{S_0}$-filtration $\cF_1$ on $\df_{\Dpik}$),\;the Hodge filtration $\fil_{\bullet}^{H}D_{\dr}(\Dpik)$ is parameterized by an element $\underline{\sL}(\Dpik)$ in $\bZ_{S_0}\backslash\GLN_{n}/\bB$,\;which we call the $p$-adic parameter of $\Dpik$.\;If $\Dpik$ is \textit{non-critical},\;then
	\[(\underline{\sL}(\Dpik) \bB(E),u\bP_{S_0}(E))\in\GLN_n(E)(1,w_0)(\bB\times \bP_{S_0})(E)\subset (\GLN_n/\bB\times\GLN_n/\bP_{S_0})(E).\]
	Therefore,\;we get that $\underline{\sL}(\Dpik)$ belongs to   $\bZ_{S_0}\backslash u\bB w_0\bP_{S_0}/\bP_{S_0}$ for all $u\in \sW_s$.\;Thus we get that
	\[\underline{\sL}(\Dpik)\in \bigcap_{u\in \sW_s} \bZ_{S_0}\backslash u\bP_{S_0}  \bB 
	w_0/\bB\subseteq \bZ_{S_0}\backslash\GLN_{n}/\bB.\]
\end{rmk}


By Schraen's spectral sequence and \cite[Lemma 2.26]{2019DINGSimple},\;we have  isomorphisms
\begin{equation}
	\begin{aligned}
		\homo\big(\bZ_{S^u_0}(L),E\big)&\xrightarrow{\sim}\ext^1_{G}(\pi_{\lalg}(\underline{x},\bh),\mathrm{PS}_{u}(\underline{x},\bh))\xleftarrow{\sim}\ext^1_{G}(\pi_{\lalg}(\underline{x},\bh),\mathrm{PS}_{u,1}(\underline{x},\bh)).
	\end{aligned}
\end{equation}
For $j\in\Delta_s$ and $\sigma\in \Sigma_L$,\;we have injections 
\[C_{j}:=\oplus_{\sigma\in\Sigma_L}\cF^{G}_{\bP_{\widehat{t_j^u}}}(\overline{L}(-s_{t_j^u,\sigma}\cdot{\lambda}),\mathrm{PS}^{\infty}_{\widehat{t_j^u},u}(\underline{x}))\hookrightarrow \mathrm{PS}_{u,1}(\underline{x},\bh)/\mathrm{PS}^{\lalg}_{u,1}(\underline{x},\bh).\]
For any irreducible consistent $V\neq C_{t_j^u,[u]_j,\sigma}$ of $\cF^{G}_{\bP_{\widehat{t_j^u}}}(\overline{L}(-s_{t_j^u,\sigma}\cdot{\lambda}),\mathrm{PS}^{\infty}_{\widehat{t_j^u},u}(\underline{x}))$,\;i.e.,\;$V=\cF^{G}_{\bP_{\widehat{t_j^u}}}(\overline{L}(-s_{t_j^u,\sigma}\cdot{\lambda}),V')$ for $V'\neq \pi_{\emptyset,\widehat{j},u}$,\;since $\homo_{G}\Big(\pi_{\lalg}(\underline{x},\bh),\big(\ind^G_{\bP_{\widehat{t_j^u}}(L)}V'\big)^{\infty}\Big)=0$,\;we see that
$\ext^1_G(\pi_{\lalg}(\underline{x},\bh),V)=0$.\;We show that $\dim_E\ext^1_{G}\left(\pi_{\lalg}(\underline{x},\bh),\mathrm{PS}^{\lalg}_{u,1}(\underline{x},\bh)\right)=s+1$.\;At first,\;
\[\dim_E\ext^1_{G,Z}\left(\pi_{\lalg}(\underline{x},\bh),\mathrm{PS}^{\lalg}_{u,1}(\underline{x},\bh)\right)=\dim_E\ext^{1,\lalg}_{G,Z}\left(\pi_{\lalg}(\underline{x},\bh),\mathrm{PS}^{\lalg}_{u,1}(\underline{x},\bh)\right)=s-1.\]
By comparing dimensions,\;we have the following exact sequence:
\begin{equation}
	\begin{aligned}
		0\rightarrow \ext^1_{G}\left(\pi_{\lalg}(\underline{x},\bh),\mathrm{PS}^{\lalg}_{u,1}(\underline{x},\bh)\right)\rightarrow &\;\ext^1_{G}\left(\pi_{\lalg}(\underline{x},\bh),\mathrm{PS}_{u,1}(\underline{x},\bh)\right)\\&\rightarrow \oplus_{j\in\Delta_s}  \ext^1_{G}\left(\pi_{\lalg}(\underline{x},\bh),C_j\right)\rightarrow 0 .
	\end{aligned}
\end{equation}
But the last map factors through $\ext^1_{G}\left(\pi_{\lalg}(\underline{x},\bh),\mathrm{PS}_{u,1}(\underline{x},\bh)\right)\rightarrow \ext^1_{G}\left(\pi_{\lalg}(\underline{x},\bh),\mathrm{PS}_{u,1}(\underline{x},\bh)\right)$.\;The result follows.\;

%Write (for $u\in\sW_s$ and  $?\in\{\flat,0\}$)
%\begin{equation}
%	\begin{aligned}
	%		\nu'_{\sigma}:&\;\ext^{1}(\Dpik,\Dpik)\rightarrow \homo_{\fil}(D_{\sigma},D_{\sigma}),\nu':=\oplus_{\sigma\in\Sigma_L}\nu'_{\sigma}:	\;\ext^{1}(\Dpik,\Dpik)\rightarrow \bigoplus_{\sigma\in\Sigma_L}\homo_{\fil}(D_{\sigma},D_{\sigma}),\\
	%			\nu^{',?}_{\sigma,u}:&\;\ext^{1,?}_u(\Dpik,\Dpik)\rightarrow\homo^{\flat}_{\fil,\bF_u}(D_{\sigma},D_{\sigma}),\nu^{',?}_u:=\oplus_{\sigma\in\Sigma_L}\nu^{',?}_{\sigma,u}:\ext^{1,?}_u(\Dpik,\Dpik)\rightarrow\bigoplus_{\sigma\in\Sigma_L}\homo^{\flat}_{\fil,\bF_u}(D_{\sigma},D_{\sigma}).\;
	%	\end{aligned}
%\end{equation}


\begin{lem}\label{generalimagechangelem} For $?\in\{\emptyset,\flat\}$ and $\star\in\{\emptyset,\sigma\}$,\;we have $\coker\nu^{?}_{\flat,1}\xrightarrow{\sim }\coker\nu^{?}_{\flat}$.\;As a direct consequence,\;we have:
	\[\sum_{1\neq u\in \sW_s}\mathrm{Im}(\nu^{?}_{\flat,u})+\mathrm{Im}(\nu^{?}_{\flat,1})=\sum_{1\neq u\in \sW_s}\homo^{?}_{\fil,\bF_u}(D_{\star},D_{\star})+\mathrm{Im}(\nu^{?}_{\flat,1}),\]	
	where $\homo^{?}_{\fil,\bF_u}(D,D)$ means $\oplus_{\sigma\in \Sigma_L}\homo^{?}_{\fil,\bF_u}(D_{\sigma},D_{\sigma})$.\;
\end{lem}
\begin{proof}
	We have the natural inclusion $\coker\nu^{?}_{\flat,1}\hookrightarrow\coker\nu^{?}_{\flat}$,\;but these two spaces have the same dimension,\;thus this inclusion is an isomorphism.\;
\end{proof}

It suffices to prove that
\[\sum_{1\neq u\in \sW_s}\mathrm{Ad}_{(u^{\sharp})^{-1}}(\fp_{S_0,L})\cap \mathrm{Ad}_{g_{L}}(\fb_{L})+\mathrm{Im}(\nu_{1})\cap \mathrm{Ad}_{g_{L}}(\fb_{L})=\mathrm{Im}(\nu).\]
But the inclusion $\mathrm{Ad}_{1}(\fp_{S_0,L})\cap \mathrm{Ad}_{g_{L}}(\fb_{L})/\mathrm{Im}(\nu_{1})\cap \mathrm{Ad}_{g_{L}}(\fb_{L})\hookrightarrow\mathrm{Ad}_{g_{L}}(\fb_{L})/\mathrm{Im}(\nu)$ is an isomorphism by using Lemma \ref{generalimagechangelem},\;now the result is direct consequence of the above discussion.\;



\begin{rmk}
	We have two induced filtrations on $D^{\flat}_{\sigma,(i)}$.\;Let  $\bd_i:=\dim_ED_{\sigma}^{\flat,(i)}\leq i$.\;
	\begin{itemize}
		\item[(1)] (Subspace filtration) We consider the induced subspace filtration $\fil_{\mathrm{sub}}^{\bullet}(D_{\sigma,(i)}^{\flat})$ of $\fil_H^{\bullet}(D_{\sigma})$ on $D_{\sigma,(i)}^{\flat}$.\;In an adapted basis of $D_{\sigma}=  D_{\sigma}^{\flat,(i)} \oplus D^{\flat}_{\sigma,(i)}$,\;by the non-critical assumption,\;we see that $\gr_{\mathrm{sub}}^{-\bh_j}(D^{\flat}_{\sigma,(i)})$ if and only if $1\leq j\leq n-\bd_i$ and 
		\begin{equation}
			\begin{aligned}
				\fil_{\mathrm{sub}}^{-\bh_j}(D^{\flat}_{\sigma,(i)})\cong \fil_H^{-\bh_{j}}(D_{\sigma})/\fil_H^{-\bh_{n-\bd_i+1}}(D_{\sigma})
			\end{aligned}
		\end{equation}
		when $1\leq j\leq n-d_{(i)}$.\;
		\item[(2)](Quotient filtration) We consider the induced subspace filtration $\fil_H^{\bullet}(D_{\sigma}^{\flat,(i)})$ of $\fil_H^{\bullet}(D_{\sigma})$ on $D_{\sigma}^{\flat,(i)}$ and the quotient filtration $\fil_{\mathrm{quot}}^{\bullet}(D^{\flat}_{\sigma,(i)})$ on $D^{\flat}_{\sigma,(i)}\cong D_{\sigma}/D_{\sigma}^{\flat,(i)}$.\;In an adapted basis of $D_{\sigma}=D^{\flat}_{\sigma,(i)}\oplus D_{\sigma}^{\flat,(i)} $,\;by the non-critical assumption,\;we see that 
		\begin{equation}
			\begin{aligned}
				\gr_H^{-\bh_{j}}(D_{\sigma}^{\flat,(i)})\neq 0,\quad (\text{resp.,\;}\gr_{\mathrm{quot}}^{\bullet}(D^{\flat}_{\sigma,(i)}))
			\end{aligned}
		\end{equation}
		if and only if $1\leq j\leq d_{(i)}$ (resp.,\;$d_{(i)}+1\leq j\leq n$).\;In particular,\;we have a strict exact sequence of filtered vector spaces:
		\[0\rightarrow \fil_H^{\bullet}(D_{\sigma}^{\flat,(i)})\rightarrow \fil_H^{\bullet}(D_{\sigma})\rightarrow \fil_{\mathrm{quot}}^{\bullet}(D^{\flat}_{\sigma,(i)})\rightarrow 0,\]
		such that 
		\begin{equation}\label{comparefilwithflat}
			\begin{aligned}
				\fil_H^{-\bh_{j}}(D_{\sigma}^{\flat,(i)})/\cong \fil_H^{-\bh_{j}}(D_{\sigma})/\fil_H^{-\bh_{d_{(i)}+1}}(D_{\sigma}),\quad (\text{resp.,\;}\fil_{\mathrm{quot}}^{\bullet}(D^{\flat}_{\sigma,(i)})\cong\fil_H^{-\bh_{j}}(D_{\sigma}))
			\end{aligned}
		\end{equation}
		when	$1\leq j\leq d_{(i)}$ (resp.,\;$d_{(i)}+1\leq j\leq n$).\;
	\end{itemize}
\end{rmk}

%\begin{rmk}
%In general,\;there is no neat fomula for the image in $(2)$ when $\flat=\flat$.\;Furthermore,\;we suspect that  $\ker(t^{\flat}_{D_{\sigma}}|_{\ext^1_{G,\mathrm{inf}}(\pi_{\lalg}(\underline{x},\bh),\pi^{\flat}_{S,\sigma}(\underline{x},\bh))})$  only presents numerous relationships between $p$-adic Hodge parameters.\;	
%\end{rmk}On the other hand,\;recall in (\ref{comparefilwithflat}) that
%\[\fil_H^{-\bh_l}(D^{\flat}_{\sigma,(j)})\cong\fil_H^{-\bh_{l}}(D_{\sigma})\]
%when $d_{(j)}+1\leq l\leq n$ and 
%\[\fil_H^{-\bh_l}(D^{\flat}_{\sigma,(j)})=\fil_H^{-\bh_{d_{(j)}}}(D^{\flat}_{\sigma,(j)})\]
%when $l<d_{(j)}$.\;Thus,\;we obtain that
%\begin{equation}
%	\left\{
%	\begin{array}{ll}
	%		f_{|S|}\in \homo_E\Big(D_{\sigma,(j)}^{\flat},\fil_H^{-\bh_{i_{l_{|S|}}+1}}(D_{\sigma})\Big):f_j\Big(\fil_H^{-\bh_{i_{l_{|S|-1}}+1}}/\fil_H^{-\bh_{i_{l_{|S|}}+1}}(D_{\sigma,(j)}^{\flat})\Big)=0,\;&j=|S|\\
	%		f_j\in\homo_E\Big(D_{\sigma,(j)}^{\flat},\fil_H^{-\bh_{i_{l_{j}}+1}}(D_{\sigma})\Big):f_j\Big(\fil_H^{-\bh_{i_{l_{j-1}}+1}}/\fil_H^{-\bh_{i_{l_{j}}+1}}(D_{\sigma,(j)}^{\flat})\Big)\subseteq\fil_H^{-\bh_{i_{l_{j+1}}+1}}(D_{\sigma,(j)}^{\flat}),\;&j<|S|
	%	\end{array}
%	\right.
%\end{equation}

%\begin{rmk}
%	
%	For example,\;if $\Delta\backslash S_0=\{a,a+b\}$ such that  $n=a+b+c$.\;Write $D_{\sigma}=D_{\sigma}^a\oplus D_{\sigma}^b\oplus D_{\sigma}^c,$Suppose that $a<b<c$ and $c\neq a+b$.\;Then $\Delta'=\{a,b,c,a+b,a+c,b+c\}$.\;If $S=\{a,b+c\}$,\;then $D_{\sigma,(a)}^{\flat}\cap D_{\sigma,(b+c)}^{\flat}=0$,\;so that $\ker(g_{\sigma})_S=0$.\;If $S=\{a,a+b\}$,\;then 
%	\begin{equation}
	%		D_{\sigma,(a)}^{\flat}=D_{\sigma}^b\oplus D_{\sigma}^c,\;D_{\sigma,(a+b)}^{\flat}=D_{\sigma}^c.\;
	%	\end{equation}
%Then we have
%\begin{equation}
%	\begin{aligned}
	%		&\Big\{f_{2}\in \homo_E\Big(D_{\sigma}/\big(D_{\sigma}^{\flat,(a+b)}+\fil_H^{-\bh_{a+b+1}}\big),\fil_H^{-\bh_{a+b+1}}\Big):\;f_{2}\Big(\fil_H^{-\bh_{a+1}}/\fil_H^{-\bh_{a+b+1}}\Big)=0\Big\}\\
	%	&\Big\{f_1\in\homo_E\Big(D_{\sigma}/\big(D_{\sigma}^{\flat,(a)}+\fil_H^{-\bh_{a+1}}\big),\fil_H^{-\bh_{a+1}}\Big):f_1\Big(D_{\sigma}/\fil_H^{-\bh_{a+1}}\Big)\subseteq\fil_H^{-\bh_{a+b+1}}\Big\},\;\\
	%	\end{aligned}
%\end{equation}
%
%\end{rmk}


%For $1\leq q\leq |S|$,\;let $V_q\subseteq \cS_{i_{l_{q}}}$ such that $V_{|S|}\subseteq V_{|S|-1}\subseteq\cdots \subseteq V_1$.\;Let 
%\begin{equation}
%	\begin{aligned}
	%		    &\cH_{V_j}:=g_{j,\sigma}^{-1}\Big(\left\{f\in\homo_E\big(D_{\sigma}[V_j],\fil_H^{-\bh_{i+1,\sigma}}(D_{\sigma})\big):f|_{\fil_H^{-\bh_{i+1,\sigma}}(D_{\sigma})} \text{scalar}\right\}\Big),\\
	%			&\cH_{\{V_q\}_{1\leq q\leq |S|}}:=g_{\sigma}^{-1}\Big(\sum_{j\in S}\left\{f\in\homo_E\big(D_{\sigma}[V_j],\fil_H^{-\bh_{i+1,\sigma}}(D_{\sigma})\big):f|_{\fil_H^{-\bh_{i+1,\sigma}}(D_{\sigma})} \text{scalar}\right\}\Big)=\oplus_{1\leq q\leq |S|}\cH_{V_j},\\
	%			&	\ext^1_{G,\sigma}\Big(\pi_{\lalg}(\underline{x},\bh),\pi^{\flat}_{S,\sigma}(\underline{x},\bh)\Big)_{\{V_q\}_{1\leq q\leq |S|}}:=\beta_S^{-1}\big(\cH_{\{V_q\}_{1\leq q\leq |S|}}\big).\;
	%	\end{aligned}
%\end{equation}

We draw the matrix of each $f_j\in \homo_E(D_{\sigma},D_{\sigma})$ (note that $f_j\in \homo_E(D_{\sigma,(j)}^{\flat},D_{\sigma})$ in our situation) in an adapted basis of $D_{\sigma}$ for the filtration $\fil^{-\bh_{i}}_H(D_{\sigma})$,\;then we get the matrix form of $f_{i_{l_1}}+\cdots+f_{i_{l_{|S|}}}$.\;If $f_{i_{l_1}}+\cdots+f_{i_{l_{|S|}}}=0$,\;then $f_{j}$ also satisfies that  (see \cite[(102)]{BDcritical25})
\begin{equation}
	\begin{aligned}
		&\Big\{f_{|S|}\in \homo_E\Big(D_{\sigma}/\big(D_{\sigma}^{\flat,(i_{l_{|S|}})}+\fil_H^{-\bh_{i_{l_{|S|}}+1}}\big),\fil_H^{-\bh_{i_{l_{|S|}}+1}}\Big):\;f_{|S|}\Big(\fil_H^{-\bh_{i_{l_{|S|-1}}+1}}/\fil_H^{-\bh_{i_{l_{|S|}}+1}}\Big)=0\Big\}\\
		%			&\Big\{f_{|S|}\in \homo_E\Big(D_{\sigma,(j)}^{\flat}/\fil_H^{-\bh_{i_{l_{|S|}}+1}}(D_{\sigma,(j)}^{\flat})\big),\fil_H^{-\bh_{i_{l_{|S|}}+1}}\Big):\;f_{|S|}\Big(\fil_H^{-\bh_{i_{l_{|S|-1}}+1}}/\fil_H^{-\bh_{i_{l_{|S|}}+1}}(D_{\sigma,(j)}^{\flat})\Big)=0\Big\}\\
	\end{aligned}
\end{equation}
when $j=|S|$,\;and 
\begin{equation}\label{landsinsubspace}
	\begin{aligned}
		&\Big\{f_j\in\homo_E\Big(D_{\sigma}/\big(D_{\sigma}^{\flat,(i_{l_{j}})}+\fil_H^{-\bh_{i_{l_{j}}+1}}\big),\fil_H^{-\bh_{i_{l_{j}}+1}}\Big):f_j\Big(\fil_H^{-\bh_{i_{l_{j-1}}+1}}/\fil_H^{-\bh_{i_{l_{j}}+1}}\Big)\subseteq\fil_H^{-\bh_{i_{l_{j+1}}+1}}\Big\},\;\\
		%		&\Big\{f_j\in\homo_E\Big(D_{\sigma,(j)}^{\flat}/\fil_H^{-\bh_{i_{l_{j}+1}}}(D_{\sigma,(j)}^{\flat}),\fil_H^{-\bh_{i_{l_{j}}+1}}\Big):f_j\Big(\fil_H^{-\bh_{i_{l_{j-1}}+1}}/\fil_H^{-\bh_{i_{l_{j}}+1}}(D_{\sigma,(j)}^{\flat})\Big)\subseteq\fil_H^{-\bh_{i_{l_{j+1}}+1}}(D_{\sigma,(j)}^{\flat})\Big\},\;
	\end{aligned}
\end{equation}
when $j<|S|$.\;Then the same proof of \cite[(102)]{BDcritical25} shows that the image of $\ker(g_{\sigma})$ must lands in the subspace of $(2)$.\;

\begin{pro}
	The  image  under (\ref{compositionfinalSver}) of the subspace $\ker(t^{\flat}_{D_{\sigma}}|_{\ext^1_{G,\mathrm{inf}}(\pi_{\lalg}(\underline{x},\bh),\pi^{\flat}_{S,\sigma}(\underline{x},\bh))_{\{V_q\}_{1\leq q\leq |S|}}})$ is the image of the composition
	\[\homo_E\Big(\big(\bigwedge\nolimits_{E}^{\!n-i}\!D_{\sigma}\big)^{\flat}/\fil_{S,i}^{2^{\mathrm{nd}}-\max}(D_{\sigma})^{\flat},\fil_i^{\max}(D_{\sigma})^{\flat}\Big)\]
\end{pro}


\begin{rmk}
	Applying $g_{i,\sigma}$ to $D_{\sigma,(i)}^{\flat}$,\;we have 
	\begin{equation*}
		\begin{aligned}
			g_{i,\sigma}:&\;\homo_E\Big(\bigwedge\nolimits_E^{\!n-i}\!D_{\sigma,(i)}^{\flat},\bigwedge\nolimits_E^{\!n-i}\!\fil_{\mathrm{quot}}^{-\bh_{i+1,\sigma}}(D_{\sigma,(i)}^{\flat})\Big)\\&\twoheadrightarrow \homo_E\Big(\bigwedge\nolimits_E^{\!n-i-1}\!\fil_{\mathrm{quot}}^{-\bh_{i+1,\sigma}}(D_{\sigma,(i)}^{\flat})\wedge D_{\sigma,(i)}^{\flat} ,\bigwedge\nolimits_E^{\!n-i}\!\fil_{\mathrm{quot}}^{-\bh_{i+1,\sigma}}(D_{\sigma,(i)}^{\flat})\Big)\\
			&\xrightarrow{\sim}\left\{f\in\homo_E\big(D_{\sigma,(i)}^{\flat},\fil_{\mathrm{quot}}^{-\bh_{i+1,\sigma}}(D_{\sigma,(i)}^{\flat})\big):f|_{\fil_{\mathrm{quot}}^{-\bh_{i+1,\sigma}}(D_{\sigma,(i)}^{\flat})} \text{scalar}\right\}\hookrightarrow \homo_{\fil}\big(D_{\sigma,(i)}^{\flat},D_{\sigma,(i)}^{\flat}\big).\;
		\end{aligned}
	\end{equation*}
	
\end{rmk}


\subsubsection{Relation to  the functors constructed in \cite{breuil2020probleme}}
This section gives a relationship between our discussion with the functors constructed in \cite{breuil2020probleme}.\;Keep the notation in \cite[Conjecture 5.3.1]{breuil2020probleme} and assume that $L=\bQ_p$.\;Put $\cR_E:=\cR_{E,\bQ_p}$.\;For $i\in \Delta'$,\;set $\alpha_i=e_{i}-e_{i+1}$ and 
\[(D(\wedge_E^{n-i}\wdre_{\Dpik}))^{\flat}:=\bigoplus_{u\in \cI_{i}^{\flat}}\wedge_{i\in I_i(u)^c}(\wedge_E^{r_i}\Delta_{E_i}),\]
which is a $(\varphi,\Gamma)$-submodule of the $p$-adic differential equation $D(\wedge_E^{n-i}\wdre_{\Dpik})$ associated to the Weil-Deligne representation $\wedge_E^{n-i}\wdre_{\Dpik}$.\;By \cite[Theorem 1.3]{breuil2020probleme},\;for $u\in \cI_{i}^{\flat}$,\;we obtain (our $\lambda=(\lambda_1,\cdots,\lambda_n)$ is the $``-\lambda"$ in the definition before \cite[(216)]{breuil2020probleme})
\begin{equation}
	\begin{aligned}
		F_{\alpha_i}(C_{i,u})&\simeq E_{\infty}(\chi_{\lambda})\otimes_E\homo_{(\varphi,\Gamma)}\Big(\cR_{E}((-\lambda)\lambda_{\alpha^{\vee}_i})/t^{1-\langle-\lambda,\alpha^{\vee}_i\rangle},-\Big),\\
		&\simeq E_{\infty}(\chi_{\lambda})\otimes_E\homo_{(\varphi,\Gamma)}\Big(\cR_{E}(z^{-\sum_{j=1}^i\lambda_i})/t^{\bh_i-\bh_{i+1}},-\Big)
	\end{aligned}
\end{equation}
where $\lambda_{\alpha^{\vee}_i}(x)=\diag(x,\cdots,x,1,\cdots,1)$ (the number of $x$ is $i$) for $x\in \BG_m$.\;Moreover,\;put (we refer to \cite[(217)]{breuil2020probleme})
\begin{equation}
	\begin{aligned}
		D_{\alpha_i}(\underline{x},\bh)^{\star}:=(D(\wedge_E^{n-i}\wdre_{\Dpik}))^{\star}\otimes_{\cR_E}\cR_E((-s_i\cdot(-\lambda))\lambda_{\alpha^{\vee}_i} )\otimes_{\cR_E}\cR_E(|\cdot|^{\sum_{j=n-i}^{n-1}}\chi^{-1}_{\wdre_{\Dpik}})
	\end{aligned}
\end{equation}
for $\star\in\{\emptyset,\flat\}$.\;Note that $D_{\alpha_i}(\underline{x},\bh)^{\flat}\hookrightarrow D_{\alpha_i}(\underline{x},\bh)$,\;which induces an injection 
\[\ext^1_{(\varphi,\Gamma)}\Big(\cR_{E}(z^{-\sum_{j=1}^i\lambda_i})/t^{\bh_i-\bh_{i+1}},D_{\alpha_i}(\underline{x},\bh)^{\flat}\Big)\hookrightarrow \ext^1_{(\varphi,\Gamma)}\Big(\cR_{E}(z^{-\sum_{j=1}^i\lambda_i})/t^{\bh_i-\bh_{i+1}},D_{\alpha_i}(\underline{x},\bh)\Big).\]
We have a natural isomorphism 
$\ext^1_{(\varphi,\Gamma)}\Big(\cR_{E}(z^{-\sum_{j=1}^i\lambda_i})/t^{\bh_i-\bh_{i+1}},D_{\alpha_i}(\underline{x},\bh)\Big)\xrightarrow{\sim}\bigwedge\nolimits_{E}^{\!n-i}\!D$ by \cite[Remark 5.1.3,\;(218)]{breuil2020probleme}.\;This map identities  $\ext^1_{(\varphi,\Gamma)}\Big(\cR_{E}(z^{-\sum_{j=1}^i\lambda_i})/t^{\bh_i-\bh_{i+1}},D_{\alpha_i}(\underline{x},\bh)^{\flat}\Big)$ with the subspace $\big(\bigwedge\nolimits_{E}^{\!n-i}\!D\big)^{\flat}$ of $\bigwedge\nolimits_{E}^{\!n-i}\!D$ exactly.\;Put $\pi^{\alpha_i,\flat}(\underline{x},\bh):=\bigoplus_{u\in \cI_{i}^{\flat}}C_{i,u}$,\;we deduce from \cite[Theorem 1.3]{breuil2020probleme} that
\begin{equation}
	\begin{aligned}
		F_{\alpha_i}(\pi^{\alpha_i,\flat}(\underline{x},\bh))&\simeq E_{\infty}(\chi_{\lambda})\otimes_E\homo_{(\varphi,\Gamma)}\Big(D_{\alpha_i}(\underline{x},\bh)^{\flat},-\Big)
	\end{aligned}
\end{equation}
and obtain an isomorphism:
\begin{equation}
	\cE_{\alpha_i}:\ext^1_{G}\Big(\pi^{\alpha_i,\flat}(\underline{x},\bh),\pi_{\lalg}(\underline{x},\bh)\Big)\xrightarrow{\sim}\ext^1_{(\varphi,\Gamma)}\Big(\cR_{E}(z^{-\sum_{j=1}^i\lambda_i})/t^{\bh_i-\bh_{i+1}},D_{\alpha_i}(\underline{x},\bh)^{\flat}\Big).\;
\end{equation}

